\documentclass[11pt,a4paper]{article}
\usepackage[T1]{fontenc}
\usepackage{lmodern}
\usepackage[a4paper,margin=27mm,headheight=14pt]{geometry}
\usepackage{amsmath,amssymb,amsthm,mathrsfs,mathtools}
\usepackage{array,booktabs,longtable,tabularx,enumitem,graphicx}
\usepackage{xcolor,microtype}
\usepackage{hyperref}
\hypersetup{colorlinks=true,linkcolor=blue!45!black,citecolor=blue!45!black,
 urlcolor=blue!45!black,pdftitle={LM7: The Fermat tower in the metabelian approximation},
 pdfauthor={GPT-6 Astra, OpenAI, ChatGPT Pro mode}}
\urlstyle{same}
\setlength{\parindent}{1.3em}
\setlength{\parskip}{.22em}
\setlength{\emergencystretch}{2em}
\setlength{\LTpre}{.6em}
\setlength{\LTpost}{.6em}
\setlist{nosep}
\setcounter{secnumdepth}{0}
\setcounter{tocdepth}{1}
\allowdisplaybreaks[2]
\newenvironment{statement}[3]{%
 \par\addvspace{1.1em}\phantomsection\label{st:#1}%
 \noindent\textbf{#1. #3}\quad\textsf{\small[#2]}\par\nobreak\smallskip\ignorespaces
}{\par\smallskip}
\newcommand{\statementref}[1]{\hyperref[st:#1]{#1}}
\newcommand{\Sha}{\mathop{\mathrm{Sh}}}
\begin{document}
\begin{center}
 {\LARGE\bfseries The Fermat tower in the metabelian approximation}\par
 \vspace{.6em}
 {\large Grothendieck, \emph{La Longue Marche}, \S51 and its refinements}\par
 \vspace{.5em}
 LM7\_FERMAT\_TOWER\quad ---\quad 13 November 2026\par
 \vspace{.4em}
 AI-produced by GPT-6 Astra (OpenAI, ChatGPT Pro mode).\par
 No priority is claimed.\par
 \vspace{.6em}
 \textbf{INCOMPLETE: the exact residual questions are retained in the paper and ledger.}
\end{center}

\noindent We give proofs and reproducible calculations for the integral schemes,
the Fermat modules and splitting, the cuspidal motive, the arithmetic
Kummer images, and the covering comparisons. We preserve the arithmetic
plus-Fermat model of the manuscript. Full profinite realization questions
and the remaining cubic section calculation are explicitly unresolved;
the finite and rational calculations below do not replace those questions.

\medskip
\noindent\textbf{Notation fixed on this page.} C4 means cote 140-4. A.1
gives the expanded dictionary, including the changing early and final versions.
\smallskip

{\small\renewcommand{\arraystretch}{1.18}
\noindent\begin{tabularx}{\linewidth}{@{}>{\raggedright\arraybackslash}p{.19\linewidth}>{\raggedright\arraybackslash}X >{\raggedright\arraybackslash}p{.20\linewidth}@{}}
\toprule
Our notation & Meaning and manuscript symbol & C4 scan; folio\\
\midrule
$\Pi_0,F,M$ & $F\xrightarrow{\sim}\Pi_0$, $x\mapsto\lambda_0$, $y\mapsto\lambda_1$; $M=\Pi_0^{\rm ab}$ & 34--38; 664--666\\
$\mathfrak S_{\rm d},\mathcal S$ & $\mathrm{SL}_2(\mathbb Z)/\Pi_0'$ and its indicated finite schematic stage; his $\mathfrak S$ and underlined $\mathcal G^+$ & 34--38; 664--666\\
$\gamma=\mathbb O_{12}^{*}$ & Units of his ring scheme $O_{12}$, with multiplier $\mu$ & 51--58; 673--676v\\
$N,\Gamma,\mathcal M$ & His corrected $N_{\rho,\sigma}$ and final $\Gamma_{\rho,\sigma},M_{\rho,\sigma}$ & 62, 85--88; 678v, 690--691v\\
$X_n,U_n$ & $X_n:x_0^n+x_1^n+x_\infty^n=0$; $U_n$ deletes its coordinate cusps & 97--99; 696--697\\
$\Lambda,c$ & $\widehat{\mathbb Z}[[\widehat{\mathbb Z}^{\,2}]]$; a specified generator of $\widehat F'/\widehat F''$ & comparison with 95--99; 695--697\\
$G_{\mathbb Q},\chi,\kappa_a$ & Arithmetic $\mathbb I\Gamma_{\mathbb Q}$; cyclotomic multiplier; Kummer cocycle with chosen roots & 156--170; 725v--732v\\
\bottomrule
\end{tabularx}}

\newpage
\tableofcontents

\section{Scope, choices, and verification}
We use the supplied second-pass transcriptions and inspect the available
scans at the mathematical decision points recorded in the readings ledger.
A transcription is evidence for a reading, and a proved identity is evidence
for the mathematical meaning; the latter does not settle unread handwriting.
The manuscript passage for each numbered statement is also recorded in
\texttt{LM7\_FERMAT\_TOWER\_results.csv}. All named groups in an argument
carry the discrete, profinite, pro-$\ell$, or algebraic interpretation stated
there. Prime-to-characteristic and fppf assertions are specified separately.

We label a proved calculation \emph{Proved}, a false proposed assertion
together with its proved replacement \emph{Refuted}, and a reduction to an
opened theorem \emph{Known}. An \emph{Open} label records an uncompleted step
of this return, with its proved cases and attempted approaches. Such a label
does not assert that the literature has no answer. The report has a row for
each of the three items, three leads, lettered tasks, and follow-ups in the
request. The bibliography records which original texts were opened; a
bibliographic identification alone is labelled unverified.

The programs use integer arithmetic, exact rational linear algebra, finite
permutation groups, and finite-field arithmetic. Their saved standard output
and machine-readable witnesses accompany the source. Finite verification is
used with the displayed general proofs, and a Lie-algebra computation is
transferred to a group only where the requisite comparison is proved.
The source and complete reproduction instructions are in \texttt{code/}.

\section{A. Integral structure, level twelve, and the congruence actions}

We use constant finite group schemes, including in residue characteristic
dividing their order. Thus $C_2$ below means the disjoint union of two
copies of $\operatorname{Spec}\mathbb Z$; the morphism
$C_2\to\mu_2$, $i\mapsto(-1)^i$, will always be written explicitly when
it is used. All scheme exact sequences in this section are exact for the
fppf topology. Our finite-stage notation distinguishes a group from its
schematic envelope by specifying their comparison morphism.

\begin{statement}{A.1}{Proved}{Dictionary and the finite-stage comparison maps}
Fix the following matrices, over $\mathbb Z$:
\[
 \rho=\begin{pmatrix}0&1\\-1&1\end{pmatrix},\qquad
 \sigma=\begin{pmatrix}0&-1\\1&0\end{pmatrix},\qquad
 \tau=\begin{pmatrix}-1&0\\0&1\end{pmatrix},\qquad \omega=-I.
\]
Put
\[
 \lambda_0=\begin{pmatrix}-1&2\\0&-1\end{pmatrix},\quad
 \lambda_1=\begin{pmatrix}-1&0\\-2&-1\end{pmatrix},\quad
 \Pi_0=\langle\lambda_0,\lambda_1\rangle,
 \quad M=\Pi_0^{\mathrm{ab}}=\mathbb Ze_0\oplus\mathbb Ze_1.
\]
The following dictionary fixes the successive versions in \S51.
\end{statement}

\begin{longtable}{@{}p{.18\linewidth}p{.56\linewidth}p{.18\linewidth}@{}}
\toprule
Our symbol & Definition and manuscript symbol & C4 scan; folio\\
\midrule\endhead
$\mathfrak S_{\rm d}$ & $\mathrm{SL}_2(\mathbb Z)/\Pi_0'$, the determinant-one part of the discrete $\mathfrak S$; $\Pi_0'=[\Pi_0,\Pi_0]$ & 33--34; 664--664v\\
$\mathcal S$ & Pushout of $1\to M\to\mathfrak S_{\rm d}\to H\to1$ along $M\to\underline M=\mathbb G_a^2$; his $\underline{\mathcal G}^{+}$ & 34--38; 664--666\\
$H$ & $\mathrm{SL}_2(\mathbb Z)/\Pi_0$, of order $12$; the full determinant-sign quotient has order $24$ & 33, 66; 664, 680v\\
$\mathbb O_N$ & Additive pushout $\mathbb Z_{\rm const}\amalg_{N\mathbb Z}\underline{N\mathbb Z}$, with its ring law in A.5 & 51--54; 673--674v\\
$\gamma$ & $\mathbb O_{12}^{*}$; $\mu:\gamma\to\mathbb G_m$ is evaluation; $\varepsilon_2,\varepsilon_3:\gamma\to C_2$ are the two residue characters & 55--58; 675--676v\\
$M_0$ & $\{c_0e_0+c_1e_1:c_0+c_1\equiv0\pmod3\}$, a lattice with its own vector group $\underline M_0$ & 60--62; 677v--678v\\
$N$ & The corrected $\underline N_{\rho,\sigma}$, with the loop equation in $\underline M_0$ & 62; 678v\\
$Z,G$ & $\underline Z_{\rho,\sigma}$ and $\underline G_{\rho,\sigma}$, namely $\mathbb G_a\rtimes\gamma$ and $\mathcal S\rtimes\gamma$ & 75--76; 685--685v\\
$\Gamma_{\rm e},\mathcal M_{\rm e}$ & The versions of $\Gamma_{\rho,\sigma},\underline M_{\rho,\sigma}$ in (123)--(127) & 70--72; 682v--683v\\
$N_\rho,N_\rho^{?}$ & The smooth closure of the $\rho$-normalizer, and its final bisection (212); A.13 records the smaller literal version (171) & 78--82, 94; 686v--688v, 694v\\
$N_\sigma,N_\sigma^\sharp$ & $\sigma$-normalizer and its intersection with $G^\sharp=\underline M\rtimes\gamma$ & 83--84; 689--689v\\
$\Gamma,\mathcal M$ & The final fibre products (183)--(184), retaining the central $C_2$ & 85--88; 690--691v\\
$\mathcal M(0)$,\newline $\mathcal M(1/2)$,\newline $\mathcal M(-\bar j)$ & The three graph subgroups in A.15 & 88--94; 691v--694v\\
\bottomrule
\end{longtable}

\begin{proof}
The assertions about the discrete groups are proved in A.2. A section
$h\mapsto s_h$ of $\mathfrak S_{\rm d}\to H$ gives the finite-stage
scheme $\coprod_{h\in H}\mathbb A^2_{\mathbb Z}$, with multiplication
\[
 (m,h)(n,k)=(m+a_hn+c(h,k),hk),\qquad
 c(h,k)=s_hs_ks_{hk}^{-1}\in M.
\]
Changing the section to $d_hs_h$ gives the explicit isomorphism
$(m,h)\mapsto(m-d_h,h)$. Thus this construction has a specified,
section-independent meaning. An integral cocycle is obtained from
\[
 \rho\longmapsto(R,0),\quad
 \sigma\longmapsto(S,(1/3,-1/3)),\qquad
 R=\begin{pmatrix}0&-1\\1&-1\end{pmatrix},\quad
 S=\begin{pmatrix}0&1\\1&0\end{pmatrix}.
\]
Choose words $s_h$ lifting the twelve elements of $H$ and let $t_h$ be
their affine translations. Then
$c(h,k)=t_h+a_ht_k-t_{hk}\in\mathbb Z^2$, since the indicated word lies
in $M$. The affine representation kills $\omega$; the retained component
$h$ records it. The pair consisting of the component and affine image is
faithful: on its possible kernel in $M$, the affine image is translation
by the lattice vector itself.

The envelope $E_{\mathfrak S_{\rm d}}$ of \S50 is the inverse limit of
the analogous stages obtained from $nM$. The transition from any of
these stages to the stage above is induced by the inclusion
$nM\hookrightarrow M$ and is a morphism of affine group schemes.
It gives the comparison $E_{\mathfrak S_{\rm d}}\to\mathcal S$; the
formulas below concern this specified finite stage. On profinite
points the continuous map
$\widehat{\mathfrak S_{\rm d}}\to\mathcal S(\widehat{\mathbb Z})$
uses one common $H$-component and the coordinate maps
$\widehat M\to M\otimes\mathbb Z_p$ at every prime. This describes its
image and the meaning of the diagonal points used later in the
manuscript. Likewise $\widehat{\mathbb Z}^{*}\to\gamma(\widehat{\mathbb Z})$
sends $\mu$ to its single residue class $r\bmod12$ and the unique
$k\in\widehat{\mathbb Z}$ with $\mu=r+12k$. The proof of A.5 shows
directly that this is a continuous homomorphism. These comparison maps
retain the integral and profinite data rather than identifying all
points of the schemes with one profinite group.
\end{proof}

\begin{statement}{A.2}{Proved}{The subgroup, its free basis, and the action of the quotient}
The subgroup $\Pi_0$ is normal and free on the displayed
$\lambda_0,\lambda_1$. Its indices in $\mathrm{SL}_2(\mathbb Z)$ and
$\mathrm{GL}_2(\mathbb Z)$ are $12$ and $24$. Moreover
\[
 H\simeq\mathrm{SL}_2(\mathbb Z/4)/V_4\simeq\mathrm{Dic}_3,
 \qquad V_4=\langle\lambda_0\bmod4,\lambda_1\bmod4\rangle.
\]
Conjugation on $M$ sends $\rho,\sigma,\tau$ respectively to $R,S,-I$;
in particular $\tau'=\omega\sigma\tau$ acts as $-S$.
For every divisor $N$ of $24$, the action of
$\mathrm{GL}_2(\mathbb Z)/\Pi_0$ on $M/NM$ has the following orders.
\[
\begin{array}{c|rrrrrrrr}
N&1&2&3&4&6&8&12&24\\\hline
|\operatorname{im}|&1&6&12&12&12&12&12&12\\
|\ker|&24&4&2&2&2&2&2&2
\end{array}
\]
This proves the index and lattice-action portion of Lead 2
(\S51.1--4; C4 scans 34--38, 64--67; fols. 664--666, 679v--681).
\end{statement}
\begin{proof}
Matrix multiplication gives
\[
 \lambda_0=\omega\sigma\rho\sigma\rho,\quad
 \lambda_1=\omega\rho\sigma\rho\sigma,
\]
\[
 \rho\lambda_0\rho^{-1}=\lambda_1,\quad
 \rho\lambda_1\rho^{-1}=(\lambda_1\lambda_0)^{-1},\quad
 \sigma\lambda_0\sigma^{-1}=\lambda_1,\quad
 \sigma\lambda_1\sigma^{-1}=\lambda_0.
\]
Conjugation by $\tau$ inverts each $\lambda_i$. These identities prove
normality and the stated action after abelianization. We use the
amalgam presentation
$\mathrm{SL}_2(\mathbb Z)=C_4*_{C_2}C_6$, with factors generated by
$\sigma,\rho$ and $\sigma^2=\rho^3=\omega$; it is recorded in
\cite[\S4, equation (1)]{A_KirschmerLeedhamGreen}. Adding $\lambda_0=1$
forces $\sigma\rho\sigma^{-1}=\rho^{-1}$; $\lambda_1=1$ is its
conjugate. The resulting group has the twelve normal forms
$\rho^i\sigma^j$, $0\leq i<6$, $0\leq j<2$.
Modulo four, $V_4$ has four elements, is normal, and the images of
$\rho,\sigma$ in $\mathrm{SL}_2(\mathbb Z/4)/V_4$ retain orders six
and four. Since $|\mathrm{SL}_2(\mathbb Z/4)|=48$, this realizes all
twelve normal forms and proves the index claim and the stated quotient.
Surjectivity of the reduction map follows also from elementary row
operations over $\mathbb Z/4$.

The Bass--Serre tree of the amalgam has trivial $\Pi_0$-stabilizers:
each factor, and therefore every conjugate of a factor, injects in
$H$. Its quotient by $\Pi_0$ has $12/4+12/6=5$ vertices and
$12/2=6$ edges. It is connected. Choosing a maximal tree shows directly
that its fundamental group is free of rank $6-5+1=2$, and lifting paths
identifies this group with $\Pi_0$. Since $\lambda_0,\lambda_1$ generate
it, they are a free basis. For completeness, a surjective endomorphism
of a finitely generated free group is injective: free groups are
residually finite by completing the finite labelled path of a nonempty
reduced word to permutations; on the finite set of subgroups of any
fixed finite index, inverse image under a surjection is an injective
permutation. A nontrivial element in its kernel would then lie in every
finite-index subgroup, a contradiction. This applies to the map from
a chosen free basis to the displayed generating pair.

The matrices $R,S$ satisfy $R^3=S^2=(RS)^2=I$ and give the six
distinct matrices of $S_3$ modulo every $N\geq2$. The additional
$-I$ is already $I$ modulo two. For $N\geq3$, $-I$ is outside those
six matrices: direct inspection of
$I,R,R^2,S,RS,R^2S$ proves this over $\mathbb Z/N$.
Thus the image has order twelve for each displayed $N\geq3$.
The central element $\omega$ is the order-two kernel integrally, and
the table follows. The program \texttt{a\_finite.py} constructs the
twenty-four cosets modulo four and propagates the matrix action by a
coupled breadth-first traversal; it also constructs the above quotient
graph and checks its connectivity and Euler characteristic. This is a
finite certificate for the presentation argument, not a finite test
used in place of freeness.
\end{proof}

\begin{statement}{A.3}{Proved}{The precise source of level twelve and the group of order thirty-two}
Let $\Pi_{00}$ be the inverse image in $\Pi_0$ of $M_0$, and put
$\delta=\lambda_1\lambda_0^{-1}$. Then
\[
 \Pi_{00}=\langle\!\langle\delta\rangle\!\rangle_{\mathrm{SL}_2(\mathbb Z)}
 =\Pi_0\cap\mathrm{SL}_2(\mathbb Z)',\qquad
 [\mathrm{SL}_2(\mathbb Z):\Pi_{00}]=36.
\]
Its congruence level is exactly $12$, and the reduction map gives
\[
 K_{32}:=\Pi_{00}/\Gamma(12)
 =V_4\times Q_8
 \ \subset\ \mathrm{SL}_2(\mathbb Z/4)\times\mathrm{SL}_2(\mathbb F_3)
 =\mathrm{SL}_2(\mathbb Z/12).
\]
Here $Q_8$ is the quaternion subgroup of $\mathrm{SL}_2(\mathbb F_3)$.
The subgroup $K_{32}$ is normal in $\mathrm{GL}_2(\mathbb Z/12)$.
This is the group of order $32$ asked for at C4 scan 64, folio 679v;
its exact congruence origin is the additional index-three lattice
$M_0\subset M$, together with the level-four realization of $\Pi_0$.
\end{statement}
\begin{proof}
With $e_i=[\lambda_i]$ we have
\[
 M_0=\mathbb Z(e_1-e_0)+\mathbb Z(3e_1)
     =(1-R)M+(1-S)M.
\]
The normal closure $J$ of $\delta$ is contained in $\Pi_{00}$.
In $\Pi_0/J$, the relation $\lambda_0=\lambda_1$ and its $\rho$-conjugate
give $\lambda_1=(\lambda_1\lambda_0)^{-1}$, so the quotient is cyclic
of order at most three. The map $M\to\mathbb Z/3$, $(c_0,c_1)\mapsto
c_0+c_1$, shows that its order is exactly three. Hence $J=\Pi_{00}$.
Abelianizing the amalgam gives $\mathrm{SL}_2(\mathbb Z)^{\rm ab}
\simeq\mathbb Z/12$, with $[\rho]=2$ and $[\sigma]=3$ after choosing
a generator. Each $[\lambda_i]=4$. The restriction to $\Pi_0$ therefore
has kernel precisely $\Pi_{00}$, proving the intersection formula.
In particular the two occurrences of $12$ in Lead 2 are related by the
exact sequence $0\to C_3\to C_{12}\to H^{\rm ab}=C_4\to0$.

Chinese remainders give the displayed product decomposition. The
matrix of $\delta$ is
\[
 \delta=\begin{pmatrix}1&2\\2&5\end{pmatrix},\qquad
 \delta^2\equiv5I\pmod{12}.
\]
Modulo four its normal closure is $V_4$: conjugation cyclically
permutes the three nonzero elements. Modulo three its normal closure
is the eight-element group consisting of $\pm I$ and the six
trace-zero determinant-one matrices; each of those six squares to
$-I$. It is closed, as one checks on two anticommuting generators,
and is the quaternion group. To see that the normal closure in the
product is the whole product, take commutators of $\delta$ with
$(g,1)$: these generate $V_4\times1$. We can then isolate
$(1,\delta\bmod3)$, whose normal closure is $1\times Q_8$.
Consequently the image of $\Pi_{00}$ is $V_4\times Q_8$.
Its full inverse image has index
$48\cdot24/32=36$, equal to the already proved index of $\Pi_{00}$;
the two subgroups coincide. This proves $\Gamma(12)\subset\Pi_{00}$.

For minimality, the upper transvection
$T^m=\left(\begin{smallmatrix}1&m\\0&1\end{smallmatrix}\right)$
belongs to $\Pi_{00}$ only when $4\mid m$, because the only such
element of $V_4$ is $I$, and only when $3\mid m$, because $Q_8$ has
no nonidentity element of order three. Every $\Gamma(m)$ contains
$T^m$, so $\Gamma(m)\subset\Pi_{00}$ implies $12\mid m$.
Finally $V_4$ is normalized by $\mathrm{GL}_2(\mathbb Z/4)$, and the
trace description of $Q_8$ is invariant under every invertible
conjugator. Their product is normal in the full level-twelve group.
All thirty-two matrices, their orders, their conjugates, and the
normal-closure assertions are exhaustively checked and printed by
\texttt{a\_finite.py}.
\end{proof}

\begin{statement}{A.4}{Proved}{All congruence arrows and the scope of Lead 2}
The exact finite quotients obtained from A.3 have orders
\[
\begin{array}{c|r|r}
\text{group}&\text{order}&\text{quotient by }K_{32}\\\hline
\mathrm{SL}_2(\mathbb Z/12)&1152&36\\
\operatorname{im}(\mathrm{GL}_2(\mathbb Z)\to\mathrm{GL}_2(\mathbb Z/12))&2304&72\\
\mathrm{GL}_2(\mathbb Z/12)&4608&144
\end{array}
\]
The middle image is $\det^{-1}\{1,-1\}$. Thus reduction induces
isomorphisms
\[
 \mathrm{SL}_2(\mathbb Z)/\Pi_{00}\simeq
 \mathrm{SL}_2(\mathbb Z/12)/K_{32},\qquad
 \mathrm{GL}_2(\mathbb Z)/\Pi_{00}\simeq
 \det^{-1}\{1,-1\}/K_{32}.
\]
There is also a surjective homomorphism
$\mathfrak S_{\rm d}\to\mathrm{SL}_2(\mathbb Z/12)/K_{32}$ with kernel
$M_0$. Reduction to $\mathrm{SL}_2(\mathbb Z/12)$ itself does not
factor through $\mathfrak S_{\rm d}$.
These are the actual morphisms relating the action on $M$ to the
order-$32$ subgroup. The action images in A.2 have order at most $12$;
none is the order-$32$ group. They determine the lattice $M_0$ by
$M_0=(1-R)M+(1-S)M$, whose inverse image is used in A.3.
This proves the strongest precise form of Lead 2 and the relation
requested in item 1 (C4 scans 64--67; fols. 679v--681).
\end{statement}
\begin{proof}
The determinant map on $\mathrm{GL}_2(\mathbb Z/12)$ is onto its four
units, with kernel of order $1152$. Integral invertible matrices have
determinant $\pm1$, and the elementary generators of
$\mathrm{SL}_2(\mathbb Z/12)$ lift to integral matrices. Adjoining
$\tau$ therefore gives exactly the middle group. Taking inverse
images of $K_{32}$ proves each isomorphism and each kernel claim.
For the last assertion, direct multiplication gives
\[
 [\lambda_0,\lambda_1]\equiv
 \begin{pmatrix}1&4\\4&5\end{pmatrix}\not\equiv I\pmod{12}.
\]
This element is killed in $\mathfrak S_{\rm d}$ but survives under
reduction, proving the asserted failure of factorization by an
explicit element. It dies after quotient by $K_{32}$, as the exact
kernel calculation requires. This also verifies the qualification in
the margin of scan 67. The scan 66 notation for the Sylow-two group
is transcribed as $D_4$; the six order-four elements and unique
involution prove that its abstract group is $Q_8$, independently of
the notation intended by that handwriting.
\end{proof}

\begin{statement}{A.5}{Proved}{The ring scheme and its exact unit sequence}
For a positive integer $N$, the ring scheme $\mathbb O_N$ is
$\coprod_{r\in\mathbb Z/N}\mathbb A^1_{\mathbb Z}$. A section is locally
written $n+N[x]$, with the equivalence
\[
 (n,x)\sim(n+Nk,x-k)\quad(k\in\mathbb Z).
\]
Addition is the additive pushout law, and multiplication is
\[
 (n+N[x])(m+N[y])=nm+N[ny+mx+Nxy].                    \tag{A.5.1}
\]
Its unit scheme and unit exact sequence are
\[
 \mathbb O_N^*=\coprod_{r\in(\mathbb Z/N)^*}
 \operatorname{Spec}\mathbb Z[t,(r+Nt)^{-1}],
\]
\[
 1\longrightarrow D(1+Nt)\longrightarrow\mathbb O_N^*
 \longrightarrow(\mathbb Z/N)^*_{\rm const}\longrightarrow1. \tag{A.5.2}
\]
The right arrow is smooth and surjective. The corresponding additive
sequence is $0\to\mathbb G_a\to\mathbb O_N\to(\mathbb Z/N)_{\rm const}
\to0$. These prove (64)--(66), (74), and (78)--(79), with the stated
choice of constant finite schemes (C4 scans 51--55; fols. 673--675).

More generally, for a unital associative ring $A$ and a two-sided
ideal $J$ that is finite free as an abelian group, the additive pushout
$A[J]$ has a ring structure and exact sequences
\[
 0\to\underline J\to A[J]\to(A/J)_{\rm const}\to0,
 \quad
 1\to(1+\underline J)^*\to A[J]^*
       \to((A/J)^*)_{\rm const}\to1.                  \tag{A.5.3}
\]
The second surjection is smooth. This supplies the general version of
(69)--(73), including a proof of its local surjectivity.
\end{statement}
\begin{proof}
The pushout is the quotient of
$\mathbb Z_{\rm const}\times\mathbb G_a$ by the constant subgroup
$k\mapsto(Nk,-k)$. Choosing one integer representative of every
residue identifies its components as claimed. Formula (A.5.1) is
unchanged under either replacement $(n,x)\mapsto(n+Nk,x-k)$; direct
expansion proves associativity and distributivity. Equivalently it is
the multiplication inherited from the ring with square-zero-free
ideal coordinate $N\mathbb Z\otimes R$, whose internal product is
$N[x]\,N[y]=N[Nxy]$.

A unit must have unit residue and unit evaluation $n+Nx$.
Conversely choose integers $r,r'$ with $rr'=1+Nh$ and set
\[
 y=-\frac{h+r'x}{r+Nx}.
\]
Then $(r+N[x])(r'+N[y])=1$, proving both sufficiency and the displayed
localized coordinate ring. In residue class one the coordinate is
$1+N[t]$, with multiplication $t\star t'=t+t'+Ntt'$; this is the
kernel in (A.5.2). Every geometric fibre of every component is
nonempty: when the characteristic divides $N$, its representative is
nonzero; otherwise at most one value of $t$ is excluded. The arrow is
an open subscheme of an affine-line bundle over the constant base,
so it is smooth and surjective, and therefore fppf locally admits
sections. This proves exactness without asserting surjectivity on
points over every ring.

For the general construction, take
$(A_{\rm const}\times\underline J)/J_{\rm const}$, where
$j\mapsto(j,-j)$, with multiplication
$(a,x)(b,y)=(ab,ay+xb+xy)$. Each component is an affine space of
rank $\operatorname{rank}J$. Fix a component $b\in(A/J)^*$ and a
lift $a\in A$. On this component multiplication on $J$ on the left
and on the right is a pair of square matrices. The unit locus is
the open set where both determinants are invertible. Indeed, a
chosen lift $a'$ of $b^{-1}$ leaves both inverse errors in $J$;
invertibility of left and right multiplication on $J$ corrects the
two errors, producing a left and a right inverse, which coincide.

To prove that this open set has a point in every geometric fibre,
work over an algebraically closed field $k$. In the finite-dimensional
algebra of pairs of endomorphisms of $J_k$, take the subalgebra
generated by the left-right multiplication pairs from $A$ and $J_k$.
The image of $J_k$ is an ideal, and the image of $a$ is a unit modulo
this ideal. In any finite-dimensional algebra, units lift modulo every
ideal: reduce by the nilpotent Jacobson radical, decompose the
semisimple quotient into matrix blocks, and choose arbitrary invertible
entries in the blocks killed by the ideal; units then lift across the
nilpotent radical by the finite geometric series. Lifting the
adjustment back to $J_k$ gives $a+j$ with both multiplication
operators invertible. The preceding correction supplies its inverse.
Thus every component-unit fibre is nonempty. It is open in affine
space, so the map in (A.5.3) is smooth and surjective. The kernel is
exactly the units of the residue-one component. This also proves the
pair description (73): in the finite free unitization of $J$, a right
inverse is a two-sided inverse, by the determinant criterion for
multiplication on a finite free module.
\end{proof}

\begin{statement}{A.6}{Refuted}{Three literal ring formulas and their exact replacements}
The evaluation criterion in (67b), the multiplication printed in (68),
and the full-kernel assertion (77) are false for the pushout (66).
Their correct replacements are A.5 and
\[
 \ker(\mathbb O_N^*\xrightarrow{\mu}\mathbb G_m)
 =\coprod_{r\in(\mathbb Z/N)^*}
   \operatorname{Spec}\mathbb Z[t]/(r+Nt-1).           \tag{A.6.1}
\]
The residue-one component alone is ${}_N\mathbb G_a$.
After inverting $N$, evaluation and residue give an isomorphism of
ring schemes
$\mathbb O_N\simeq\mathbb G_a\times(\mathbb Z/N)_{\rm const}$,
and of unit schemes
$\mathbb O_N^*\simeq\mathbb G_m\times(\mathbb Z/N)^*_{\rm const}$.
The putative global section of the unit-residue map using the four
integers $\pm1,\pm5$ in (83) is also false over $\mathbb Z$.
These are mathematical corrections to C4 scans 51--56, fols. 673--675v;
they do not settle an unattached handwritten reading.
\end{statement}
\begin{proof}
Over $\mathbb F_2$, the sections $0+2[0]$ and $0+2[1]$ are distinct
points of the same affine component of $\mathbb O_2$, while their
evaluations coincide. Thus evaluation does not detect equality.
For (68), the product $12[1]\,12[1]$ evaluates to $144$ over
$\mathbb Z$; the printed formula evaluates to $12$. Formula (A.5.1)
inserts exactly the missing factor $N$.

Evaluation on the $r$-component is $r+Nt$, giving (A.6.1).
For $N=12$, the section $5+12[-1/3]$ over $\mathbb Q$ has evaluation
one and residue five. It is a nonidentity element of the full kernel,
although ${}_{12}\mathbb G_{a,\mathbb Q}=0$. After $N$ is inverted,
$(\mu,r)\mapsto r+N[(\mu-r)/N]$ is the inverse of the displayed
comparison, proving the stronger structural assertion.
Finally a $\mathbb Z$-point of $\mathbb O_{12}^*$ evaluates to
$\pm1$, and hence has residue $1$ or $11$. The other two residue
components have no integral point. Thus there is no global section
of the four-component residue map over $\mathbb Z$; the numbers
$\pm5$ can be used as affine coordinate representatives only.
All these counterexamples are checked by \texttt{a\_schemes.py}.
\end{proof}

\begin{statement}{A.7}{Proved}{Integral coordinates on the multiplier scheme}
Write $a=\nu_2,b=\nu_3\in\{0,1\}_{\rm const}$, so that
$\varepsilon_2=(-1)^a$ and $\varepsilon_3=(-1)^b$. The four charts of
$\gamma$ have coordinates
\[
 r_0=1-6a+4b,\qquad \mu=r_0+12k\in R^*.
\]
The regular functions used in the manuscript are
\[
 \nu=-3a+2b+6k,\quad
 \nu'=-a+b+3k,\quad \nu''=-a+b+2k,                 \tag{A.7.1}
\]
and satisfy, integrally,
\[
 \mu=1+2\nu=1-2a+4\nu'=1-2b+6\nu'',\qquad
 2\nu'-a=3\nu''-b.                                \tag{A.7.2}
\]
For the alternative representative
$\mu_0=(-5)^a5^b$ used in (84)--(92), with
$\mu=\mu_0+12\nu_0$, the required formulas instead are
\[
 \nu'=-a+b-6ab+3\nu_0,\qquad
 \nu''=-a+b-4ab+2\nu_0.                            \tag{A.7.3}
\]
The transcribed coefficients $-3ab,-2ab$ are false. The substitutions
(A.7.3) correct the corresponding expansions in (138)--(141) as well.
The map $\gamma\to\operatorname{Spec}\mathbb Z$ is smooth, flat, of
relative dimension one, with four geometric components at every prime.
\end{statement}
\begin{proof}
The representatives $1,-5,5,-1$ are congruent to
$(-1)^a\bmod4$ and $(-1)^b\bmod3$. Each equation in (A.7.2) follows
by expansion. Conversely the simultaneous data in (63) give
$k=\nu'-\nu''$, and substitution recovers (A.7.1) without division
by two or three. Thus these are integral isomorphisms of schemes.
Expanding $(-5)^a5^b$ on the four constant components gives
$1-6a+4b-24ab$, whence (A.7.3). The discrepancy is detected over
$\mathbb F_7$ by $a=b=1,\nu_0=0$: $\mu=-25\equiv3$ is a unit,
whereas either of the transcribed formulas reconstructs $\mu=1$.
The exact polynomial checks are in \texttt{a\_schemes.py}.
Smoothness and the component claim follow from the four open
affine-line charts of A.5. In characteristics two and three, each
chart is the entire affine line since its evaluation is a constant
nonzero element. These calculations concern C4 scans 51, 57--58,
74--75; fols. 673, 676--676v, 684v--685.
\end{proof}

\begin{statement}{A.8}{Proved}{The first requested exact sequence, over every base ring}
Set $\delta=e_1-e_0$ and $e=3e_1$, a basis $(\delta,e)$ of $M_0$.
For $\alpha=a_0e_0+a_1e_1$ and $\beta=b_0e_0+b_1e_1$, put
\[
 F(\alpha,\beta)=(1-R)\alpha-(1-S)\beta
   =(-a_0-a_1+b_0-b_1)\delta+a_1e.                 \tag{A.8.1}
\]
Define $N$ by $F(\alpha,\beta)=\nu''e$ \emph{in}
$\underline M_0$, with multiplier in $\gamma$. Then
\[
 (t,s,\underline\mu)\longmapsto
 \left(\underline\mu,
       \alpha=t e_0+\nu''\delta,
       \beta=t e_0+s(e_0+e_1)\right)                \tag{A.8.2}
\]
is an isomorphism of schemes $\mathbb G_a^2\times\gamma\simeq N$,
and its group law is
\[
 (t,s,\underline\mu)(t',s',\underline\mu')
 =(t+\mu t',\ s+\mu s',\ \underline\mu\underline\mu'). \tag{A.8.3}
\]
Consequently the sequence requested at the beginning of \S51.4 is
split exact over $\mathbb Z$:
\[
 1\longrightarrow N^{0+}=\mathbb G_a^2
 \longrightarrow N\longrightarrow\mathbb O_{12}^*
 \longrightarrow1.                               \tag{A.8.4}
\]
The middle group is smooth and flat of relative dimension three,
including at two and three. This proves (104)--(105) after the
author's correction in scans 60--62, fols. 677v--678v.
\end{statement}
\begin{proof}
Multiplication by the displayed $R,S$ gives (A.8.1). The two
coordinates of its required equality say exactly
$a_1=\nu''$ and $a_0+a_1=b_0-b_1$. Therefore the inverse to (A.8.2)
is $t=a_0+a_1$, $s=b_1$, with no division by an integer. In
particular $F:M^2\to M_0$ is a split surjection of free lattices.

We verify the group law and its automorphism meaning. For
$\underline\mu\in\gamma(R)$, define an automorphism $u_{\underline\mu}$
of $\mathcal S_R$ by
\[
 u_{\underline\mu}(\rho)=\lambda_1^{\nu}\rho,\qquad
 u_{\underline\mu}(\sigma)=\omega^a\sigma,
 \qquad u_{\underline\mu}(m)=\mu m\quad(m\in\underline M). \tag{A.8.5}
\]
Exponents on the lattice denote additive vector-group coordinates.
These formulas define the morphism on every component of the
finite-stage scheme. Indeed $(1+R+R^2)e_1=0$ gives
$(\lambda_1^\nu\rho)^3=\omega$, and
$(\omega^a\sigma)^2=\omega$. Substituting in
$\lambda_0=\omega\sigma\rho\sigma\rho$ and
$\lambda_1=\omega\rho\sigma\rho\sigma$ gives multiplication by
$1+2\nu=\mu$ on both lattice generators. The commutator subgroup of
the lattice is therefore killed. These checks construct a compatible
map on the twelve charts of $\mathcal S$, using the scalar map on
their affine coordinates; they do not infer uniqueness merely from
the values of two sections in a nonreduced fibre. The identity
\[
 \nu(\underline\mu\underline\mu')
 =\nu(\underline\mu)+\mu\nu(\underline\mu')
\]
holds as an integral polynomial on every pair of components, by
(A.7.1). The sign $a$ adds modulo two. Thus
$u_{\underline\mu}u_{\underline\mu'}=u_{\underline\mu\underline\mu'}$;
the formula for the inverse follows, so these are automorphisms.

The pair of elements
\[
 \left(\operatorname{Int}(t e_0)u_{\underline\mu},\quad
       (-s(e_0+e_1))u_{\underline\mu}\right)        \tag{A.8.6}
\]
records the effective automorphism and the $\sigma$-normalizer
conjugator of (A.8.2). To check this directly,
$(1-R)\delta=3e_1$ and $\nu=3\nu''-b$ give the first generator
formula after the prescribed $\tau^b$ correction; the second follows
from $t e_0-\beta=-s(e_0+e_1)$ and
$(1-S)(e_0+e_1)=0$. Each factor in (A.8.6) multiplies with scalar
action $\mu$ on its one-dimensional translation coordinate. Their
fibre product over $\gamma$ therefore has exactly (A.8.3), which
proves the transported group law on the manuscript's data. Its
section is $t=s=0$. Smoothness follows from (A.8.2) and A.7.
\end{proof}

\begin{statement}{A.9}{Refuted}{The premature exact sequences and their flat replacement}
The surjection onto $\mathbb G_m\times C_2^2$ in (54), and the
surjection $M^2\to M$ in (59), are false as written. The latter has
cokernel $\mathbb Z/3$, and the correct integral lattice sequence is
\[
 0\to P\to M^2\xrightarrow F M\to\mathbb Z/3\to0,
 \quad P=\mathbb Z(e_0,e_0)\oplus
               \mathbb Z(0,e_0+e_1).              \tag{A.9.1}
\]
The raw scheme obtained by imposing the loop equation only in
$\underline M$ is
\[
 \widetilde Z\simeq
 \operatorname{Spec}_{\gamma}
     \mathcal O_\gamma[t,s,z]/(3z).                \tag{A.9.2}
\]
Its projection to $\gamma$ is split surjective, with kernel
$\mathbb G_a^2\times{}_3\mathbb G_a$ as a scheme. The closed subscheme
$z=0$ is exactly $N$ of A.8 and is the schematic closure of
$\widetilde Z_{\mathbb Z[1/3]}$.
The fine kernel sequence (56), after the sign in (55) is corrected
to $1-R$, is exact with raw kernel
$\mathbb G_a^2\times{}_3\mathbb G_a$ and quotient
${}_4\mathbb G_a$. These are the full corrections to (54)--(59),
(94)--(103) (C4 scans 48--49, 59--62; fols. 671v--672, 677--678v).
\end{statement}
\begin{proof}
The sum of the two coordinates of $F(\alpha,\beta)$ is $3a_1$.
Its image is therefore $M_0$, and (A.8.1) proves surjectivity onto
that lattice and the stated kernel. Tensoring this map with a ring
does not preserve the monomorphism $\underline M_0\to\underline M$:
in its basis $(\delta,e)$ the determinant is $-3$. Imposing the
equation in $\underline M$ instead yields $3(a_1-\nu'')=0$;
the other equation determines $b_0$ from $a_0,a_1,b_1$.
Take $z=a_1-\nu''$, $t=a_0+a_1$, $s=b_1$ to obtain (A.9.2).
Its ideal of functions vanishing after inverting three is $(z)$,
which proves the closure assertion. In particular the characteristic
three raw fibre has dimension four, while the corrected fibre has
dimension three. The kernel version has dimensions three and two.
Finite point enumeration over $\mathbb F_2,\mathbb F_3,\mathbb F_5,
\mathbb F_7$ verifies these counts in \texttt{a\_schemes.py}.

For (54), over the algebraic closure of $\mathbb F_2$, equation (63)
forces $\mu=1$ on every component, so any other nonzero element of
that field has no lift, even fppf locally. Equivalently, the already
specified map to the naive character base factors through $\gamma$;
the correct exact target is the one in A.8. For (56), put
$a=b=0$ and $w=\nu'$. Its equations are $4w=0$ and
$F(\alpha,\beta)=2w e_1$. A section of the projection to
${}_4\mathbb G_a$ is
$\alpha=2w e_0-2w e_1$, $\beta=0$:
$F(\alpha,0)=-6w e_1=2w e_1$ since $4w=0$.
The kernel is $F=0$, already computed. The literal plus sign in
$(1+R)\alpha$ in (55) is inconsistent with (57), for example at
$\alpha=\beta=e_0,w=0$ over $\mathbb Z$. We have stated which
replacement makes the sequence exact.

Scan 61 (100)--(101) contains another sign discrepancy visible in
the attached image: the zero equation (96) requires
$F(\alpha,\beta)-\nu''e=0$, or equivalently
$(1-R)(\alpha-\nu''\delta)-(1-S)\beta=0$.
The plus sign in the shifted $\alpha$ gives $6e_1$ on
$\alpha=\delta,\beta=0,\nu''=1$ over $\mathbb Z$.
Thus this replacement is forced by (94)--(96); it is not an
unannounced choice between readings.
\end{proof}

\begin{statement}{A.10}{Proved}{The integral comparison with the fine scheme of LM1}
The fine scheme of scan 41, with component bits $a,b$, coordinate
$w=\nu'$, and loop equation $3a_1=2w-a+b$, is isomorphic over
$\mathbb Z$ to the corrected $N$ of A.8. The comparison and its
inverse are
\[
 k=w-a_1,\quad \nu''=a_1,\quad
 t=a_0+a_1,\quad s=b_1,
 \qquad
 w=3k-a+b,\quad a_1=2k-a+b,                       \tag{A.10.1}
\]
with $\alpha,\beta$ recovered by (A.8.2). In particular it is an
isomorphism on fibres at two and three and on points over every ring.
The alternate LM1 coordinate using $\nu=2w+a$ is carried to the
present coordinate using $\nu=2w-a$ by $w_{\rm here}=w_{\rm LM1}+a$.
This settles the finite-stage comparison left in LM1 E.20 for this
specific parameter scheme (C4 scans 40--41, 62, 71--72;
fols. 667v--668, 678v, 683--683v).
\end{statement}
\begin{proof}
The equation $2w-3a_1=a-b$ implies
$w=3(w-a_1)-a+b$ and $a_1=2(w-a_1)-a+b$.
Thus (A.10.1) is integral in both directions. Substituting gives
$\mu=1-2a+4w=1-6a+4b+12k$, the same inverted element in both
coordinate rings. The remaining loop coordinate equation is exactly
$a_0+a_1=b_0-b_1$, whose solutions are (A.8.2).
This proves the claimed scheme isomorphism, including all nilpotent
test rings. Both groups use the automorphism and conjugator
multiplication (A.8.6), so the comparison also respects their group
structures. The final coordinate translation follows by equating
$2w_{\rm here}-a=2w_{\rm LM1}+a$ as integral functions; substitution,
not cancellation in a characteristic-two ring, proves the equality.
\end{proof}

\begin{statement}{A.11}{Proved}{Effective automorphisms and the early quotient sequences}
The effective group, its geometric extension, and the early quotient
groups are
\[
 Z=\mathbb G_a\rtimes_\mu\gamma,
 \quad G=\mathcal S\rtimes_u\gamma,
 \quad \Gamma_{\rm e}=\mathbb G_a\rtimes_\mu\gamma,
\]
\[
 \mathcal M_{\rm e}\simeq(\mathbb G_a\times\mathcal S)\rtimes\gamma,
 \qquad
 \mathcal M_{\rm e}^{\sharp}\simeq\mathbb G_a^3\rtimes_\mu\gamma.
                                                               \tag{A.11.1}
\]
The $\gamma$-action on $\mathcal S$ is (A.8.5). In $N$, the subgroup
$L_0$ is the $t$-axis and the quotient is the $s$-coordinate. We obtain
the exact sequences
\[
\begin{gathered}
1\to L_0\to N\to\Gamma_{\rm e}\to1,\qquad
1\to\mathbb G_a\to\Gamma_{\rm e}\to\gamma\to1,\\
1\to\mathcal S\to\mathcal M_{\rm e}\to\Gamma_{\rm e}\to1,
\quad
1\to\underline M\to\mathcal M_{\rm e}^{\sharp}\to\Gamma_{\rm e}\to1,\\
1\to L_0\to Z\to\gamma\to1,
\qquad1\to\mathcal S\to G\to\gamma\to1.
\end{gathered}                                                     \tag{A.11.2}
\]
All are split and exact on points over every ring. The lattice kernel
sequences (127)--(132) are
$0\to\mathbb Z\to P\to\mathbb Z\to0$ and their vector-group
base changes; the maps are $(t,s)\mapsto s$ and $t\mapsto(t,0)$.
These prove (123)--(127), (131)--(132), (142), (145)--(150)
(C4 scans 70--76; fols. 682v--685v).
\end{statement}
\begin{proof}
Under (A.8.6), the first factor is
$\operatorname{Int}(t e_0)u_{\underline\mu}$ and the second has
coordinate $-s$. The kernel of the effective projection is the
$s$-axis; its image is the displayed $Z$. Inner conjugation by
$t e_0$ embeds $L_0\simeq\mathbb G_a$ into $Z$, and (A.8.5)
normalizes it by the scalar $\mu$. Thus the quotients are represented
by the displayed affine schemes, with the stated sections.

For the crossed quotient defining $\mathcal M_{\rm e}$, take
$\mathcal S\rtimes N$ with the action of $N$ through $Z$ and quotient
by $t\mapsto(-t e_0,(t,0,1))$. This is a normal subgroup because
its action on $\mathcal S$ equals inner conjugation by $t e_0$.
Every coset has exactly one representative with $t=0$; its remaining
coordinates are $(x,s,\underline\mu)$ and multiply as in (A.11.1).
The same calculation with $\underline M$ in place of $\mathcal S$
gives the sharp subgroup. All kernel and section assertions now
follow by coordinate projection, so they are exact on arbitrary
ring points and hence as fppf sheaves. The kernel lattice $P$ was
computed in A.9; its quotient by the diagonal $e_0$-line has the
free generator $(0,e_0+e_1)$. This proves every lattice and
vector-group sequence asserted here.

The ambiguous exponent in scan 76 (151) is fixed by the verified
formula $\lambda_1^{\nu}\rho$ in scan 93, and independently by
$u(\lambda_i)=\mu\lambda_i$ in A.8. The candidate exponent
$6\nu''$ fails this test, since the required relation is
$\nu=3\nu''-b$. We retain this correction explicitly in the
readings ledger.
\end{proof}

\begin{statement}{A.12}{Proved}{The two normalizers and the central extension of the multiplier scheme}
Let $d=\nu''$ and set, in $G=\mathcal S\rtimes\gamma$,
\[
 v_{\underline\mu}=\sigma^b\delta^{-d}u_{\underline\mu}.
\]
For the raw normalizer of (152), before the closure in (159),
the literal sequence (155) is
\[
 1\to C_6\times{}_3\mathbb G_a\to N_\rho^{\rm raw}
   \to\gamma\to1,
 \qquad N_\rho^{\rm raw}\simeq{}_3\mathbb G_a\times N_\rho.
                                                               \tag{A.12.0}
\]
The smooth closure of the $\rho$-normalizer is represented by
$\{\rho^k v_{\underline\mu}:k\in\mathbb Z/6\}$; hence it has the
exact sequence
\[
 1\to C_6=L_\rho\to N_\rho\to\gamma\to1.          \tag{A.12.1}
\]
The final subgroup specified in (212), scans 93--94, is
\[
 N_\rho^?=\{\omega^i v_{\underline\mu}:i\in\mathbb Z/2\},
 \qquad 1\to C_2=\mu_\rho\to N_\rho^?\to\gamma\to1. \tag{A.12.2}
\]
Its multiplication on $C_2\times\gamma$ is
\[
 (i,\underline\mu)(i',\underline\mu')
  =(i+i'+c(\underline\mu,\underline\mu'),
       \underline\mu\underline\mu'),\qquad
 c=(a+b)b'\pmod2.                                 \tag{A.12.3}
\]
Thus (A.12.2) is the pullback along
$(\varepsilon_2,\varepsilon_3):\gamma\to C_2^2$ of the order-eight
dihedral central extension $1\to C_2\to D_8\to C_2^2\to1$.
Moreover
\[
 N_\sigma^\sharp=\mathbb G_a\rtimes_\mu\gamma,\qquad
 N_\sigma=(C_4\times\mathbb G_a)\rtimes\gamma,       \tag{A.12.4}
\]
where $\gamma$ acts by $\varepsilon_2$ on $C_4$ and by $\mu$ on
$\mathbb G_a$. These assertions prove (155), (159), and (173)--(181),
and compute the extension asked for after (196)
(C4 scans 78--84, 87, 93--94; fols. 686v--689v, 691, 694--694v).
\end{statement}
\begin{proof}
The identities $(1-R)\delta=3e_1$ and $\nu-3d=-b$ give
\[
 \delta^{-d}u_{\underline\mu}(\rho)\delta^d
   =\lambda_1^{-b}\rho.
\]
For $b=1$, direct multiplication in the modular group gives
$\sigma\lambda_1^{-1}\rho\sigma^{-1}=\rho^{-1}$.
For $b=0$ the displayed expression is $\rho$.
Thus $v_{\underline\mu}$ has the required normalizing action
$\rho\mapsto\rho^{\varepsilon_3}$.
In the finite quotient $H$, the centralizer of $\rho$ is its cyclic
subgroup of order six. Within each of these six components, the
remaining translation condition is $(1-R)m=0$. Its kernel has
equations equivalent to $3z=0$, because $\det(1-R)=3$.
Explicitly $m=z\delta$, with $3z=0$. The cyclic component
$C_6$ fixes this vector. Conjugation by $v_{\underline\mu}$
acts on it as $\mu\varepsilon_3$, and
$\mu\varepsilon_3=1+6\varepsilon_3d$; hence it also fixes
$z\delta$. This proves the central direct product (A.12.0)
and the original sequence (155), including its extra fibre
in characteristic three.
The closure of its characteristic-not-three part is the zero
translation. Translation by $v_{\underline\mu}$ therefore gives
exactly the smooth normalizer stated above, including all its
components, with kernel $L_\rho$.

The identity
\[
 d(\underline\mu\underline\mu')
   =\varepsilon_3(\underline\mu')d(\underline\mu)
      +\mu d(\underline\mu')                     \tag{A.12.5}
\]
is an integral polynomial identity on the sixteen pairs of charts;
it follows by multiplying $\mu=\varepsilon_3+6d$ and using (A.7.1).
In particular its proof does not cancel six in a test ring. Moving
$\sigma^{b'}$ past $\delta^{-d}$ uses $S\delta=-\delta$, and
$u_{\underline\mu}(\sigma)=\omega^a\sigma$. The two central factors
are $\omega^{ab'}$ and $\omega^{bb'}$. Equation (A.12.5) then gives
\[
 v_{\underline\mu}v_{\underline\mu'}
     =\omega^{(a+b)b'}v_{\underline\mu\underline\mu'}.
\]
Consequently the multiplication in $N_\rho$ is
$(k,\underline\mu)(k',\underline\mu')
=(k+\varepsilon_3k'+3c,\underline\mu\underline\mu')$.
Restricting $k$ to $0,3$ proves (A.12.2)--(A.12.3).
The bilinear function $((a,b),(a',b'))\mapsto(a+b)b'$ is a
two-cocycle. The resulting eight-element group has five involutions
and two elements of order four, and is generated by two involutions
with noncommuting product of order four. This proves it is $D_8$;
the full cocycle and element-order checks are in
\texttt{a\_finite.py}. Exactness of (A.12.2) holds even on ring
points, because $v_{\underline\mu}$ is a global section as a scheme,
although it is not a group section.

For $\sigma$, the centralizer in $H$ is $C_4$, generated by
$\sigma$. Its translation kernel is
$\ker(1-S)=\mathbb Z(e_0+e_1)$, also after every base change.
These translations commute with $\sigma$, so the geometric
centralizer is $C_4\times\mathbb G_a$. Equation (A.8.5) supplies
a group section over $\gamma$ and its action on both factors,
which proves (A.12.4) and all the stated exact sequences. The sharp
intersection retains just the translation component. Similarly
$N_\rho\cap G^\sharp$ is the section over
$\gamma'=\ker\varepsilon_3$, since the component
$\rho^k\sigma^b$ is the identity precisely when $k=b=0$.
Put $\gamma''=\ker\varepsilon_2$, $\gamma^0=\gamma'\cap\gamma''$,
and $\gamma'''=\{a=b\}$. The neutral partial section of (166)--(167)
is $v(\gamma^0)$. The partial section of (168)--(169) is exactly
\[
 \underline\mu\longmapsto\omega^b v_{\underline\mu}
 \quad(\underline\mu\in\gamma'''),
\]
because its cocycle vanishes on $a=b$, and at $\underline\mu=-1$
it is $\tau'$. Conjugation by $\tau'$ sends
$v_{\underline\mu}$ on $\gamma'$ to
$\omega^a v_{\underline\mu}$, as (A.12.3) shows. It preserves
the partial section precisely on $\gamma^0$.
This proves the prime and double-prime intersections and the
normalization assertions of (160)--(170) with their exact domains.
\end{proof}

\begin{statement}{A.13}{Refuted}{The literal definition (171) has only half the stated image}
The group defined in (171) from $Z_\rho^{0\sharp}$, $\omega$, and
$\tau'$ maps onto
\[
 \gamma'''=\{\varepsilon_2=\varepsilon_3\}\subset\gamma,
\]
and is the inverse image of this subgroup in the final $N_\rho^?$
of (212). Thus (172) with target all of $\gamma$ is false for
the literal definition (171). Its exact replacement is
$1\to C_2\to N_{\rho,\mathrm{literal}}^?\to\gamma'''\to1$;
the final fibres (212) yield the larger exact sequence (A.12.2).
The surviving last line of (194), with $N_\rho$ and no question mark,
requires the replacement
$1\to C_3\to N_\rho/\mu_\rho\to\gamma\to1$;
the isomorphism to $\gamma$ holds for $N_\rho^?/\mu_\rho$.
The final $\Gamma$ in (184) also satisfies
$\Gamma/C_2\simeq\mathbb G_a\rtimes\gamma$; the unquotiented
version of this isomorphism in the third row of (195) is false.
These corrections concern C4 scans 82, 85--87, 94;
fols. 688v, 690--691, 694v.
\end{statement}
\begin{proof}
The identity component $Z_\rho^{0\sharp}$ maps isomorphically to
$\gamma^0=\{a=b=0\}$ by A.12. The central $\omega$ has trivial
character, and $\tau'$ has multiplier $-1$, hence $(a,b)=(1,1)$.
The literal generated subgroup therefore has image in $a=b$.
Conversely multiplication by $-1$ carries the whole $00$ chart to
the whole $11$ chart, so the image is exactly $\gamma'''$.
At $\underline\mu=-1$, $d=0$ and
$\tau'=\omega v_{-1}$; hence both lifts over each of these charts
occur. This proves the asserted inclusion and its exact image and
kernel. Over $\mathbb F_2$, the $\gamma$-point with $a=1,b=0,k=0$
has evaluation one and has no lift to this literal subgroup. This
finite counterexample is recorded in \texttt{a\_schemes.py}.

Quotienting (A.12.1) by its central order-two subgroup leaves
kernel $C_6/C_2=C_3$, which proves the correction to (194).
The final fibre product defining $\Gamma$ has kernel
$C_2\times\mathbb G_a$ over $\gamma$, by A.12. Its quotient by
the indicated central $C_2$ is $N_\sigma^\sharp$, giving the claimed
isomorphism. On $\mathbb F_2$ points the two groups have $32$ and
$16$ elements, respectively; the explicit projection has kernel of
order two. Thus the missing quotient in (195) is detected by an
actual map and a finite counterexample.
\end{proof}

\begin{statement}{A.14}{Proved}{The final fibre products and every surviving extension sequence}
Using the final (212), define
\[
 \Gamma=N_\rho^?\times_\gamma N_\sigma^\sharp,
 \quad \mathcal M=\Gamma\times_\gamma G,
 \quad \mathcal M^\sharp=\Gamma\times_\gamma G^\sharp.
                                                               \tag{A.14.1}
\]
The coordinates on $\Gamma$ are $(i,s,\underline\mu)$ and its law is
\[
 (i,s,\underline\mu)(i',s',\underline\mu')
  =(i+i'+(a+b)b',\ s+\mu s',\
       \underline\mu\underline\mu').              \tag{A.14.2}
\]
There are central quotient morphisms, with the indicated kernels,
\[
 \Gamma\twoheadrightarrow\Gamma_{\rm e},\qquad
 \mathcal M\twoheadrightarrow\mathcal M_{\rm e},\qquad
 \mathcal M^\sharp\twoheadrightarrow\mathcal M_{\rm e}^\sharp,
 \qquad\ker=C_2=\mu_\rho.                         \tag{A.14.3}
\]
In addition to A.8, A.11, and A.12, the final exact sequences are
\[
\begin{gathered}
1\to\mathcal S\to\mathcal M\to\Gamma\to1,\qquad
1\to\underline M\to\mathcal M^\sharp\to\Gamma\to1,\\
1\to C_2\times\mathbb G_a\times\mathcal S
 \to\mathcal M\to\gamma\to1,\\
1\to C_2\times\mathbb G_a\times\underline M
 \to\mathcal M^\sharp\to\gamma\to1,\qquad
1\to C_2\times\mathbb G_a\to\Gamma\to\gamma\to1.
\end{gathered}                                                     \tag{A.14.4}
\]
For $X=\Gamma,\mathcal M,\mathcal M^\sharp$, the auxiliary group
$S_0X=Z\times_\gamma X$ has split exact sequence
$1\to L_0\to S_0X\to X\to1$.
These prove (182)--(193) and (196), with the corrected quotient
versions of (194)--(195). The central extensions are precisely the
pullbacks of the explicit $D_8$ extension in A.12.
\end{statement}
\begin{proof}
Every map to $\gamma$ in (A.14.1) is a homomorphism, and the
multiplication of the fibre product is componentwise. Insert
(A.12.3)--(A.12.4) to obtain (A.14.2). Forgetting $i$ is a
homomorphism, since its cocycle is central, and its kernel is the
constant $C_2$. The resulting three groups have exactly the
coordinates and actions in (A.11.1), proving the actual comparisons
(A.14.3) between the early and final revisions.
The projection $\mathcal M\to\Gamma$ has the group section
$(a,b)\mapsto(a,b,u_{\underline\mu})$; its kernel is $\mathcal S$.
Replacing $\mathcal S$ by $\underline M$ gives the sharp row.
Kernels over $\gamma$ are the products of the three displayed
geometric kernels, proving (A.14.4). These projections are onto
on ring points: use $v_{\underline\mu}$, $u_{\underline\mu}$, and
zero translation as a scheme section. Finally $Z\to\gamma$ has
the group section (A.8.5), so the auxiliary projections are split
exact. This proves exactness and representability directly for the
entire final diagram.

The central factor $\mu_\rho$ here is $(\omega,1,1)$ in
$(a,b,U)$-coordinates. The central subgroup $\mu_G$ is
$(1,1,\omega)$ and has a separate inclusion. The quotient map
in (A.14.3) removes the first of these, exactly as its kernel states;
no central factors have been combined without a map.
\end{proof}

\begin{statement}{A.15}{Proved}{The three distinguished subgroups: all requested points and fibres}
The final subgroups of scans 88--94 are the graph subgroups
\[
 \mathcal M(0)=\{U=u_{\underline\mu}\},\qquad
 \mathcal M(1/2)=\{U=b_\sigma\},\qquad
 \mathcal M(-\bar j)=\{U=a_\rho\}                 \tag{A.15.1}
\]
inside $\mathcal M=\{(a_\rho,b_\sigma,U)\}$.
Projection on $(a_\rho,b_\sigma)$ identifies each with $\Gamma$ as
a group scheme over $\mathbb Z$. They are smooth and flat of relative
dimension two, with eight geometric components at every prime.

For any connected ring $R$, their common point set is exactly
\[
 \left\{(a,b,k,i,s): a,b,i\in\{0,1\},\ k,s\in R,
           \ 1-6a+4b+12k\in R^*\right\}.          \tag{A.15.2}
\]
For a general ring the bits are locally constant functions on its
spectrum. The group law is (A.14.2), with multiplier product
\[
\begin{split}
 a''&=a+a'\pmod2,\qquad b''=b+b'\pmod2,\\
 k''&=\frac{r_0r_0'-r_0''}{12}
             +r_0k'+r_0'k+12kk'.
\end{split}                                                       \tag{A.15.3}
\]
The numerator is an integer divisible by twelve on every component.
In particular the requested rings have the following complete
parametrizations, with $s\in R$ and $i\in C_2$ always free.
\begin{center}\small
\begin{tabular}{p{.19\linewidth}p{.71\linewidth}}
\toprule
$R$ & Allowed multipliers and component bits; $k=(\mu-r_0)/12$\\\midrule
$\mathbb Z$ & $\mu=1:(a,b,k)=(0,0,0)$; $\mu=-1:(a,b,k)=(1,1,0)$.\\
$\mathbb Z[1/2]$ & $\mu=\pm2^m$ ($m\in\mathbb Z$); $a$ arbitrary,
                   $(-1)^b=\mu\bmod3$.\\
$\mathbb Z[1/3]$ & $\mu=\pm3^m$ ($m\in\mathbb Z$); $b$ arbitrary,
                   $(-1)^a=\mu\bmod4$.\\
$\mathbb Z[1/6]$ & $\mu=\pm2^m3^n$ ($m,n\in\mathbb Z$); $a,b$ arbitrary.\\
$\mathbb F_2$ & $a,b,k$ arbitrary; $\mu=1$;
              $|\mathcal M(p)(\mathbb F_2)|=32$.\\
$\mathbb F_3$ & $a,b,k$ arbitrary; $\mu=(-1)^b$;
              $|\mathcal M(p)(\mathbb F_3)|=72$.\\
\bottomrule
\end{tabular}
\end{center}
Here $p$ denotes any of the three displayed points, not a prime.
As abstract groups, the integral points are
$C_2\times(\mathbb Z\rtimes\{\pm1\})$; the two finite groups are
$D_8\times C_2^2$ over $\mathbb F_2$ and
$C_3\times(C_3\rtimes D_8)$ over $\mathbb F_3$, where $D_8$ acts
on the second $C_3$ through $\varepsilon_3$.
\end{statement}
\begin{proof}
The maps $\Gamma\to G$ given respectively by
$(a_\rho,b_\sigma)\mapsto u_{\underline\mu},b_\sigma,a_\rho$
are homomorphisms. Their graphs have the stated inverse projection,
proving all three group-scheme isomorphisms, not just equality of
dimensions or of point counts. Formula (A.15.2) follows from
$\Gamma\simeq C_2\times\mathbb G_a\times\gamma$ as a scheme.
Equations (A.14.2) and multiplication of
$r_0+12k$ give (A.15.3); its carry term is integral because residue
classes multiply according to the two sign bits.

For $\mathbb Z$, invertibility forces $\mu=\pm1$ and the congruences
in (A.7.2) force exactly the displayed bits. The cocycle then vanishes
since $a=b$. For $\mathbb Z[1/2]$, the unit group is
$\{\pm2^m\}$, and divisibility by twelve in that ring is equivalent
to the displayed congruence modulo three; modulo four imposes no
condition because two is inverted. The argument with the primes
reversed proves the next row, and inverting six proves the fourth.
Uniqueness of $k$ in these torsion-free rings follows directly.
In the two finite fields, twelve vanishes and all representatives
$r_0$ are units, giving $4\cdot2\cdot2\cdot2=32$ and
$4\cdot3\cdot2\cdot3=72$ points. These were also counted directly
from the original lattice equations by \texttt{a\_finite.py}.

There is an explicit ring isomorphism
$\mathbb O_N(\mathbb Z/q)\to\mathbb Z/Nq$,
$r+N[t]\mapsto r+Nt$. Therefore
$\gamma(\mathbb F_2)=(\mathbb Z/24)^*\simeq C_2^3$ and
$\gamma(\mathbb F_3)=(\mathbb Z/36)^*\simeq C_3\times C_2^2$.
The residue maps split as maps of these finite abelian groups, with
kernels $C_2$ and $C_3$. Pulling back the $D_8$ extension yields the
claimed finite-group descriptions, including its action on $s$.
Finally each of the eight charts is an open subset of affine
two-space, and over characteristics two and three it is the whole
affine two-space. This proves smoothness, flatness, and the component
count on the actual fibres.

Under (A.14.3), the first two graphs contain $\mu_\rho$, while the
third intersects it trivially. Indeed $U=a_\rho$ forces
$U=\omega$ when $a_\rho=\omega$. Its image in
$\mathcal M/\mu_\rho$ is a bisection over $\Gamma_{\rm e}$, with
kernel $\mu_G$. This gives explicitly the change of descriptions in
scans 91--94 rather than identifying the central quotients with
all three unquotiented graphs.
\end{proof}

\begin{statement}{A.16}{Proved}{An actual comparison with the groups of \S48}
Let $H_{48}$ be the group of pairs of quadratic-normalizer matrices
\[
 A=c\rho+dI,\quad
 B=\begin{pmatrix}u&t\\-et&eu\end{pmatrix},\quad
 Q_3=c^2+cd+d^2,\quad Q_2=t^2+u^2,
\]
with $Q_2,Q_3$ units and $e^2=(e')^2=1$, satisfying
$\det B=eQ_2=e'Q_3$. Its character map is
\[
 H_{48}\to\gamma_{48}=\mathbb G_m\times\mu_2^2,
 \quad(A,B)\mapsto(\det B,e,e').
\]
There is a specified comparison of character schemes
\[
 \iota:\gamma\longrightarrow\gamma_{48},\qquad
 \underline\mu\longmapsto(\mu,(-1)^a,(-1)^b),      \tag{A.16.1}
\]
which is an isomorphism over $\mathbb Z[1/6]$. Define the actual
base change $H_{48}^\dagger=H_{48}\times_{\gamma_{48}}\gamma$.
Projection followed by the zero-translation section is a morphism
of affine group schemes
\[
 H_{48}^\dagger\longrightarrow\Gamma/\mu_\rho,
 \quad(A,B,\underline\mu)\longmapsto(0,\underline\mu). \tag{A.16.2}
\]
Its kernel is $ST_\rho\times ST_\sigma$, the product of the two
norm-one quadratic groups, and its image is exactly the displayed
copy of $\gamma$. For the full groups the analogous map
$\mathcal M_{48}\times_{\gamma_{48}}\gamma\to\mathcal M/\mu_\rho$
has kernel $ST_\rho\times ST_\sigma\times\mathrm{SL}_2$ and image
the canonical $\gamma$ section.

The \S48 identifications
$\mathcal M_{48}(0)\simeq H_{48}$ and
$\mathcal M_{48}(1/2)\simeq H_{48}$ hold by their graph maps.
The graph $U=A$ is defined there over $e'=1$, and is identified with
$H_{48}|_{e'=1}$. Its comparison is the restriction of (A.16.2)
over $b=0$. Thus all three \S51 groups are related to the earlier
groups by these explicit base changes, character maps, and central
quotients. The resulting schemes are not isomorphic: their generic
relative dimensions are respectively two and three, and their bad
fibres have the distinct descriptions proved below.
\end{statement}
\begin{proof}
Solving $A\rho=\rho A$ and $B\sigma=e\sigma B$ gives the two
displayed matrix forms, and their determinants give exactly the
stated equations. The character map is fppf surjective: given its
three coordinates, adjoining square roots of the required unit norms
and taking scalar elements in the quadratic algebras supplies lifts.
The extensions adjoining square roots are finite free and faithfully
flat, including in characteristic two. Its kernel is precisely the
two norm-one groups. Formula (A.16.1) is a homomorphism by the ring
and sign formulas in A.5--A.7. Over $\mathbb Z[1/6]$, its inverse
recovers the constant signs and sets $k=(\mu-r_0)/12$, proving the
asserted isomorphism.

The base-changed character map retains its exact kernel, so its
composition with the group section of $\Gamma/\mu_\rho$ has the
kernel and image claimed in (A.16.2). For the full matrix triple,
the additional determinant-prescribed $U$ coordinate has kernel
$\mathrm{SL}_2$. The section $U=\operatorname{diag}(\det B,1)$
gives the graph at zero, and $U=B$ gives the graph at one-half.
For $U=A$, the determinant condition reads $e'=1$, proving exactly
the specified restriction. This derives the comparison from the
equations used in LM1 B.2--B.5 instead of importing their conclusion.

Here is also a direct fibre computation. Eliminating the common
determinant gives
\[
 H_{48}=\operatorname{Spec}
 \mathbb Z[c,d,t,u,e,e',Q_2^{-1}]/
 (Q_3-ee'Q_2,e^2-1,(e')^2-1).
\]
Relative to the sign scheme $\mu_2^2$ its hypersurface is smooth:
the gradient in $c,d,t,u$ is
$(2c+d,c+2d,-2ee't,-2ee'u)$ and cannot vanish on $Q_2\ne0$.
For characteristic different from two, use the last two entries;
in characteristic two the first two become $d,c$ and cannot both
vanish because $Q_3$ is a unit. The geometric fibres are nonempty,
so this is a smooth surjective map of relative dimension three.
Thus the generic dimension is three. At three,
$Q_3=(c-d)^2$; setting $v=c-d\ne0$ shows four smooth geometric
components $\mathbb G_m^2\times\mathbb G_a$.
At two put $x=e-1,y=e'-1$. The sign algebra is
$\mathbb F_2[x,y]/(x^2,y^2)$; after reduction, $Q_2=(t+u)^2$ and
the remaining equation is the norm-one equation
$a^2+ab+b^2=1$. Hence the reduced fibre is
\[
 \mathbb G_a\times\mathbb G_m\times
 R^1_{\mathbb F_4/\mathbb F_2}\mathbb G_m.
\]
It is geometrically connected, and smoothness over the sign algebra
gives generic length four before reduction. Restricting $e'=1$
replaces that length by two. The full $\mathcal M_{48}$ is, as a
scheme, the product of this scheme with $\mathrm{SL}_2$, by
$U=\operatorname{diag}(\det B,1)h$.
These fibre calculations and the explicit morphisms above prove
the asserted comparison with A.15, whose characteristic-two fibre
has eight reduced affine-plane components. They also prove flatness
of the \S48 models and their failure of smoothness exactly at two.
No nonexistence of further mathematical comparisons is asserted.
\end{proof}

\begin{statement}{A.17}{Proved}{Linear congruence descent from the actual braid automorphism}
Use the standard tangential and braid normalization
\[
 B_3=\langle\beta_1,\beta_2:
              \beta_1\beta_2\beta_1=\beta_2\beta_1\beta_2\rangle,
 \quad x=\beta_1^2,\quad y=\beta_2^2.
\]
For $(\mu,f)\in\widehat{\mathrm{GT}}_0$, where
$f\in\widehat{\langle x,y\rangle}'$, let
\[
 \Phi(\beta_1)=\beta_1^\mu,\qquad
 \Phi(\beta_2)=f^{-1}\beta_2^\mu f.                \tag{A.17.1}
\]
For every $n\geq2$, define the continuous surjection
\[
 q_n:\widehat B_3\longrightarrow\mathrm{SL}_2(\mathbb Z/n),
 \quad\beta_1\mapsto U=\begin{pmatrix}1&1\\0&1\end{pmatrix},
 \quad\beta_2\mapsto W=\begin{pmatrix}1&0\\-1&1\end{pmatrix}.
\]
There is a unique matrix
\[
 V_n(\Phi)=\begin{pmatrix}\mu&t_n(\Phi)\\0&1\end{pmatrix}
 \quad\text{such that}\quad
 q_n\Phi=\operatorname{Int}(V_n(\Phi))q_n.          \tag{A.17.2}
\]
The maps $\Phi\mapsto V_n(\Phi)$ are continuous homomorphisms,
compatible under reduction. In particular all linear congruence
kernels are invariant. They give a specified action on the full
$\mathrm{GL}_2(\mathbb Z/n)$ by the same conjugation formula.
This supplies a direct proof of the descent needed in LM2 C.d1 and
its requested independent check at $n=12$.
\end{statement}
\begin{proof}
The group-theoretic input is that (A.17.1) is an automorphism of
$\widehat B_3$ for an actual $\widehat{\mathrm{GT}}_0$ pair.
Lochak--Schneps prove the braid identity in Lemma 5 of
\cite{A_LochakSchneps}; Harbater--Schneps, \S1.3, Proposition 4,
proof of (i)$\Rightarrow$(iii), explicitly applies that lemma to the
full profinite group, \cite{A_HarbaterSchneps}. Both proofs were
opened and checked. The input uses the two cycle relations and the
automorphism requirement; no pentagon is used here.
The following reduction from that theorem to (A.17.2) is elementary
and also proves the requisite generation condition cannot be omitted.

Write $A=q_n\Phi(\beta_1)=U^\mu$ and
$B=q_n\Phi(\beta_2)=\left(\begin{smallmatrix}a&b\\c&d\end{smallmatrix}\right)$.
Conjugacy to $W^\mu$ gives $a+d=2$, $ad-bc=1$.
The lower-left entry of $ABA=BAB$ gives
\[
 c(1+\mu c)=0.                                    \tag{A.17.3}
\]
The pair $A,B$ generates all of $\mathrm{SL}_2(\mathbb Z/n)$,
because $q_n$ is onto and $\Phi$ is an automorphism. For each prime
$p\mid n$, the element $c\bmod p$ cannot vanish: otherwise both
matrices would be upper triangular, and could not generate
$\mathrm{SL}_2(\mathbb F_p)$. Thus $c$ is a unit modulo $n$.
Equation (A.17.3) yields $c=-\mu^{-1}$; the determinant and trace
then give
\[
 d=2-a,\quad b=\mu(a-1)^2,\quad t_n=\mu(1-a).
\]
Direct multiplication proves
$V_nUV_n^{-1}=A$, $V_nWV_n^{-1}=B$. Agreement on the topological
generators proves (A.17.2) on the full profinite group.
For surjectivity of $q_n$, elementary row operations over each local
ring $\mathbb Z/p^r$ reduce a determinant-one matrix to diagonal
form, and diagonal determinant-one matrices are products of
elementary matrices. Chinese remainders combine these operations.
Powers of $U,W$ give all the elementary matrices.

A matrix commuting with $U,W$ is scalar, by comparing its four
entries. A scalar relating two matrices with lower-right entry one
is the identity. This proves uniqueness. The product $V_n(\Phi)
V_n(\Psi)$ has the same normalization and implements the composite
automorphism, which proves the homomorphism assertion. The formula
for $a$ in terms of $q_n(f)$ proves continuity; uniqueness proves
compatibility in $n$. Taking the same inner conjugation on
$\mathrm{GL}_2(\mathbb Z/n)$ is now a well-typed extension of the
action on its determinant-one subgroup.

For comparison with the modular group, the central braid
$z=(\beta_1\beta_2)^3$ maps to $\omega$, and
$\widehat B_3/\overline{\langle z^2\rangle}
\simeq\widehat{\mathrm{SL}_2(\mathbb Z)}$.
The cited proof also gives $\Phi(z)=z^\mu$, so this quotient is
preserved. The subgroup $\widehat\Pi_0$ is the inverse image of
$V_4$ modulo four, preserved by the matrices in A.18.
Alternatively its generators are
$\lambda_0=\omega U^{-2}$, $\lambda_1=\omega W^{-2}$; their images
are respectively $\lambda_0^\mu$ and a conjugate of
$\lambda_1^\mu$ by an element of $\widehat\Pi_0'$. To verify the
last assertion, evaluate the commutator word $f$ in
$U^2=\omega\lambda_0^{-1}$ and
$W^2=\omega\lambda_1^{-1}$: its central factors cancel. This
constructs the comparison to the action on $M/12M$ used below.

The program \texttt{a\_finite.py} independently enumerates every
signed matrix-coboundary solution used in LM2 C.d1 at level twelve:
there are twelve for each of the four multipliers and no failed
intertwiner. The program \texttt{a\_braid.py} additionally
enumerates the eight possible values of $f$ in the finite
commutator image, for every multiplier. All thirty-two pairs
satisfy the braid relation; eight pairs fail generation and are
explicitly rejected. The remaining twenty-four pairs give the
twelve distinct matrices in A.18. This records both the computation
and the exact reason the discarded matrix solutions cannot be
automorphisms of the profinite source.
\end{proof}

\begin{statement}{A.18}{Proved}{Exact images at level twelve, with the genuine normalization stated}
For the normalization (A.17.1)--(A.17.2), the images of
$G_{\mathbb Q}$, $\widehat{\mathrm{GT}}$, and
$\widehat{\mathrm{GT}}_0$ in $\mathrm{GL}_2(\mathbb Z/12)$ are all
\[
 \mathcal V_{12}=
 \left\{\begin{pmatrix}\mu&t\\0&1\end{pmatrix}:
     \mu\in\{1,5,7,11\},\quad t\in\{0,4,8\}\right\},
 \qquad |\mathcal V_{12}|=12.                      \tag{A.18.1}
\]
Their specified actions in $\operatorname{Aut}(\mathrm{GL}_2(\mathbb Z/12))$
therefore have image $\operatorname{Ad}(\mathcal V_{12})$ of order
$12$. Their images in
$\operatorname{Out}(\mathrm{GL}_2(\mathbb Z/12))$ are trivial.
The canonical outer images on $\mathrm{SL}_2(\mathbb Z/12)$ are all
$(\mathbb Z/12)^*/((\mathbb Z/12)^*)^2=C_2^2$, of order four.
On $M/12M$ the genuine image of each group is the scalar subgroup
$\{\mu I:\mu\in(\mathbb Z/12)^*\}$, also of order four.

For $G_{\mathbb Q}$, choose the tangential base point at $\lambda=0$
with tangent $+1$ and the real interval path to the tangent at one;
this is the Belyi normalization in (A.17.1).
The kernel of its normalized level-twelve matrix representation is
exactly
\[
 G_{\mathbb Q(\zeta_{12},\sqrt[3]{2})}.           \tag{A.18.2}
\]
The extended group generated by the geometric modular group and
any of the three arithmetic groups acts on the same linear
congruence quotient with image
$\operatorname{Ad}(\mathrm{GL}_2(\mathbb Z/12))$ of order $1152$.
These statements give the precise actions requested after LM2;
the order $32$ remains the normal subgroup $K_{32}$ of A.3.
\end{statement}
\begin{proof}
Modulo four, $q_4(x)=U^2$ and $q_4(y)=W^2$ commute, so
$q_4(f)=I$. Hence $q_4\Phi(\beta_2)=W^\mu$, whose upper-left
entry is one. The formula in A.17 gives $t\equiv0\pmod4$.
Thus the whole $\widehat{\mathrm{GT}}_0$ image lies in
$\mathcal V_{12}$. This set is a group under matrix multiplication
and has twelve elements. To prove equality for the arithmetic
subgroup, we now compute a compatible finite geometric cover,
including its base-point comparison.

Consider the Legendre curve
$E_\lambda:y^2=x(x-1)(x-\lambda)$ over
$\mathbb P^1-\{0,1,\infty\}$, and its finite etale cover of
cyclic order-three subgroups. Its four points are the four
$x$-coordinates of nonzero three-torsion modulo sign. With the
oriented vanishing cycles above the intervals $[0,\lambda]$ and
$[\lambda,1]$, a positive loop around either colliding pair of
branch points is a full twist. On the double cover it acts by the
square of the transvection in the corresponding vanishing cycle.
Taking their intersection to be one gives the matrices
$U^2$ and $W^2$. This constructs the geometric monodromy comparison
from $x,y$ in (A.17.1) to the projectivization of $E[3]$.
Changing the simultaneous homology basis conjugates both sides;
the four-point permutation representation and its coefficient
field are unchanged. The central signs used to define $\lambda_i$
disappear under this specified projectivization.

We also specify the arithmetic fibre functor. At the tangent
$\lambda=0,+1$, the fibre is obtained by adjoining formal roots
$\lambda^{1/m}$ and letting $G_{\mathbb Q}$ act on their coefficients,
fixing the symbols $\lambda^{1/m}$. The real interval path gives the
usual second generator. In every abelian Kummer covering generated
by $\lambda^{1/m}$ and $(1-\lambda)^{1/m}$, the leading coefficients
at the two tangents are one. Thus the path cocycle has zero
abelianization; the induced based Galois action has precisely
$x\mapsto x^\chi$, $y\mapsto f^{-1}y^\chi f$ with $f\in\widehat F_2'$.
The arithmetic compatibility and the inclusions
$G_{\mathbb Q}\subset\widehat{\mathrm{GT}}\subset
\widehat{\mathrm{GT}}_0$ are the tangential actions described in
\cite[\S0.3 and \S1.2--1.3]{A_HarbaterSchneps}.
Consequently the arithmetic conjugation of this four-point
Legendre monodromy is the projective reduction of $V_{12}$.
There is no extra projective conjugator ambiguity here: a matrix
centralizing the projective images of $U^2,W^2$ modulo three
centralizes their actual matrices, since their trace two is
nonzero; comparing entries then makes it scalar.

The three-division polynomial, obtained from the tangent-and-chord
addition formula on the displayed cubic, is
\[
 \psi_3(x)=3x^4-4(1+\lambda)x^3+6\lambda x^2-\lambda^2. \tag{A.18.3}
\]
At zero it is $x^3(3x-4)$. The root with constant term $4/3$
lies in $\mathbb Q[[\lambda]]$, by the formal implicit-function
theorem. For the other roots put $\lambda=s^3$, $x=s^2z$.
Dividing the polynomial by $s^6$ gives
\[
 3s^2z^4-4(1+s^3)z^3+6sz^2-1.                    \tag{A.18.4}
\]
Each root $c$ of $c^3=-1/4$ is simple at $s=0$, and uniquely
lifts to $z\in\mathbb Q(c)[[s]]$. Hence the exact field of all
Puiseux coefficients is
$L=\mathbb Q(\zeta_3,\sqrt[3]{2})$: its containment follows by
the unique recursion for the coefficients, and the reverse
containment follows from the three leading coefficients, because
$( -2c)^3=2$. The polynomial $X^3-2$ is irreducible by Eisenstein
at two, and its discriminant $-108$ is not a square. Therefore
$[L:\mathbb Q]=6$ and the arithmetic action on the four lines
has image $S_3$, fixing the one rational branch.

Modulo four, $V_{12}$ is $\operatorname{diag}(\chi,1)$, so its
remaining field is $\mathbb Q(i)$. The field $L$ has unique
quadratic subfield $\mathbb Q(\sqrt{-3})$, by the subgroup lattice
of $S_3$, and therefore $L\cap\mathbb Q(i)=\mathbb Q$.
The pair consisting of $V\bmod4$ and $[V\bmod3]$ determines
every $V\in\mathcal V_{12}$: a scalar with lower-right entry one
is the identity, and Chinese remainders recover both entries.
Thus the arithmetic kernel fixes exactly the degree-twelve field
$L(i)=\mathbb Q(\zeta_{12},\sqrt[3]{2})$, proving (A.18.1)--(A.18.2).
The two containing GT groups have the same image by the already
proved upper bound. The formal polynomial calculations (A.18.3)--
(A.18.4), the leading field, and all twelve matrices are checked
in \texttt{a\_braid.py}.

The centralizer of $\mathrm{SL}_2(\mathbb Z/12)$ in
$\mathrm{GL}_2(\mathbb Z/12)$ consists of its four scalar units.
The normalization $V_{22}=1$ makes their intersection with
$\mathcal V_{12}$ trivial. This proves the genuine automorphism
image count and the trivial outer image on the full GL group.
A conjugation by $V$ is inner on SL exactly when
$V\in\mathrm{SL}_2(\mathbb Z/12)\cdot
((\mathbb Z/12)^*I)$; taking determinants gives its unit square
class. All four units modulo twelve square to one, proving the
outer SL count. The scalar action on $M/12M$ follows from the
generator comparison at the end of A.17, and cyclotomic
surjectivity realizes all four scalars. Finally the geometric
image is SL and the determinants of $\mathcal V_{12}$ are all
units, so $\langle\mathrm{SL}_2(\mathbb Z/12),\mathcal V_{12}\rangle
=\mathrm{GL}_2(\mathbb Z/12)$. Quotienting its order $4608$ by
the scalar centralizer of order four gives $1152$.

This specifies which choices give a genuine action and which
quotients remove them. The congruence action is obtained from
$\widehat B_3$ and its modular quotient; A.4 proves that the
reduction to the unquotiented GL or SL group cannot instead be
defined by first passing through $\mathfrak S_{\rm d}$.
No equality between the normalized upper-right entry and the
Kummer cocycle based at $1/2$ is used in this proof.
\end{proof}

\begin{statement}{A.19}{Proved}{The normalized cyclotomic arithmetic image in the minimum approximation}
In the tangential normalization of A.17--A.18, the arithmetic
homomorphism into the final \S51 minimum approximation has image in
$\mathcal M(0)$: its effective coordinate is $U=u_{\underline\mu}$.
Its restriction to $G_{\mathbb Q^{\rm cyc}}$ has image in the abelian
group
\[
 \ker(\mathcal M(0)(\widehat{\mathbb Z})^{\Delta}
                   \longrightarrow\widehat{\mathbb Z}^{*})
       =C_2\times\widehat{\mathbb Z}.              \tag{A.19.1}
\]
After the central quotient by $\mu_\rho$, this kernel is
$\widehat{\mathbb Z}$ under addition. In particular the closed
commutator subgroup of $G_{\mathbb Q^{\rm cyc}}$ maps to the identity
in this normalized minimum approximation. Here the superscript
$\Delta$ requires the single global component choice described in
A.1; it is essential to (A.19.1).
\end{statement}
\begin{proof}
The congruence-normalized braid action has conjugator
$f\in\widehat F_2'$. In the modular quotient its image is in
$\widehat\Pi_0'$ by the evaluation calculation in A.17, so it
vanishes in the minimum approximation. Thus both elementary
generators $U,W$ are sent to their $\mu$-th powers.
The automorphism (A.8.5) has exactly these images on diagonal
profinite points: from $U=(\sigma\rho)^{-1}$,
$W=(\rho\sigma)^{-1}$ and
$\lambda_0=\omega U^{-2}$, $\lambda_1=\omega W^{-2}$,
the possible central factor is $\omega^{\nu+a}=1$, since
$\nu=(\mu-1)/2\equiv a\pmod2$.
This proves $U=u_{\underline\mu}$ for the effective coordinate,
which is precisely the graph defining $\mathcal M(0)$.
The retained normalizer conjugators give its other two coordinates
as in (183), so its projection to those coordinates is the
isomorphism of A.15.

For a diagonal cyclotomic point with $\mu=1$, its residue is one
and the unique quotient $(\mu-1)/12$ is zero in
$\widehat{\mathbb Z}$: hence $a=b=k=0$. Formula (A.14.2) on this
fibre is $(i,s)(i',s')=(i+i',s+s')$. This proves (A.19.1), its
central quotient, and the commutator assertion. It also explains
why replacing diagonal profinite points by arbitrary
$\gamma(\mathbb F_2)$-points would give the wrong kernel.
\end{proof}

\section{B. The Fermat tower and its homology}

We use the arithmetic model actually written in (220):
\[
 X_n^+:\ x_0^n+x_1^n+x_\infty^n=0,\qquad
 D_n=\{x_0x_1x_\infty=0\},\qquad U_n^+=X_n^+\setminus D_n.
\]
The notation \(X_n:x^n+y^n=z^n\), used in our geometric chain
calculations, is related to this model explicitly in B.1. Throughout this
part, \(F=\langle x,y\rangle\), \(c=[x,y]=xyx^{-1}y^{-1}\), and
\(F''=[F',F']\). In a profinite group, a commutator subgroup means its
closed version. Put
\[
 A_n=(\mathbb Z/n)^2,\quad
 R_n=\mathbb Z[A_n]=
 \mathbb Z[X^{\pm1},Y^{\pm1}]/(X^n-1,Y^n-1),
 \quad K_n=\ker(F\longrightarrow A_n).
\]
Capital letters \(X,Y\) in a group ring denote the images of \(x,y\).
We write
\(M=\mathbb Z^3/\mathbb Z(1,1,1)\), with basis the classes \(e_0,e_1\)
and \(e_\infty=-e_0-e_1\).

\begin{statement}{B.1}{Proved}{Arithmetic models and the precise problem}
The morphism
\[
 q_n:U_n^+\longrightarrow U_1^+,\qquad
 [x_0:x_1:x_\infty]\longmapsto[x_0^n:x_1^n:x_\infty^n]
\]
is defined over \(\mathbb Q\), is equivariant for coordinate permutations,
and has deck group \(\mu_n^3/\mu_n\). The coordinate
\[
 z=-x_0/x_\infty,\qquad
 \psi(z)=[-z:z-1:1]
\]
identifies \(U_1^+\) with \(\mathbb P^1-\{0,1,\infty\}\).
For \(\alpha_n^n=-1\), the geometric comparison
\[
 \iota_n:X_n^+\longrightarrow X_n,\qquad
 [x_0:x_1:x_\infty]\longmapsto[x_0:x_1:\alpha_n x_\infty]
\]
is an isomorphism over \(\mathbb Q(\alpha_n)\).
At odd levels one can take \(\alpha_n=-1\) over \(\mathbb Q\);
at level four the displayed comparison uses \(\alpha_4=\zeta_8\).
In particular \(X_4^+(\mathbb R)=\varnothing\), whereas
\([0:1:1]\in X_4(\mathbb Q)\).

The \emph{position du probl\`eme} in \S51.11 is to identify the affine
representation
\[
 \widetilde{\mathcal M}\longrightarrow
 G_{\rho,\sigma}/\mu_G
 \longrightarrow \widehat{\mathbb Z}^{\,*}\times S_3
\]
with the action on a fibre of the universal abelian Fermat covering,
and separately to interpret its specialization at \(z=1/2\).
Its translation group is \(\widehat M(1)\).
\emph{Manuscript: \S51.11, folios 695--697, scans 95--99, (218)--(230);
\S51.14, folio 723, scan 151, (315)--(317).}
\end{statement}
\begin{proof}
The asserted equations follow on taking \(n\)-th powers. On the chart
\(x_\infty\ne0\), division by \(\alpha_n\) gives
\((x_0/\alpha_nx_\infty)^n+(x_1/\alpha_nx_\infty)^n=1\);
the induced coordinate on the base is \(z\). Compatible choices, for
example \(\alpha_n=\exp(\pi i/n)\) after fixing an embedding in
\(\mathbb C\), commute with the power maps. Thus all later geometric
comparisons have an actual map behind them. The even-level real assertion
follows because a sum of three real fourth powers vanishes only when
all three coordinates vanish. Equations (228)--(230) explicitly identify
the deck lattice and its cyclotomic action; (219) singles out the remaining
rank-one affine specialization. These are the data just stated.
\end{proof}

\begin{statement}{B.2}{Proved}{The universal abelian covering}
Over \(\mathbb C\), \(q_n\) corresponds to the subgroup \(K_n\) of
\(\pi_1(U_1)=F\). Its monodromy is \(x\mapsto(1,0)\),
\(y\mapsto(0,1)\). Consequently the pro-\'etale tower
\(\varprojlim_n U_n^+\) is the universal abelian geometric covering, with
deck module \(\widehat M(1)\). Its deck and coordinate-permutation
actions form \((\mu_n^3/\mu_n)\rtimes S_3\) at finite level.
\emph{Manuscript: \S51.11, folios 696--697, scans 97--99, (220)--(230);
\S51.13, folio 715, scan 135.}
\end{statement}
\begin{proof}
In the standard geometric chart, \(u^n=z\) and \(v^n=1-z\). Analytic
continuation around \(0\) multiplies \(u\) by \(\zeta_n\) and leaves \(v\)
unchanged; continuation around \(1\) does the reverse. These two elements
generate \(A_n\), so the covering is connected and its subgroup is exactly
\(K_n\). Every finite quotient of \(\widehat{\mathbb Z}^2\) factors through
some \(A_n\); this proves the universal assertion, rather than only the
existence of a collection of abelian coverings. A change of compatible
geometric roots identifies the deck group with \(\widehat{\mathbb Z}^2\)
by a change of basis; intrinsically it is
\((\widehat{\mathbb Z}^3/\widehat{\mathbb Z})\otimes
\widehat{\mathbb Z}(1)\). Permutations act on the three coordinates,
commute with power maps, and conjugate the diagonal multipliers by
permuting them. This proves the stated semidirect product and the
compatibility of the action with the inverse system.
\end{proof}

\begin{statement}{B.3}{Refuted}{The withdrawn freeness and its exact replacement}
For \(n>1\), the action of \(\mu_n^3\) on \(X_n^+\) is not free:
its diagonal \(\mu_n\) fixes the whole curve. The quotient group scheme
\(A_n^{\rm geom}=\mu_n^3/\mu_n\) acts freely on \(U_n^+\), and
\(U_n^+\to U_1^+\) is an \(A_n^{\rm geom}\)-torsor.
At a coordinate cusp its stabilizer in \(A_n^{\rm geom}\) is a copy of
\(\mu_n\). Thus the quotient action on the compact curve is not free.
These are assertions about actions on schemes; the corresponding
integral homology modules are computed in B.5--B.7.
\emph{Manuscript: \S51.11, folio \(696'\), scan 98, deletion before (224)
and the inserted \(U_n\) in the paragraph following (225).}
\end{statement}
\begin{proof}
A multiplier \((a_0,a_1,a_\infty)\) fixes a point with all coordinates
invertible precisely when the three multipliers agree projectively.
Its class in the quotient is then the identity. At a point
\([0:b:c]\), with \(bc\ne0\), the equality required is
\(a_1=a_\infty\); \(a_0\) is unrestricted. Modulo the diagonal this is
exactly \(\mu_n\). This calculation also describes the stabilizer
functors, and hence applies to group schemes, including their
nonreduced fibres.
On a torus chart the power map is given by adjoining independent
\(n\)-th roots of two units. Its coordinate algebra is locally free of
rank \(n^2\), and the standard root-of-unit torsor identity proves
faithful flatness and the torsor assertion. Over \(\mathbb Z[1/n]\)
it is finite \'etale. At residue characteristics dividing \(n\) it is
still a torsor for the fppf topology; this does not assert \'etaleness or
smoothness of those Fermat fibres.
The finite checks in \texttt{code/b\_freeness.py} enumerate the
full and quotient stabilizers over finite fields for \(1\le n\le8\).
They give full-group stabilizer order \(n\) on the open part,
quotient stabilizer order one there, and quotient stabilizer
order \(n\) at a cusp.
\end{proof}

\begin{statement}{B.4}{Proved}{The discrete Alexander module}
Conjugation induces a free rank-one
\(R=\mathbb Z[X^{\pm1},Y^{\pm1}]\)-module structure
\[
 F'/F''=R\,c.
\]
With the convention for \(c\) fixed above, the injective Fox-chain map is
\[
 R\,c\longrightarrow R^2,\qquad
 f c\longmapsto ((1-Y)f,(X-1)f),
\]
and its image is
\(\ker\bigl((X-1,\;Y-1):R^2\to R\bigr)\).
\end{statement}
\begin{proof}
The graph with one vertex and oriented edges \(x,y\) has fundamental
group \(F\). Its universal abelian covering has vertex module \(R\),
edge module \(R^2\), and the displayed boundary. Its first homology is
the abelianization of its fundamental group \(F'\). The chain of
\(xyx^{-1}y^{-1}\) is \((1-Y,X-1)\).
For a cycle \((a,b)\), reduction of
\((X-1)a+(Y-1)b=0\) modulo \(X-1\) gives
\((Y-1)\bar b=0\) in the domain \(\mathbb Z[Y^{\pm1}]\).
Thus \(b=(X-1)f\). Substitution and cancellation of the nonzero
divisor \(X-1\) give \(a=(1-Y)f\). The same cancellation proves
uniqueness. This supplies both the generator and the proof of freeness.
\end{proof}

\begin{statement}{B.5}{Proved}{All-cusp-deleted integral homology}
There are canonical \(R_n\)-linear descriptions
\[
 H_1(U_n,\mathbb Z)=
 \ker\left(R_n^2\xrightarrow{(X-1,\;Y-1)}R_n\right),             \tag{B.1}
\]
\[
 0\longrightarrow
 R_n/\mathbb Z\,\mathcal N_n
 \xrightarrow{\ f\mapsto f c\ }H_1(U_n,\mathbb Z)
 \xrightarrow{q_{n*}} n\mathbb Z^2\longrightarrow0,\qquad
 \mathcal N_n=\sum_{a,b=0}^{n-1}X^aY^b.                       \tag{B.2}
\]
The deck action on the right-hand term is trivial.
The group \(H_1(U_n,\mathbb Z)\) has rank \(n^2+1\), and over
\(\mathbb C\) the trivial deck character occurs twice and every other
character occurs once.
\end{statement}
\begin{proof}
The finite covering of the two-circle graph has \(n^2\) vertices and
\(2n^2\) edges, with boundary as in (B.1). A spanning tree proves that
its cycle group is free of rank \(2n^2-n^2+1\).
The inclusion \(K_n\subset F\) induces \(K_n^{\rm ab}\to F^{\rm ab}\).
Its image is \(n\mathbb Z^2\); its kernel is the image of \(F'\),
because \(F'\subset K_n\) and \(K_n/F'=n\mathbb Z^2\).
B.4 identifies that image with the orbit module generated by \(c\).
The kernel of \(R_n\to R_n^2,\ f\mapsto ((1-Y)f,(X-1)f)\)
consists of elements fixed by both translations. Equality of all
coefficients says precisely that it is \(\mathbb Z\mathcal N_n\).
This proves (B.2), including its exactness over \(\mathbb Z\).
On a deck character \((\zeta_n^a,\zeta_n^b)\), the boundary in
(B.1) has rank one unless \(a=b=0\); in that case it has rank zero.
This proves the character formula directly.
\end{proof}

\begin{statement}{B.6}{Proved}{The compact integral homology module}
Define
\[
 S_X=\sum_{i=0}^{n-1}X^i,\quad S_Y=\sum_{j=0}^{n-1}Y^j,\quad
 P_n=\sum_{0\le j\le i<n}X^iY^j,\quad
 I_n=(S_X,S_Y,P_n)\subset R_n.
\]
The compact homology has the exact integral presentation
\[
 H_1(X_n,\mathbb Z)\simeq R_n/I_n,\qquad 1\longmapsto c.        \tag{B.3}
\]
Equivalently, (B.1) is divided by the \(3n\) cusp cycles
\[
 Y^j(S_X,0),\quad X^i(0,S_Y),\quad
 X^i(S_{XY},X S_{XY})\quad(0\le i,j<n),
 \qquad S_{XY}=\sum_{i=0}^{n-1}(XY)^i.                        \tag{B.4}
\]
The cusp subgroup is primitive of rank \(3n-1\).
The rank of (B.3) is \((n-1)(n-2)\), and its complex characters
are exactly those \((a,b)\) with \(a,b,a+b\ne0\pmod n\), each once.
\end{statement}
\begin{proof}
Filling a puncture kills its oriented peripheral loop in first homology.
Lifting the peripheral words \(x^n,y^n,(xy)^n\) gives precisely (B.4).
The homology sequence of the pair consisting of the compact surface and
its punctured surface identifies the subgroup of peripheral classes
with \(\mathbb Z[D_n]/\mathbb Z\sum_{P\in D_n}P\).
The quotient is the integral homology of an oriented compact surface,
so is torsion-free. This proves primitivity and rank \(3n-1\).

We give an integral derivation of the ring presentation. First fill
the \(2n\) cusps over \(0,1\). The resulting affine curve
\(u^n+v^n=1\) retracts onto the inverse image of the segment \([0,1]\).
To see this, take a closed regular neighbourhood of that segment
in the \(z\)-plane. Outside it the finite map is unramified and its
ends are collars. Inside it the local charts at the only branch
points are \(w\mapsto w^n\); their lifted discs and the intervening
strips are a regular neighbourhood of the lifted segment.
Shrinking these collars, discs and strips gives the stated retraction.
That graph has vertices \(A_n/\langle X\rangle\) and
\(A_n/\langle Y\rangle\), and one edge for each element of \(A_n\);
it is the complete bipartite graph on two sets of \(n\) vertices.
Its cycles are the integral matrices whose row and column sums vanish.
The map
\[
 R_n\longrightarrow\{\text{such matrices}\},\qquad
 f\longmapsto(1-X)(1-Y)f
\]
is surjective: subtracting the last row and column expresses every such
matrix as an integral sum of elementary rectangles. Its kernel is
\((S_X,S_Y)\), as one sees by successively taking differences of rows
and columns. The elementary rectangle is the commutator cycle, up to
the fixed choice of its orientation.

Now fill the \(n\) remaining cusps. In the original two-edge chain
module the identities
\[
 (1-Y)P_n=S_X-Y S_{XY},\qquad
 (X-1)P_n=S_Y-S_{XY}
\]
give
\[
 (S_{XY},X S_{XY})
 =-Y^{-1}P_n\,c+Y^{-1}(S_X,0)+Y^{-1}(0,S_Y).
\]
All infinity cusp cycles are translates of this one. Thus their
classes in the bipartite graph kill exactly the ideal generated by
\(P_n\), proving (B.3).
For a character with both \(a,b\ne0\), the last two displayed identities
make \(P_n\) zero precisely when \(a+b\ne0\); when \(a+b=0\),
\((X-1)P_n=-n\) is nonzero. The character statement follows. Its
cardinality is \((n-1)(n-2)\), agreeing with the topological rank.
\end{proof}

\begin{statement}{B.7}{Proved}{Finite integral computations and module freeness}
The exact calculations for \(1\le n\le8\) give
\[
\begin{array}{c|rrrrrrrr}
n&1&2&3&4&5&6&7&8\\ \hline
g(X_n)&0&0&1&3&6&10&15&21\\
\operatorname{rk}H_1(U_n)&2&5&10&17&26&37&50&65\\
\operatorname{rk}H_1(X_n)&0&0&2&6&12&20&30&42\\
\operatorname{rk}\langle\text{cusps}\rangle&2&5&8&11&14&17&20&23
\end{array}
\]
For every \(n>1\), \(H_1(U_n,\mathbb Z)\) is not a free \(R_n\)-module.
For every \(n\ge3\), \(H_1(X_n,\mathbb Z)\) is cyclic but not free
as an \(R_n\)-module. At \(n=1,2\) the compact module is zero.
The exact comparison maps, deck actions and \(S_3\)-actions are supplied
as integer matrices in \texttt{code/b\_fermat\_modules.json}.
\end{statement}
\begin{proof}
The program builds the lifted graph, chooses the tree obtained from the
paths \(x^i y^j\), and uses each remaining edge as a fundamental cycle.
It checks the Fox boundary, computes the deck transformations and the
two based substitutions in B.9, and expresses every cusp cycle in
this integral basis. It checks that the nonzero Smith entries of the
cusp matrix and of the compact annihilator presentation are all one.
It also checks that the map to \(F^{\rm ab}\) has Smith entries \(n,n\).
These are full integer presentations, rather than character counts.
All output rows displayed here were generated and checked by that program.

A free \(R_n\)-module has integral rank divisible by \(n^2\).
The affine rank \(n^2+1\) fails this test for \(n>1\).
The positive compact rank \((n-1)(n-2)\) is strictly below \(n^2\).
These rank conclusions apply through the explicit presentations
(B.1) and (B.3); they do not replace the scheme-action assertion of B.3.
\end{proof}

\begin{statement}{B.8}{Proved}{The full profinite rank-one theorem}
Let \(\widehat F\) be the free profinite group on \(x,y\), and put
\[
 \Lambda=\widehat{\mathbb Z}[[\widehat{\mathbb Z}^{\,2}]]
 =\varprojlim_{m,n}(\mathbb Z/m)[A_n].
\]
There are continuous \(\Lambda\)-linear isomorphisms
\[
 \widehat F'/\widehat F''
 \xrightarrow{\ \sim\ }\varprojlim_n H_1(U_n,\widehat{\mathbb Z})
 \xrightarrow{\ \sim\ }\varprojlim_n H_1(X_n,\widehat{\mathbb Z})
 \simeq\Lambda\,c.                                          \tag{B.5}
\]
The transition maps on homology are pushforwards of the covering maps.
The same limits can be formed diagonally with coefficients
\(\mathbb Z/n\). This is the full profinite assertion, including
the prime-to-\(\ell\) factors of the completed group ring.
\end{statement}
\begin{proof}
We first prove the group-to-tower comparison. Let
\(\widehat K_n=\ker(\widehat F\to A_n)\) and
\(H=\bigcap_n\widehat K_n=\widehat F'\).
The map \(H\to\varprojlim_n\widehat K_n^{\rm ab}\) is surjective:
compatible abelianized cosets are nonempty compact subsets of
\(\widehat K_n\), nested along a cofinal divisibility chain; their
intersection supplies a lift in \(H\).
Its kernel contains \(H'\). To prove the reverse inclusion, take a
continuous homomorphism \(f:H\to C\) to a finite abelian group. Choose
an open normal subgroup \(V\) of \(\widehat F\) with
\(V\cap H\subset\ker f\). Since \(H\) is normal,
\([V,H]\subset V\cap H\), and \(f\) extends to
\(HV\) by \(f(hv)=f(h)\). Compactness gives
\(\widehat K_n\subset HV\) for a sufficiently divisible \(n\).
It follows that every element in
\(\bigcap_n\widehat K_n'\) is annihilated by \(f\). Finite abelian
quotients separate \(H/H'\), proving the first isomorphism.
The covering graph identifies
\(\widehat K_n^{\rm ab}\) with \(H_1(U_n,\widehat{\mathbb Z})\):
the finite-index topology on \(K_n\subset F\) is its full profinite
topology, because the intersection of the finitely many \(F\)-conjugates
of a finite-index subgroup of \(K_n\) is finite-index normal in \(F\).
The extra quotient \(n\widehat{\mathbb Z}^2\) in (B.2) has transition
maps given by inclusion and has inverse limit
\(\bigcap_n n\widehat{\mathbb Z}^2=0\). Thus its rank-two contribution
at every finite level does not survive in (B.5).

For the second isomorphism, work first modulo \(m\).
The kernel of \(H_1(U_n,\mathbb Z/m)\to H_1(X_n,\mathbb Z/m)\)
is the peripheral module of B.6. Under \(U_{dn}\to U_n\), a small
loop around a cusp maps to \(d\) times a small loop. Thus the image of
that kernel is zero when \(m\mid d\). The inverse system is pro-zero.
All groups are finite, so taking the inverse limit preserves the short
exact sequences. This gives the second isomorphism modulo \(m\),
and passage to the inverse limit over \(m\) gives the assertion.

Finally use (B.3) modulo \(m\). It stays exact because its integral
quotient is free. Under \(R_{dn}\to R_n\),
\[
 S_{X,dn}\mapsto dS_{X,n},\quad
 S_{Y,dn}\mapsto dS_{Y,n},\quad
 P_{dn}\mapsto dP_n+\binom d2\mathcal N_n.
\]
The last formula counts pairs \(0\le j\le i<dn\) with each prescribed
pair of residues modulo \(n\). A value of \(d\) divisible by \(2m\)
kills every image generator of \(I_{dn}\) modulo \(m\).
Hence \(\varprojlim_n I_n(\mathbb Z/m)=0\), and the finite inverse-limit
exact sequence identifies \(\varprojlim_n H_1(X_n,\mathbb Z/m)\)
with \(\varprojlim_n(\mathbb Z/m)[A_n]\), by the compatible generator
\(c\). Passing to \(m\) proves freeness. This argument agrees with
\cite[Theorem 6.1]{Otsubo2016}, whose full profinite statement was
opened and checked. It does not use a power-series-domain argument.
The diagonal \((m,n)=(r,r)\) is cofinal in the divisibility product,
which proves the coefficient formulation in the statement.
\end{proof}

\begin{statement}{B.9}{Proved}{The exact permutation comparison}
On \(\Lambda c\) the selected \(S_3\)-action is semilinear:
\[
 r(f(X,Y)c)=f(Y,X^{-1}Y^{-1})c,\qquad
 s(f(X,Y)c)=-f(Y,X)c,                                       \tag{B.6}
\]
where \(r^3=s^2=(rs)^2=1\). It agrees with the coordinate-permutation
action on the Fermat tower after a compatible deck translation.
Changing that translation conjugates the action by the corresponding
deck operator. The same comparison is represented integrally at every
finite level by the matrices of B.7.
\end{statement}
\begin{proof}
Conjugation by the modular generators acts on \(F\) by
\(\rho:x\mapsto y,\ y\mapsto x^{-1}y^{-1}\) and
\(\sigma:x\mapsto y,\ y\mapsto x\).
Replace \(\rho\) by \(x\rho\); its based substitution is
\[
 r:x\mapsto xyx^{-1},\quad y\mapsto y^{-1}x^{-1}.
\]
Direct free reduction gives \(r(c)=c\) and \(s(c)=c^{-1}\).
Their actions on \(F^{\rm ab}\) are
\[
 R=\begin{pmatrix}0&-1\\1&-1\end{pmatrix},\qquad
 S=\begin{pmatrix}0&1\\1&0\end{pmatrix}.
\]
The defects in \(r^3\) and \((rs)^2\) are inner automorphisms by
elements of \(F'\). Such automorphisms act trivially on \(F'/F''\)
and on \(K_n^{\rm ab}\), since \(F'\subset K_n\). Formula (B.6)
therefore gives a genuine \(S_3\)-action.
At the level of deck extensions this is the complement constructed
in B.11. The coordinate permutations provide another complement
of the same \(A_n\)-extension. B.10 proves that there is a unique
deck translation conjugating the two. Uniqueness makes these
translations compatible as \(n\) varies. Passing to homology and
then to (B.5) gives the claimed comparison, with its ambiguity stated.
\end{proof}

\begin{statement}{B.10}{Proved}{The two cohomology groups governing uniqueness}
For every abelian group \(I\), with trivial \(S_3\)-action on \(I\),
put \(M_I=I^3/I(1,1,1)\). Then
\[
 H^0(S_3,M_I)=H^1(S_3,M_I)=0.                               \tag{B.7}
\]
The same assertions hold for continuous cohomology with a profinite
\(I\). For any profinite group \(B\) acting on \(I\) and commuting
with the permutation action,
\[
 H^0(B\times S_3,M_I)=H^1(B\times S_3,M_I)=0.
\]
\emph{Manuscript: \S51.12--13, folios 702, 715--716, scans 109,
136--138, (295)--(297); \S51.14, folio \(721'\), scan 148, (309).}
\end{statement}
\begin{proof}
Use coordinates in which \(r,s\) act by \(R,S\) above.
An invariant \((u,v)\) satisfies \(u=v\), \(u=-v\), and
\(u=2v\). Thus \(2v=3v=0\), so \(u=v=0\), over every \(I\).
A crossed homomorphism is specified by \(f(r)=a=(a_0,a_1)\),
\(f(s)=b\). The relation \(s^2=1\) gives \(b=(t,-t)\).
The relation \(r^3=1\) imposes no condition, because \(1+R+R^2=0\).
The relation \((rs)^2=1\) gives
\[
 a_0=2a_1+3t.
\]
The element \(m=(-a_1-2t,-a_1-t)\) satisfies
\((R-1)m=a\) and \((S-1)m=b\).
Hence the cocycle is a coboundary, and the conjugating element is
unique by the invariant calculation. Every formula uses only integral
addition, so it also proves the continuous profinite assertion.
For \(B\times S_3\), first subtract the unique coboundary that kills
the restriction to \(S_3\). The cocycle identity for commuting
\(b\in B,s\in S_3\) then says that \(f(b)\) is \(S_3\)-invariant.
It is zero. This proves the last assertion without a spectral-sequence
convergence assumption.
\end{proof}

\begin{statement}{B.11}{Proved}{An explicit integral splitting}
The extension \(\mathfrak S^+/\mu_{\mathfrak S}\) of \(S_3\) by \(M\)
splits. In the affine realization
\[
 \rho=(R,0),\qquad \sigma=(S,t),\qquad
 t=(1/3,-1/3),
\]
a complement is generated by
\[
 \rho'=\lambda_0\rho=(R,(1,0)),\qquad \sigma'=\sigma.
\]
It fixes \(p=(2/3,1/3)\). In particular its integral extension
class, and the profinite class after completion, are zero.
\emph{Manuscript: \S51.12, folios 700--702, scans 105--109;
\S51.13, folio \(715'\), scan 136.}
\end{statement}
\begin{proof}
For affine multiplication \((A,u)(B,v)=(AB,u+Av)\),
\(1+R+R^2=0\) gives \((\rho')^3=1\),
\((1+S)t=0\) gives \((\sigma')^2=1\), and
\[
 RS=\begin{pmatrix}-1&0\\-1&1\end{pmatrix},\qquad
 (1,0)+Rt=(4/3,2/3),\qquad
 (1+RS)(4/3,2/3)=0
\]
gives \((\rho'\sigma')^2=1\). Their linear parts generate all of
\(S_3\), so this is a section. Its translations lie in the original
affine cosets, since the only changed generator was multiplied by
\(\lambda_0\in M\). Solving \((1-R)p=(1,0)\) and
\((1-S)p=t\) gives the stated fixed point.

For complete cocycle data choose \(h=r^i s^j\) and put
\(t_h=R^i(jt)\). The original integral cocycle is
\[
 c(h,k)=t_h+A_h t_k-t_{hk}.
\]
The section just found has translation \(p-A_hp\).
Thus the integral cochain \(b_h=p-A_hp-t_h\) satisfies
\(c+\delta b=0\). This proves vanishing as a cohomology class,
not merely satisfaction of a few relations. All formulas commute
with reduction modulo \(N\) and with profinite completion.
\end{proof}

\begin{statement}{B.12}{Proved}{Splitting and uniqueness for the full quotient}
The continuous extension
\[
 1\longrightarrow\widehat M\longrightarrow
 G_{\rho,\sigma}/\mu_G\longrightarrow
 \widehat{\mathbb Z}^{\,*}\times S_3\longrightarrow1
\]
has zero class in continuous \(H^2\). Its set of continuous splittings
is a torsor under \(\widehat M\), by conjugation.
For the canonical scalar lift \(u_\mu\) of the manuscript, put
\(\nu=(\mu-1)/2\in\widehat{\mathbb Z}\). A section commuting with the
complement of B.11 is
\[
 s(\mu)=(-\nu,-\nu)\,u_\mu.
\]
\emph{Manuscript: \S51.12, folios 704--\(704'\), scans 113--114,
(264)--(266); \S51.13, folios 715--\(716'\), scans 135--138.}
\end{statement}
\begin{proof}
The manuscript's scalar automorphism sends
\(\rho\) to \(\lambda_1^\nu\rho\) and fixes \(\sigma\) modulo the
central involution, and it acts by multiplication by \(\mu\) on \(M\).
In the affine realization it is conjugation by
\[
 (\mu I,w_\mu),\qquad w_\mu=(-\nu/3,\nu/3).
\]
Indeed \((1-R)w_\mu=\nu(0,1)\) and
\(\mu t+(1-S)w_\mu=t\). The product law follows from
\(\nu(\mu\mu')=\nu(\mu)+\mu\nu(\mu')\).
After multiplying by the integral translation \((-\nu,-\nu)\),
its affine translation becomes \((1-\mu)p\).
These are precisely the scalar homotheties about \(p\); they commute
with the finite complement fixing \(p\), and multiply according to
\(\mu\mu'\). This gives a continuous section of the full product.
The thirds in this verification are computed in
\(\widehat{\mathbb Z}\otimes_{\mathbb Z}\mathbb Q\); the correcting
cochain is integral at every prime, including \(3\).
The half is integral because every profinite unit is odd at \(2\).

Two sections differ by a continuous \(1\)-cocycle with coefficients
\(\widehat M\), acted on by \((\mu,h)\) as \(\mu A_h\).
B.10 makes that cocycle a coboundary and makes the conjugating
element unique. Thus the asserted torsor and the vanishing of the
actual \(H^2\)-class both hold.
\end{proof}

\begin{statement}{B.13}{Proved}{The finite \(2\)- and \(3\)-primary checks}
For \(N\mid24\), the actual extension class of B.12 in coefficients
\(M/NM\) is zero, with trivial \(2\)- and \(3\)-primary parts.
The independent normalized-cochain computation gives
\[
 \left|H^0(S_3,M/NM)\right|=\left|H^1(S_3,M/NM)\right|=1,
 \qquad \left|H^2(S_3,M/NM)\right|=\gcd(N,3).
\]
In particular the zero answer for the actual class at \(N=3\)
is a genuine cocycle computation in a nonzero obstruction group.
\end{statement}
\begin{proof}
The integral cochain of B.11 is included in
\texttt{code/b\_splitting.json}. For the full product we use the
reference affine lift
\[
 t_{\mu,h}=w_\mu+\mu t_h,\qquad
 b_{\mu,h}=(-\nu,-\nu)+\mu b_h.
\]
The program verifies \(c+\delta b=0\) modulo \(N\) for every pair of
elements of \((\mathbb Z/6N)^*\times S_3\). The modulus \(6N\)
retains the halves and thirds in these reference formulas.
For \(N=1,2,3,4,6,8,12,24\) the numbers of checked pairs are,
respectively,
\[
 144,\ 576,\ 1296,\ 2304,\ 5184,\ 9216,\ 20736,\ 82944.
\]
The normalized differentials \(C^0\to C^1\to C^2\to C^3\)
are integer matrices. Their nonzero Smith diagonals are
\(1^2,\ 1^8,\ 1^{41},3\).
For an integer matrix with nonzero Smith entries \(d_i\), its kernel
modulo \(N\) has cardinality \(N^{b-r}\prod_i\gcd(N,d_i)\),
where \(b\) is its column count and \(r\) its rank.
Applying this formula to those three matrices gives the displayed
cohomology orders. The code checks consecutive differentials compose
to zero and verifies the full-product cocycle by exact rational
arithmetic before reduction.
\end{proof}

\begin{statement}{B.14}{Proved}{Identification and rigidity of the arithmetic tower}
Let
\[
 E=\pi_1^{\mathrm{et}}([U_1^+/S_3],\bar b),\qquad
 E_{\rm ab}=E/\overline{[\pi_1(U_{1,\overline{\mathbb Q}}),
                              \pi_1(U_{1,\overline{\mathbb Q}})]}.
\]
Its exact extension
\[
 1\to\widehat M(1)\to E_{\rm ab}\to\Gamma_{\mathbb Q}\times S_3\to1
\]
splits, with one conjugacy class of splittings and a unique
\(\widehat M(1)\)-conjugator between any two. The affine representation
on a Fermat fibre has target
\(\widehat M\rtimes(\widehat{\mathbb Z}^{\,*}\times S_3)\).
The representation supplied by the quotient
\(G_{\rho,\sigma}/\mu_G\), with the same geometric translation map
and cyclotomic character, agrees with it after a unique translation
of the fibre. This proves the identification posed in B.1.
\emph{Manuscript: \S51.13, folios 715--\(718'\), scans 135--142,
(294)--(298).}
\end{statement}
\begin{proof}
The coordinate-permutation action is defined over \(\mathbb Q\).
Consequently the arithmetic quotient of the orbifold group by the
geometric fundamental group is
\(\Gamma_{\mathbb Q}\times S_3\).
The fibres of \(U_n^+\to U_1^+\) form compatible torsors under
\(M/nM(1)\). Choose one point in their inverse-limit fibre.
Monodromy then acts by affine transformations with linear action
\(\chi(\gamma)A_h\), and its restriction to the geometric group is
the abelianization map of B.2. Thus it factors through \(E_{\rm ab}\).
The map from \(E_{\rm ab}\) to the pullback of this affine group over
\(\Gamma_{\mathbb Q}\times S_3\) is an isomorphism: it is the identity
on its kernel \(\widehat M(1)\) and on its quotient, so injectivity
and surjectivity follow by lifting a quotient element and correcting
in the kernel. The pullback of a semidirect product is split.

B.10 with \(B=\Gamma_{\mathbb Q}\), acting through \(\chi\), proves
uniqueness of the conjugacy class and of its conjugator.
The algebraically constructed representation and the Fermat
representation have the same restriction to \(\widehat M(1)\)
and the same quotient action. Their translation difference is
therefore a cocycle on \(\Gamma_{\mathbb Q}\times S_3\).
It is the unique coboundary given by B.10, which proves the
claimed comparison. This identifies the resulting unbased torsor
canonically; a based homomorphism still requires the fibre point
chosen above. The group \(E\) in this argument is the arithmetic
fundamental group, with quotient the actual Galois group.
No equality of that group with a group of all loop-preserving
outer automorphisms is used.

The source changes its tilde notation here. At scan 136 it describes
the geometric-abelian quotient, while at scans 139--140, in (298),
it uses \(\widetilde{\mathcal M}=\mathcal M/\mu_{\mathcal M}\)
for the orbifold group. Our notation retains the successive maps
\(\mathcal M\to E\to E_{\rm ab}\). Thus the argument uses the
displayed kernels, rather than identifying the two quotients
because they share a symbol in the manuscript.
\end{proof}

\paragraph{Scope of the two ranks.}
The map in B.14 first uses the quotient
\(\widehat F\to\widehat F^{\rm ab}=\widehat M(1)\).
The \(\Lambda\)-module in B.8 is obtained by taking the homology
\emph{of the kernels} of its finite quotients. This is the exact
passage from the universal abelian covering to the commutator
module seen by the metabelian fundamental group. It is this map
that connects the rank-two construction in \S51 to the universal
homology module.

\section{C. The relevant 1-motives and the cuspidal extension}

We keep the arithmetic plus model \(X_n^+\) and its reduced divisor
\(D_n\) of \(3n\) coordinate cusps. All groups and motives in this
part are over characteristic zero. The geometric comparison of B.1
transports the computations on \(x^n+y^n=z^n\) to this model; its
field and its effect on the arithmetic descent are retained.
Write \(J_n=\operatorname{Pic}^0(X_n^+)\) and
\(\mathcal J_n=\operatorname{Pic}^0(X_n^+,D_n)\), the degree-zero
generalized Jacobian with reduced modulus \(D_n\).

\begin{statement}{C.1}{Known}{The actual Fermat 1-motive and its realizations}
Put
\[
 L_n=\operatorname{Div}^0_{D_n}(X_{n,\overline{\mathbb Q}}^+),\qquad
 \mathcal P_n=[L_n\xrightarrow{D\mapsto\mathcal O(D)}J_n].
\]
This is a Deligne 1-motive, with its natural descent over \(\mathbb Q\).
Its Cartier dual, using the principal polarization of \(J_n\), is
\[
 \mathcal A_n=[0\longrightarrow\mathcal J_n].
\]
There is an exact sequence of semiabelian varieties
\[
 1\longrightarrow
 T_{D_n}:=(\operatorname{Res}_{D_n/\mathbb Q}\mathbb G_m)/\mathbb G_m
 \longrightarrow\mathcal J_n\longrightarrow J_n\longrightarrow1.
                                                               \tag{C.1}
\]
The character lattice of \(T_{D_n}\) is \(L_n\), and its cocharacter
lattice is \(\mathbb Z[D_n]/\mathbb Z\sum_{P\in D_n}P\).
The Betti and Tate realizations of \(\mathcal A_n\) are
\[
 T_{\mathbb Z}(\mathcal A_n)=H_1(U_n^+(\mathbb C),\mathbb Z),\qquad
 T_\ell(\mathcal A_n)=
 H_1^{\mathrm{et}}(U_{n,\overline{\mathbb Q}}^+,\mathbb Z_\ell).
                                                               \tag{C.2}
\]
Its de Rham realization is the dual of
\(H^1_{\mathrm{dR}}(U_n^+/\mathbb Q)\).
Its weights are \(-2\), with rank \(3n-1\), and \(-1\), with
rank \((n-1)(n-2)\); its total rank is \(n^2+1\).
\emph{Source: Deligne, \cite[10.1.3--10.1.8, 10.2.10--10.2.11,
10.3.4--10.3.8]{DeligneHodgeIII}.}
\end{statement}
\begin{proof}
We spell out the reduction to the cited construction. A point of the
generalized Picard group is a line bundle on \(X_n^+\), together with
a trivialization along \(D_n\), modulo scalar automorphisms of the
line bundle. The fibre over a fixed degree-zero line bundle is thus
the product of the nonzero fibres at the cusps modulo simultaneous
scaling. This proves (C.1) and both of its lattice descriptions.

Deligne's definition 10.3.4 assigns to a smooth open curve
\(X-D\) the complex
\([\ker(\mathbb Z[D]\to\mathbb Z)\to\operatorname{Pic}^0(X)]\).
For \(X=X_n^+\), this is precisely \(\mathcal P_n\).
Construction 10.3.6 identifies its \(m\)-torsion realization with
\(H^1(U_n^+,\mu_m)\): a representative is a pair
\((\mathcal L,D)\) with \(\mathcal L^{\otimes m}\simeq\mathcal O(D)\),
modulo pairs \((\mathcal O(E),mE)\). This is the same complex
defining the finite realization of \([L_n\to J_n]\).
Thus the comparison includes the cusp lattice and has no missing
toric summand. Construction 10.2.11 dualizes this complex to the
extension by \(\operatorname{Hom}(L_n,\mathbb G_m)\) whose class
evaluated at \(D\) is \(\mathcal O(D)\); the line-bundle description
above identifies this extension with \(\mathcal J_n\).
Dualizing the finite realizations with values in \(\mu_m\)
gives \(H_1(U_n^+,\mathbb Z/m)\), and their inverse limits give (C.2).

The Betti construction 10.1.3 and the comparison 10.3.8 give the
integral homology realization; 10.1.7--10.1.8 and 10.3.9 give the
de Rham comparison. Concretely, its dual Hodge filtration has
\[
 F^1H^1_{\mathrm{dR}}(U_n^+)
   =H^0(X_n^+,\Omega^1_{X_n^+}(\log D_n)),
 \quad \dim F^1=g_n+3n-1.
\]
The torus in (C.1) contributes weight \(-2\) and the abelian
variety weight \(-1\). Their ranks are \(3n-1\) and \(2g_n\),
respectively. B.6 gives \(2g_n=(n-1)(n-2)\).
Every construction used here is functorial for the deck and
permutation automorphisms, which proves compatibility of all
the realizations with those actions.
\end{proof}

\begin{statement}{C.2}{Refuted}{The rank-two identification in the lead}
The assertion that the rank-two module \(\widehat M(1)\) of \S51,
or its reduction modulo a nontrivial integer \(m\), is the full
Tate realization of the generalized Jacobian of \(X_n^+\) with its
cusps is false for \(n>1\).
Its exact replacement is
\[
 \widehat M(1)=T_\bullet(T),\qquad
 T=(\mathbb G_m^3)/\mathbb G_m=\mathcal J_1,
                                                               \tag{C.3}
\]
and the covering pushforward has the exact realization
\[
 0\longrightarrow
 (R_n/\mathbb Z\mathcal N_n)\otimes\widehat{\mathbb Z}
 \longrightarrow T_\bullet(\mathcal J_n)
 \xrightarrow{q_{n*}}n\widehat M(1)\longrightarrow0.             \tag{C.4}
\]
The Galois structure on the first term is the one transported by its
inclusion into the geometric homology in (C.2).
At level \(n=2\), the ranks in (C.4) are \(3,5,2\).
\emph{Manuscript: \S51.11, folio \(697\), scan 99, (228);
\S51.13--14, folios 719--720, scans 143--146, (299)--(305).}
\end{statement}
\begin{proof}
At \(n=1\) the compact curve is a projective line, so its ordinary
Jacobian vanishes. Formula (C.1) then gives (C.3), with cocharacter
lattice \(M\). Abel--Jacobi sends a peripheral loop to the corresponding
cocharacter; hence its Tate realization is the geometric abelianization
used in B.2 and B.14.
For arbitrary \(n\), functoriality of the generalized Jacobian gives
the norm map induced by \(q_n\). On the Betti realization it is
the map \(K_n^{\rm ab}\to F^{\rm ab}\). B.5 computes its kernel
and its image exactly. Tensoring (B.2) with the flat ring
\(\widehat{\mathbb Z}\) and using (C.2) proves (C.4).
This is an explicit comparison morphism; at level two its matrix
and its Smith entries \(2,2\) are supplied in B.7.

The geometric module in the middle has rank \(n^2+1\); the proposed
module has rank two. Modulo \(m>1\), their orders are
\(m^{n^2+1}\) and \(m^2\), respectively. Thus even an abstract
isomorphism of the proposed finite realizations fails, with
\(n=m=2\) giving the finite orders \(32\) and \(4\).
The finite order comparison is checked in
\texttt{code/b\_freeness.py}. The norm comparison (C.4) gives the actual relationship that
replaces that proposed identification.
\end{proof}

\begin{statement}{C.3}{Proved}{Complex multiplication with its actual ring map}
For \(n\ge3\), over a field containing \(\mu_n\), the deck action gives a homomorphism
\[
 R_n\longrightarrow\operatorname{End}(J_n)
\]
whose kernel is \(I_n\) of B.6. The induced commutative semisimple
algebra
\[
 \mathcal E_n=\mathbb Q[A_n]/(I_n\otimes\mathbb Q)
 \ \subset\ \operatorname{End}^0(J_n)
\]
has dimension \((n-1)(n-2)=2\dim J_n\). Its factors are cyclotomic
fields indexed by the Galois orbits of characters
\((a,b)\) with \(a,b,a+b\ne0\pmod n\). This is the precise
complex-multiplication algebra supplied by the Fermat automorphisms.
The action on \(\mathcal J_n\) restricts to the permutation action
on the torus lattice in (C.1).
\end{statement}
\begin{proof}
Automorphisms of the curve act on divisor classes and therefore on
the Jacobian. An endomorphism of a complex abelian variety is zero
as soon as it is zero on integral first homology: after lifting to
the complex universal cover it is a complex linear map, and the
integral lattice spans that cover over \(\mathbb R\).
B.6 identifies this homology with \(R_n/I_n\), so the kernel is
exactly \(I_n\). The rational group algebra of a finite abelian
group is a product of number fields, and its character
decomposition selects exactly the characters in B.6. Each has
multiplicity one. The dimension assertion follows by counting them.
Complex conjugation pairs \((a,b)\) with \((-a,-b)\);
none of the selected characters is real: a real character would
have \(a,b\) of order at most two and would force one of
\(a,b,a+b\) to vanish. Thus the selected cyclotomic factors are
CM fields. This proves the full degree required for complex
multiplication.
The line-bundle-with-trivialization description of \(\mathcal J_n\)
shows directly that deck and coordinate automorphisms permute the
cuspidal fibres, giving the torus action claimed. The notation
\(\mathbb Z[\mu_n]^2\) in Lead 3 is therefore replaced by this
explicit group-ring homomorphism and its kernel, with no
interpretation of that ambiguous expression left implicit.
\end{proof}

\begin{statement}{C.4}{Proved}{Exact order of cuspidal differences}
For every \(n\ge3\) and every two distinct geometric coordinate
cusps \(P,Q\) of \(X_n^+\), the class \([P-Q]\in J_n\) has exact
order \(n\). Consequently the subgroup generated by all
cuspidal degree-zero divisor classes has exponent exactly \(n\).
At \(n=1,2\) that subgroup is zero.
\end{statement}
\begin{proof}
Let \(\ell_Q=0\) be the tangent line at a coordinate cusp \(Q\).
For example, at \(Q=[0:q:1]\), with \(q^n=-1\), its equation is
\(y-qz=0\). Substitution into \(x^n+y^n+z^n=0\) gives \(x^n=0\);
the intersection divisor is exactly \(nQ\). The same holds at
every coordinate cusp. Therefore
\[
 \operatorname{div}(\ell_P/\ell_Q)=n(P-Q),                    \tag{C.5}
\]
and the order divides \(n\).

We prove exactness, including the even levels. At \(Q=[0:q:1]\),
the functions \(z/\ell_Q\) and \(x/\ell_Q\) have their only poles
at \(Q\), of orders \(n\) and \(n-1\), respectively. Thus the pole
semigroup \(H(Q)\) contains \(\langle n-1,n\rangle\).
The latter numerical semigroup has exactly
\((n-1)(n-2)/2\) positive gaps: the least element in residue
class \(r\pmod{n-1}\) is \(rn\), for \(0\le r<n-1\), and counting
the smaller positive members of those residue classes gives this
number. Riemann--Roch gives exactly \(g_n\) gaps for \(H(Q)\):
for \(a\ge2g_n\), \(\ell(aQ)=a+1-g_n\), and every successive
dimension increase is zero or one. B.6 gives
\(g_n=(n-1)(n-2)/2\). Since the contained semigroup already has
all \(g_n\) gaps, it equals \(H(Q)\).

A proper positive divisor \(d\) of \(n\) satisfies
\(d\le n/2<n-1\), hence is a gap. A principal divisor
\(d(P-Q)\) would supply a nonconstant function with its only
pole at \(Q\), of order \(d\), a contradiction. The exact order
is therefore \(n\). The genus-zero cases have zero Jacobian.
This gives a proof for every \(n\), without a parity restriction.
\end{proof}

\begin{statement}{C.5}{Proved}{The extension class and its exact isogeny exponent}
The class of (C.1) is the homomorphism
\[
 L_n\longrightarrow J_n^\vee,\qquad D\longmapsto[D],
                                                               \tag{C.6}
\]
under the principal polarization and the extension--character
correspondence
\(\operatorname{Ext}^1(J_n,T_{D_n})
 \simeq\operatorname{Hom}(L_n,J_n^\vee)\).
For every \(n\ge3\), this class has exact order \(n\).
Its pullback by \([n]:J_n\to J_n\) is split, and the resulting
isogeny
\[
 T_{D_n}\times J_n\longrightarrow\mathcal J_n
\]
has kernel isomorphic to \(J_n[n]\) and degree \(n^{2g_n}\).
The splitting of that pullback descends to \(\mathbb Q\).
There is no single positive integer that kills these extension
classes at every Fermat level.
\end{statement}
\begin{proof}
Push (C.1) out by the torus character associated with
\(D=\sum a_PP\). The fibre of the resulting \(\mathbb G_m\)-extension
over a line bundle \(\mathcal L\) is the nonzero part of
\(\bigotimes_P\mathcal L(P)^{\otimes a_P}\).
This is the restriction of the Poincare bundle corresponding
to \([D]\in J_n^\vee\); it proves (C.6).
Conversely, these character pushouts determine a torus extension,
because a split torus is a product of copies of \(\mathbb G_m\).
Descent retains exactly the Galois-equivariant homomorphisms in
(C.6); morphisms from an abelian variety to a torus vanish,
so there is no ambiguity in descending a splitting.
This is also the Cartier-duality description in
\cite[10.2.11--10.2.12]{DeligneHodgeIII}.

C.4 says exactly that (C.6) is killed by \(n\), and that no
proper divisor of \(n\) kills it. Pullback by the multiplication
map \([m]\) multiplies the associated point of \(J_n^\vee\)
by \(m\), as is immediate from the Poincare bundle. Thus
\([n]^*(\text{C.1})\) is the trivial torus extension.
Projection from that pullback to \(\mathcal J_n\) has kernel
\(J_n[n]\), proving the degree and the isogeny statement.
A splitting over \(\overline{\mathbb Q}\) is unique because
\(\operatorname{Hom}(J_n,T_{D_n})=0\), so it is Galois-invariant
and descends. Finally C.4 gives exponents \(n\) for unbounded
\(n\), proving the last assertion.

This statement describes the order in the algebraic extension
group. Its realizations give a splitting of the weight
filtration in the isogeny category, and an integral comparison
whose finite defect is controlled by \(n\). It does not replace
that specified extension group by an unspecified Galois
cohomology group.
\end{proof}

\begin{statement}{C.6}{Known}{Rohrlich's theorem and the requested computations}
For \(n\ge3\), the geometric cuspidal subgroup is
\[
 C_n\simeq
 \begin{cases}
  (\mathbb Z/n)^{3n-7},&n\text{ odd},\\
  (\mathbb Z/n)^{3n-7}\times\mathbb Z/2,&n\text{ even}.
 \end{cases}
\]
This is \cite[Theorem 1, p.~101]{Rohrlich1977},
\emph{Points at infinity on the Fermat curves}. We opened the original
theorem and its definitions on p.~95. The odd case is also reproduced
with an independent proof in \cite[Theorem 1, p.~28-02]{Velu1979}.
For the requested levels the exact orders of the extension class
are
\[
\begin{array}{c|rrrr}
n&3&4&5&7\\ \hline
g_n&1&3&6&15\\
\operatorname{rk}T_{D_n}&8&11&14&20\\
\operatorname{rk}T_\bullet(\mathcal J_n)&10&17&26&50\\
\operatorname{ord}[(\mathrm{C.1})]&3&4&5&7
\end{array}
\]
At \(n=3,5,7\), the computed nontrivial Smith factors of the
checked cuspidal presentation are respectively
\(3^2,\ 5^8,\ 7^{14}\).
\end{statement}
\begin{proof}
Rohrlich defines \(\mathcal D^\infty\) to be the degree-zero divisors
supported on the \(3n\) coordinate cusps of \(x^n+y^n=z^n\), and
\(\mathcal F^\infty\) to be the principal divisors in this lattice.
The geometric isomorphism \(\iota_n\) of B.1 identifies their quotient
with the image of \(L_n\to J_n\). By (C.6), that is precisely the image
of the homomorphism classifying the generalized-Jacobian extension.
The exponent of the displayed group is \(n\), also in the even case
because \(2\mid n\). Therefore multiplication by \(n\) kills the
extension class, and no smaller positive integer does. C.5 then
constructs the splitting isogeny, over \(\mathbb Q\) itself.
This is the reduction from the original theorem to the asserted
weight-extension statement.

For independent integral computations at the requested odd levels,
choose a cusp
\(P_*\), and identify the degree-zero cusp lattice with the
free lattice on \(P-P_*\), for \(P\ne P_*\).
Velu's proof determines its principal-divisor sublattice modulo
\(n\), with six relation classes. In the standard ordering
\(a_i,b_i,c_i\), these can be written
\[
\begin{gathered}
 \sum_i(a_i-P_*),\quad \sum_i(b_i-P_*),\quad
 \sum_i(c_i-P_*),\\
 \sum_i i(a_i-b_i),\quad \sum_i i(a_i-c_i),\quad
 \sum_i i^2(a_i+b_i+c_i-3P_*).
\end{gathered}
\]
The same displayed presentation appears in
\cite[Theorem 29]{BanerjeeMerel2022}, which was also opened.
Together with \(n\) times each basis vector it is the relation
matrix used by \texttt{code/c\_cuspidal\_exponent.py}.
Its exact integer Smith form gives the three listed odd-level
groups, hence their exponents. Formula (C.6) identifies the image
of the extension-class homomorphism with that same cuspidal group,
so the exponent is again the order of the extension class.

For \(n=4\) the program does not apply an odd-level presentation.
It checks instead the independent exact certificate of C.4:
the tangent divisor is \(4(P-Q)\), the pole semigroup
\(\langle3,4\rangle\) has gaps \(1,2,5\), and the proper divisors
\(1,2\) of \(4\) are gaps. Therefore the class has exact order
four. The same certificate is checked for all four requested
levels and, as an additional arithmetic verification, for every
\(3\le n\le32\). The proof of C.4 supplies its validity for all
levels.
\end{proof}

\begin{statement}{C.7}{Proved}{The toric motive actually constructing the tower}
Let \(T=\mathbb G_m^3/\mathbb G_m\), and let
\(\psi:U_1^+\hookrightarrow T\) be the map of B.1. There is a
Cartesian square over \(\mathbb Q\)
\[
\begin{array}{ccc}
 U_n^+&\longrightarrow&T\\
 {\scriptstyle q_n}\downarrow&&\downarrow{\scriptstyle[n]}\\
 U_1^+&\xrightarrow{\ \psi\ }&T.
\end{array}                                                   \tag{C.7}
\]
For a point \(e\in U_1^+(k)\), the 1-motive
\[
 \mathcal K_e=[\mathbb Z\xrightarrow{1\mapsto\psi(e)}T]
\]
has finite realizations
\[
 0\to T[n]\to T_{\mathbb Z/n}(\mathcal K_e)\to\mathbb Z/n\to0,
                                                               \tag{C.8}
\]
whose fibre over \(1\) is \(q_n^{-1}(e)\).
The compatible fibres give exactly the Fermat torsor and its
affine Galois action of B.14. Thus the rank-two module of \S51
is the weight-\(-2\) part of this toric Kummer construction.
\emph{Manuscript: \S51.13--14, folios 719--724, scans 143--153,
(299)--(323), especially (317) and (323).}
\end{statement}
\begin{proof}
A point \([a_0:a_1:a_\infty]\) of \(T\) lies over
\([-e:e-1:1]\) under \([n]\) precisely when its \(n\)-th powers
have that projective ratio. Their sum is zero, which is the
equation of \(U_n^+\). This proves the Cartesian square
scheme-theoretically by its equations on a torus chart.
Deligne's construction 10.1.5 represents
\(T_{\mathbb Z/n}([\mathbb Z\to T])\) by pairs
\((a,t)\), with \(t^n=\psi(e)^a\), modulo
\((a,t)\sim(a+nb,t\psi(e)^b)\). Projection to \(a\bmod n\)
gives (C.8), and the fibre at \(a=1\) is the torsor of
\(n\)-th roots of \(\psi(e)\). By (C.7) it is the stated
Fermat fibre. All formulas commute with automorphisms of \(k\)
and with coordinate permutations, proving the claimed action
and its compatibility with the representation in B.14.

This construction has weights \(0,-2\), with quotient
\(\mathbb Z\) and toric Tate kernel of rank two.
The generalized-Jacobian construction in C.1 is connected to it
by the norm morphism (C.4); that map records the full additional
kernel at level \(n\).
\end{proof}

\begin{statement}{C.8}{Refuted}{Lead 3 and the reading of the manuscript}
Lead 3's proposed identification of \S51.13--14 with a splitting
of the generalized Jacobian of every Fermat curve is false.
The source's explicit construction is (C.7) and the Kummer
fibre (C.8). Its further elliptic-curve construction starts
with an elliptic curve \(E/S\) and the three nonzero
two-torsion points \(D=E[2]\setminus\{0\}\), and has exact sequence
\[
 1\to \mathbb G_m\otimes
       (p_*\mathbb Z_D/\mathbb Z_S)
   \to\mathcal J_{E,D}\to E\to1.
                                                               \tag{C.9}
\]
The Tate kernel of (C.9) has rank two, as written in (303).
It is identified with the torus of the three-point quotient
\(E/\{\pm1\}\) punctured at the images of \(D\).
The weight-extension result for the Fermat Jacobians is C.5,
with exact and unbounded exponents \(n\); the rank-two group
extension of B.12 has zero class and unique splitting conjugacy
class.
\emph{Manuscript: \S51.13, folios 719--720, scans 143--145,
(299)--(304); \S51.14, folios \(720'\)--724, scans 146--153,
(305)--(323).}
\end{statement}
\begin{proof}
The divisor used in scan 144 is explicitly the degree-three
subscheme of nonzero two-torsion points on the elliptic curve,
and (302) has quotient that elliptic curve. Applying the
line-bundle construction of C.1 to \((E,D)\) proves (C.9).
The three-point quotient map \(E\to E/\{\pm1\}\) identifies the
two degree-three etale divisors. Hence it identifies their
cocharacter lattices \(p_*\mathbb Z_D/\mathbb Z_S\) and their
associated tori. On the generalized Jacobians this torus isomorphism
is the pullback of a trivial line bundle with prescribed nonzero
values at the three marked points; the three values are carried
to the corresponding points of \(D\), modulo common scaling.
This is the explicit comparison explaining
the repeated rank-two Tate module. The degree-three divisor
is etale over bases on which \(2\) is invertible, precisely
the qualification made at scan 152.

The other torus comparison written at scan 147 is also explicit:
\[
 [x_0:x_1:x_\infty]\longmapsto
 (x_1/x_\infty,x_\infty/x_0,x_0/x_1)
 \in\ker(\mathbb G_m^3\xrightarrow{\mathrm{product}}\mathbb G_m).
\]
It is an isomorphism: its first two coordinates recover
\(x_1/x_\infty\) and \(x_0/x_\infty\).
For an even permutation it commutes with coordinate permutation;
for an odd permutation its value is the inverse of the permuted
value. This proves the manuscript's marginal qualification that
it is not equivariant for the two untwisted full \(S_3\)-actions.
The \(S_3\)-module used in B.10 is always the quotient lattice,
so no product-one torus action is substituted in the vanishing proof.

At scans 151 and 153 the manuscript supplies \(\psi\) and the
fibre of the universal multiplication tower at the marked
section \(e\). They are exactly the equations proved in C.7.
For the full Fermat generalized Jacobian, the torus has rank
\(3n-1\) and the comparison with this rank-two torus is the
noninjective norm morphism (C.4). Thus the proposed equality
fails by that explicit morphism and its nonzero kernel already
at \(n=2\). The exponents in C.5 further show that a fixed
bounded torsion correction cannot describe all levels.

The supplied table of contents retains the unread initial
letters in \emph{[??]tivisation}; the body heading at folio 708, scan 121,
reads \emph{Retour sur la tour de Fermat}. The mathematics
settles the interpretation just described. It does not settle
the missing letters of the table-of-contents word.
\end{proof}

\section{D. The arithmetic action, its Kummer coordinates, and the two CM factors}
\label{sec:D}

We use $G_k=\operatorname{Gal}(\overline k/k)$, $K=\mathbb Q^{\rm cyc}=\mathbb Q^{\rm ab}$, and the cyclotomic character $\chi:G_{\mathbb Q}\to\widehat{\mathbb Z}^{\times}$. Compatible roots of unity identify $T\mathbb G_m$ with $\widehat{\mathbb Z}(1)$. A coordinate in this Tate module depends on that choice. The group in scans156--170 is the arithmetic $G_{\mathbb Q}$; the larger loop-automorphism groups enter only in the explicitly designated comparisons below. The relevant scans157,166--170 have been inspected as images, independently of their transcription. A positive-weight calculation always refers to the kernel of the multiplier and to its lower-central filtration.

\subsection{The homomorphism of scan157 and the rational-prime questions}

\begin{statement}{D.1}{Proved}{The actual affine Kummer homomorphism}
Choose compatible roots $a^{1/n}$ of $a=-1/2$. The homomorphism of \S51.15, (326), (329)--(335), scans156--160, is, after a choice of origin in its torsor,
\[
 \rho_a:G_{\mathbb Q}\longrightarrow
 \widehat{\mathbb Z}^{\times}\ltimes\widehat{\mathbb Z}(1),
 \qquad g\longmapsto(\chi(g),\kappa_a(g)),\qquad
 g(a^{1/n})/a^{1/n}=\zeta_n^{\kappa_a(g)\bmod n}.
\]
Its image is the closed subgroup of index two
\[
 \{(u,v):(-1)^v=\chi_{-8}(u)\},                         \tag{D.1}
\]
where $\chi_{-8}$ is the quadratic character of $\mathbb Q(\sqrt{-2})$, with values $+1$ at $1,3\pmod8$ and $-1$ at $5,7\pmod8$. Changing the root system translates the origin and changes $\kappa_a$ by $(\chi-1)b$. On $G_K$ its image is exactly $2\widehat{\mathbb Z}(1)$.
\end{statement}
\begin{proof}
The cocycle identity is obtained by applying $g$ to $h(a^{1/n})$:
$\kappa_a(gh)=\kappa_a(g)+\chi(g)\kappa_a(h)$. This is precisely the multiplication in the displayed affine group. In the homogeneous Fermat model of the manuscript, the point $1/2$ has coordinates $(-1/2,-1/2,1)$; its fibre is the diagonal pair of root torsors of $-1/2$. Dividing the two coordinates removes the diagonal ambiguity and gives the rank-one torsor of (329)--(335). Thus this construction supplies the stated morphism, rather than merely a group of the same abstract form. The unread signs in (332), scan159, do not change this argument: the homogeneous equation and the specified point determine the two $n$th powers to be $-1/2$.

The square root $\sqrt{-2}/2$ lies in $K$, and hence the parity of $\kappa_a(g)$ is exactly the quadratic action on $\sqrt{-2}$. On $G_K$ the roots of $-1$ are fixed, so $\kappa_a=-\kappa_2$. D.2 below proves that this has image $2\widehat{\mathbb Z}(1)$. Finally $\chi$ is onto. Every fibre over $u$ is a coset of that translation image; the parity constraint therefore describes the entire fibre, proving (D.1).
\end{proof}

\begin{statement}{D.2}{Proved}{Exact joint images and independence for distinct rational primes}
For distinct positive rational primes $p_1,\ldots,p_r$, the continuous homomorphism
\[
 (\kappa_{p_1},\ldots,\kappa_{p_r}):G_K\longrightarrow
 \widehat{\mathbb Z}(1)^r
\]
has image $2\widehat{\mathbb Z}(1)^r$, of index $2^r$. Consequently every odd-primary projection is onto, the $2$-primary image is $2\mathbb Z_2^r$, and the characters are linearly independent over $\widehat{\mathbb Z}$. On $G_{\mathbb Q}$ the full affine image consists exactly of the pairs $(u,(v_i))$ satisfying
$(-1)^{v_i}=\chi_{\mathbb Q(\sqrt{p_i})}(u)$ for every $i$.
These assertions answer the individual and simultaneous questions of scans168--169; integral surjectivity as written there is false.
\end{statement}
\begin{proof}
Let $a=\prod p_i^{e_i}>0$ and $n\geq1$. We first compute the exact relation
\[
 a\in K^{\times n}\quad\Longleftrightarrow\quad n\mid2e_i\quad\hbox{for every }i.       \tag{D.2}
\]
Choose the positive real root $\alpha=a^{1/n}$. Any other root differs by an element of $K$, so membership in $K$ is independent of this choice. A finite subfield $\mathbb Q(\alpha)$ of the abelian extension $K/\mathbb Q$ is Galois and totally real. Every conjugate of $\alpha$ is a real root of $X^n-a$, of which there are at most two. Thus $\alpha^2\in\mathbb Q$ and valuation at $p_i$ gives $n\mid2e_i$. Conversely this divisibility makes $\alpha$ a product of rational numbers and square roots of rational primes. Such square roots are cyclotomic: $\sqrt2=\zeta_8+\zeta_8^{-1}$, and the quadratic Gauss sum for an odd prime has square $(-1)^{(p-1)/2}p$. Multiplication by $i$ removes the latter sign. This proves (D.2).

For completeness, the finite Kummer pairing used here needs no rank estimate. Let $B$ be the subgroup generated by the classes of the $p_i$ in $K^\times/K^{\times n}$ and let $L=K(p_i^{1/n}:1\leq i\leq r)$. The pairing
$\operatorname{Gal}(L/K)\times B\to\mu_n$, $(g,a)\mapsto g(a^{1/n})/a^{1/n}$,
is nondegenerate on both sides: a Galois element fixing every root is trivial, and a root fixed by every element belongs to $K$. Both groups have exponent dividing $n$; injections into their finite character groups therefore have equal orders and are isomorphisms. By (D.2) the annihilator in $(\mathbb Z/n)^r$ of the relation subgroup is exactly $2(\mathbb Z/n)^r$. These are all finite images. Compactness gives the inverse-limit image stated above. A linear relation vanishes on $2\widehat{\mathbb Z}^r$ only when all coefficients vanish, since multiplication by two on $\widehat{\mathbb Z}$ is injective. Parity and surjectivity of $\chi$ give the last assertion just as in D.1.

The program \texttt{d\_kummer\_soule.py} enumerates these annihilators for every $n\leq24$ and records the indices for $r=1,2,3,4$. This checks, in particular, the finite obstruction to the proposed surjectivity.
\end{proof}

\begin{statement}{D.3}{Proved}{General openness and two corrections to the manuscript}
Let $k$ be a field finitely generated over $\mathbb Q$ and $A\subset k^\times$ a finitely generated subgroup of torsion-free rank $r$. The Kummer homomorphism on $G_{k(\mu_\infty)}$ has open image in $\operatorname{Hom}(A,\widehat{\mathbb Z}(1))\simeq\widehat{\mathbb Z}(1)^r$. This proves the openness requested under (355), scans166--167, also for normal connected schemes essentially of finite type over $\mathbb Q$ and their groups of units. The stronger surjectivity proposed in scan167 is false even for the full unit group of a normal arithmetic scheme of dimension two. The criterion for cyclotomic surjectivity proposed in scan163 is also false.
\end{statement}
\begin{proof}
Torsion in $A$ contributes no homomorphism to the torsion-free Tate module, so choose a free complement. Write $H=\operatorname{Gal}(k(\mu_\infty)/k)$, an open subgroup of $\widehat{\mathbb Z}^{\times}$: the algebraic closure of $\mathbb Q$ inside $k$ is a number field. Choose an integer $m$ such that $H$ contains all units congruent to one modulo $m$. There is $u\in H$ with $u-1=dv$, where $d=\operatorname{lcm}(2,m)$ and $v\in\widehat{\mathbb Z}^{\times}$: take $u_p=1+p^{v_p(m)}$ at $p\mid m$, $u_p=2$ at odd $p\nmid m$, and $u_2=3$ when $2\nmid m$.

Suppose $\alpha^n=a\in A$ with $\alpha\in k(\mu_\infty)$. Put $\tau(\alpha)/\alpha=\zeta_n^{c_\tau}$. Since $H$ is abelian, comparison of its two cocycle products gives
$(u-1)c_\tau=(\chi(\tau)-1)c_\sigma$ for an element $\sigma$ with character $u$. Setting $b=v^{-1}c_\sigma$ modulo $n$ shows that $\alpha^d\zeta_n^{-b}$ is fixed by $H$. Thus
\[
 a^d\in k^{\times n}.                                      \tag{D.3}
\]
Choose a normal finitely generated $\mathbb Z$-algebra $R$ with fraction field $k$ containing $A$ as units. Its unit group is finitely generated. One proof is to take a normal projective compactification over $\operatorname{Spec}\mathbb Z$: a unit has divisor supported on the finitely many codimension-one boundary components, while the kernel of this divisor map is the unit group of the ring of global functions of the compactification. That ring is finite over $\mathbb Z$; its units are finitely generated by the number-field unit theorem. Normalization is finite here, so this compactification argument applies to $R$.

If $b^n=a^d$ in $k$, both $b$ and $b^{-1}$ are integral over $R$, hence belong to $R$. In the free lattice $U=R^\times/\mu(R)$, the saturation of $A$ modulo $A$ is finite; let $e$ annihilate it. Equation (D.3) then implies $de\,a\in nA$ in additive lattice notation. The kernel of $A/nA\to k(\mu_\infty)^\times/k(\mu_\infty)^{\times n}$ has order at most $(de)^r$. The perfect finite pairing in D.2 therefore bounds all finite indices of the Kummer image by this same number. A closed subgroup of $\widehat{\mathbb Z}^r$ with uniformly bounded finite indices has finite index: any collection of more distinct cosets is separated in one finite quotient. This proves openness.

For a normal connected scheme, the Kummer covers of units are finite etale and the generic-point map is onto their monodromy groups; the preceding field computation consequently gives the same image. For the counterexample take $S_1=\operatorname{Spec}\mathbb Z[1/2][t]$ and $A=\mathcal O(S_1)^\times=\{\pm2^j\}$. Its generic fibre has the Kummer image $2\widehat{\mathbb Z}$ of D.2. Moreover $A$ is saturated in $\mathbb Q[t]^\times$ modulo its torsion, so the proposed saturated variant fails too. Finally $k=\mathbb Q(\sqrt5)$ has only the roots of unity $\pm1$, but its cyclotomic image has index two because $k\subset\mathbb Q(\mu_5)$.
\end{proof}

\begin{statement}{D.4}{Proved}{Exact local inertia images}
Embed $G_{\mathbb Q_\ell}$ into $G_{\mathbb Q}$ and put $I_\ell^+=I_\ell\cap\ker\chi$. For a positive rational prime $p$, the image of $\kappa_{p,\ell}$ on $I_\ell^+$ is $\ell^{e_{p,\ell}}\mathbb Z_\ell(1)$, where
\[
 e_{p,\ell}=\begin{cases}
 0&\ell>2, p=\ell,\\
 v_\ell(p^{\ell-1}-1)-1&\ell>2, p\ne\ell,\\
 1&\ell=p=2,\\
 v_2(p^2-1)-2&\ell=2, p\text{ odd}.
 \end{cases}                                               \tag{D.4}
\]
Every image in scan170 is therefore open, and the formula determines exactly which are onto. The program \texttt{d\_local\_cm.py} records every pair $p,\ell\leq50$; for example the image for $(p,\ell)=(3,11)$ has index $11$, and for $(19,7)$ index $49$.
\end{statement}
\begin{proof}
The fixed field of $I_\ell^+$ is $L=\mathbb Q_\ell^{\rm ur}(\mu_{\ell^\infty})$, abelian over $\mathbb Q_\ell$. For odd $\ell$, choose a cyclotomic multiplier with $u-1\in\mathbb Z_\ell^\times$. The cocycle argument in D.3, now with $d=1$, proves that an $\ell^k$th root of $a\in\mathbb Q_\ell^\times$ belongs to $L$ precisely when $a\in\mathbb Q_\ell^{\times\ell^k}$. The decomposition
$\mathbb Q_\ell^\times=\ell^{\mathbb Z}\times\mu_{\ell-1}\times(1+\ell\mathbb Z_\ell)$ and the $\ell$-adic logarithm give the first two cases.

For $\ell=2$ use multiplier $-1$. The same argument proves that $a$ having a $2^k$th root in $L$ implies $a^2\in\mathbb Q_2^{\times2^k}$, equivalently $a=\pm b^{2^{k-1}}$ for $b\in\mathbb Q_2^\times$. Conversely every quadratic extension of $\mathbb Q_2$ lies in $L$: the square classes are generated by $-1,2,5$, the first two have cyclotomic square roots and $\mathbb Q_2(\sqrt5)$ is unramified. Thus $\pm b^{2^{k-1}}$ indeed has a $2^k$th root in $L$. For $a=2$ the largest such $k$ is one. For odd $p$, write $p=\epsilon u$, $\epsilon\in\{\pm1\}$, $u\in1+4\mathbb Z_2$; logarithms give the largest $k$ as $v_2(u-1)-1=v_2(p^2-1)-2$. Finite Kummer duality converts these root heights into the displayed closed images.
\end{proof}

\subsection{The universal Fermat cocycle and its actual comparison maps}

\begin{statement}{D.5}{Known}{Full profinite universal measure, and the precise correction to Lead 1}
Put $\Lambda=\widehat{\mathbb Z}[[\widehat{\mathbb Z}^2]]$. The compatible projective Fermat homology of Otsubo's Theorem6.1 is a free cyclic $\Lambda$-module. For its specified generator $c$, write $g(c)=F_g c$; then
$F_{gh}=g(F_h)F_g$. On $G_K$ this is a homomorphism into $\Lambda^\times$. The comparison with Anderson's relative-path measure $B_g$ is the identity
\[
 (1-s^r)(1-t^m)F_g=(1-s^{\chi(g)r})(1-t^{\chi(g)m})B_g.       \tag{D.5}
\]
Here $r,m$ are integers and $s,t$ are the two deck generators. This is Otsubo's (7.2). In particular $F_g=B_g$ on $G_K$. Its $\ell$-power specialization is related to Ihara's power series by the explicit derived-module basis comparison below.

The morphisms relating this invariant to D.1 are the exact quotient
\[
 1\longrightarrow\widehat\pi'/\widehat\pi''
 \longrightarrow\widehat\pi/\widehat\pi''
 \longrightarrow\widehat\pi^{\rm ab}\longrightarrow1
                                                               \tag{D.6}
\]
and the root torsors obtained by evaluating the abelian covering at $1/2$. For rational primes the additional exact comparison is
\[
 p=\prod_{a=1}^{p-1}(1-\zeta_p^a),\qquad
 \kappa_p\big|_{G_K}=\sum_{a=1}^{p-1}\kappa_{1-\zeta_p^a}\big|_{G_K}.       \tag{D.7}
\]
Lead1's proposed Kummer comparison is made exact by (D.5)--(D.7) and the coefficient extraction in D.23 below. In particular the rank-one affine coordinate is obtained from the full universal measure by a specified continuous map. This comparison retains the factor two in D.2 and the prime-to-$\ell$ deck components that a pure pro-$\ell$ series discards.
\end{statement}
\begin{proof}
We use \cite[Theorem6.1, Corollary6.2, equation(7.2)]{Otsubo2016}, whose full text was opened. The non-zero-divisor assertion in Corollary6.2 permits cancellation in (D.5); when $\chi=1$ the two measures coincide. The maps to $\ell$-power deck quotients and $\mathbb Z_\ell$ coefficients are continuous quotient maps of completed group rings; hence this comparison is profinite before any pro-$\ell$ specialization. Part B constructs the derived-module comparison in (D.6); evaluation of its abelian quotient at the displayed point gives D.1 by the explicit root equations. Finally (D.7) is $\Phi_p(1)=p$ and functoriality of roots under multiplication. Applying valuations at the distinct primes to its rational norm products, followed by the exact pairing of D.2, proves independence through these cyclotomic elements. The possible quadratic relations, including the factor two, are retained; an assertion of unrestricted integral independence of all circular units would not give the required exact image.

Otsubo's Remark8.2 identifies the original results precisely as Ihara1986, Theorem7 for prime-power level, and Anderson1989, Theorem6 in general. Anderson's complete original has now been retrieved from the official GDZ scan and checked. Its \S\S1.3--1.4 use the full ring
$\widehat{\mathbb Z}[[\widehat{\mathbb Z}(1)^2]]$;
Proposition5.5.1, p.88, constructs the all-degree torsion comparison from relative interval homology and states
$e_N(B_\sigma\,\sigma(v))=\sigma(e_N(v))$.
Theorem6, pp.118--124, then interpolates Gauss sums using a specified Frobenius lift and the hyperadelic gamma function; it is not by itself the torsion comparison.

The original comparison to Ihara is \cite[Proposition13.3.2 and \S\S13.4--13.6]{dAnderson1989}.
For odd $\ell$ it constructs a Galois-equivariant map
$\theta:\pi_\ell'/\pi_\ell''\simeq
\varprojlim_n T_\ell J_{\ell^n}$ and a basis-change unit
$Q(u,v)$, with exact formula
\[
\begin{aligned}
 F^{\rm Ih}_\sigma(u,v)&=S_\sigma(u)S_\sigma(v)
 B_{\sigma,\ell}(u,v)
 \frac{Q((1+u)^{\chi_\ell(\sigma)}-1,
          (1+v)^{\chi_\ell(\sigma)}-1)}{Q(u,v)},\\
 S_\sigma(T)&=\frac{(1+T)^{\chi_\ell(\sigma)}-1}{T}.
\end{aligned}
\]
The map $\theta$ sends a deck transformation to its translation
on the compatible division points of the divisor
$(\delta_n)-(\infty_n)$, where
$\delta_n=(t_n,-t_n(1-t^{-1})_n)$ in \S13.5.
That section records the plus-model twist for $\ell=2$;
RemarkII in \S13.3 proves $F^{\rm Ih}_\sigma=B_{\sigma,\ell}$
on the cyclotomic kernel for every $\ell$.
These maps specify both the basis coboundary and the arithmetic twist.

The original Ihara1986 and Coleman1989 proofs remain to be obtained; the separate coefficient formula used below was checked in Ihara--Kaneko--Yukinari's primary paper and now also in Anderson's original. The retrieved original \cite[Supplement]{dIhara2000} corrects Ihara1986, Theorem7 by replacing the Jacobi indices $(a,b,c)$ with $(-a,-b,-c)$, and corrects the signs in Theorem10. It also instructs replacing Anderson1989, p.125 line11's group-ring identification by its inverse, and records the inverse deck convention on p.85. These corrections explain why a convention must be attached to the original theorem numbers; D.6 uses the explicitly checked Otsubo convention.
\end{proof}

\begin{statement}{D.6}{Known}{Frobenius specialization with nonsplit primes and the plus model}
Let $p\nmid n$, $f=\operatorname{ord}_n(p)$, $q=p^f$, and let $\theta:\mathbb F_q^\times\to\mu_n$ be fixed. For $a,b,a+b\ne0$ modulo $n$, put
\[
 j_n^{a,b}=-\sum_{t\in\mathbb F_q}\theta(t)^a\theta(1-t)^b,
\]
with nontrivial characters extended by zero. Otsubo's Theorem8.1 identifies this number with the character specialization of the product of $F_{\operatorname{Frob}_p}$ along the orbit
$(a,b),(pa,pb),\ldots,(p^{f-1}a,p^{f-1}b)$. For $x^n+y^n=z^n$ the sum of these $j_n^{a,b}$ is the trace of geometric Frobenius on $H^1$ over $\mathbb F_q$. For the manuscript's $x_0^n+x_1^n+x_\infty^n=0$ it is instead the sum of
\[
 j_{n,+}^{a,b}=\theta(-1)^{a+b}j_n^{a,b}.                    \tag{D.8}
\]
\end{statement}
\begin{proof}
The cited theorem \cite[Theorem8.1]{Otsubo2016} is formulated for an arbitrary residue degree; its product, rather than a single value at $p$ on a non-fixed character, is essential. The comparison with $H^1$ uses the ordinary arithmetic Frobenius action on the Tate module and its dual geometric action on cohomology; these have the same characteristic polynomial in this convention. The change of models is defined over $\mathbb Q(\epsilon_n)$ by
$(x_0:x_1:x_\infty)\mapsto(x_0:x_1:\epsilon_n x_\infty)$, $\epsilon_n^n=-1$. Its Frobenius discrepancy is $\epsilon_n^{q-1}=(-1)^{(q-1)/n}$. Projectively this multiplies the first two coordinates by the same element. The $(a,b)$ character is therefore multiplied by $\theta(-1)^{a+b}$, which proves (D.8). This retains the arithmetic descent of scan171, including even $n$.
\end{proof}

\begin{statement}{D.7}{Proved}{Complete finite-field computation in the requested range}
The files \texttt{d\_jacobi.py}, \texttt{d\_jacobi\_field.cpp}, \texttt{d\_jacobi.out} and \texttt{d\_jacobi.json} compute every nondegenerate character value for every prime $p\leq50$ and $3\leq n\leq12$ with $p\nmid n$, over the exact field $\mathbb F_{p^{\operatorname{ord}_n(p)}}$. Every value satisfies $j\overline j=q$. Character sums and independent point counts agree for both models. Levels1 and2 have zero projective $H^1$, and the output separately retains the bad-reduction pairs $p\mid n$.
\end{statement}
\begin{proof}
For non-semiprimitive pairs the code finds a monic polynomial for which multiplication by its residue class has order $q-1$. This certifies that all nonzero elements of the quotient ring form a cyclic group, hence that the quotient is the required finite field. It enumerates the cyclotomic-number array
$C_{ij}=\#\{t\ne0,1:\theta(t)=\zeta_n^i,\theta(1-t)=\zeta_n^j\}$.
Discrete Fourier summation gives each $j$ as an integral polynomial modulo $\Phi_n$. All norm identities are checked in this exact quotient ring. Independently, the standard curve has
$n^2C_{00}+2n+n\,[\theta(-1)=1]$ points, and the plus curve has
$n^2C_{ee}+3n\,[e=0]$, where $\theta(-1)=\zeta_n^e$.

For the larger semiprimitive fields, $p^h\equiv-1\pmod n$, $q=p^{2h}$, the code uses and checks the exact identity $j=-p^h$. Indeed the characters are trivial on $\mathbb F_{p^h}^\times$. In the quadratic extension, summing the additive character over trace fibres gives the Gauss sum $G(\theta^a)=p^h\theta^a(\tau)$ for a nonzero trace-zero element $\tau$; the quotient $G(\theta^a)G(\theta^b)/G(\theta^{a+b})$ is $p^h$. This proves the shortcut. The smaller semiprimitive cases are also enumerated directly. The largest direct field is $\mathbb F_{47^5}$; it was actually enumerated. For an explicit arithmetic-twist check, $(p,n)=(43,6)$ gives standard trace $32$ and plus-model trace $-64$, with respective point counts $12$ and $108$. These computations verify the finite-field side of D.6 and its character conventions; they do not claim a numerical oracle evaluating an arbitrary abstract Galois element independently of the theorem.
\end{proof}

\subsection{Circular Kummer characters, regularity, and the first irregular prime}

\begin{statement}{D.8}{Known}{The checked Ihara--Anderson--Coleman coefficient formula}
For odd $\ell$, let $U=\log(1+u)$, $V=\log(1+v)$, $W=-U-V$. On $G_{\mathbb Q(\mu_{\ell^\infty})}$ the universal series has the expansion
\[
 F_g(u,v)=\exp\left(\sum_{m\geq3\ {\rm odd}}
 \frac{\kappa_m(g)}{(1-\ell^{m-1})m!}
 (U^m+V^m+W^m)\right),                                  \tag{D.9}
\]
where $\kappa_m$ is the compatible Kummer character of
$\epsilon_n^{(m)}=\prod_{a\in(\mathbb Z/\ell^n)^\times}(\zeta_{\ell^n}^a-1)^{\langle a^{m-1}\rangle_n}$.
The integral factor $1-\ell^{m-1}$ is a unit. Vandiver at $\ell$ is equivalent to surjectivity of all these $\kappa_m$ for $m=3,5,\ldots,\ell-2$.
\end{statement}
\begin{proof}
This is the reduction of \cite[Theorems A$_2$ and B, equations(9)--(11)]{dIKY1987} to the root convention of D.1. Their character is defined by the same quotient of a chosen $\ell^n$th root and its Galois translate, so no unnamed rescaling is involved. Their Theorem B identifies $\beta_m=(1-\ell^{m-1})^{-1}\kappa_m$, and Theorem A$_2$ gives the exponential with $\beta_m/m!$. The paragraph immediately following Theorem B proves the stated Vandiver equivalence. Both pages were inspected visually. Reducing modulo $\ell$ permits replacing $\zeta^a-1$ by $1-\zeta^a$: the resulting sign is an $\ell$th power. These are the exact characters computed next.

Anderson's original \S11.6, Proposition11.6.1 and formula(13.7.3),
pp.117--118 and130, give an independent normalization check.
His $\beta_m$ sums $i^{m-1}$ times the Kummer cocycle of
$1-\zeta_{\ell^n}^{-i}$ over every $0<i<\ell^n$.
On the cyclotomic kernel this is the cocycle of
$\zeta_{\ell^n}^{i}-1$, since their ratio is
$\zeta_{\ell^n}^{-i}$ and all of its $\ell$-power roots are fixed.
Splitting the sum into $\ell\nmid i$ and $i=\ell j$ gives
$\beta_m=\kappa_m+\ell^{m-1}\beta_m$.
Thus his exponential coefficient is exactly the one in (D.9).
\end{proof}

\begin{statement}{D.9}{Known}{An exact obstruction module, not an irregular-prime guess}
Let $\mathcal F_\ell\subset\mathbb Z_\ell[[u,v]]^\times$ be the group defined by the five conditions of Ichimura--Kaneko1989, \S1.2: congruence $F\equiv1\pmod{uvw}$, inversion $F\overline F=1$, symmetry in $u,v,w$, symmetry of $F*F$ in four variables with product of $1+$variables equal to one, and
\[
 \prod_{\zeta^\ell=1}F(\zeta(1+u)-1,\zeta(1+v)-1)
       =F((1+u)^\ell-1,(1+v)^\ell-1).
\]
Their Theorem1 identifies equality $\operatorname{Im}F=\mathcal F_\ell$ with Vandiver at $\ell$. Their Theorem2 computes the characteristic power series of the defect $(\mathcal F_\ell/\operatorname{Im}F)(-1)$ as that of
$\operatorname{Hom}_{\mathbb Z_\ell}(A_\infty^+,\mathbb Q_\ell/\mathbb Z_\ell)$ over the cyclotomic Iwasawa algebra. Thus the relevant obstruction is the plus class-group module, not merely divisibility of a Bernoulli numerator. At every regular odd prime this defect is zero. D.10 proves the same equality at every odd prime at most50, including37.
\end{statement}
\begin{proof}
We use \cite[Theorems1--2]{dIchimuraKaneko1989}; the defining conditions and both theorem statements were checked on the original pages314--315. Equation(D.9) identifies their $F$ and Kummer characters with ours. The plus part $A_\infty^+$ is the direct limit of the $\ell$-Sylow parts of the ideal class groups of $\mathbb Q(\mu_{\ell^n})$, with the authors' transition maps. The stated characteristic-series equality records exactly the size invariant proved there; it is not promoted to an isomorphism of modules. Regularity implies that the full class number at the first level is prime to $\ell$, hence its plus part is, and so Vandiver holds. For the bounded range, the certificates below establish the equivalent surjectivities directly.
\end{proof}

\begin{statement}{D.10}{Proved}{Soulé characters modulo every prime at most50}
For every odd prime $\ell\leq50$ and each odd $3\leq m\leq\ell-2$, the program \texttt{d\_kummer\_soule.py} exhibits an explicit nonzero value of $\kappa_m\bmod\ell$. Consequently all these individual characters are onto $\mathbb Z_\ell$, their distinct Teichmuller-weight components are independent, and Vandiver holds throughout this range. At $\ell=37$ all seventeen characters are certified by the single prime $q=5477$. The Bernoulli computation finds the irregular pair $(37,32)$ and no other irregular pair below50; it causes no missing plus-part image here.
\end{statement}
\begin{proof}
For a prime $q\equiv1\pmod{\ell^2}$, a primitive root $g\bmod q$, and $r=g^{(q-1)/\ell}$, the code evaluates
$c_m(r)=\prod_{a=1}^{\ell-1}(1-r^a)^{a^{m-1}\bmod\ell}$ and finds the unique $k\bmod\ell$ such that
$c_m(r)^{(q-1)/\ell}=r^k$. A nonzero $k$ proves that the corresponding cyclotomic element is not an $\ell$th power. All choices, exponents and residues are retained in the JSON, so each witness is checked by modular exponentiation alone.

This is also a certificate on the infinite cyclotomic field, not only at its first layer. In $\mathbb Q(\zeta_\ell)^\times/\ell$ the class of $c_m$ has Teichmuller character $b\mapsto b^{1-m}$. The only order-$\ell$ subextension of the cyclotomic $\mathbb Z_\ell$-extension has Kummer class $\zeta_\ell$, of character $b\mapsto b$. These characters differ in the indicated range. Thus the nonzero class cannot disappear in that extension. The character takes a unit value in $\mathbb Z_\ell$, proving surjectivity; projection by the distinct idempotents of $\mathbb F_\ell^\times$ proves joint independence of these weight components. D.8 now proves Vandiver without inferring it from the irregularity test.

The witness primes, grouped by $\ell$, are
\[
\begin{array}{c|rrrrrrr}
\ell&5&7&11&13&17&19&23\\\hline
q&101&197,491&727&677,2029&3469&10831&12697
\end{array}
\]
\[
\begin{array}{c|rrrrrr}
\ell&29&31&37&41&43&47\\\hline
q&10093,13457&15377&5477&16811&22189&8837.
\end{array}
\]
For $\ell=3$ the specified weight list is empty; for $\ell=2$ there is likewise no weight in the requested range. No omitted character is hidden by these conventions.
\end{proof}

\subsection{Metabelian leading terms and the lifting issue}

\begin{statement}{D.11}{Proved}{Exact special-derivation images through weight10}
Let $L$ be the free Lie algebra over $\mathbb Q$ on $x,y$, with the three peripheral elements $x,y,-x-y$. Modulo inner derivations normalize a positive-degree special derivation by $D(x)=0$, $D(y)=[f,y]$. For weight $n\geq2$, it is special exactly when some $g\in L_n$ satisfies
\[
 [f,y]=[g,x+y].                                           \tag{D.10}
\]
The metabelian quotient is
$L/L''=\mathbb Qx\oplus\mathbb Qy\oplus\mathbb Q[X,Y]c$, where $c=[x,y]$, $X=\operatorname{ad}x$, $Y=\operatorname{ad}y$; the last summand is abelian. The following exact rational ranks were computed, with source witnesses, by \texttt{d\_metabelian.py}:
\[
\begin{array}{c|rrrrrrrrrr}
n&1&2&3&4&5&6&7&8&9&10\\\hline
\text{special outer source}&0&0&1&0&3&0&6&4&13&12\\
\text{its metabelian image}&0&0&1&0&3&0&5&0&7&0\\
S_3\text{-invariant source}&0&0&1&0&1&0&2&1&4&2\\
\text{its metabelian image}&0&0&1&0&1&0&1&0&2&0
\end{array}                                               \tag{D.11}
\]
Thus all three weight-five directions survive. The extra invariant direction in weight seven dies. In weight nine two invariant directions survive. These are images of derivations of the free Lie algebra; special derivations defined only on $L/L''$ have dimension $n-2$ for every $n\geq3$, and are a larger calculation.
\end{statement}
\begin{proof}
Subtracting the inner derivation that conjugates $x$ gives the normalization. For $n\geq2$, the centralizer of $x$ has no component of degree $n$, so it is unique. The condition at the third peripheral element gives (D.10), with a harmless simultaneous sign determined by the convention for $g$. The free metabelian description follows by Jacobi: the commuting operators $X,Y$ generate $L'/L''$ from $c$, and its monomials are independent, for instance in the usual upper-triangular matrix realization over $\mathbb Q[X,Y]$.

The program uses the Lyndon basis in each degree and expands its brackets in the free associative algebra. It forms the exact matrix of $(f,g)\mapsto[f,y]-[g,x+y]$ and computes its rational nullspace. The invariant subspace is separately defined by the checked LM2 equations
\[
 f(x,y)+f(y,x)=0,\qquad
 [x,f(x,y)]-[x+y,f(-x-y,y)]=0.                             \tag{D.12}
\]
These follow by conjugating the normalized outer derivation by the transposition and the three-cycle of the three peripheral elements and then renormalizing $D(x)=0$. Thus the matrix computes the actual infinitesimal centralizer, rather than the two linear cycle equations. Projection of each Lyndon bracket to $\mathbb Q[X,Y]c$ gives the image ranks. Every pivot-image vector retains an exact preimage in the JSON.

In the metabelian algebra (D.10) reads $Yp=(X+Y)q$, hence $(X+Y)\mid p$. This gives the stated $n-2$ comparison dimension. At weight five the full image consists of $(X+Y)$ times all degree-two polynomials; the invariant image is spanned by
$X^3+2X^2Y+2XY^2+Y^3$.
At weight seven its invariant image is spanned by
$X^5+3X^4Y+5X^3Y^2+5X^2Y^3+3XY^4+Y^5$.
At weight nine one extra independent vector is $X^2Y^2(X+Y)^3$; the output gives the other vector and both preimages. The zero even-weight images, including weights eight and ten with nonzero sources, are computed by this projection, not assumed from parity.
\end{proof}

\begin{statement}{D.12}{Proved}{The reviewed pro-$\ell$ lifting assertion and its precise scope}
For a free pro-$\ell$ group on finitely many peripheral generators, every integral special derivation is the leading term of a special automorphism. For type $(0,3)$ and $\ell\geq5$, every $S_3$-invariant integral leading term lifts to an element of the centralizer of the geometric $S_3$ in the pro-$\ell$ outer special-automorphism group. Thus the image ranks in (D.11) are ranks of actual pro-$\ell$ special-automorphism images, and of their centralizer for $\ell\geq5$.
\end{statement}
\begin{proof}
Write $\Gamma_m$ for the closed lower-central series. Its associated graded is the free Lie algebra over $\mathbb Z_\ell$. Start with conjugators representing the given degree-$n$ Lie elements. The product-of-peripheral-generators equation holds modulo $\Gamma_{n+2}$ because its leading error is the special-derivation equation. At stage $m$ the error lies in $\Gamma_m/\Gamma_{m+1}$; changing the conjugators by elements of $\Gamma_{m-1}$ changes it by
$\sum_i[a_i,x_i]$. This map is onto $L_m$: iterated Jacobi rewrites every Lie monomial of degree $m$ as a sum $[x_i,b_i]$. In the two-generator normalization fixing $x$, the two available operators are $[\ ,y]$ and $[\ ,x+y]$; their sum still spans $[x,L_{m-1}]+[y,L_{m-1}]=L_m$. We can therefore cancel each error without altering the prescribed leading term.

The corrections converge in the free pro-$\ell$ group. Indeed every finite $\ell$-group quotient is nilpotent, so $\Gamma_m$ eventually lies in every open normal subgroup. The resulting endomorphism is the identity on the Frattini quotient and hence is onto by the pro-$\ell$ basis theorem; a finitely generated profinite group is Hopfian, so it is an automorphism. This independently proves the reviewed assertion.

For the invariant assertion, the group of IA automorphisms in this pro-$\ell$ setting is pro-$\ell$, as are the successive kernels. A finite group of order prime to $\ell$ acts on it. For a quotient by a stable closed normal subgroup, fixed elements lift to fixed elements: for finite $\ell$-groups this follows by induction on a central series and averaging the additive obstruction, and the profinite assertion follows by compactness. Apply this to the order-six group $S_3$ and successive filtration quotients. This proves the claimed lift. No such convergence argument applies to the full free profinite group: a finite nonnilpotent quotient need not kill any $\Gamma_m$.
\end{proof}

\begin{statement}{D.13}{Open}{The remaining full profinite image comparison in D(e)}
The maps from actual full profinite loop-preserving outer automorphisms, and from $\widehat{\rm GT}_0$, to their pro-$\ell$ leading terms land in the second and fourth rows of (D.11), respectively. D.12 proves the reverse inclusions for the pro-$\ell$ groups, not for those full profinite source groups. The full profinite realization of the extra weight-five metabelian directions, and of the extra invariant weight-nine direction, remains unproved here. Consequently D(e), as asked for the three full profinite groups, remains undecided in this return.
\end{statement}
\begin{proof}[Established comparisons and attempted lifting]
Taking the maximal pro-$\ell$ quotient is characteristic, so it defines the displayed continuous homomorphisms of outer automorphism groups. Taking a leading term is injective on each filtration quotient by its definition; the peripheral and centralizer equations give (D.10)--(D.12). D.12 realizes every admissible leading vector after the change of source to the pro-$\ell$ group. A proposed full profinite lift would have to solve the same conjugacy and generation equations in every finite quotient, including nonnilpotent quotients. The lower-central correction sequence controls none of the residual errors in such a quotient. We also checked the weaker route of solving just the metabelian Lie equations: its larger dimension $n-2$ proves that this route discards further lifting conditions already at weight four. These two failed transfers are exact gaps, not counterexamples to a full profinite lifting theorem.

For the Galois source, (D.9) gives at most one rational direction in each odd weight and zero in even weights in this metabelian image. D.10 supplies nonzero directions through weight ten at $\ell=37$; E.3 supplies the rational general-prime comparison. The weight-five pro-$\ell$ image therefore distinguishes the whole special group from its $S_3$ centralizer and from Galois. D.20--D.21 below add the full profinite beta/gamma compatibility and prove equality of the rational metabelian graded images of Galois and full $\widehat{\rm GT}$ through weight ten. The pentagon used there is an additional equation on $\widehat{\rm GT}_0$. The exact missing lift test for the latter is D.22; the pro-$\ell$ construction alone supplies only the subfamily of its tests with nilpotent quotients. D.24 gives an explicit test of the general automorphism-lifting route and its peripheral limitation. D.25 constructs nonarithmetical finite GT shadows with generation verified; it leaves their genuineness, and hence the indicated full profinite comparison, undecided.
\end{proof}

\subsection{The integral signs, CM coordinates, and the genuine congruence action}

\begin{statement}{D.14}{Proved}{What the first integral sign congruences say}
The norm equations give $\epsilon_2=\chi_4(\mu)$ and $\epsilon_3=\chi_{-3}(\mu)$, already at moduli4 and3. On the Galois image $\mu=\chi$. Thus these signs add no independent condition to the full universal datum $(\chi,F_g)$: the exact comparison is evaluation of $\chi$ modulo4 and3. The unadorned series restricted to $G_K$ has $\chi=1$, and consequently both signs there are one.
\end{statement}
\begin{proof}
For a unit Gaussian norm $a^2+b^2$, exactly one of $a,b$ is odd; its value is one modulo4. For the Eisenstein norm $c^2+cd+d^2$, reduction modulo3 is $(c-d)^2$, hence is one for a unit. The signed norm equalities in LM2 C.d2 therefore force the stated signs. Conversely the local norm images are $1+4\mathbb Z_2$ and $1+3\mathbb Z_3$: the norm on sufficiently deep principal units is the trace under logarithm, and inspection at moduli8 and9 supplies the first step. These finite images are also enumerated inside \texttt{d\_local\_cm.py}. Evaluating the already specified multiplier proves the asserted factorization; it does not equate an unmarked beta series with its missing multiplier.
\end{proof}

\begin{statement}{D.15}{Known}{The two normalizer coordinates are the projective CM representations}
Use the rigidified modular stack $\mathcal M_{1,1}/\{\pm1\}\simeq[U_{0,3}/S_3]$ and the positive tangential point $\lambda=0$ with tangent1 and the standard branch-cut paths on $0<\lambda<1$. Choose the same etale paths and homology basis in the modular congruence quotient and in LM2 C.d4. Then its coordinates
$T=A^{-1}V$ and $W=\beta^{-1}V$, each modulo its single global sign, are respectively the projective adelic Tate representations of
\[
 E_3:y^2=x^3+1,\qquad E_4:y^2=x^3-x.                       \tag{D.13}
\]
Their determinants are $\chi$; their actions on the CM orders are complex conjugation according to $\chi_{-3}$ and $\chi_4$. Changing the paths conjugates the indicated representations by their transported bases. Quadratic twists of the displayed curves make exactly the same representations modulo the global sign.
\end{statement}
\begin{proof}
The quotient of an elliptic curve by $[-1]$ is a four-pointed projective line, with the identity among the four branch points. Ordering the other three points gives the Legendre parameter and is a torsor under $S_3$. Conversely the double cover of this four-pointed line gives the elliptic curve; its two possible covering automorphisms are exactly $\pm1$. These two constructions are inverse on families and their isomorphisms after rigidification, proving the displayed stack equivalence over $\mathbb Q$.

The projective Tate local system consequently has geometric monodromy $\operatorname{SL}_2(\widehat{\mathbb Z})/\{\pm I\}$. At the elliptic points its inertia generators are the matrices $R,S$ of LM2, satisfying $R^3=S^2=-I$. Transport the Galois action at either elliptic point to the cusp by a fixed path. The resulting path cocycle is a conjugator $A$, respectively $\beta$, with
$u_0(R)=AR^{\chi_{-3}}A^{-1}$ and
$u_0(S)=\epsilon_2\beta S\beta^{-1}$. These are precisely the conjugators in C.d4: their prescribed reductions modulo2 remove the finite inertia ambiguity. More explicitly, the centralizer of either torsion generator in the free profinite product $C_2\amalg C_3$ is its cyclic factor, whose intersection with the level-two free subgroup is trivial. Hence two such normalized path conjugators coincide. Transport of the Tate action now gives exactly $A^{-1}V$ and $\beta^{-1}V$, proving the formulas.

At $\lambda=1/2$, the change $X=4x-2$, $Y=8y$ gives $Y^2=X^3-4X$, the quadratic twist by2 of $E_4$. At $\lambda^2-\lambda+1=0$, shifting $x$ by $(1+\lambda)/3$ gives $y^2=x^3+(2\lambda-1)/9$; this constant is the cube of $-(2\lambda-1)/3$, so it is a quadratic twist of $E_3$ over $\mathbb Q(\sqrt{-3})$. The descent to $\mathbb Q$ is fixed by its action on the three nonzero two-torsion points: for $j=0$ the splitting field $\mathbb Q(\sqrt{-3})$ forces the cubic constant to be a rational cube up to a quadratic twist, and for $j=1728$ rational two-torsion forces the coefficient to be a negative square. Thus the descended projective representations are those in (D.13). This is a typed arithmetic comparison, not an inference from the shapes of the two centralizer rings. The CM reciprocity used next is \cite[Corollary3.4]{dLozano2022}, or equivalently \cite[Theorem3.10]{dMilne2007}, both opened.
\end{proof}

\begin{statement}{D.16}{Proved}{Exact joint CM image: one additional quadratic relation}
Let $K_3=\mathbb Q(\sqrt{-3})$, $K_4=\mathbb Q(i)$ and $L=K_3K_4$. In the product of the two normalizer groups of D.15 with common determinant and the modulo-two conditions of LM2 C.d4, the joint Galois image has index two. Its exact additional condition is equality of the two restrictions to
\[
 M=\mathbb Q(i,\sqrt[4]{3});\qquad
 K_3^{\rm ab}\cap K_4^{\rm ab}=\mathbb Q^{\rm ab}M.         \tag{D.14}
\]
Thus the two-tori lead has an open image, but the Galois image does not fill the proposed full pair of equal norms. The missing coset is a single $C_2$; the existence of an element of $\widehat{\rm GT}_0$ realizing that coset is a separate unresolved realization question.
\end{statement}
\begin{proof}
Both quadratic rings have class number one, as follows already from their Euclidean norms. CM reciprocity identifies the Weber-function torsion fields with the ray class fields; taking all levels in \cite[Corollary3.4]{dLozano2022} gives $K_j^{\rm ab}$. The field defined by the Tate representation modulo $\pm1$ over $K_j$ contains these Weber values ($x^3$ for $E_3$, $x^2$ for $E_4$), and is contained in $K_j^{\rm ab}$ since the Tate action over $K_j$ is linear over its commutative CM order. It is therefore exactly $K_j^{\rm ab}$. An element outside $G_{K_j}$ cannot act as $\pm I$, because it conjugates a nonreal CM endomorphism to its distinct complex conjugate. Hence the field of the projective representation over $\mathbb Q$ is also exactly $K_j^{\rm ab}$.

More precisely the image fills the specified individual factor. Reciprocity identifies $\operatorname{Gal}(K_j^{\rm ab}/K_j)$ with $\widehat{\mathcal O}_{K_j}^{\times}/\mathcal O_{K_j}^{\times}$. Reduction modulo2 maps $\mu_6/\{\pm1\}$ isomorphically to $(\mathcal O_{K_3}/2)^\times$, and $\mu_4/\{\pm1\}$ isomorphically to $(\mathcal O_{K_4}/2)^\times$. Thus removing that finite root of unity gives an isomorphism with the units congruent to one modulo2, modulo the single global sign. The curves in (D.13) have their full two-torsion over the stated fields, so their projective representations give exactly this isomorphism. The determinant is restriction to the cyclotomic subfield by the Weil pairing. Outside $G_L$, the two conjugation signs are already forced by that determinant and are all realized by $G_{\mathbb Q}$.

It remains to compute the intersection; this step determines the joint image, not merely its openness. Put $E=K_3^{\rm ab}\cap K_4^{\rm ab}$. It contains $\mathbb Q^{\rm ab}$. In $G=\operatorname{Gal}(E/\mathbb Q)$ the two distinct index-two subgroups fixing $K_3,K_4$ are abelian. Their intersection is central, and together they generate $G$. Choosing representatives $x,y$ of their two nontrivial cosets, the closed commutator group is generated by the central element $[x,y]$. Since $x^2$ and $y^2$ are central, this element has order at most two. Therefore $[E:\mathbb Q^{\rm ab}]\leq2$.

The splitting field $M$ of $X^4-3$ has dihedral group of order eight: its automorphisms are $(a,b)$, $0\leq a<4$, $0\leq b<2$, acting by $\sqrt[4]3\mapsto i^a\sqrt[4]3$, $i\mapsto(-1)^b i$, with product $(a,b)(c,d)=(a+(-1)^bc,b+d)$. The subgroups fixing $K_4$ and $K_3$ are respectively $b=0$ and $a+b=0\pmod2$, both abelian. Thus $M\subset E$, and $M$ is nonabelian over $\mathbb Q$. This realizes the bound and proves (D.14). Galois theory now identifies the joint image with the fibre product over this exact intersection. The common-determinant target is the fibre product over $\mathbb Q^{\rm ab}$, so its index is two. The program checks the dihedral group and this finite index explicitly.
\end{proof}

\begin{statement}{D.17}{Proved}{The local indices at two and three}
For $\mathcal O_L=\mathbb Z[i,z]$, $z^2-z+1=0$, the paired local norm map to the equal-norm units of $\mathcal O_{K_3}\otimes\mathbb Z_p$ and $\mathcal O_{K_4}\otimes\mathbb Z_p$ has index two at each $p=2,3$, and is onto at every other prime. Its two local obstruction classes are identified by the global root-of-unity pairs, leaving the single global relation of D.16. Exact finite enumeration gives
\[
\begin{array}{c|rrrrrrrr}
 p^k&2&4&8&16&32&3&9&27\\\hline
 \text{equal-norm pairs}&6&48&384&3072&24576&24&648&17496\\
 \text{norm image}&6&24&192&1536&12288&12&324&8748.
\end{array}                                               \tag{D.15}
\]
\end{statement}
\begin{proof}
At either bad prime one quadratic subfield is unramified and the other ramified. Given a compatible pair $(a,b)$, first lift its coordinate in the unramified subfield through the ramified quadratic norm. It is a norm because the quadratic Hilbert-symbol obstruction equals the Hilbert symbol of its rational norm, and that rational norm is already a norm from the ramified quadratic subfield. Its valuation is zero, so a lift is a unit. The remaining discrepancy has norm one in the ramified subfield. Hilbert90 writes it as $c/\bar c$. It can be corrected by a unit from the biquadratic field exactly when the valuation of $c$ in that ramified subfield is even: norms from its unramified quadratic extension have even valuation, and multiplying $c$ by a rational scalar changes its valuation by an even integer. This gives precisely one $C_2$ obstruction. The unit norm in an unramified extension is onto, so no further obstruction remains.

The pairs $(-1,1)$ and $(1,i)$ represent the nonzero class at each prime. The identity
$N_{L/K_3}(z+i\bar z)=-1$, $N_{L/K_4}(z+i\bar z)=-i$
shows that their sum is zero. At other primes the integral norm torus has connected smooth kernel; norm surjectivity on residue fields and on successive principal-unit quotients proves surjectivity. The code enumerates the rings directly: for $a+bz+i(c+dz)$ its two norms are
$(a+bz)^2+(c+dz)^2$ and
$(a+ci)^2+(a+ci)(b+di)+(b+di)^2$.
It tests invertibility by the rational norm and records the entire image size. The signs above are also tested at every displayed level. The global roots of unity contribute the same nonzero class at both bad primes, hence remove only the diagonal $C_2$. This agrees with, and does not replace, the field-intersection proof of D.16.
\end{proof}

\begin{statement}{D.18}{Known}{The genuine cusp congruence action retains the full rank-one Kummer coordinate}
In the positive tangential normalization $\lambda=0$, tangent1, take $q^{1/n}$ with positive real leading coefficient $256^{-1/n}$ and $\zeta_n=\exp(2\pi i/n)$. Use the resulting Tate basis, matched to the Legendre cycles over $[0,\lambda]$ and $[\lambda,1]$, with positive local monodromy $\left(\begin{smallmatrix}1&2\\0&1\end{smallmatrix}\right)$. The genuine normalized congruence representation is
\[
 V_g=\begin{pmatrix}\chi(g)&-8\kappa_2(g)\\0&1\end{pmatrix}.            \tag{D.16}
\]
Compatible roots may be chosen so that
$\kappa_{-1/2}=(\chi-1)/2-\kappa_2$. Consequently the continuous inverse comparison on this image is
\[
 (\chi,V_{12})\longmapsto
 \left(\chi,\frac{\chi-1}{2}+\frac{V_{12}}8\right).          \tag{D.17}
\]
Multiplication by eight is a homeomorphism from $\widehat{\mathbb Z}$ onto $8\widehat{\mathbb Z}$, so no part of this Kummer coordinate is lost in the genuine normalized congruence action. In the basis of the manuscript with local monodromy $\left(\begin{smallmatrix}1&-2\\0&1\end{smallmatrix}\right)$, conjugation by $\operatorname{diag}(-1,1)$ reverses the off-diagonal sign. On \emph{outer} congruence actions this translation is forgotten: the comparison retains only $\chi$ modulo squares.
\end{statement}
\begin{proof}
The Legendre invariant is
$j(\lambda)=256(1-\lambda+\lambda^2)^3/(\lambda^2(1-\lambda)^2)$.
Inverting $j(q)=q^{-1}+744+\cdots$ gives
$q=\lambda^2/256\,(1+\lambda\mathbb Q[[\lambda]])$.
Tate uniformization gives the exact torsion sequence with basis $\zeta_n,q^{1/n}$ and matrix $\left(\begin{smallmatrix}\chi&\kappa_q\\0&1\end{smallmatrix}\right)$; see \cite[\S11, Theorem11.17 and the preceding uniformization]{dJoshiTate}. For the formal field used here, the comparison follows by base change of the universal Tate equation and its torsion representatives over $\mathbb Z((q))$: the convergent uniformization identities are identities of Laurent series with integral coefficients, and remain identities after $q=q(\lambda)$. Thus the torsion representatives are $\zeta_n^a q^{b/n}$ over $\overline{\mathbb Q((\lambda))}$. The tangential section fixes compatible $\lambda^{1/n}$ and the unique roots with constant coefficient one of the last power-series unit. Its only remaining Kummer action on $q^{1/n}$ is therefore that on $256^{-1/n}$, namely $-8\kappa_2$. The Legendre curve is a quadratic twist of the split Tate curve at this cusp; that twist multiplies the entire matrix by a single sign and disappears in the projective congruence local system. Normalizing its lower-right entry to one yields (D.16). To fix the residual shear, let $K(m)=\int_0^{\pi/2}(1-m\sin^2\theta)^{-1/2}d\theta$. The branch-cut periods have ratio $\tau=iK(1-\lambda)/K(\lambda)$, purely imaginary and positive for $0<\lambda<1$. Hence $q=\exp(2\pi i\tau)$ and its chosen roots are positive real. A shear would replace $q^{1/n}$ by $\zeta_n^b q^{1/n}$ and change the off-diagonal entry by $b(1-\chi)$; positivity for every $n$ forces $b=0$. This proves that the Tate basis is the specified global Legendre branch-cut basis, rather than merely a basis with the same local monodromy.

Choose $(-1)^{1/n}=\zeta_{2n}$. Its Kummer cocycle is $(\chi-1)/2$, which proves (D.17) by multiplication of root systems. The stated change of basis explicitly handles the orientation convention. Finally conjugation by $V_g$ on the special linear congruence quotient has outer class determined by its determinant modulo squares, as established in Part A; its unipotent upper-triangular factor is inner. Thus the often proposed claim that the rank-one translation is visible only metabelianly is false for genuine normalized actions, while the precise outer quotient does discard it.
\end{proof}

\begin{statement}{D.19}{Open}{The residual GT and full profinite realization questions}
D.1--D.4 settle item3, including general openness, the quadratic failure of surjectivity, joint independence, and the local inertia images. D.15--D.18 settle the stated Galois comparison of the two CM coordinates and the rank-one congruence coordinate. The unresolved parts of D(e) and D-after-LM2(1) concern the full profinite GT beta/gamma presentation and actual realization outside Galois: whether the extra leading metabelian directions of D.11 lift from $\widehat{\rm GT}_0$ or the full profinite loop-preserving group, and whether $\widehat{\rm GT}_0$ supplies the missing index-two CM coset of D.16. These assertions remain on the coverage ledger; none follows from solving only matrix or Lie equations.
\end{statement}
\begin{proof}[Concrete residual instances]
The first extra metabelian special vector occurs at weight five, where the free special source has dimension three and its invariant part dimension one; exact preimages are in \texttt{d\_metabelian.json}. The first extra invariant metabelian vector occurs at weight nine. D.12 supplies their pro-$\ell$ realizations and identifies the nonnilpotent-quotient obstruction to transferring that construction. For the CM problem, D.17 supplies an explicit finite matrix coset, represented locally by $(-1,1)$ at one bad prime and the identity at the other, whose two obstruction classes disagree. D.16 proves that this coset has no Galois lift. Its possible lift to an actual full profinite $\widehat{\rm GT}_0$ element must solve the conjugator and generation equations before congruence reduction; the matrix reconstruction of LM2 C.d4 does not do so. The source search reached the original Ihara Part II through the RIMS1224 PostScript file, and the exact modern profinite lift criterion through Dolgushev--Guynee. Part I's RIMS1177 catalogue entry has no download; the PDF, PostScript and DVI locations returned404. D.20--D.22 record the additional results actually obtained. The unresolved steps are the complete lift tests for the two indicated metabelian vectors and for the missing CM coset.
\end{proof}

\subsection{The original profinite beta/gamma relations and their graded consequence}

\begin{statement}{D.20}{Known}{Ihara's exact cyclic-cover compatibility and gamma multiplication}
Write a full profinite GT element as $(\lambda,f)$ acting by
$x\mapsto x^\lambda$, $y\mapsto f^{-1}y^\lambda f$.
For $n\geq1$ let $H_n=\ker(\widehat F_2\to\mathbb Z/n)$, where $x\mapsto1$, $y\mapsto0$, and define the continuous epimorphism
\[
 \phi_n:H_n\longrightarrow\widehat F_2,\qquad
 x^n\mapsto x,\quad y\mapsto y,\quad x^{-j}yx^j\mapsto1\quad(0<j<n).
\]
For every $g\in G_{\mathbb Q}$ one has
\[
 \phi_n(f_g)=y^{-\kappa_n(g)}f_g,\qquad
 \phi_n\circ g|_{H_n}=g\circ\phi_n .                         \tag{D.18}
\]
Here the roots of the positive integer $n$ use the positive real branches of D.2. These are \cite[Theorem1 and Proposition1]{dIhara2000}.

The full profinite beta function of $\sigma\in\widehat{\rm GT}$ has Ihara's factorization
\[
 B_\sigma(s,t)=\frac{\Gamma_\sigma(s)\Gamma_\sigma(t)}{\Gamma_\sigma(s+t)},
 \qquad s,t\in\mathbb Q/\mathbb Z.                          \tag{D.19}
\]
The coefficient ring and completed group-ring interpretation are those of the hyperadelic gamma function; equality precedes every prime-power specialization. For Galois elements Anderson's multiplication formula, in Ihara's corrected convention, is
\[
 \frac{\prod_{nc=0}\Gamma_g(s+c)/\Gamma_g(c)}{\Gamma_g(ns)}
       =1\otimes\exp(-2\pi i\,ns\kappa_n(g)).                 \tag{D.20}
\]
Ihara Part II, Introduction and \S2, states (D.19) and (D.20), attributes (D.20) to Anderson's Corollary8.6.3, and gives the following exact relations among its intermediate subgroups:
$G_{\mathbb Q}\subset\widehat{\rm GT}_K\subset\widehat{\rm GT}_{A'}$ and
$G_{\mathbb Q}\subset\widehat{\rm GT}_A\subset\widehat{\rm GT}_{A'}$.
Here $K$ means all the commuting squares in (D.18), $A$ means (D.20), and $A'$ means the two distribution identities in the proof below.
\end{statement}
\begin{proof}
The original RIMS1224 version was downloaded, decompressed, converted from PostScript, and read. Its proof of Theorem1 identifies $H_n$ with the fundamental group of $\mathbb P^1-\{0,\mu_n,\infty\}$, and $\phi_n$ with the map induced by filling the punctures different from1, after the change $t\mapsto t^n$. Coefficientwise Galois action on Puiseux series commutes with this change, proving the second equality of (D.18). Applying the two sides to $y$ shows that $\phi_n(f_g)f_g^{-1}$ centralizes $y^{\chi(g)}$, hence is a power of $y$. Its exponent is the logarithmic Kummer quasi-cocycle $\Psi_n^{(0)}(g)$, characterized in the original Theorem1 by
$(n^{-r})^{g-1}=\exp(2\pi i r\Psi_n^{(0)}(g))$ for $r\in\mathbb Q$.
Thus $\Psi_n^{(0)}=-\kappa_n$, proving the first equality with its sign.

For completeness the identities defining $A'$ are
\[
 \Psi_{mn}^{(0)}=\Psi_m^{(0)}+\Psi_n^{(0)},\qquad
 \Psi_{mn}^{(in)}=\Psi_m^{(i)}\quad(i\ne0\bmod m).
\]
The functions $\Psi_n^{(i)}$ are defined by the action of $f$ on
$(1-\zeta_n^{-i}t^{1/n})^r$; Proposition2 proves these identities directly from (D.18). Proposition3 proves that they are exactly the logarithmic derivative of (D.20), using Part I's Theorem2, whose formula is reproduced there. The original Introduction explicitly says that the proof of (D.19) uses the five-cycle relation. These statements specify which equation is used and which group it constrains.

Anderson's original Corollary8.6.3, p.103, was also visually
checked. With $\mathcal W=\prod_pW(\overline{\mathbb F}_p)$
from his \S7.1, it states, in
$(\overline{\mathbb Q}\otimes\mathcal W)^\times$,
\[
 \prod_{nb=a}\Gamma_g(b)=
 \bigl((n^{\langle-a\rangle})^{g-1}\otimes1\bigr)
 \Gamma_g(a)\prod_{nc=0}\Gamma_g(c).
\]
The roots of $n$ are positive real. Put $a=ns$ and divide
by the two displayed products; their Kummer factor is
$\exp(-2\pi i\,ns\kappa_n(g))\otimes1$.
Ihara's coefficient convention places $\mathcal W$ first;
the tensor flip $z\otimes w\mapsto w\otimes z$ gives exactly
(D.20), including its sign.

There is a precise obstruction to reversing the logarithmic derivative in the hyperadelic ring. For complementary nonzero idempotents $e,e'\in\widehat{\mathbb Z}$, $ee'=0$, the nonconstant element $e x^{e'}$ has derivative zero. It is nonconstant already after projection to a coefficient prime supported on $e$ and a finite deck quotient supported on $e'$. Thus the kernel of the derivation contains more than constants. Ihara's final Remark explicitly records that $\widehat{\rm GT}_K\subset\widehat{\rm GT}_A$ was not proved there. We use only the two inclusions actually proved.
\end{proof}

\begin{statement}{D.21}{Proved}{Galois and full GT have the same rational metabelian graded image through weight ten}
For every prime $\ell$, the rational graded images of the actual groups $G_{\mathbb Q}$ and $\widehat{\rm GT}$ in the metabelian outer action, in weights1 through10 and after removing the cyclotomic multiplier, both have dimensions
\[
                 0,0,1,0,1,0,1,0,1,0.                     \tag{D.21}
\]
In weight $n\geq3$ odd the common line, expressed in the $p(X,Y)[x,y]$ coordinate of D.11, is
\[
 \mathbb Q_\ell\,
 \frac{(X+Y)^n-X^n-Y^n}{XY}.                                \tag{D.22}
\]
The inclusion of this line in the $\widehat{\rm GT}_0$ upper-bound space is the explicit polynomial inclusion in D.11. The latter upper bound has an additional line in weight nine.
\end{statement}
\begin{proof}
We give the graded comparison explicitly. Map the free profinite group to its pro-$\ell$ quotient, take logarithms in each rational nilpotent quotient, and then take the first nonzero homogeneous term of a filtered GT element. Expanding the two-cycle, three-cycle and five-cycle equations gives their infinitesimal equations for this term. The infinitesimal pure-braid generators are the initial classes of the logarithms of the pure-braid generators. Their relations are
$[t_{ij},t_{kl}]=0$ for four distinct indices and
$[t_{ij},t_{ik}+t_{jk}]=0$; the successive split fibrations forgetting a strand identify their free Lie kernels, so these relations give the associated graded pure-braid Lie algebra used in the expansion. Thus the displayed map sends the leading term into $\mathfrak{grt}_1(\mathbb Q_\ell)$.

Furusho's \cite[Corollary0.3, Theorem0.5, Proposition5.1 and Lemma5.2]{dFurusho2011} gives the gamma constraint on its metabelian image. Here is the comparison with our coordinate, including the factor that matters. For a Lie element $h=h_xx+h_yy$, put $M(h)=(h_yy)^{\rm ab}$. Direct expansion gives
\[
 M([x,y])=XY,\quad M([x,h])=XM(h),\quad M([y,h])=YM(h)
 \quad(h\in L').
\]
It follows that $M(p(X,Y)[x,y])=XYp(X,Y)$ and that $M$ kills $L''$. Furusho's gamma constraint is
$M(h)=\gamma(X)+\gamma(Y)-\gamma(X+Y)$.
In homogeneous degree $n$ it is a scalar multiple of
$X^n+Y^n-(X+Y)^n$. Since multiplication by $XY$ is injective, this bounds the image by the line (D.22). Equivalently the actual module action has leading multiplier $-XYp$, because
$D([x,y])=[x,[p[x,y],y]]=-XYp[x,y]$.
This is the leading term of the full profinite beta multiplier in (D.19), rather than an unlabelled identification of the two series.

The program \texttt{d\_gamma.py} checks $M(h)=XY\bar h$ on all224 Lyndon basis words of degrees2 through10, and solves for each gamma polynomial in the exact invariant image basis of \texttt{d\_metabelian.json}. In weight nine its coordinates in that basis are $(15/2,9/2)$. The first basis vector $X^2Y^2(X+Y)^3$ is independent of the gamma line, verified by an exact rank-two calculation. Thus the comparison with the extra vector is computed in the same coordinate system as D.11.

The special equation forces $X+Y$ to divide $p$, as proved in D.11. At $Y=-X$ the numerator of (D.22) is $-(1+(-1)^n)X^n$, so every positive even-degree scalar must vanish in characteristic zero. Weight one vanishes directly. For each odd weight3,5,7,9, E.3 supplies a nonzero Galois class for every $\ell$, using \cite[Theorem8.4, \S8.4, Lemma8.5 and Remark8.6]{E_HainMatsumoto}; at $\ell=37$ D.10 also supplies direct finite-field certificates. Since $G_{\mathbb Q}\subset\widehat{\rm GT}$, these lower bounds equal the upper bounds just proved. This proves (D.21) for actual profinite source groups with rational graded target. It does not replace either source by its pro-unipotent automorphism group.
\end{proof}

\begin{statement}{D.22}{Known}{The exact finite-quotient criterion for the remaining GT0 lifts}
Let $N\triangleleft B_3$ have finite index and lie in $PB_3$, and let $\operatorname{GT}(N)$ be the finite set of pairs satisfying the two hexagon equations, the commutator condition on $f$, the unit condition on $\lambda$, and surjectivity onto $B_3/N$. Write $R_{K,N}$ for reduction from a finer such quotient. Then the image of the actual full profinite group is exactly
\[
 \operatorname{Im}\bigl(\widehat{\rm GT}_0\longrightarrow\operatorname{GT}(N)\bigr)
   =\bigcap_{K\triangleleft B_3,\ [B_3:K]<\infty,\ K\subset N}
          R_{K,N}\bigl(\operatorname{GT}(K)\bigr).             \tag{D.23}
\]
This is \cite[Theorem5.2, Corollary5.4 and Proposition5.5]{dDolgushevGuynee2024}. It includes all finite quotients of the braid group, including the nonnilpotent ones. Consequently the finite Lie and CM solutions in D.11 and D.17 specify candidates and concrete finite tests, but their realization in $\widehat{\rm GT}_0$ requires the intersection in (D.23).
\end{statement}
\begin{proof}
One inclusion is reduction of a genuine element. For the other, the fibres over a fixed member of the right side form an inverse system of nonempty finite sets indexed by all finer quotients. Any finite collection is dominated by their intersection, so the system has the finite intersection property. Compactness gives a compatible family. Theorem5.2 identifies these compatible pairs with the profinite group of hexagon-preserving invertible automorphisms, and Proposition5.5 identifies that group with Harbater--Schneps' $\widehat{\rm GT}_0$. Surjectivity in each finite quotient supplies the invertibility condition; it cannot be omitted. This proves the equality, including the exact hypothesis that the pro-$\ell$ lift of D.12 alone does not supply.
\end{proof}

\begin{statement}{D.23}{Proved}{The exact coefficient map from the profinite universal measure to every rational Kummer character}
Use the positive interval paths and the Ihara convention of D.20. Write
$c=[x,y]$, $\Lambda=\widehat{\mathbb Z}[[\widehat{\mathbb Z}^2]]$,
and let $s,t$ act on $\widehat F_2'/\widehat F_2''=\Lambda c$ by conjugation by $x,y$.
For $g\in G_{\mathbb Q}$ put $f_g\bmod\widehat F_2''=u_g c$ and
$A_\lambda(s)=(s^\lambda-1)/(s-1)$, where $\lambda=\chi(g)$.
Then the universal measures of D.5 satisfy the exact identities
\[
 F_g=A_\lambda(s)A_\lambda(t)
          \bigl(1+(s-1)(t-1)u_g\bigr),\qquad
 B_g=1+(s-1)(t-1)u_g.                                    \tag{D.24}
\]
For a positive integer $n$, reduce $u_g$ by $t=1$, $s^n=1$, and write
$u_{g,n}=\sum_{j=0}^{n-1}a_j(g)s^j\in\widehat{\mathbb Z}[\mathbb Z/n]$.
The coefficient map is
\[
 \kappa_n(g)=a_0(g)-a_{n-1}(g).                            \tag{D.25}
\]
Consequently the homomorphism of scan157 is recovered from $(\chi,F)$ by
\[
 g\longmapsto
 \left(\chi(g),\frac{\chi(g)-1}{2}
                 -\bigl(a_0(g)-a_1(g)\bigr)\right),        \tag{D.26}
\]
where the two coefficients in (D.26) use $n=2$.
For $p\ne\ell$, the extraction map in (D.25) on the $\mathbb Z_\ell$ coefficient component does not factor through the quotient with only pro-$\ell$ deck groups; the proof gives an explicit element of its kernel on which extraction is1. For every odd $\ell$, this failure already occurs on actual Galois elements: a pro-$\ell$ tame-inertia subgroup at2 has trivial full cyclotomic character and trivial normalized pro-$\ell$ action on $\pi_1(\mathbb P^1-\{0,1,\infty\})$, but its $\kappa_{2,\ell}$ image is all of $\mathbb Z_\ell$.
\end{statement}
\begin{proof}
All divisions in (D.24) are well-defined continuous operations on their stated images. The product $(s-1)(t-1)$ is a non-zero-divisor by Otsubo's Corollary6.2, hence its individual factors are also non-zero-divisors. The finite geometric sums defining $A_m$ for integers $m$ extend by continuity to $\lambda\in\widehat{\mathbb Z}$. For a unit $\lambda$,
$A_\lambda(s)A_{\lambda^{-1}}(s^\lambda)=1$, by multiplication by $s-1$ and cancellation, so $A_\lambda$ is a unit. The inverse on a closed divisible image is continuous because the completed group ring is compact and Hausdorff.

In additive notation for the derived module,
$f_g^{-1}y^\lambda f_g=(t^\lambda-1)u_gc\cdot y^\lambda$.
The commutator identities give
\[
 [x^\lambda,y^\lambda]=A_\lambda(s)A_\lambda(t)c,
 \qquad
 [x^\lambda,(t^\lambda-1)u_gc]=(s^\lambda-1)(t^\lambda-1)u_gc.
\]
This proves the first identity of (D.24); Otsubo's comparison (D.5) proves the second. To fix the path convention, the lift of $[x,y]$ along the positive interval has homology chain
$(s-1)(1-t)\delta=-(1-s)(1-t)\delta$. Thus it is the negative of Otsubo's Pochhammer generator. The factor is the constant $-1$, so their Galois multipliers agree for every $\chi$, not only on the cyclotomic kernel. The leading derivation convention in D.11 uses $D(y)=[f,y]$, so its $f$ is the negative of the leading term of this Ihara conjugator; this accounts for the sign $-XYp$ in D.21.

Consider the Schreier generators $e_j=x^{-j}yx^j$ of $H_n$, besides $x^n$. In $H_n^{\rm ab}$,
$c=e_{n-1}-e_0$. Conjugation by $x^a$ sends this to
$e_{-a-1}-e_{-a}$, with indices modulo $n$; conjugation by $y$ acts trivially on this abelianization. Therefore the coefficient of $y$ in $\phi_n(u_gc)$ is $a_{n-1}-a_0$.
By Ihara's exact equality (D.18), that coefficient is $-\kappa_n(g)$, proving (D.25). Formula (D.26) now follows from the chosen roots
$(-1/2)^{1/m}=\zeta_{2m}/2^{1/m}$ and
$\kappa_{-1}=(\chi-1)/2$. The formulas give a continuous map from the actual metabelian action to the affine image of D.1, with no assertion that arbitrary coefficients satisfy Galois relations.

Here is the promised prime-to-$\ell$ kernel element. In
$\mathbb Z_\ell[[\widehat{\mathbb Z}^2]]$, take the normalized Haar idempotent $e_p$ of the $\mathbb Z_p$ factor of the first deck group. Its images at finite $p^r$-levels are
$p^{-r}\sum_{j=0}^{p^r-1}s^j$; these are compatible since $p$ is a unit in $\mathbb Z_\ell$. Projection to the pro-$\ell$ deck quotient sends $e_p$ to1. At $t=1,s^p=1$ its image is $p^{-1}\sum_{j=0}^{p-1}s^j$. Hence $1-e_p$ lies in the kernel of the pro-$\ell$ deck projection but its coefficient extraction is1. This proves the exact failure of that factorization on the ambient module.

For the asserted arithmetic witness, take an odd $\ell$ and a pro-$\ell$ Sylow subgroup $J$ of inertia at2. Its tame quotient is $\mathbb Z_\ell(1)$. Every root of unity of odd order lies in the maximal unramified extension of $\mathbb Q_2$, and the action on the $2$-power roots has pro-$2$ image. Therefore $\chi(J)=1$. The extensions obtained by adjoining $2^{1/\ell^r}$ over the maximal unramified field are totally tamely ramified of degree $\ell^r$; thus $\kappa_{2,\ell}(J)=\mathbb Z_\ell$. The pro-$\ell$ outer action of the thrice-punctured line is unramified at2, as in E.3 and \cite[\S8.4]{E_HainMatsumoto}. It follows that the normalized Ihara action of any $g\in J$ is inner. An inner automorphism fixing $x$ has conjugator in the centralizer $\langle x\rangle$. Comparing its action on $y$ with $f_g^{-1}yf_g$, where $f_g\in F_\ell'$, shows that $f_g$ differs from a power of $x$ by a power of $y$. Abelianization forces both powers to vanish, so this normalized action is the identity. Hence its pure pro-$\ell$ universal measure is1, while the full profinite coefficient extraction ranges over $\mathbb Z_\ell$. This proves the asserted failure on actual Galois elements, even when the whole cyclotomic character is retained.

The program \texttt{d\_kummer\_measure.py} independently computes abelian Fox derivatives and Schreier coefficients on160 explicit free words at all24 levels $1\leq n\leq24$, giving3840 exact checks of (D.25)'s coefficient identity. It also checks210 commutator-action identities for positive and negative integer exponents, which extend continuously as above. Its output does not substitute computation for the Galois theorem (D.18); that theorem supplies the arithmetic meaning of the checked coefficients.
\end{proof}

\subsection{An additional full profinite lifting test}

\begin{statement}{D.24}{Known}{What unrestricted profinite lifting supplies, with an exact peripheral failure}
Put $P=\widehat F_2/\widehat F_2''$. The canonical homomorphism
\[
 \operatorname{Aut}(\widehat F_2)\longrightarrow
 \operatorname{Aut}(P)
\]
is surjective. Its kernel contains an automorphism which does not
preserve the conjugacy class of the peripheral subgroup generator
$x$, although it acts identically on the entire metabelian quotient.
Thus the unrestricted lifting theorem supplies lifts in the displayed
source; it supplies no peripheral conclusion for the chosen lift.
The existence of peripheral lifts for the additional directions in
D.13 remains a separate unsolved step of that calculation.
\end{statement}
\begin{proof}
We use the profinite Gasch\"utz lemma in
Bleher--Chinburg--Lubotzky, Lemma~2.3 and Corollary~2.4
\cite{D_BleherChinburgLubotzky}. Applied to a quotient of a
two-generated free profinite group, it lifts any generating pair
of the quotient to a generating pair of the free profinite group.
For an automorphism of the quotient, lift the images of $x,y$ in
this way. The resulting endomorphism of $\widehat F_2$ is
surjective and hence invertible, by Hopficity of finitely generated
profinite groups. This proves the asserted surjectivity and explains
the imported theorem in our precise setting.

Here is a lift on which the peripheral condition fails. Let
\[
 q:\widehat F_2\longrightarrow A_5,\qquad
 q(x)=(12345),\quad q(y)=(123),
\]
and let $t=(12)\in S_5$. The map
\[
 (p,q):\widehat F_2\longrightarrow P\times A_5
\]
is onto. Both projections are onto, and Goursat's argument identifies
any proper subdirect product with a nontrivial common quotient of
$P$ and $A_5$. The latter group is simple and perfect, whereas every
quotient of $P$ has trivial second derived subgroup. There is no such
common quotient, proving the assertion for this actual map.

The continuous automorphism
$\operatorname{id}_P\times\operatorname{Ad}(t)$ of $P\times A_5$
lifts by the same lemma to an automorphism $\alpha$ of
$\widehat F_2$. It lies in the kernel of the displayed map on
automorphism groups. The two elements $q(x)$ and $tq(x)t^{-1}$
belong to the two different conjugacy classes of five-cycles in
$A_5$. Consequently $\alpha(x)$ cannot be conjugate to $x$ in
$\widehat F_2$. A loop-preserving lift here would have multiplier
one, because its action on $P^{\rm ab}=\widehat{\mathbb Z}^{\,2}$
is the identity. It would therefore have to send $x$ to a conjugate
of $x$, which this lift does not do.

The program \texttt{d\_lifting.py} verifies the order sixty quotient,
generation by the specified pair, simplicity by the normal closures
of all five conjugacy classes, perfectness, and the two disjoint
classes of twelve five-cycles. It also checks all $3600$
multiplication identities for $\operatorname{Ad}(t)$.
This example diagnoses the missing condition in the attempted
unrestricted-lifting route. It makes no claim that a different lift
of a given metabelian automorphism cannot preserve all three loops.
\end{proof}

\subsection{The CM field and explicit nonarithmetical finite shadows}

\begin{statement}{D.25}{Proved}{A finite CM entanglement field and nonarithmetical GT0 shadows}
Put $r_0=(3/4)^{1/4}$ and $M'=\mathbb Q(i,r_0)$. There are exact equalities
\[
 \mathbb Q(E_3[4])=M'=\mathbb Q\bigl(x(E_4[3])\bigr),\qquad
 M'\mathbb Q^{\rm ab}=M\mathbb Q^{\rm ab},                  \tag{D.25.1}
\]
where $E_3:y^2=x^3+1$, $E_4:y^2=x^3-x$, and $M$ is D.16's field.
The projective $E_3[12]$ representation recovers the full $E_3[4]$
representation: choose its unique global sign fixing the rational
point $(0,1)$ of order three. Thus the extra relation of D.16 is
already detected by $T\bmod12$ and $W\bmod3$, with this sign rule.

Here is a finite refinement of the matrix problem, including
generation and the two hexagons. Write
\[
 B=\langle s,r:s^2=r^3=1\rangle,\quad
 Q=\mathrm{SL}_2(\mathbb Z/12)/\{\pm I\},\quad
 N=\ker(B\to Q),\quad G_j=B/[N,N]N^j.
\]
Use $s\mapsto\left(\begin{smallmatrix}0&1\\-1&0\end{smallmatrix}\right)$
and $r\mapsto\left(\begin{smallmatrix}0&-1\\1&-1\end{smallmatrix}\right)$
projectively. Then $|Q|=576$, $N$ is free of rank97, and
$|G_j|=576j^{97}$ for $j=2,3,6$.
Among the48 pairs of matrix conjugators with multiplier one and
$V_{12}=I$,24 satisfy the CM relation and24 violate it.
Exact computation in $G_2$ and $G_3$ gives
\[
\begin{array}{c|rr}
\text{condition}&\text{CM-compatible}&\text{extra CM coset}\\\hline
\text{matrix conjugators}&24&24\\
\text{cusp relation in }G_2&12&12\\
\text{cusp relation in }G_3&24&24\\
\text{automorphisms of }G_6&12&12.
\end{array}                                               \tag{D.25.2}
\]
Each surviving row has an explicit certificate whose induced
matrix on $N^{\rm ab}/p$ has rank97 for $p=2,3$.
The twelve extra pairs yield twelve distinct nonarithmetical
finite $\widehat{\rm GT}_0$ shadows for the pullback of
$[N,N]N^6$ to $B_3$. One representative has
\[
 A_{12}=I,\qquad
 \beta_{12}=\begin{pmatrix}9&8\\4&9\end{pmatrix},\qquad
 V_{12}=I.                                                \tag{D.25.3}
\]
Here ``nonarithmetical'' means that no element of $G_{\mathbb Q}$
has that finite shadow. Whether any of these shadows is genuine,
that is, comes from an actual full profinite
$\widehat{\rm GT}_0$ element, is not established.
\end{statement}
\begin{proof}
For a nonzero two-torsion point of $E_3$ with $x$-coordinate
$e\in\{-1,-\zeta_3,-\zeta_3^2\}$, the coordinates of its halves
satisfy $x=e(1\pm\sqrt3)$. Their $y$-coordinates have squares
$-9\pm6\sqrt3$. For the plus sign choose $y$ with
\[
 y=\sqrt3(\sqrt3-1)r_0.
\]
The other square root is obtained using the product
$3\sqrt{-3}$, up to sign. These formulas prove
$\mathbb Q(E_3[4])=\mathbb Q(i,r_0)$ in both directions.
For $E_4$ the three-division polynomial is
$3x^4-6x^2-1$. Choose two roots $a,b$ with
$a^2=1+2/\sqrt3$, $b^2=1-2/\sqrt3$ and $ab=i/\sqrt3$.
Then $w=a+ib$ satisfies $w^4=4/3$; conversely
\[
 a=(w+2/w)/2,\qquad b=(w-2/w)/(2i).
\]
The splitting field is therefore the same $M'$. Since
$r_0=3^{1/4}/\sqrt2$, adjoining the abelian field
$\mathbb Q(\sqrt2)$ proves the second equality in (D.25.1).

The single global sign in this comparison is essential.
For $P\in E_3[4]$ and $Q_0=(0,1)\in E_3[3]$, the addition law gives
\[
 y(P)=\frac{x(P)^2}{4}
       \bigl(x(P-Q_0)-x(P+Q_0)\bigr).
\]
The points on the right belong to $E_3[12]$, proving directly
that its projective coordinate field contains the full
four-torsion field. Equivalently every Galois matrix on $E_3[3]$
fixes $Q_0$, so exactly one of the two simultaneous signs on
$E_3[12]$ has that property. The separate projective
$E_3[4]$ field is only $\mathbb Q(\zeta_{12})$; the displayed
addition formula is the precise map supplying the missing sign.

We describe the finite calculation so that its certificates
can be verified without enumerating $576\cdot6^{97}$ elements.
Let $F=\ker(B\to S_3)$ and $H=q(F)\subset Q$; then $|H|=96$.
The graph $N\backslash\mathcal T$ of the Bass--Serre tree of
$C_2*C_3$ has edges $Q$ and vertices
$Q/\langle s\rangle\amalg Q/\langle r\rangle$.
There are576 edges and480 vertices. A spanning tree gives97
integral fundamental cycles, a basis of $N^{\rm ab}$.
Write $e_q$ for the oriented edge from its $s$-vertex to its
$r$-vertex. The path cocycle into $\mathbb Z[Q]$ is
\[
 P(s)=0,\quad P(r)=e_1-e_r,\quad
 P(gh)=P(g)+q(g)P(h).                                    \tag{D.25.4}
\]
Its restriction identifies $N^{\rm ab}$ with the cycle group.
Reduction modulo $j$ realizes $G_j$ faithfully as the image
of $B$ in $(\mathbb Z/j)[Q]\rtimes Q$.

For $A\in C_H(r)$ and $\beta\in C_H(s)$ choose the recorded
paths $P_A,P_\beta$. There are respectively six and eight such
choices. A lift of the cusp-preserving relation
\[
 \beta s\beta^{-1}A rA^{-1}=sr
\]
exists modulo $p$ exactly when
\[
 (1-s)P_\beta+
 s\bigl((1-r)P_A+(A-1)P(r)\bigr)
 \in (1-s)C_p+s(1-r)C_p,\quad
 C_p=H_1(N\backslash\mathcal T,\mathbb F_p).               \tag{D.25.5}
\]
This is an explicit linear system. The right side has rank96
at $p=2$ and97 at $p=3$. Its solution counts are those in
(D.25.2). For every surviving pair the code gives corrections
$a,b\in C_p$ and checks generation as follows. The induced
equivariant map on edges sends
\[
 e_q\longmapsto q\bigl(P_A+a-P_\beta-b+e_A\bigr).
\]
Its restriction to the97 fundamental cycles has rank97 in
each characteristic. The homomorphism already induces the
identity on $Q$, so it is an automorphism of $G_p$.
Chinese remainders combine the two recorded corrections into
single conjugators in $G_6$, giving an automorphism there.
The code also multiplies the resulting group elements in
the semidirect model and checks both hexagon identities
directly. All spanning-tree choices, matrix pairs, cycle
coefficients and ranks are reproducible from
\texttt{d\_cm\_shadow.py}; its JSON retains the certificates.

The CM test in these coordinates is particularly small.
Use the unique representative of an element of $H$ whose
modulo-four image lies in the four-group generated by
$U^2,W^2$. For every one of the six $A$ choices, $A_4=I$.
The rational order-three line is $\ker(R-I)$ modulo3, spanned
by $(1,-1)$. If $A_3^{-1}$ acts on this line by $-1$, the
global-sign normalization changes $T_4=I$ to $-I$, and the
restriction to $M'$ sends $r_0$ to $-r_0$; otherwise it fixes
$r_0$. On the other factor $\beta_3^{-1}$ is in
$\{\pm I,\pm S\}$; its class is nontrivial exactly for
$\pm S$, again sending $r_0$ to $-r_0$. Comparing these two
bits gives the24/24 split. In (D.25.3) the bits are zero and
one, respectively.

The conjugator bits are invariants of the finite action,
not arbitrary choices in a matrix centralizer. The same
linear computation of (D.25.4) gives
\[
 q\bigl(C_{G_2}(s)\cap\operatorname{im}F\bigr)=1,
 \qquad
 q\bigl(C_{G_3}(r)\cap\operatorname{im}F\bigr)=1.          \tag{D.25.6}
\]
For example the first test asks whether
$(1-s)P_\beta$ lies in $(1-s)C_2$; only $q(\beta)=1$
survives. The second uses
$(1-r)P_A+(A-1)P(r)$ and $(1-r)C_3$, again with only the
identity surviving. Thus $G_6$ uniquely determines the
reductions of both normalized conjugators in $Q$.

Finally we give the exact passage to finite GT shadows.
Put $u=(sr)^{-1}$, $v=sus$, $x=u^2$, $y=v^2$; these are
the images of the standard braid generators and their
squares under $B_3\to B_3/Z(B_3)=B$.
Let $\alpha$ be one of the finite automorphisms just constructed.
It fixes $u$ and satisfies
$\alpha(s)=\beta s\beta^{-1}=sf$, where
$f=s\beta s\beta^{-1}$. Choose a lift of $\beta$ to the
discrete free group $F=\langle x,y\rangle$ and write its
abelianization as $a_0x+b_0y$. Since $s$ swaps $x,y$,
$f^{\rm ab}=(b_0-a_0)(x-y)$.
Set $k=b_0-a_0$, $\beta'=x^k\beta$, and
$f'=s\beta's(\beta')^{-1}$. Then $f'\in[F,F]$.
The gauge $\alpha'=\operatorname{Int}(x^k)\alpha$ still
fixes $u$ and satisfies
\[
 \alpha'(s)=sf',\qquad\alpha'(r)=(f')^{-1}r,
 \qquad\alpha'(v)=(f')^{-1}vf'.                           \tag{D.25.7}
\]
The last two equalities, with $r=uv$ and $s=uvu$, are
exactly the two $m=0$ hexagons (2.1)--(2.2) of
\cite[Definitions2.1 and2.3]{dBDHP2026}.
The pullback $\widetilde N$ of $[N,N]N^6$ to $B_3$ is
normal, finite index, and contained in $PB_3$, since
$G_6\to S_3$ is the standard permutation quotient.
Here $\widetilde N_{\rm ord}=36$: the cusp $sr$ has orders
24,36,72 in $G_2,G_3,G_6$, respectively, its square has
order36 in $G_6$, and the braid center is killed.
Thus $m=0$ gives the required unit multiplier one,
$f'$ has the required commutator image, and $\alpha'$ being
an automorphism gives surjectivity. This proves that
$(0,f'\bmod\widetilde N)$ is a finite GT shadow.

The gauge multiplies $A,\beta,V$ on the left by the same
$x^k$, so it fixes $T=A^{-1}V$ and $W=\beta^{-1}V$.
Equation(D.25.6) makes their reductions recoverable from
the finite shadow. The twelve different extra pairs have
different such coordinates, hence give distinct shadows.
No Galois element can give one of them: its two restrictions
to the single field $M'$ in (D.25.1) must agree. This proves
nonarithmeticality without assuming that an arbitrary
finite shadow lifts to $\widehat{\rm GT}_0$.
For these specific shadows the remaining genuine-lift
test is precisely the intersection over all finer finite
quotients in D.22.
\end{proof}

\subsection{Isolated shadows and a complete tower of homology refinements}

\begin{statement}{D.26}{Proved}{The CM witnesses survive every abelian-kernel refinement}
Keep $B,Q,N,F$ and $G_j$ from D.25, and let
$\widetilde N_j\triangleleft B_3$ be the inverse image of
$[N,N]N^j$. Every $\widetilde N_j$ is an isolated object of
the gentle GT-shadow groupoid: every shadow with that target
induces an automorphism of $G_j$.

Put $C=N^{\rm ab}\simeq\mathbb Z^{97}$, and let $s,r$ act
on $C$ by conjugation. For $t=sr$ define $c=[t^{12}]\in C$.
The integral maps
\[
 L:C\oplus C\longrightarrow C,\quad
 L(b,a)=(1-s)b+s(1-r)a,\qquad
 \widetilde L(b,a,h)=L(b,a)-hc                         \tag{D.26.1}
\]
have Smith invariant factors $(1^{96},2)$ and $(1^{97})$,
respectively. In particular $\operatorname{coker}L=C_2$
and $\widetilde L$ is onto over $\mathbb Z$.

There is an exact sequence and a specified quotient map
\[
 1\longrightarrow\widehat{\mathbb Z}^{97}
 \longrightarrow\mathcal G:=\varprojlim_j G_j
 \longrightarrow Q\longrightarrow1,\qquad
 \widehat B\twoheadrightarrow\mathcal G
   =\widehat B/\overline{[\widehat N,\widehat N]} .       \tag{D.26.2}
\]
Each of D.25's twenty-four certified automorphisms of $G_6$
lifts to a continuous automorphism of $\mathcal G$ inducing
the identity on $Q$ and sending $t$ to $t^\lambda$, for some
$\lambda\in\widehat{\mathbb Z}^{\times}$ with
$\lambda\equiv1\pmod{72}$. A single cusp gauge supplies a
compatible family of GT shadows at all $\widetilde N_j$.
For each of the twelve extra CM pairs, its member at level six
can be chosen to be one of D.25's nonarithmetical shadows,
with $m=(\lambda-1)/2\equiv0\pmod{36}$.
Thus all twelve extra CM pairs survive the entire family of
abelian-kernel refinements of $N$.

This family does not exhaust the finite quotients in D.22:
there is an explicit finite quotient of $B$ below $G_6$ whose
kernel does not contain $[N,N]$. The existence of a genuine
full profinite GT$_0$ lift of these shadows remains undecided.
\end{statement}
\begin{proof}
First we verify isolation directly from finite shadow equations.
Let $(m,f)$ be a shadow at $\widetilde N_j$, put
$\lambda=2m+1$, and choose a representative $f\in[F,F]$.
The hexagons give a surjection $\phi:B_3\to G_j$ and send
the braid center to its $\lambda$th power, which is trivial
in $G_j$. Hence $\phi$ factors through $B$.
The cusp $t$ has order $12j$ in $G_j$: its image in $Q$ has
order twelve, and the cycle $c$ is primitive in $C$.
The last assertion also follows from the integral cycle
coordinates in the calculation below, one of which is one
and all the others are zero. Thus
$(\widetilde N_j)_{\rm ord}=6j$ and the shadow's multiplier
is a unit modulo twelve.

Evaluate the same word $f$ in the linear matrices $U^2,W^2$
of A.17 modulo twelve. Its reduction modulo four is $I$,
since $U^2,W^2$ generate an abelian four-group modulo four.
Consequently the linear matrices
$A=U^\lambda$ and $D=f^{-1}W^\lambda f$ satisfy the braid
identity with positive sign. Indeed the projective identity
first gives $ADA=\pm DAD$, and modulo four the choice
$\lambda\equiv\pm1$ gives the positive identity for
$U^\lambda,W^\lambda$; an invertible matrix cannot equal
its negative modulo four. The pair $A,D$ generates all of
$\mathrm{SL}_2(\mathbb Z/12)$: its generated subgroup
surjects onto $Q$, and a proper such subgroup would be a
complement to $\{\pm I\}$. That complement cannot exist,
because both lifts $\pm S$ of the order-two element $s\in Q$
have square $-I$.

Write $D=(\begin{smallmatrix}a&b\\c_0&d\end{smallmatrix})$.
The trace, determinant, braid identity and generation now give
\[
 a+d=2,\quad ad-bc_0=1,\quad
 c_0(1+\lambda c_0)=0,\quad c_0\in(\mathbb Z/12)^\times.
\]
Thus $c_0=-\lambda^{-1}$ and the matrix
$V=(\begin{smallmatrix}\lambda&\lambda(1-a)\\0&1\end{smallmatrix})$
conjugates $U,W$ to $A,D$, respectively. This repeats the
elementary matrix argument of A.17 using only the finite
shadow hypotheses. It proves $q\phi=\operatorname{Int}(V)q$.
In particular $\phi(N)\subset C/jC$. Since this subgroup
is abelian of exponent $j$, $\phi$ kills $[N,N]N^j$.
It therefore induces a surjective endomorphism of the finite
group $G_j$, hence an automorphism. Its kernel is precisely
$\widetilde N_j$. This is isolation in
\cite[Definition3.13]{dDolgushevGuynee2024}.

For the integral assertion use the same oriented graph and
spanning tree as in D.25. The ninety-seven fundamental cycles
form a matrix with the identity matrix in the non-tree-edge
coordinates, so projection to those coordinates is an exact
integral inverse on $C$. The columns of $L$ are
$(1-s)C$ and $s(1-r)C$ in this basis. The program
\texttt{d\_cm\_refine.py} computes their integral Smith form
as $(1^{96},2)$. Appending the column $-P(t^{12})$ gives
Smith form $(1^{97})$; an independent Hermite-normal-form
calculation of this augmented matrix returns $I_{97}$.
The program retains the complete integer coordinates and
reproduces the spanning tree. Thus these are integer-image
computations, not deductions from ranks modulo two and three.

We construct the lift, retaining the multiplier's full
residue. Fix one of the twenty-four certified matrix pairs
$A_0,\beta_0$ and choose the fixed integral paths $P_A,P_\beta$
of D.25. Let $E\in C$ be its integral cusp error, the left
side of (D.25.5). With cycle corrections $(b,a)$ the exact
cusp equation for multiplier $\lambda=1+12h$ is
\[
       E+L(b,a)-hc=0.                                  \tag{D.26.3}
\]
To see the exponent, $P(t^{1+12h})=P(t)+hc$, since
$q(t)c=c$. This proves (D.26.3) in every finite quotient,
and by continuity for $h\in\mathbb Z_p$.

At $p=2,3$ begin with D.25's correction modulo $p$ and
$h=0\pmod p$. Suppose coefficients have been constructed
modulo $p^k$; their error in (D.26.3) is divisible by $p^k$.
Surjectivity of $\widetilde L$ supplies a correction
$(\delta b,\delta a,\delta h)$ modulo $p$ equal to the
negative error divided by $p^k$. Add $p^k$ times that
correction. This constructs compatible solutions over
$\mathbb Z_p$. Their induced matrix on $C\otimes\mathbb Z_p$
reduces to D.25's invertible matrix modulo $p$, hence is
invertible. The associated map on the quotient $Q$ is the
identity, so it is an automorphism of the whole extension.
The initial values $h=0\pmod2$ and $h=0\pmod3$ are
preserved. Accordingly the combined multiplier satisfies
$\lambda\equiv1\pmod{72}$, not just modulo twelve.
As a check of the power and sign conventions, the program
carries every one of the twenty-four solutions through
exponents $2,4,8,16$ and $3,9,27,81$, and directly verifies
both general-$m$ hexagons in their faithful group models:
there are384 successful identity checks.

For $p>3$ take the identity automorphism on
$\mathcal G_p=\varprojlim_k G_{p^k}$, with multiplier one,
but choose conjugators with the required images
$A_0,\beta_0$ in $Q$. Such choices exist. For example,
an arbitrary lift $a_0$ of $A_0\in C_Q(r)$ has
$a_0ra_0^{-1}r^{-1}=z\in C\otimes\mathbb Z_p$ and
$(1+r+r^2)z=0$. Multiplication by three is invertible,
and the identity
$(1-r)(2+r)=3-(1+r+r^2)$ shows that
$z=(1-r)((2+r)z/3)$. Adjusting $a_0$ by the negative of
this cycle gives a lift centralizing $r$. The same argument
for $s$ uses $(1-s)z=2z$ when $(1+s)z=0$ and division by
two. These are compatible $p$-adic choices, and both belong
to the preimage of $H=q(F)$.

The extensions $\mathcal G_p$ have the same quotient $Q$,
and their fibre product over $Q$ is precisely $\mathcal G$.
This follows either from the path model or from Chinese
remainders on $C/jC$. The constructed automorphisms and
conjugators therefore assemble to $\alpha\in
\operatorname{Aut}(\mathcal G)$ and $A,\beta\in\overline F$,
with $\alpha(s)=\beta s\beta^{-1}$,
$\alpha(r)=ArA^{-1}$ and $\alpha(u)=u^\lambda$.
Here $\overline F$ denotes the image of $\widehat F$ in
$\mathcal G$. Its abelianization is
$\overline F^{\rm ab}=\widehat{\mathbb Z}x\oplus
\widehat{\mathbb Z}y$, because the kernel
$\overline{[\widehat N,\widehat N]}$ is contained in
$\widehat F'$.

Write $\beta^{\rm ab}=ax+by$, set $k=b-a$, and conjugate
$\alpha$ by $x^k$. With $\beta'=x^k\beta$ and
$f'=s\beta's(\beta')^{-1}$, the same abelianization
calculation as in D.25 gives $f'\in\overline F'$.
The resulting automorphism satisfies
\[
 \alpha'(u)=u^\lambda,\qquad
 \alpha'(v)=(f')^{-1}v^\lambda f',\qquad
 \alpha'(s)=sf'=(f')^{-1}s.                             \tag{D.26.4}
\]
Since $r=uv$, $s=uvu=vuv$, and $m=(\lambda-1)/2$, its
two hexagons are exactly
\[
 u^\lambda(f')^{-1}v^\lambda f'=(f')^{-1}r x^{-m},\qquad
 (f')^{-1}v^\lambda f'u^\lambda=vu y^{-m}f'.             \tag{D.26.5}
\]
Reduction to every $G_j$ gives a GT shadow: the center is
killed, the unit condition holds, $f'$ has commutator image,
and the induced map is an automorphism. The single global
pair makes these reductions compatible.

For precision concerning D.25's finite gauge, choose its
discrete representative of $\beta$ to agree with the
constructed profinite $\beta$ both in $G_6$ and in
$\overline F^{\rm ab}/36$. Such a representative exists
because $F$ is dense in $\overline F$ and this is a finite
joint quotient. Then its integer $b-a$ agrees with $k$
modulo36, the order of $x$ in $G_6$. Thus its D.25 gauge
is exactly the reduction of (D.26.4), with $m=0\pmod{36}$.
The gauge preserves $T=A^{-1}V$ and $W=\beta^{-1}V$.
Consequently all twelve extra CM pairs remain
nonarithmetical at level six in the compatible tower.
This specifies the finite gauges that extend; it does not
make an assertion about an independently chosen gauge.

Finally we locate the remaining finite tests by an actual
quotient. Choose a free basis of the rank97 group $N$ and
map its first two generators to $(12),(23)\in S_3$ and
the others to the identity. Let $J$ be the normal core in
$B$ of the intersection of this map's kernel with
$[N,N]N^6$. Then $J\triangleleft B$ has finite index,
$J\subset[N,N]N^6$, but $[N,N]\not\subset J$, since
the commutator of the first two basis elements has nontrivial
image in $S_3$. The quotient map $B\to B/J$ therefore
cannot factor through any $G_j$: every such factorization
would kill $[N,N]$. This also proves the kernel in
(D.26.2) is nontrivial and that the homology tower is not
cofinal in the system of D.22.

The source search supplies no additional lifting assertion
for these quotients. In particular isolation gives reduction
\emph{homomorphisms}, as stated in
\cite[Remark3.16]{dDolgushevGuynee2024}; the all-finer-quotient
criterion is its Theorem5.2 and Corollary5.4.
The abelian sufficient criterion in
\cite[TheoremB.2]{dDolgushevLeLorenz2024} requires the
pure quotient to be abelian. Here its image already
surjects onto
$H\simeq C_2^2\times\mathrm{SL}_2(\mathbb F_3)$:
the isomorphism sends the unique lift with modulo-four
image in $V_4$ to its modulo-four and modulo-three components.
The latter group's second derived subgroup has order two.
Similarly the dihedral quotients of
\cite[Remark3.2]{dBDHP2026} embed their $F$ images into
products of three dihedral groups and are metabelian.
No such pure quotient can surject onto this $H$, since
surjections preserve the derived-series images.
These explicit maps explain why those sufficient families
do not supply the remaining lift.
\end{proof}

\section{E. The next approximations and exactness on coverings}
\label{sec:E}

We write $F=\langle x,y\rangle$ for the discrete group, $F^{(i)}$
for its derived series, and $\gamma_iF$ for its lower central series,
with $F^{(0)}=\gamma_1F=F$. All commutators in this section are
$[a,b]=aba^{-1}b^{-1}$. The puncture used to define depth is the
puncture labelled $1$; it is additional data. The manuscript anchors
are \S51.16, folios 733v--736v, scans 172--178, especially (371)--(379).
The arithmetic model in those passages is
$X_n:x_0^n+x_1^n+x_\infty^n=0$, with $U_n$ obtained by deleting the
coordinate divisors. We retain that model over $\mathbb Q$.

\subsection{Derived quotients and the depth filtration}

\begin{statement}{E.1}{Proved}{The third derived quotient and its exact relation module}
Put $Q=F/F''$ and $R=\mathbb Z[Q]$. There is an exact sequence of
discrete groups and an exact sequence of left $R$-modules
\[
 1\longrightarrow F''/F'''\longrightarrow F/F'''
       \longrightarrow Q\longrightarrow1,
 \qquad
 0\longrightarrow F''/F'''\longrightarrow R^2
 \xrightarrow{d}R\xrightarrow{\epsilon}\mathbb Z\longrightarrow0,
\]
where $d(a,b)=a(\bar x-1)+b(\bar y-1)$, with right factors in this
formula and the regular left action on each term. This is an explicit
module description over the noncommutative ring $R$.
\end{statement}
\begin{proof}
The cover of the two-loop graph corresponding to the normal subgroup
$F''$ has vertices $Q$ and oriented edges $(q,x),(q,y)$ from $q$ to
$q\bar x,q\bar y$. Its cellular differential is the displayed map.
The graph is connected, so its reduced zeroth homology vanishes.
Its fundamental group is $F''$, and its first homology is
$F''/[F'',F'']=F''/F'''$. Lifting paths identifies conjugation by
$F$ with the left deck action of $Q$. These identifications give every
arrow, its kernel and its image. The group sequence is the quotient
sequence for the nested normal subgroups $F'''\subset F''$.

For the algebraic constructions below we use the rational relative
unipotent completion over a specified finite quotient, constructed in
E.14. For an unbased finite-dimensional representation problem one
can also use the affine proalgebraic group with coordinate Hopf algebra
the matrix coefficients of all finite-dimensional rational
representations of $F/F'''$. Evaluation is the canonical group map
into its rational points, and factoring matrix coefficients proves
its universal property. These constructions specify an envelope and
its comparison maps; they do not require a finite-index abelian
subgroup of $F/F'''$.
\end{proof}

\begin{statement}{E.2}{Proved}{The depth-two Lie algebra and its comparison with the derived quotient}
Let $L=\mathbb L_{\mathbb Q}(x,y)$, graded by total degree, and let
$D^rL$ be the span of Lie monomials with at least $r$ occurrences of
$y$. Put $v_i=\operatorname{ad}_x^i(y)$ and
$V=\bigoplus_{i\geq0}\mathbb Qv_i$. There is an isomorphism of graded
Lie algebras
\[
 L/D^3L\simeq \mathbb Qx\ltimes(V\oplus\Lambda^2V).
\]
Here $[v_i,v_j]=v_i\wedge v_j$, the bracket of $V$ with
$\Lambda^2V$ vanishes, and $[x,v_i]=v_{i+1}$; the action of $x$
on $\Lambda^2V$ is the induced derivation. In weight $n\geq2$ the
depth-two dimension is $\lfloor(n-1)/2\rfloor$.
The identity on $x,y$ induces the surjection
$L/L'''\twoheadrightarrow L/D^3L$ with kernel $D^3L/L'''$.
\end{statement}
\begin{proof}
We give the elimination argument. Jacobi moves every occurrence of
$x$ in a bracket in the ideal generated by $y$ onto a bracket entry.
Consequently that ideal is generated as a Lie algebra by the $v_i$.
In the tensor algebra, the leading words of the $v_i$ are $x^iy$.
The words obtained by concatenating these blocks have unique
factorizations at the occurrences of $y$. Thus the homomorphism
from the tensor algebra on the symbols $v_i$ into
$\mathbb Q\langle x,y\rangle$ is injective, by its triangular leading
words. Its primitive Lie subalgebra is the free Lie algebra on the
$v_i$. Keeping at most two such symbols yields $V\oplus\Lambda^2V$.
The derivation formula follows from Jacobi. A basis in weight $n$
is $v_i\wedge v_j$ with $i<j$ and $i+j=n-2$, which has the stated
cardinality.

Every element of $L'$ has depth at least one, hence
$L''\subset D^2L$ and $L'''\subset D^4L\subset D^3L$.
This proves the displayed quotient map and its kernel.
The element $[y,[y,[y,x]]]$ has weight four and belongs to $D^3L$;
its tensor expansion is nonzero, whereas $L'''$ has no weight-four
part. Thus the kernel is nonzero. Also
$[[x,y],[x,[x,y]]]$ survives in $\Lambda^2V$, so the identity on
$x,y$ does not descend from $L/L''$ to $L/D^3L$: its defining
kernel $L''$ contains this explicitly nonzero image.
Completion for the total-degree filtration gives the same maps of
pronilpotent Lie algebras, and exponentiation gives maps of rational
prounipotent groups. No integral exponential is used.
\end{proof}

\begin{statement}{E.3}{Known}{The rational Galois calculation in depth two through weight fourteen}
Normalize the depth-one generators by
$\bar\sigma_{2m+1}=\operatorname{ad}_x^{2m}(y)$.
For the positive-weight Lie algebra of the outer pro-$\ell$ Galois
image, tensored with $\mathbb Q_\ell$, the depth-one dimensions are
one in weights $3,5,7,9,11,13$ and zero in the other weights at most
fourteen. Its depth-two dimensions in weights $1,\ldots,14$ are
\[
 0,0,0,0,0,0,0,1,0,1,0,1,0,2.
\]
The first relation among the depth-one brackets is
\[
 \{\bar\sigma_3,\bar\sigma_9\}
       -3\{\bar\sigma_5,\bar\sigma_7\}=0
 \quad\text{in depth two}.
\]
This is a statement after tensoring with $\mathbb Q_\ell$ and taking
the specified filtrations; it asserts no integral saturation index.
\end{statement}
\begin{proof}
We use Hain--Matsumoto, Theorem 8.4, its application in \S8.4,
Lemma 8.5 and Remark 8.6 of the opened version
\cite{E_HainMatsumoto}, and Brown's depth-two bracket formula and
exact sequence (7.8), Example 7.2 \cite{E_BrownDepth}.
Here $U_1$ has $H^1_{\mathrm{et}}(U_{1,\overline{\mathbb Q}},
\mathbb Q_\ell(1))=\mathbb Q_\ell^2$ with trivial Galois action:
the two Kummer covers are those of $t$ and $1-t$.
Its outer pro-$\ell$ action is unramified away from $\ell$.
The indicated theorem therefore identifies the weight-graded Lie
algebra of its Zariski closure with the rationalized graded Galois
image. The weight-one generator maps to zero by Lemma 8.5;
the odd generators in weights at least three have nonzero
depth-one coefficient, as recorded in Remark 8.6. This supplies
actual Galois elements after rationalization, rather than only
solutions of a linear equation.

The polynomial form of the depth-two bracket is
\begin{align*}
 B_{a,b}(X,Y)={}&X^a\bigl(Y^b-(Y-X)^b\bigr)
 +(Y-X)^a\bigl(X^b-Y^b\bigr)\\
 &+Y^a\bigl((Y-X)^b-X^b\bigr),
 \qquad a,b\text{ positive and even}.
\end{align*}
One can derive this formula directly by substituting
$D_f(x)=0$, $D_f(y)=[y,f]$ into
$[f,g]+D_f(g)-D_g(f)$ in the tensor algebra, and reading the
two $y$-blocks from the right. The program \texttt{e\_depth.py}
does both expansions independently over the integers and compares
every coefficient. The only pairs in weights $8,10,12,14$ are,
respectively, $(3,5)$, $(3,7)$, $(3,9),(5,7)$, and
$(3,11),(5,9)$. Their exact ranks are $1,1,1,2$.
The other weights contain no pair.

For completeness of this span in depth two, every bracket word of
length at least three in the Galois generators has depth at least
three. A generator itself has a nonzero depth-one image in its
weight; different generator weights cannot cancel in a weight-graded
piece. Thus the generation theorem reduces the calculation exactly
to the displayed pairs. The polynomial
\[
 X^2Y^2(X-Y)^3(X+Y)^3
 =X^8Y^2-3X^6Y^4+3X^4Y^6-X^2Y^8
\]
satisfies the antisymmetry and the three-term period equation.
Substitution in $B_{a,b}$ gives the displayed relation, also checked
word by word in the free tensor algebra. Brown's exact sequence
identifies its line with the weight-twelve cusp-form period line.
This verifies the requested modular relation without making a
freeness assertion for an integral Galois Lie algebra.
\end{proof}

\begin{statement}{E.4}{Proved}{What the metabelian representation forgets about these brackets}
Under the canonical Lie homomorphism $L\to L/L''$, every Ihara
bracket of two elements of $L'$ maps to zero. The depth-two image
of $\{\bar\sigma_3,\bar\sigma_5\}$ is nonzero. Thus the specified
comparison maps exhibit information already added in weight eight;
the first dependence between these new bracket directions occurs
in weight twelve as in E.3.
\end{statement}
\begin{proof}
Write the derived module of $L/L''$ as
$\mathbb Q[X,Y]c$, where $c=[x,y]$ and $X,Y$ act by
$\operatorname{ad}_x,\operatorname{ad}_y$. For $f=pc$ one has
$D_f(c)=[x,[y,pc]]=XYpc$. The derivation $D_f$ vanishes on the
operators $X,Y$ when applied to the derived module: its correction
there is a bracket of two derived elements. Hence
$D_f(qc)=XYpq\,c=D_g(pc)$ for $g=qc$, while $[f,g]=0$.
The Ihara bracket therefore vanishes. In E.3 the polynomial
$B_{2,4}=2X^5Y-5X^4Y^2+5X^2Y^4-2XY^5$ is explicitly nonzero.
Both claims concern the images of the same derivations through the
two quotient maps just constructed.
\end{proof}

\subsection{Localizers and finite coverings}

\begin{statement}{E.5}{Proved}{Two-out-of-three and relative Theorem A}
For each of abelianization, the metabelian quotient, a fixed
nilpotent quotient $G/\gamma_{c+1}G$ with $c\geq1$, pro-$p$
completion, and rational prounipotent completion, the class
$W_{\mathfrak a}$ in the question satisfies two-out-of-three,
retract closure, the terminal-object axiom and relative Theorem A
for all small categories.
\end{statement}
\begin{proof}
We first specify the tests determining each class. They are the
classifying groupoids $BH$ for, respectively, all abelian groups,
all metabelian groups, all groups of nilpotence class at most $c$,
all finite $p$-groups, and all groups $U(\mathbb Q)$ of rational
points of finite-dimensional unipotent rational algebraic groups.
We include all discrete sets as tests to detect components.
For a category $C$, functors $C\to BH$ invert every arrow and
therefore factor uniquely through its fundamental groupoid.
For a connected pointed category the homotopy fiber of evaluation
$\operatorname{Fun}(C,BH)\to BH$ is equivalent to the discrete groupoid
$\operatorname{Hom}(\pi_1C,H)$. The universal properties of the five
approximations, followed by Yoneda, prove that a component-bijective
functor is in $W_{\mathfrak a}$ precisely when it induces an
equivalence of these functor groupoids for every indicated test.
Conjugations cause no loss here: the evaluation fiber restores the
pointed homomorphisms. Two-out-of-three and retract closure follow
from equivalences of groupoids. A category with a terminal object
has contractible nerve and passes every test.

Here is the relative argument. For $C\to S$, the Grothendieck
construction of $s\mapsto C/s$ has objects
$(s,c,u:f(c)\to s)$. Its forgetful functor to $C$ has a section
$c\mapsto(f(c),c,1)$ and a natural transformation from section
followed by forgetful to the identity, furnished by $u$.
It therefore has the same fundamental groupoid as $C$.
Functors from this construction to a groupoid $T$ consist exactly
of functors $C/s\to T$ with invertible coherent transition
transformations along the arrows of $S$. This is the pseudolimit
of the groupoids $\operatorname{Fun}(C/s,T)$. A pointwise
equivalence of such diagrams induces an equivalence of pseudolimits:
choose objectwise lifts, transport the transitions through full
faithfulness, and use faithfulness to verify their coherence;
the same argument applies to morphisms. A map $C\to D$ over $S$
which passes every test on every comma category consequently passes
the test on $C\to D$. This is the stated relative Theorem A.
\end{proof}

\begin{statement}{E.6}{Proved}{The exact comparison maps for quotient approximations}
Let $q:F\twoheadrightarrow G$ be any finite quotient, $N=\ker q$,
and $V(H)$ one of $H'$, $H''$, or $\gamma_{c+1}H$.
There is a natural short exact sequence
\[
 1\longrightarrow\frac{N\cap V(F)}{V(N)}
 \longrightarrow\frac{N}{V(N)}
 \longrightarrow
 \ker\left(\frac{F}{V(F)}\longrightarrow\frac{G}{V(G)}\right)
 \longrightarrow1.
\]
The comparison requested with target $G$ itself is the specialization
$V(G)=1$ of this exact sequence. The map $F/V(F)\to G$ with
composite $q$ exists exactly when $V(G)=1$.
\end{statement}
\begin{proof}
Surjectivity of $q$ gives $q(V(F))=V(G)$, by the definitions of
the verbal subgroups. The inverse image of $V(G)$ is $NV(F)$,
so the kernel in the display is $NV(F)/V(F)$. The map from $N$
onto it has kernel $N\cap V(F)$ and kills $V(N)$.
This proves exactness and identifies every map. A factorization
of $q$ through $F/V(F)$ kills $q(V(F))=V(G)$; the converse is
the same calculation. The formulas thus also give an explicit
replacement target for every quotient for which the arrow to $G$
written in the question does not descend.
\end{proof}

\begin{statement}{E.7}{Proved}{All nontrivial finite quotients have a nonzero comparison kernel}
For every finite quotient $F\twoheadrightarrow G$ of order $d>1$,
each of the three comparison kernels in E.6 is nonzero.
For abelianization it is a free abelian group of rank $d-1$;
the image is a rank-two lattice of index $|G^{\mathrm{ab}}|$ in
$F^{\mathrm{ab}}$. For the nilpotent quotient of class $c$ its
rational Hirsch rank is
\[
 \sum_{j=1}^c\bigl(w_j(d+1)-w_j(2)\bigr),\qquad
 w_j(r)=\frac1j\sum_{e\mid j}\mu(e)r^{j/e}.
\]
\end{statement}
\begin{proof}
The connected Cayley graph of $G$ has $d$ vertices and $2d$ edges.
A maximal tree gives $d+1$ free generators of $N$, hence
$N^{\mathrm{ab}}=\mathbb Z^{d+1}$. The abelian image and its index
follow from E.6. A subgroup of a free abelian group is free, and
the rank of the kernel is $(d+1)-2=d-1$.

For the nilpotent case the domain is the free nilpotent group on
$d+1$ generators and the image is a finite-index subgroup of the
free nilpotent group on two generators. Both are torsion-free.
Their lower central ranks are the Witt numbers: the tensor algebra
identity $1/(1-rt)=\prod_{j\geq1}(1-t^j)^{-w_j(r)}$, obtained by
ordering a homogeneous Lie basis in the universal enveloping
algebra, gives the displayed formula by logarithms and M\"obius
inversion. Hirsch rank is additive in a short exact sequence of
torsion-free nilpotent groups. The first summand is already $d-1$.

For the metabelian case choose
$a\in N\cap F'$ with nonzero image $A$ in $N^{\mathrm{ab}}$;
such an element exists by the abelian calculation. Choose $t\in N$
whose abelian class $T$ is linearly independent from $A$ over
$\mathbb Q$. In the free metabelian quotient of $N$ the module
$N'/N''$ injects into $\mathbb Z[N^{\mathrm{ab}}]^{d+1}$ by the
covering-graph differential. It is therefore torsion-free over this
integral domain. The image of $[t,a]$ is nonzero: its degree-two
term is $T\wedge A$. The element
\[
 w=[a,tat^{-1}]\in N\cap F''
\]
equals $[a,[t,a]]$ and has image $(A-1)[t,a]$ in that module
for our convention $[u,v]=uvu^{-1}v^{-1}$. This image is nonzero.
Thus $w\notin N''$, giving a specific construction of a nontrivial
element of $(N\cap F'')/N''$ for every such quotient.
For $d=1$, $N=F$ and every comparison is the identity.
\end{proof}

\begin{statement}{E.8}{Proved}{The pro-$p$ comparison and the precise exactness criterion}
Write $G^{(p)}$ for the maximal $p$-group quotient of the finite
group $G$. There is a continuous surjection
\[
 \widehat N_p\longrightarrow
 \ker(\widehat F_p\longrightarrow G^{(p)}).
\]
Its source is free pro-$p$ of rank $|G|+1$ and its target is free
pro-$p$ of rank $|G^{(p)}|+1$. It is an isomorphism exactly for
$p$-groups $G$. For $p$-groups the target is the open subgroup
corresponding to the given cover.
More explicitly, put $r=|G|+1$ and $s=|G^{(p)}|+1$.
There are free pro-$p$ bases $a_1,\ldots,a_s,b_1,\ldots,b_{r-s}$
in the source and $\bar a_1,\ldots,\bar a_s$ in the target
for which the comparison sends $a_i\mapsto\bar a_i$ and
$b_j\mapsto1$. Its kernel is exactly the closed normal subgroup
generated by those $r-s$ elements $b_j$. This description is
noncanonical; the original continuous comparison is canonical.
\end{statement}
\begin{proof}
The image is the closure of $N$. The quotient by that closure is
the maximal pro-$p$ quotient of the finite group $F/N$, namely
$G^{(p)}$. This proves surjectivity. For a normal subgroup of
$p$-power index the induced pro-$p$ topology is its full topology:
the core in the ambient group of a normal subgroup of finite
$p$-power index is the intersection of finitely many conjugates,
and its quotient is an extension of finite $p$-groups.
Completing therefore identifies $\widehat N_p$ with the open
subgroup when $G$ is a $p$-group. The Schreier rank calculation
applied to finite $p$-quotients, or the covering tree, gives the
two ranks. An isomorphism would identify their Frattini quotients,
whose dimensions are these ranks, and hence forces
$|G|=|G^{(p)}|$. This equality says precisely that $G$ is a
$p$-group. All arrows used here are continuous homomorphisms of
compact pro-$p$ groups.

For the kernel description, lift a target basis to elements
$a_1,\ldots,a_s$. Their classes in the source Frattini quotient
are independent. Extend them to a basis of that vector space and
choose lifts $b_j$. The pro-$p$ basis theorem and Hopficity show
that the resulting $r$ elements form a free basis. The subgroup
generated by the $a_i$ maps isomorphically onto the target:
the induced surjection of free groups of rank $s$ is an
isomorphism by Hopficity. For each $b_j$, choose the unique
word $w_j(a_1,\ldots,a_s)$ with its target image and replace
$b_j$ by $b_jw_j^{-1}$. The Frattini classes remain a basis,
so these are again free generators. The quotient by their closed
normal closure is now the free group on the $a_i$, proving the
stated kernel and image, not merely their ranks.
\end{proof}

\begin{statement}{E.9}{Proved}{Ordinary rational completion of a finite-index subgroup}
For $|G|=d$, the canonical morphism of rational prounipotent groups
$N^{\mathrm{un}}\to F^{\mathrm{un}}$ is surjective. Its map on
abelianized Lie algebras is a surjection
$\mathbb Q^{d+1}\to\mathbb Q^2$ with kernel of dimension $d-1$.
It is an isomorphism exactly for $d=1$.
In suitable noncanonical free pronilpotent coordinates, its Lie
algebra map sends $X,Y$ to the two target generators and
$z_1,\ldots,z_{d-1}$ to zero. Its kernel is the closed ideal
generated by the $z_j$; by free-Lie elimination it is the completed
free Lie algebra on
$\mathbb Q\langle X,Y\rangle\otimes
\langle z_1,\ldots,z_{d-1}\rangle$, where a word acts by iterated
adjoints. Exponentiation is the exact kernel group.
Every morphism from a connected unipotent rational algebraic group
to the constant finite group $G$ is trivial; thus the specified
map to $G$ in a kernel formula descends only for the trivial
quotient. The preceding surjection is the comparison obtained
from the trivial completion of $G$.
\end{statement}
\begin{proof}
A homomorphism from $F$ to a unipotent group has the same Zariski
closure on $N$ as on $F$: the quotient of the closures has finite
order, while a unipotent group in characteristic zero is connected
and admits no nontrivial finite quotient. Equivalently, every
$f\in F$ has a positive power in $N$, and taking roots in the
one-parameter unipotent subgroup containing its image recovers
that image. The resulting map of completions is surjective.
Both Lie algebras are completed free Lie algebras on their
abelianizations, of ranks $d+1$ and two. The abelian comparison
is the rationalization of E.7. A morphism from a connected scheme
to a constant finite scheme is constant, which proves the last
assertion and types the replacement map.

For the coordinate assertion, lift the target free Lie generators
and complete their independent linear terms to a basis of the
source abelianization. The homomorphism of completed free Lie
algebras determined by these lifts is invertible: its map in
each graded degree is invertible, and degreewise inversion
converges. Its first two generators give a section of the target
free Lie algebra. Subtract from every remaining generator the
section of its image. This triangular change is again invertible
and puts that image equal to zero. The kernel is consequently
the closed ideal generated by the remaining $z_j$.
Jacobi writes this ideal as the Lie algebra generated by the
iterated $X,Y$ adjoints of the $z_j$. In the free associative
algebra their leading words have unique cuts at the letters
$z_j$, proving freeness exactly as in E.2. This also computes
the whole kernel for every finite quotient, including those for
which the proposed map to the original $G$ has no descent.
\end{proof}

\begin{statement}{E.10}{Proved}{Which finite regular coverings preserve the localizers}
For a regular covering with monodromy $q:F\twoheadrightarrow G$,
membership in $\operatorname{Fib}_{W_{\mathfrak a}}$ holds only for
$G=1$ for abelianization, the metabelian quotient, every fixed
positive nilpotent class, and rational prounipotent completion.
For pro-$p$ completion it holds exactly for finite $p$-groups $G$.
These are exactly the cases of isomorphism in E.7--E.9, with the
properly descended target of E.6 or E.8.
\end{statement}
\begin{proof}
We first give uniform counterexamples over $BF$, so the universal
base-change requirement is tested, not merely the given fiber.
For a nontrivial finite $G$, choose a prime $r\nmid |G|$ and let
$V$ be the augmentation-zero summand of the regular
$\mathbb F_r[G]$-module. Averaging shows $V_G=0$, and $V\ne0$.
Let $A=V\rtimes F$ for the action through $q$, with retraction
$A\to F$. The inclusion $F\to A$ is an abelian and every
nilpotent equivalence: $[F,V]=V$ puts $V$ in every term
$\gamma_jA$, $j\geq2$. It is also a rational unipotent
equivalence, since the finite normal subgroup $V$ maps trivially
to every rational unipotent group. Pullback along the $G$-cover
replaces the inclusion by $N\to V\times N$. Its abelian and
nilpotent quotients retain $V$, so they fail those tests.
For the rational test use instead the integral augmentation ideal
$V=\ker(\mathbb Z[G]\to\mathbb Z)$. Every homomorphism
$A=V\rtimes F\to U(\mathbb Q)$, with $U$ a rational unipotent
group, kills $V$. To prove this, take the Zariski closure of the
image of $V$. It is an abelian unipotent group. Conjugation by
each image of an element of $F$ acts unipotently on its Lie
algebra and has finite order, since the action on $V$ factors
through $G$. A finite-order unipotent linear operator in
characteristic zero is the identity. Hence the image factors
through $V_G$, which is finite because its rationalization is
zero. A unipotent group in characteristic zero has no nontrivial
finite subgroup, so that image is trivial. The universal property
therefore identifies the rational prounipotent completions of
$F$ and $A$. After pullback, $N\to V\times N$ retains the
nonzero vector quotient $V\otimes\mathbb Q$, of dimension $d-1$.

For the metabelian test form a finite quotient
$H=F/(N'N^r)$, with $r$ any prime. The inclusion
$N\cap F'\hookrightarrow N$ has nonzero primitive abelian
image of rank $d-1$, by E.7. Its reduction modulo $r$ remains
nonzero. Thus $H'$ is nontrivial and the image of $N'$ is trivial.
Choose a prime $s\nmid |H|$ and put
$V=(1-e_{H'})\mathbb F_s[H]$, where
$e_{H'}=|H'|^{-1}\sum_{h\in H'}h$. This is a nonzero $H$-module
with $[H',V]=V$. For $A=V\rtimes F$ one has
$A''=V\rtimes F''$, so $F\to A$ induces an isomorphism of
metabelian quotients. On the pulled-back group $V\rtimes N$,
$N'$ acts trivially on $V$; its second derived subgroup is
$\{0\}\rtimes N''$. Consequently the metabelian quotient
retains $V$, and $N\to V\rtimes N$ fails the test. These are
explicit finitely presented group constructions. Their standard
finite presentation complexes and mapping cylinders give finite
simplicial-category models with the same fundamental-group
calculation.

For a finite $p$-group cover, E.8 applies to every group occurring
after any base change, not just $F$. An isomorphism of pro-$p$
completions identifies the finite orbit sets and the appropriate
open stabilizer subgroups, whose completions are those open
subgroups. Hence pullback preserves $W_p$.

For a regular cover with $G$ not a $p$-group, let $P$ be a Sylow
$p$-subgroup. The image of
$F\to G\times(F/F'F^p)$ is the fiber product over
$G/G'G^p$: this follows by taking the two projections and their
kernels, or directly from $q(F'F^p)=G'G^p$.
The map $P\to G/G'G^p$ is surjective, because its cokernel has
order both a power of $p$ and prime to $p$. Choose words $u,v$
mapping into $P$ whose images in $F/F'F^p$ are $x,y$.
The endomorphism $x\mapsto u,y\mapsto v$ induces a pro-$p$
automorphism: it is onto the Frattini quotient, hence onto the
finitely generated free pro-$p$ group, and that group is Hopfian.
Its pulled-back regular $G$-set has image contained in the proper
subgroup $P$, and is therefore disconnected. This violates the
component condition in $W_p$. Finally, a trivial covering is a
product with a discrete set and preserves all five classes.
\end{proof}

\begin{statement}{E.11}{Proved}{The explicit three-letter obstruction in every nilpotent approximation}
The endomorphism $\alpha(x)=x$, $\alpha(y)=y[x,y]$ induces the
identity on every graded lower central quotient and an
automorphism on every nilpotent quotient, every pro-$p$
completion and the rational prounipotent completion of $F$.
Pullback of the action $x\mapsto(12),y\mapsto(13)$ changes its
orbit set from one orbit to $\{1,2\},\{3\}$.
It consequently refutes intrinsic-fibration membership for this
cover for each of these approximations, and for every functor
of their inverse system which sends isomorphisms to isomorphisms.
\end{statement}
\begin{proof}
Substituting a generator times a commutator in an iterated
commutator changes it only in the next lower central degree.
Thus the graded action is the identity. Induction in the central
extensions gives an automorphism on each nilpotent quotient.
The Frattini argument of E.10 proves the pro-$p$ statement;
the completed free rational Lie algebra gives the unipotent
statement by degreewise inversion. The inverses commute with
the transition maps, so the map of inverse systems is an
isomorphism. For $a=(12),b=(13)$, multiplication gives
$b[a,b]=a$. Both pulled-back generators therefore act by
$(12)$. The program \texttt{e\_obstructions.py} verifies these
permutations and the orbit sets. The ordinary metabelian quotient
retains the map to $S_3$, since $S_3''=1$; the map of E.11 is
not substituted for a metabelian equivalence.
\end{proof}

\begin{statement}{E.12}{Proved}{Exhaustive small-quotient and Schreier verification}
For all ordered generating pairs in every group model of order
at most twenty-four, the program \texttt{e\_smallgroups.py}
constructs an integral cycle basis of $N^{\mathrm{ab}}$, the
comparison to $F^{\mathrm{ab}}$, and the deck character; it
computes the free-metabelian relation-module generic rank and
the first four lower central ranks. Every generated kernel has
rank $|G|+1$, the comparison lattice has index
$|G^{\mathrm{ab}}|$, and the rational deck character is
$1+\operatorname{reg}_G$. All checks pass.
\end{statement}
\begin{proof}
The input is the complete order-at-most-twenty-four list in the
opened primary GroupNames database \cite{E_GroupNames}.
It contains seventy-four groups. The accompanying input file
records the permutation generators and exact source URLs.
The program closes each set of generators, checks its order,
constructs its multiplication table, and enumerates every ordered
pair. It tests generated subgroup size without a presentation
heuristic. For every generating pair it constructs the $2|G|$
edges and a spanning tree, and checks that every fundamental
cycle has zero boundary. The non-tree edge coefficients form
an integral basis, so this proves the integral calculation.
The gcd of all two-by-two minors of the pushforward matrix
checks the Smith index. Deck traces are calculated for every
group element. Independent tests exhaust the images of a
generating pair in every other two-generated model of the same
order to verify pairwise nonisomorphism of the models used.

The lower central ranks are calculated from the proved Witt
formula in E.7. For the free group $N$ on $r=|G|+1$ generators,
the graph over $N^{\mathrm{ab}}$ gives
$N'/N''=\ker(R_N^r\to R_N)$, $R_N=\mathbb Z[N^{\mathrm{ab}}]$;
localization to its fraction field has dimension $r-1=|G|$.
Thus the printed free-metabelian invariant has an exact module
presentation behind it. The output is accompanied by
\texttt{e\_all\_surjections.csv}, with a row for every ordered
epimorphism, and integral cycle witnesses for every group.
The formulas E.6--E.10 prove the assertions for all finite
quotients, so the bounded computation also has a general proof.
\end{proof}

\begin{statement}{E.13}{Proved}{Specialization to the Fermat coverings}
For $N_n=\ker(F\to(\mathbb Z/n)^2)$ the abelian comparison is
\[
 0\longrightarrow\mathbb Z^{n^2-1}
 \longrightarrow N_n^{\mathrm{ab}}
 \longrightarrow n\mathbb Z^2\longrightarrow0.
\]
The metabelian and class-$c$ comparisons are the exact sequences
of E.6 with kernels $F''/N_n''$ and
$\gamma_{c+1}F/\gamma_{c+1}N_n$, respectively; both are nonzero
for $n>1$. Pro-$p$ completion is exact exactly when $n$ is a
power of $p$. With $n=p^a m$, $(m,p)=1$, its replacement target
is the open subgroup of $\widehat F_p$ of index $p^{2a}$, and
the source has rank $n^2+1$. The rational unipotent comparison
is the surjection of E.9, with abelian kernel dimension $n^2-1$.
\end{statement}
\begin{proof}
The abelian covering map sends a closed edge cycle to its two
exponent sums. Each sum is divisible by $n$, and the loops
$x^n,y^n$ realize the two generators of $n\mathbb Z^2$.
The source rank is $n^2+1$ by the graph calculation. Since
$N_n$ contains $F'$, it contains both $F''$ and every
$\gamma_{c+1}F$ for $c\geq1$, giving exactly the two kernels
in the statement. E.7 proves their nonvanishing. The maximal
$p$-quotient of $(\mathbb Z/n)^2$ is $(\mathbb Z/p^a)^2$,
so E.8 gives the asserted ranks and the exactness criterion.
E.9 gives the last assertion. These are the maps induced by
the actual covering inclusion, and coincide with the Fox
matrices computed in section B.
\end{proof}

\begin{statement}{E.14}{Proved}{Relative completion retains the entire covering kernel}
The rational relative unipotent completion $\mathcal R_q$ of
$F\twoheadrightarrow G$, for finite $G$, fits into
\[
 1\longrightarrow N^{\mathrm{un}}
 \longrightarrow\mathcal R_q\longrightarrow G_{\mathbb Q}
 \longrightarrow1.
\]
It splits over $\mathbb Q$, with splittings conjugate by
$N^{\mathrm{un}}(\mathbb Q)$. Its quotient with abelian kernel
has vector kernel $H_1(N,\mathbb Q)$, of dimension $|G|+1$ and
character $1+\operatorname{reg}_G$.
For $q_2:F\to(\mathbb Z/2)^2$ its four character
multiplicities are $2,1,1,1$.
\end{statement}
\begin{proof}
Choose a transversal of $G$ in $F$. Its conjugation maps on $N$
extend uniquely to $N^{\mathrm{un}}$, and its products give a
factor set with values in $N\to N^{\mathrm{un}}$.
The associated multiplication on the disjoint union of $|G|$
copies of $N^{\mathrm{un}}$ is associative by the factor-set
identities. This constructs $\mathcal R_q$ and the displayed
exact sequence as rational proalgebraic group schemes.
Every map from $F$ into an extension of $G$ by a unipotent group
extends uniquely on $N$ and then on the chosen components.
It respects multiplication by the same identities. This proves
the relative universal property.

At each successive central vector quotient of the unipotent
kernel, a section lifts because the cohomology in positive
degree of the finite group $G$ with rational coefficients
vanishes. To see the vanishing, average a cochain contraction
over $G$ and divide by $|G|$ in the standard bar resolution.
The same averaging in degree one conjugates two sections.
Induction gives compatible sections and conjugators through
the central quotients, and their inverse limits give the
assertions for $N^{\mathrm{un}}$. This specifies the conjugacy
ambiguity; uniqueness of the conjugating element is not claimed.
Finally, the chain modules of the finite Cayley graph are
$\mathbb Q[G]^2$ and $\mathbb Q[G]$, with zeroth homology
$\mathbb Q$. Semisimplicity of $\mathbb Q[G]$ gives
$H_1(N,\mathbb Q)\simeq\mathbb Q\oplus\mathbb Q[G]$.
For $G=(\mathbb Z/2)^2$, expanding its regular character gives
the four stated multiplicities. This action has a $-1$
eigenvalue; it survives in the relative completion, whereas
the map of E.9 takes coinvariants and loses these characters.
\end{proof}

\begin{statement}{E.15}{Proved}{The genuine action on a relative completion}
Let $\operatorname{Aut}(F,q)$ consist of the automorphisms of $F$
carrying $N$ to itself. Its homomorphism to
$\operatorname{Aut}_{\mathbb Q}(\mathcal R_q)$ is injective,
and its image is exactly the subgroup preserving the embedded
discrete subgroup $F\subset\mathcal R_q(\mathbb Q)$ setwise.
The automorphism $\operatorname{Int}([x,y])$, at $q=q_2$, acts
nontrivially on the relative completion, although its action
on $F^{\mathrm{ab}}$ and the image of its conjugator there are
both the identity.
\end{statement}
\begin{proof}
The universal property in E.14 extends an automorphism and its
inverse, with their induced automorphism of $G$, so the
extension is an automorphism. The map $N\to N^{\mathrm{un}}$
is injective. Write a nontrivial reduced word in maximal
nonzero-power blocks $z_{i_1}^{a_1}\cdots z_{i_k}^{a_k}$,
with adjacent $i_j$ different. Under $z_i\mapsto1+Z_i$,
the coefficient of $Z_{i_1}\cdots Z_{i_k}$ with these
$k$ alternating blocks is $a_1\cdots a_k\ne0$:
omitting a block cannot yield that sequence of $k$
blocks, and each retained block must contribute its
linear term. A finite truncation of the power-series
ring detects the word in its unipotent group $1+(Z_i)$.
Since
components detect $F/N$, the map $F\to\mathcal R_q(\mathbb Q)$
is injective as well. This proves faithfulness. An algebraic
automorphism preserving this subgroup restricts to an
automorphism of $F$ preserving its identity component $N$;
the universal property identifies its extension with the
given automorphism. This proves the image description.

At $q_2$ the commutator $w=[x,y]$ has a nonzero cycle class in
$H_1(N_2,\mathbb Q)$, as can be read from the unit-square
cycle in the four-vertex graph. Choose another cycle with
independent class. Their bracket is nonzero in the degree-two
free Lie algebra of $N_2^{\mathrm{un}}$, so conjugation by $w$
is nontrivial. Passing from a genuine action to its outer
class kills this particular inner action. Thus this example
also specifies the actual quotient map responsible for the
ambiguity, rather than identifying an outer action with a
genuine one.
\end{proof}

\subsection{Sections and arithmetic obstructions}

\begin{statement}{E.16}{Proved}{All sections of the rank-two toric quotient}
Let $T=\mathbb G_m^{\{0,1,\infty\}}/\mathbb G_m$ over
$\mathbb Q$, with its rational identity. Its arithmetic
fundamental-group extension has geometric kernel
$\widehat{\mathbb Z}(1)^2$. Its continuous sections modulo
conjugacy by this kernel form the abelian group
\[
 \left(\varprojlim_m\mathbb Q^*/\mathbb Q^{*m}\right)^2.
\]
One factor has the explicit description
\[
 \{\pm1\}\times
 \left\{(a_p)_p\in\prod_{p\text{ prime}}\widehat{\mathbb Z}:
 \forall m,\ a_p\in m\widehat{\mathbb Z}
 \text{ for all but finitely many }p\right\}.
\]
This computes the cohomology in (379) for its specified
rank-two quotient at every level $n$.
\end{statement}
\begin{proof}
The section at the identity gives the semidirect product
$\widehat{\mathbb Z}(1)^2\rtimes G_{\mathbb Q}$.
A section has form $g\mapsto(c(g),g)$; multiplication says
$c(gh)=c(g)+g c(h)$. Conjugation changes $c$ by a
coboundary. At finite level the Kummer sequence and
Hilbert 90 identify $H^1(G_{\mathbb Q},\mu_m)$ with
$\mathbb Q^*/\mathbb Q^{*m}$. To verify the reduction,
a multiplicative cocycle of a finite Galois extension
is a coboundary by choosing $b$ for which
$\sum_\tau c(\tau)\tau(b)\ne0$ and translating the sum;
linear independence of field embeddings supplies such
a $b$. Every algebraic cocycle is defined over a finite
extension. Taking limits gives the formula. The possible
$\varprojlim^1 H^0(\mu_m)$ vanishes because these groups
are finite and their descending images stabilize.

Unique prime factorization gives
$\mathbb Q^*=\{\pm1\}\times\bigoplus_p p^{\mathbb Z}$.
An element modulo $m$ has finitely many nonzero valuation
residues. Compatible residues are exactly the families
in the second display. Even $m$ retain the sign and odd
$m$ kill it; their compatible limit is $\{\pm1\}$.

The map $U_n\to T$ is inclusion on the open with all
coordinates nonzero. On geometric groups it is the
covering pushforward of E.13 divided by $n$, hence
surjective onto $\widehat{\mathbb Z}(1)^2$. The map to
the base torus in coordinates is
$(u,v)\mapsto(-u^n,-v^n)$, with
$t=-u^n,\ 1-t=-v^n$. It induces multiplication by $n$
on geometric kernels and retains the arithmetic
translation $(-1,-1)$. Thus the quotient extension
is the arithmetic fundamental group of $T$ even
when $U_n$ has no rational point. This supplies
the comparison underlying (373)--(379).
\end{proof}

\begin{statement}{E.17}{Refuted}{Pairwise separation of cusp sections by the minimum quotient}
For every odd $n$, the rational cusps
$P_0=[0:1:-1]$ and $P_\infty=[1:-1:0]$ admit rational
tangent choices whose sections become conjugate in
the rank-two quotient of (373). Both toric Kummer
classes are $(1,-1)$ in E.16. For the three punctures
of $U_1$ the complete families of toric images are
$(w,1)$, $(1,-w)$ and $(w^{-1},-w^{-1})$, with
$w\in\mathbb Q^*$, in the conventional coordinates
$(t,1-t)$ and the tangent coordinates specified below.
\end{statement}
\begin{proof}
Use $u=x_0/x_\infty,v=x_1/x_\infty$. Near $P_0$
take $s=u$; the branch is
$v=-(1+s^n)^{1/n}$ with constant coefficient $-1$.
The unit tangent has leading torus coefficients
$(1,-1)$. Near $P_\infty$ take $s=x_\infty/x_0$;
then $(u,v)=(s^{-1},-s^{-1}(1+s^n)^{1/n})$,
giving the same leading coefficients. Adjoining
compatible roots of these units computes the
tangential cocycle: coefficientwise Galois action
fixes the formal parameter and acts on roots of
the leading coefficients. Both classes therefore
coincide. This also derives the applicable case of
Wickelgren's tangential formula, Lemma 4
\cite{E_Wickelgren}.

For $U_1$ in coordinate $t$, at $0$ with tangent
$w$ the units $t,1-t$ give $(w,1)$; at $1$ with
coordinate $t-1$ they give $(1,-w)$; at infinity
with coordinate $1/t$ they give
$(w^{-1},-w^{-1})$. These exhaust rational
tangential choices. E.16 determines equality of
their images. The unit tangent at $0$ and tangent
$-1$ at $1$, for example, both give $(1,1)$.
The exhibited coincidences are recorded in
\texttt{e\_obstructions.py}. They concern precisely
this quotient, with the quotient map specified.
\end{proof}

\begin{statement}{E.18}{Proved}{The real obstruction for the manuscript quartic}
The geometrically abelianized arithmetic fundamental
group of $U_4/\mathbb Q$ has no continuous section.
The section-class sets are also empty for its
geometrically metabelian, full profinite and
geometric pro-$2$ quotients. The obstruction
already lies in a group of order two in
$H^2(C_2,H_1(U_4(\mathbb C),\mathbb Z_2))$.
\end{statement}
\begin{proof}
The plus quartic has no real points, by positivity
of fourth powers. Complex conjugation acts freely
on $U_4(\mathbb C)$. Its quotient is a connected
nonorientable punctured surface of Euler
characteristic $-8$, since the double cover has
Euler characteristic $-16$. It retracts to a
graph, so its fundamental group $J$ is free of
rank nine. The geometric subgroup $N$ of index
two is free of rank seventeen. Finite complex
topological covers with conjugation descent
data identify the real arithmetic extension
with the profinite completion of
$1\to N\to J\to C_2\to1$: descend the complex
cover with its antilinear involution; in the
other direction analytify a finite etale cover.
This is the equivariant finite-cover comparison,
and finite-index completion preserves this
extension.

Free-generator changes put $J\to C_2$ in the
form $a\mapsto\bar1$, $b_i\mapsto\bar0$
for $1\leq i\leq8$. A Schreier basis of $N$ is
\[
 a^2,\quad b_i,\quad ab_i a^{-1}
 \quad(1\leq i\leq8).
\]
Conjugation by $a$ fixes the first abelian basis
vector and exchanges each remaining pair. The
factor set at $(\bar1,\bar1)$ is $a^2$.
Replacing the lift by $na$ changes its square
by $(1+\operatorname{Ad}a)[n]$. Its first
coordinate is $1+2n_0$, which cannot vanish
in $\mathbb Z_2$. Each exchanged pair has
invariant diagonal equal to the image of
the norm. Hence the indicated cohomology
group is $\mathbb Z/2$, and the extension
is its nonzero class. The finite program
tests the factor-set equation for every
modulus dividing twenty-four and for
additional odd powers; precisely the
even moduli are obstructed.

A rational section restricts to a real
decomposition group, and each quotient
in the statement projects to this 2-adic
abelian quotient. All are therefore
obstructed. Its image under
$N^{\mathrm{ab}}\to4\widehat{\mathbb Z}(1)^2$
is zero in $H^2$: conjugation acts as
$-1$ on that torsion-free rank-two module,
whose invariants vanish. This explicit
pushout explains the splitting of (373)
together with its failure to lift.
The odd-adic and rational pushouts of
the first-coordinate obstruction also
vanish because two is invertible there;
this cohomology calculation makes no
assertion of a global section there.
\end{proof}

\begin{statement}{E.19}{Proved}{The cubic cusps are separated by their own 3-adic abelianization}
The three rational cusps $[0:1:-1]$,
$[1:0:-1]$, $[1:-1:0]$ determine pairwise
nonconjugate sections of the geometric
pro-$3$ abelianization of
$\pi_1(U_3/\mathbb Q)$, for every rational
tangent choice. They remain distinct in
the geometric abelian, metabelian and
full profinite quotients.
\end{statement}
\begin{proof}
Filling the cusps maps the affine
fundamental group to the Tate module of
the elliptic curve $X_3$, based at one
cusp. A cusp section maps to the point
section there, independently of its
tangent. Each difference of two distinct
cusps has exact order three: the ratio
of their tangent lines has divisor
three times their difference, and a
genus-one curve has no function with a
single simple pole. Section C gives
the same exact-order calculation.

The cubic has good reduction at two.
Its equation over $\mathbb F_2$ has
exactly the three projective points
with two nonzero coordinates, as the
program enumerates. The kernel of
good reduction over $\mathbb Q_2$ is
pro-$2$, from its formal-parameter
congruence filtration. Odd torsion
therefore injects into this group of
three points. A rational point of
exact order three cannot be divisible
by three: a preimage would have order
nine. Its Kummer class modulo three
is nonzero. The difference of the
two point sections is this class
in $H^1(G_{\mathbb Q},X_3[3])$.
Projection from affine 3-adic homology
detects it and proves the assertion.
This map also identifies information
lost in the quotient of E.17.
\end{proof}

\begin{statement}{E.20}{Proved}{A new obstruction to lifting a toric Kummer section}
For $U_1$ based at the tangent from $0$
towards $1$, the abelian section with
Kummer coordinates $(2,3)$ does not
lift to the class-two quotient with
geometric group
\[
 \langle x,y,z\mid x^2=y^2=z^2=1,
 [x,y]=z,\ [x,z]=[y,z]=1\rangle.
\]
Its obstruction is the nonzero class
$\{2\}\cup\{3\}$ in
$H^2(G_{\mathbb Q},\mu_2)$.
\end{statement}
\begin{proof}
The eight normal forms $x^a y^b z^c$
have central multiplication term
$b a'$. The code verifies this group
law and presentation. The tangential
Galois action is trivial on it:
$x\mapsto x^\chi=x$ and
$y\mapsto f^{-1}y^\chi f=y$, since
$\chi$ is odd and the commutator
$f$ is central here. Pulling the
central factor set back along the
two Kummer characters gives their
cup product (the order of factors
does not change its class modulo
two).

Over $\mathbb Q_3$, the quadratic
extension generated by $\sqrt2$ is
unramified because $2$ is a
nonsquare modulo three. Norms have
even valuation, so $3$ is not a
norm. The crossed-product
description of the cup product
identifies its vanishing with
the norm equation
$3=u^2-2v^2$ over $\mathbb Q_3$:
represent $\sqrt2$ by its
multiplication matrix, and the
matrix equation for the second
quaternion generator is exactly
this norm equation. Thus the
local and the global class are
nonzero. A lifted section would
give a cochain whose coboundary
is this factor set, a contradiction.
\end{proof}

\begin{statement}{E.21}{Refuted}{The unmarked depth assignment fails to descend}
The quotient of $\mathbb L(x,y)$
by $D_y^3$, with its canonical
projection, is not functorial for
all automorphisms of the unmarked
free group. It is functorial for
maps preserving the ideal of the
marked puncture and its lower
central filtration. Thus E.5
tests the five specified ordinary
group approximations; the depth
construction additionally uses
this marking and these morphisms.
\end{statement}
\begin{proof}
The swap of $x,y$ sends
$[y,[y,[y,x]]]\in D_y^3$ to
$[x,[x,[x,y]]]$, a nonzero
depth-one element of E.2.
It therefore fails to preserve
the kernel of the canonical
projection, which proves the
non-descent. A map preserving
the marked ideal preserves its
lower central terms and does
induce the quotient map.
No LM8 return specifying an
additional unmarked functor is
present in the supplied archive.
\end{proof}

\begin{statement}{E.22}{Refuted}{The scheme coordinates do not carry the requested relative action}
Put $q_2:F\to(\mathbb Z/2)^2$ and extend the relative completion
of E.14 to $\mathbb Q_\ell$. There is no homomorphism from the
normalized diagonal arithmetic image in
$\mathcal M(0)(\widehat{\mathbb Z})$ to
$\operatorname{Out}(\mathcal R_{q_2,\mathbb Q_\ell})$
whose composite with $G_{\mathbb Q}$ is its actual outer action.
In particular no such genuine action descends from those coordinates.
The actual action is defined on the continuous geometric
fundamental group, and its relative universal property gives its
action on $\mathcal R_{q_2,\mathbb Q_\ell}$. E.15 gives the
faithful discrete version and its exact image.
\end{statement}
\begin{proof}
First we type the comparison with the arithmetic action. The continuous
unipotent completion of the free profinite group $\widehat N_2$
over $\mathbb Q_\ell$ is the free prounipotent group on its five
Schreier generators. Indeed each finite tuple in a unipotent
$\mathbb Q_\ell$-group generates a relatively compact subgroup,
as follows from upper-triangular coordinates and their polynomial
group law. Its homomorphism from the discrete free group consequently
extends continuously to the free profinite group. The two universal
properties identify this completion with
$N_2^{\rm un}\otimes\mathbb Q_\ell$. The transversal construction
in E.14 then gives the continuous relative completion and its
functorial arithmetic action. The normalized Galois action is the
identity on $(\mathbb Z/2)^2$, because its abelian multiplier is
$\chi\equiv1\pmod2$.

The homomorphism $F\to G\times F^{\rm un}$, $g\mapsto(q_2(g),g)$,
and the relative universal property give the canonical surjection
\[
 \mathcal R_{q_2,\mathbb Q_\ell}\longrightarrow
 F^{\rm un}_{\mathbb Q_\ell}.
\]
It is the maximal unipotent quotient, so its kernel is characteristic.
Every genuine automorphism descends through this map and inner
automorphisms descend to inner automorphisms. This gives the needed
map of outer automorphism groups. A nontrivial outer action on its
target is therefore a nontrivial outer action on its source.

A.19 proves that the normalized arithmetic image is contained in
$\mathcal M(0)\simeq\Gamma$. On $G_{\mathbb Q^{\rm cyc}}$ its
multiplier-scheme coordinate is the identity, and the exact law is
that of $C_2\times\widehat{\mathbb Z}$, an abelian group.
Every commutator of two elements of $G_{\mathbb Q^{\rm cyc}}$
thus maps to the identity in these scheme coordinates.

We now exhibit a commutator whose actual outer unipotent action is
nontrivial. E.3 gives nonzero actual graded Galois classes in
weights three and five, with depth-one coefficients proportional
to $\bar\sigma_3,\bar\sigma_5$. Choose representatives $g_3,g_5$
of nonzero such classes and choose $\tau\in G_{\mathbb Q}$ with
$\chi_\ell(\tau)=1+\ell$. The cyclotomic action on weight $w$
is multiplication by $\chi_\ell^w$. Hence
$h_w=[\tau,g_w]$, for $w=3,5$, has leading class
$((1+\ell)^w-1)$ times that of $g_w$, which is nonzero over
$\mathbb Q_\ell$. Both $h_w$ lie in
$G_{\mathbb Q^{\rm ab}}=G_{\mathbb Q^{\rm cyc}}$, since they
are commutators. Their commutator has nonzero weight-eight,
depth-two outer class: its polynomial is a nonzero scalar multiple
of $B_{2,4}$ computed in E.3. This assertion uses actual graded
Galois classes, not arbitrary special derivations.

Consequently $[h_3,h_5]$ is killed by $\mathcal M(0)$ and acts
nontrivially on $\operatorname{Out}(F^{\rm un}_{\mathbb Q_\ell})$,
and thus on the relative completion. A homomorphism agreeing with
both actions would send the identity to this nonidentity element,
which proves the claimed non-descent. For a separate genuine
geometric witness, $w=[x,y]\in\Pi_0'$ has trivial final coordinates
$(a,b,U)=(1,1,\bar w)$ but its inner action on the relative
completion is nontrivial by E.15. The arithmetic witness above
also survives passage to outer automorphisms.
\end{proof}

\begin{statement}{E.23}{Open}{The remaining cubic section calculation}
The complete set of continuous section classes for the geometric
metabelian and marked depth-two arithmetic quotients of
$U_3/\mathbb Q$ is not determined in this return. The remaining
calculation is the image and every fiber of their section maps
to E.16. The first unresolved concrete nonabelian case here is
the class-two pro-$3$ quotient of the cubic's rank-ten geometric
fundamental group, with its integral arithmetic action.
\end{statement}
\noindent\emph{Work and obstruction.}
E.18 determines the empty set for the plus quartic and all
quotients retaining its $2$-adic geometric abelianization;
E.25 extends this to every even plus-Fermat level. E.24 computes
the entire cubic abelian section set and every fiber over the
minimum torus. E.19 separates the three rational cubic cusps
inside that set, while E.17 proves their collision under the
stated toric map.

The relation-module presentation E.1, the rank-two factor-set
calculation, local cup products, the exact generalized-Jacobian
extension, and a new three-punctured-line map on the cubic were
successive attempts to compute the nonabelian fibers. They give
the following actual refinements. E.28 exhibits a whole
nonempty abelian toric fiber with no class-two lifts. E.29
classifies every rational toric pair for the inherited cubic
class-two quotient. E.32 completes that calculation for every
completed pro-$3$ Kummer pair, including the exact convergence
condition and all three classes in each nonempty fiber; it also
exhibits an actual obstruction invisible at every finite level.
E.31 computes a genuine intrinsic quotient with its
arithmetic origin term: in one toric fiber the zero abelian
class lifts, whereas an explicit other class and its entire
congruence coset have no intrinsic class-two lift. Thus the
two minimum coordinates do not determine intrinsic lifting,
as witnessed by this exact same-fiber comparison.
E.34 computes the entire intrinsic abelian family unramified
outside three and proves that every one of its classes lifts
to the full geometric pro-$3$ group.
E.35 evaluates every section of the intrinsic finite
Heisenberg quotient by a cubic norm equation and gives all
of its nonempty fibers in explicit Kummer coordinates.
E.36 proves that every abelian class lifts through the
cuspidally central quotient retaining all ten abelian
coordinates; each fiber is the specified completed group
of rational units. Its quotient map from the full
intrinsic class-two group has a central kernel of rank $44$.

These calculations use the quotient homomorphisms specified in
E.25. They do not supply every obstruction and every fiber for
the full rank-ten intrinsic class-two group in the statement.
After E.36 the next calculation is the actual obstruction
in that rank-$44$ kernel, including its integral lattice
and continuous inverse-limit conditions.
Wickelgren's opened formulas apply directly to the ambient
rank-two group; the additional intrinsic quotient in E.31
requires its separately computed arithmetic action. E.22
settles the proposed descent from the manuscript scheme by an
actual Galois counterexample. The next unresolved step is to
evaluate the full intrinsic central extension on every class
of the explicitly computed set E.24, retaining the origin
terms and the geometric directions killed by E.25's quotient
map. A name for that nonabelian cohomology is not its evaluation.

\IfFileExists{sections/E_cubic.tex}{\subsection{The complete cubic abelian calculation and the precise depth maps}

In this subsection $E=X_3^+$ is the cubic
$x^3+y^3+z^3=0$, $D=E\setminus U_3$, and
$G=J_{E,D}$ is the generalized Jacobian for the reduced divisor $D$.
We use the origin $O=[1:-1:0]$, and set
$P=[0:-1:1]$, $Q=[-1:0:1]$.
For an abelian group $A$ we write
$A^\wedge=\varprojlim_m A/mA$, with $m$ ordered by divisibility.
This is the completion for multiplication by $m$; the quotients
need not be finite. All inverse limits of cohomology groups below
carry their inverse-limit prodiscrete topology.
For a countable set $\mathcal I$ define
\[
 \mathcal R_{\mathcal I}(\widehat{\mathbb Z})
 =\left\{(a_j)\in\prod_{j\in\mathcal I}\widehat{\mathbb Z}:
       \text{for every }m,\ a_j\in m\widehat{\mathbb Z}
       \text{ for all but finitely many }j\right\}.
\]
Its topology is the one coming from
$\varprojlim_m\bigoplus_{\mathcal I}\mathbb Z/m$.
In particular this space is not asserted to be compact.

\begin{statement}{E.24}{Proved}{All cubic abelian sections and every fiber over the toric minimum}
Put $K=\mathbb Q(\zeta)$, $\zeta^2+\zeta+1=0$,
and $\lambda=1-\zeta$. There is a retraction of algebraic groups
over $\mathbb Q$
\[
 r:G\longrightarrow\mathbb G_m^2,\qquad
 G\simeq\mathbb G_m^2\times G',
 \qquad
 1\longrightarrow(\operatorname{Res}_{K/\mathbb Q}\mathbb G_m)^3
 \longrightarrow G'\longrightarrow E\longrightarrow0.
\]
The map $r$ is the map induced by the two units $x/z,y/z$ on $U_3$.
The group $E(\mathbb Q)$ is cyclic of order three, and
\[
 G'(\mathbb Q)\simeq\mu_6^3\oplus
                  \bigoplus_{j\in\mathcal I}\mathbb Z .
\]
Here one may take $\mathcal I$ to consist of three copies of the
nonzero prime ideals of $\mathbb Z[\zeta]$, with the first copy of
$(\lambda)$ replaced by one new generator. The displayed abstract
isomorphism includes explicit choices of prime generators.

Choose the unit rational tangent at $O$, with local parameter $z/x$,
as the origin of the set of sections of the geometrically abelian
arithmetic fundamental group. Up to geometric conjugation, that
entire set is then the group
\[
 H^1(G_{\mathbb Q},T_\bullet G)
 \simeq
 (\mathbb Q^{\times,\wedge})^2
 \times\mu_6^3
 \times\mathcal R_{\mathcal I}(\widehat{\mathbb Z})
 \times\widehat{\mathbb Z}.
 \tag{E24.1}
\]
The last factor and its splitting are noncanonical. The actual map
to the toric minimum of E.16 is the projection onto the first two
factors followed by translation by the Kummer class of $(1,-1)$.
Consequently it is surjective, and every fiber is a torsor under
\[
 \mu_6^3\times
 \mathcal R_{\mathcal I}(\widehat{\mathbb Z})\times\widehat{\mathbb Z}.
 \tag{E24.2}
\]
For the pro-$3$ geometric abelianization the same formulas hold
with $\mathbb Q^{\times,\wedge}$ replaced by
$\varprojlim_a\mathbb Q^\times/\mathbb Q^{\times3^a}$,
$\mu_6^3$ replaced by $\mu_3^3$, and each
$\widehat{\mathbb Z}$ replaced by $\mathbb Z_3$.
The translation by $-1$ is zero in this pro-$3$ version.
\end{statement}
\begin{proof}
\emph{The rational elliptic-curve calculation.}
The maps
\[
 X=-\frac{3z}{x+y},\qquad Y=\frac{4x-5y}{x+y},
 \qquad
 [x:y:z]=[Y+5:4-Y:-3X]
\]
identify the smooth cubic with $Y^2+Y=X^3-7$, with $O$ at infinity.
Indeed the exact identity is
\[
 Y^2+Y-X^3+7
 =\frac{27(x^3+y^3+z^3)}{(x+y)^3}.
\]
The inverse formula proves the birational assertion, which extends
to an isomorphism of the smooth projective curves.

For the rank calculation we use the explicit identification over
$\mathbb Q$ of the Fermat cubic with $X_0(27)$ in
Li--Long--Tu, \S2.4 and Remark 4.2
\cite{EC_LiLongTu}. The corresponding normalized differential has
newform
$f(\tau)=\eta(3\tau)^2\eta(9\tau)^2$.
The Mellin formula therefore gives
\[
 L(E,1)=2\pi\int_0^\infty
           \eta(3it)^2\eta(9it)^2\,dt>0 .
\]
The product formula for $\eta$ makes the integrand strictly
positive for $t>0$, and cuspidality gives convergence at both ends.
The automorphism induced by multiplication of one Fermat coordinate
by $\zeta$ gives complex multiplication by
$\mathbb Z[\zeta]$; this is the maximal order of the norm-Euclidean
field $K$. Coates--Wiles, Theorem 1 on p.~223
\cite{EC_CoatesWiles}, now gives rank zero.
The smooth cubic has exactly three points over $\mathbb F_2$,
so the prime-to-$2$ rational torsion injects into a group of order
three. An order-two point on the displayed Weierstrass model would
have $Y=-1/2$ and $X^3=27/4$, which has no rational solution.
Thus there is no $2$-primary torsion. The distinct points $O,P,Q$
already give three rational points, so $E(\mathbb Q)=\langle P\rangle
\simeq\mathbb Z/3$.

\emph{The retraction and the exact rational-point presentation.}
The cusp scheme is the disjoint union of three rational points and
three copies of $\operatorname{Spec}K$. Taking the frame at $O$ as
one identifies the torus of $G$ with
\[
 T_D=\mathbb G_m^2
       \times(\operatorname{Res}_{K/\mathbb Q}\mathbb G_m)^3.
\]
We write its coordinates as $(a,b,A,B,C)$, at $P,Q$ and the three
quadratic cusp orbits, respectively. The convention is that a frame
is a nonzero vector in the fiber of the line bundle.

Here is an explicit construction of the characters used in $r$.
For a unit $h$ on $U_3$, write
$\operatorname{div}(h)=\sum_{d\in D}m_d d$.
For a framed line bundle $(L,(e_d))$, choose a rational section
$s$ nonvanishing along $D$, and set
\[
 \chi_h(L,(e_d))
 =h(\operatorname{div}s)\,
        \prod_{d\in D}(s(d)/e_d)^{-m_d}.
 \tag{E24.3}
\]
Products over nonrational points mean the corresponding norms.
Replacing $s$ by $fs$ leaves this expression unchanged, since
Weil reciprocity gives
$h(\operatorname{div}f)=f(\operatorname{div}h)$.
Tensor products multiply (E24.3). The construction in families, or
the divisor presentation of the generalized Jacobian, consequently
gives a morphism of algebraic groups. On the torus it is
$\prod e_d^{m_d}$, and on a divisor $u-u_0$ in $U_3$ it is
$h(u)/h(u_0)$. Thus it is precisely the character induced by $h$.

For $h=x/z,y/z$ their divisors are the respective coordinate-cusp
orbits minus the orbit at infinity. Hence
\[
 r_T(a,b,A,B,C)
  =\bigl(aN(A)/N(C),\ bN(B)/N(C)\bigr).
 \tag{E24.4}
\]
The map $i(a,b)=(a,b,1,1,1)$ is an algebraic section of $r$.
The isomorphism with the direct product is explicitly
$g\mapsto(r(g),\,g\,i(r(g))^{-1})$.
The kernel torus is parametrized by
\[
 (A,B,C)\longmapsto
 \bigl(N(C)/N(A),\,N(C)/N(B),\,A,B,C\bigr).
 \tag{E24.5}
\]
Hilbert 90 gives $H^1(\mathbb Q,T_D)=0$ and
$H^1(\mathbb Q,(\operatorname{Res}_{K/\mathbb Q}\mathbb G_m)^3)=0$.

To calculate the extension on rational points, let $z_0\in G(\mathbb Q)$
be the framed line bundle $\mathcal O(P-O)$ with the local frames
specified next. Put
$f=(y+z)/(x+y)$, so
$\operatorname{div}(f)=3P-3O$.
At $P$ use $s=x/z$ and the frame $s^{-1}$ for
$\mathcal O(P-O)$; at $O$ use $s=z/x$ and the frame $s$.
At the other cusps use the canonical section $1$.
Multiplication by $f$ trivializes its third tensor power.
The two leading values of $f$ are
\[
 \lim_{P}f/s^3=1/3,\qquad
 \lim_{O}fs^3=3,\qquad f(Q)=-1 .
\]
For example $y/z=-(1+s^3)^{1/3}$ at $P$, which proves the
first equality; the same expansion at $O$ proves the second.
At the nonrational cusps use the representatives
$[0:-\zeta:1]$, $[-\zeta:0:1]$, $[1:-\zeta:0]$.
Dividing every frame value by the value $3$ at $O$ gives
\[
 3z_0=
 \left(3^{-2},-3^{-1},
             \lambda^{-1},\lambda^{-2},\zeta^2\lambda^{-3}\right)
       \in\mathbb Q^{\times2}\times K^{\times3}.
 \tag{E24.6}
\]
This calculation also follows directly from
$3=-\zeta^2\lambda^2$.
Equation (E24.4), with $N(\lambda)=3$, sends (E24.6) to $(1,-1)$.
The rational cube roots of $1$ and $-1$ are unique, so
$r(z_0)=(1,-1)$.
For $z'=z_0-i(1,-1)$ the presentation becomes
\[
 G'(\mathbb Q)=
 (K^{\times3}\oplus\mathbb Z z')/
 \left\langle 3z'-
   (\lambda^{-1},\lambda^{-2},\zeta^2\lambda^{-3})\right\rangle .
 \tag{E24.7}
\]
There are no additional relations: the quotient onto
$E(\mathbb Q)=\mathbb Z/3$ is surjective by Hilbert 90, the class of
$z'$ generates it, and (E24.7) specifies its third multiple in the
kernel.

The norm-Euclidean ring $\mathbb Z[\zeta]$ has units $\mu_6$,
so prime factorization gives
$K^\times=\mu_6\oplus\bigoplus_{\mathfrak p}\mathbb Z$ after choosing
one generator for each prime ideal. In (E24.7) the coefficient of
the first $\lambda$ generator is $-1$. Eliminate that generator
and retain $z'$ instead. This operation changes no torsion relation
and proves the asserted description of $G'(\mathbb Q)$.
As an independent finite-support check,
\texttt{e\_cubic\_abelian.py} obtains the Smith diagonal
$1,6,6,6$ for the relation involving the three $\mu_6$ generators,
the three $\lambda$ generators and $z'$.

\emph{The full cohomology calculation.}
Write $B=(\operatorname{Res}_{K/\mathbb Q}\mathbb G_m)^3$.
The extension of $E$ by $B$ in $G'$ is killed by three: the
extension of $E$ by $T_D$ is classified by the cusp-difference
map, all of whose values are killed by three by C.5, and (E24.5)
and the direct product exhibit the component defining $G'$.
It follows that the connecting map
\[
 \delta:H^1(\mathbb Q,E)\longrightarrow H^2(\mathbb Q,B)
\]
is killed by three. Since $H^1(\mathbb Q,B)=0$,
$H^1(\mathbb Q,G')$ injects into $H^1(\mathbb Q,E)$.
Every element of the maximal divisible subgroup of the latter
is three times an element of that subgroup; hence its image under
$\delta$ is zero. Consequently the injection identifies maximal
divisible subgroups:
\[
 \operatorname{Div}H^1(\mathbb Q,G')
   =\operatorname{Div}H^1(\mathbb Q,E)\simeq\mathbb Q/\mathbb Z .
 \tag{E24.8}
\]
The last equality is exactly
\c{C}iperiani--Stix, Theorem 18(2), printed p.~18
\cite{EC_CiperianiStix}, applied to the analytic-rank-zero curve
already established above. Their additional restrictions on primes
concern the separate assertion about the Tate--Shafarevich
intersection. They impose no restriction on (E24.8).

The Kummer sequences and inverse limit give
\[
 0\longrightarrow G'(\mathbb Q)^\wedge
 \longrightarrow H^1(G_{\mathbb Q},T_\bullet G')
 \longrightarrow
 \operatorname{Hom}(\mathbb Q/\mathbb Z,H^1(\mathbb Q,G'))
 \longrightarrow0 .
 \tag{E24.9}
\]
For precision, take the countable cofinal subsystem $m=a!$.
The transitions $G'(\mathbb Q)/m\to G'(\mathbb Q)/d$, $d\mid m$,
are surjective, so their first derived inverse limit vanishes.
The finite groups $H^0(\mathbb Q,G'[m])$ satisfy the
Mittag--Leffler condition, which identifies continuous $H^1$ of
$T_\bullet G'$ with the inverse limit of the finite-coefficient
$H^1$ groups. No finiteness of
$G'(\mathbb Q)/m$ or $H^1(\mathbb Q,G'[m])$ is used.
The last term in (E24.9) is the limit of
$H^1(\mathbb Q,G')[m]$ under multiplication transitions.
Every homomorphic image of $\mathbb Q/\mathbb Z$ is divisible,
so (E24.8) identifies it with
$\operatorname{End}(\mathbb Q/\mathbb Z)=\widehat{\mathbb Z}$.

Choose a lift of $1$ in (E24.9). Multiplication by
$\widehat{\mathbb Z}$ then gives a continuous splitting: modulo
$m$ its value depends only on the scalar modulo $m$.
The description (E24.7) gives
\[
 G'(\mathbb Q)^\wedge
  =\mu_6^3\times\mathcal R_{\mathcal I}(\widehat{\mathbb Z}).
 \tag{E24.10}
\]
Indeed the image at each $m$ has finite support in $\mathcal I$,
and compatibility is exactly the restricted-sequence condition
in the definition. This proves both the algebraic and
topological assertions in (E24.9).
The identical argument with the subsystem $m=3^a$ gives the
pro-$3$ version, whose divisible quotient is
$\operatorname{Hom}(\mathbb Q_3/\mathbb Z_3,
 \mathbb Q_3/\mathbb Z_3)=\mathbb Z_3$.

Finally, $T_\bullet G=H_1(U_{3,\overline{\mathbb Q}},
\widehat{\mathbb Z})$ by the generalized-Jacobian realization in
C.1. A rational tangential section splits the corresponding
arithmetic extension; its conjugacy classes of splittings are
exactly continuous one-cocycles modulo coboundaries.
Thus (E24.9), (E24.10), Hilbert 90 and the algebraic retraction
give (E24.1). At $O$, with $s=z/x$, the leading terms of
$(x/z,y/z)$ are $(s^{-1},-s^{-1})$. This makes the toric origin
the Kummer class of $(1,-1)$ and proves the claimed affine map.
The covering pushforward $q_{3*}$ is multiplication by three
after $r_*$ on geometric homology; the map used here is the
actual coordinate-unit map $r_*$, as in E.16. This fixes the
comparison and the translation, and proves the fiber assertion.
\end{proof}

\begin{statement}{E.25}{Proved}{Intrinsic depth, inherited depth, and the even-level obstruction}
Let $N_n=\pi_1(U_{n,\overline{\mathbb Q}})$, embedded in the
geometric profinite group $\widehat F$ of $U_1$ by the covering.
In this statement write $F=\widehat F$ in the profinite formulas.
Use $y$ for
inertia at the base puncture $1$.
For an intrinsic depth filtration choose a Galois-stable set of
cusps above $1$, and let $I_n$ be the closed normal subgroup
generated by their inertia. Put $D_{\mathrm{intr}}^r=\gamma_r I_n$.
One allowed choice for $n=3$ is the rational cusp $Q$ alone;
for $n=4$ take the entire four-point coordinate-cusp orbit.
There are actual quotient maps
\[
 N_n/D_{\mathrm{intr}}^3
   \longrightarrow
 \operatorname{im}(N_n\longrightarrow F/D_y^3),
 \qquad
 \ker=\frac{N_n\cap D_y^3}{D_{\mathrm{intr}}^3},
 \tag{E25.1}
\]
where $D_y^3=\gamma_3\langle\!\langle y\rangle\!\rangle$.
The same formulas apply to the corresponding Lie ideals and
their completed rational Lie algebras.

Over each field $\mathbb Q_\ell$ the inherited prounipotent
image is all of $F_{\mathbb Q_\ell}^{\mathrm{un}}/D_y^3$.
The corresponding inherited arithmetic extension is the base
arithmetic extension and has a rational tangential section,
for every $n$. On abelianized Lie algebras the map from the
intrinsic quotient has a kernel of dimension $n^2-1$.

For every even $n$, the inherited pro-$2$ depth quotient,
the inherited pro-$2$ class-two quotient, and every intrinsic
quotient retaining the curve's own pro-$2$ abelianization have
empty section sets over $\mathbb Q$.
The obstruction is the explicit central extension with class
$\kappa_2(-1)\cup\kappa_2(-1)$ described below.
\end{statement}
\begin{proof}
Inertia above $1$ maps to a conjugate of a power of $y$.
Consequently $I_n$ maps into $\langle\!\langle y\rangle\!\rangle$,
and lower central series give (E25.1), including its exact kernel.
The chosen sets of cusps are Galois stable, so their normal
closures and lower central terms are preserved by the full
arithmetic fundamental group. These are therefore arithmetic
quotients, not just quotients of unmarked geometric groups.

By E.9 the map
$N_{n,\mathbb Q_\ell}^{\mathrm{un}}\to
 F_{\mathbb Q_\ell}^{\mathrm{un}}$ is surjective: positive powers
of $x,y$ lie in $N_n$, and unique roots in characteristic-zero
unipotent groups recover $x,y$. Its inherited quotient is thus
the entire base quotient. The map on arithmetic extensions
induces this surjection on geometric kernels and the identity
on $G_{\mathbb Q}$. Quotienting by the geometric kernel of that
surjection identifies the inherited arithmetic extension with
the base one, which is split by the unit tangent at $0$ towards
$1$. Equivalently its fibers are the full fibers of
$(F_{\mathbb Q_\ell}^{\mathrm{un}}/D_y^3)\rtimes G_{\mathbb Q}$.
Since both depth ideals lie in the respective derived Lie
algebras, their abelianizations are the original ones.
The kernel dimension is therefore exactly the $n^2-1$ of B.5
and E.13.

The integral obstruction uses the quotient
\[
 F\longrightarrow
 D_8=\langle x,y,c\mid
   x^2=y^2=c^2=1,\ c\text{ central},\ [x,y]=c\rangle .
 \tag{E25.2}
\]
With the standard rational tangential section, the Galois action
on this quotient is trivial: in the normalized action of E.20
the cyclotomic exponent is odd, while the conjugating element
lies in $F'$ and hence has central image in $D_8$.
The arithmetic base quotient is $D_8\times G_{\mathbb Q}$.

In the plus model, with $t=-x^n/z^n$, the two covering roots are
\[
 (x/z)^n=-t,\qquad(y/z)^n=-(1-t).
 \tag{E25.3}
\]
At the unit tangent $t=0$, their constant root factors are
$n$th roots of $-1$. For even $n$, their reductions to quadratic
root factors both give
$\epsilon=\kappa_2(-1):G_{\mathbb Q}\to\mathbb Z/2$.
The sign in the stabilizer condition has no effect modulo two.
The geometric image of $N_n$ in $D_8$ is exactly
$\langle c\rangle$: its two exponent sums are even, and
$[x,y]\in N_n$ maps to $c$. The arithmetic image is consequently
the explicit pullback
\[
 \mathcal P=
 \{(d,g)\in D_8\times G_{\mathbb Q}:
           d^{\mathrm{ab}}=(\epsilon(g),\epsilon(g))\}.
 \tag{E25.4}
\]
For each $g$ there is a lift in $\pi_1(U_n)$, and multiplication
by the geometric element $c$ gives both lifts in (E25.4);
this proves equality of the image, not just its containment.
Using $(xy)^{\epsilon(g)}$ as a set-theoretic lift gives the
factor set $c^{\epsilon(g)\epsilon(h)}$.
Thus (E25.4) represents
$\kappa_2(-1)\cup\kappa_2(-1)$.

Complex conjugation has $\epsilon=1$. Both possible lifts,
$xy$ and $xyc$, square to $c\ne1$.
The extension cannot have a section even after restriction to
$G_{\mathbb R}$, hence cannot have one over $\mathbb Q$.
The group $D_8$ kills $\gamma_3F$ and hence $D_y^3$.
Its geometric image from $N_n$ is abelian, so it also kills
$N_n'$. These two explicit factorizations prove all three
nonexistence assertions in the statement.
They recover E.18 at level four and prove the assertion for
all even levels. Pushing to the rank-two toric lattice kills
the central element $c$, exactly as in E.18; extending scalars
to $\mathbb Q_\ell$ gives the inherited comparison already proved.
Thus neither operation is silently identified with the integral
intrinsic quotient.

For clarity the intrinsic marked-depth quotient also has a
precise comparison with the third derived quotient. Their common
quotient is
$N_n/(N_n'''D_{\mathrm{intr}}^3)$.
For the markings just specified neither kernel contains the other.
Indeed filling the marked cusps gives free groups of ranks
$9$ and $13$, respectively. On six distinct remaining generators
use the word
$w=[[[a,b],[a,c]],[[d,e],[d,f]]]$.
Its degree-eight Magnus term has the monomial $abacdedf$ with
coefficient one, and is nonzero. This third-derived word survives
the filling map,
so $N_n'''$ is not contained in $D_{\mathrm{intr}}^3$.
In the other direction choose a marked cusp generator $z_0$
and an unmarked free generator $a$. The element
$[z_0,[z_0,[z_0,a]]]$ lies in $\gamma_3 I_n$.
Its degree-four Magnus term is nonzero, whereas every element
of $N_n'''$ has zero Magnus terms below degree eight.
Thus it does not lie in $N_n'''$.
The corresponding maps from $N_n$, and the common quotient,
state the relation without imposing the base rank-two
containment from E.2 on a different free group.
\end{proof}
}{}
\subsection{The first obstruction beyond the splitting of \S51.12}

\begin{statement}{E.26}{Proved}{The class-two splitting obstruction is the generator modulo three}
Let
\[
 P=\langle\rho,\sigma\mid\rho^3=\sigma^2=1\rangle
   =\operatorname{PSL}_2(\mathbb Z),\qquad
 q:P\longrightarrow S_3,
\]
where $q(\rho)$ is a three-cycle and $q(\sigma)$ a transposition.
The kernel is the free group
\[
 F=\langle x,y\rangle,\qquad
 x=(\sigma\rho)^2,\quad y=(\rho\sigma)^2.
\]
The extension of $S_3$ by $F/F'$ splits by B.11. Its class-two
refinement
\[
 1\longrightarrow F/\gamma_3F\longrightarrow P/\gamma_3F
       \longrightarrow S_3\longrightarrow1                         \tag{E.26a}
\]
does not split. Neither does its profinite completion, its relative
pro-$3$ completion, or any of the discrete or profinite extensions
$P/F^{(j)}\to S_3$ for $j\ge2$.

Write $c=[x,y]$ and $r'=x\rho$. The inverse image in (E.26a) of
the chosen abelian complement is
\[
 1\longrightarrow\mathbb Z(\mathrm{sign})
   \xrightarrow{\,1\mapsto c\,}
   \langle r',\sigma\mid\sigma^2=1,\ \sigma r'\sigma=(r')^{-1}\rangle
   \longrightarrow S_3\longrightarrow1,
 \qquad (r')^3=c.                                             \tag{E.26b}
\]
Its cohomology class is $1$ in
$H^2(S_3,\mathbb Z(\mathrm{sign}))\simeq\mathbb Z/3$.
The finite quotient
\[
 P/(\gamma_3F\,F^3)\longrightarrow S_3                          \tag{E.26c}
\]
has order $162$, Heisenberg kernel of order $27$, and no section.
The inverse image of the chosen abelian complement is the dihedral
group of order $18$.

For the class-two relative pro-$p$ extension, the set of continuous
sections modulo conjugation by its kernel is empty at $p=3$, has
two elements at $p=2$, and has one element at every $p\ge5$.
Over $\mathbb Q$ or $\mathbb Q_\ell$, the rational class-two
relative extension does split, with one conjugacy class of sections.
An explicit section in every one of these cases where it exists is
\[
 \bar\rho\longmapsto c^{-1/3}x\rho,
 \qquad \bar\sigma\longmapsto\sigma.                          \tag{E.26d}
\]
Here $c^a$ denotes the actual central coordinate $a$ of the
class-two kernel; it is defined over $\mathbb Z_p$ for $p\ne3$
and over the indicated fields.

\emph{Manuscript: \S51.12, folios 700--704 and 715--716,
scans 105--114 and 135--138; in particular the abelian splitting
on scan 136. This answers the splitting part of E(c) with its
integral, profinite, and rational comparison maps.}
\end{statement}
\begin{proof}
The free-kernel identification and the matrices of its conjugation
action were proved in A.2 and B.9. In the present central quotient
of the modular group the exact free-group substitutions are
\[
 \rho x\rho^{-1}=y,\qquad
 \rho y\rho^{-1}=x^{-1}y^{-1},\qquad
 \sigma x\sigma=y,\quad\sigma y\sigma=x.                      \tag{E.26e}
\]
They also follow immediately by substituting the displayed words
for $x,y$ and using $\rho^3=\sigma^2=1$.
Consequently, already in $P$,
\[
 (x\rho)^3=xyx^{-1}y^{-1}=c,
 \qquad ((x\rho)\sigma)^2=1.                                 \tag{E.26f}
\]
The second identity follows alternatively by reducing
$(\sigma\rho\sigma\rho^2\sigma)^2$.

We specify the integral class-two group rather than replacing it
by a vector space. Every element of $F/\gamma_3F$ has a unique
normal form $x^a y^b c^k$, $a,b,k\in\mathbb Z$, with law
\[
 (a,b,k)(a',b',k')
   =(a+a',b+b',k+k'-a'b).                                    \tag{E.26g}
\]
Existence is collection using the central commutator $c$.
Uniqueness follows from the injective matrix model
\[
 x^a y^b c^k\longmapsto
 \begin{pmatrix}1&a&ab+k\\0&1&b\\0&0&1\end{pmatrix}.
\]
Thus $c$ has infinite order. Conjugation by $r'$ fixes $c$ and
conjugation by $\sigma$ sends $c$ to $c^{-1}$. The inverse image
of the complement generated by $\overline{x\rho},\bar\sigma$
in $P/F'$ has elements $c^k(r')^i\sigma^j$, with
$0\le i<3$, $0\le j<2$. Since $c=(r')^3$, these are exactly
the normal forms of the infinite dihedral group in (E.26b).
There are no additional relations: the image in $S_3$ first
determines $i,j$, and $c^k=1$ then forces $k=0$.

For clarity we compute the cohomology class. Write elements of
$S_3$ as $(i,j)=\bar\rho^i\bar\sigma^j$, with $0\le i<3$ and
$0\le j<2$. The section $(i,j)\mapsto(r')^i\sigma^j$ gives the
normalized integral two-cocycle
\[
 f((i,j),(i',j'))
   =\left\lfloor\frac{i+(-1)^j i'}{3}\right\rfloor,
 \qquad (i,j)\cdot a=(-1)^j a.                               \tag{E.26h}
\]
Its restriction to the cyclic subgroup has carry $1$, namely
$(r')^3=c$. We can compute the whole cohomology group directly.
In an arbitrary extension by $A(\mathrm{sign})$, for
$A=\mathbb Z$ or $\mathbb Z_3$, choose lifts $r,s$. The elements
$s^2$ and $(rs)^2$ lie in $A$ and are fixed by conjugation by
$s$ and $rs$, respectively. Those conjugations act as $-1$ on
$A$; because $A$ has no $2$-torsion, both elements vanish.
The only remaining parameter is $r^3=a\in A$. Replacing $r$
by $b r$, with additive notation on $A$, changes $a$ to $a+3b$.
Conversely every $a\in A$ defines an extension: take
$(A\oplus\mathbb Zr)/(3r-a)$ and adjoin the involution acting
by negation. This proves $H^2(S_3,A(\mathrm{sign}))=A/3A$ and
identifies (E.26b) with $1$. Equivalently,
the one-cochain $(i,j)\mapsto i$ has coboundary $3f$.

To exclude sections of the whole group, we must also exclude
other abelian complements. B.10 proves, over every coefficient
group used here,
\[
 H^0(S_3,M)=H^1(S_3,M)=0,
 \qquad M=F^{\mathrm{ab}}.
\]
Thus conjugation by a unique element of $M$ carries any abelian
section to the B.11 section. The surjection
$F/\gamma_3F\to M$ lifts this conjugating element.
After this conjugation any proposed section of (E.26a) must send
$\bar\rho$ to $c^a r'$ and $\bar\sigma$ to $c^b\sigma$.
Its order-three relation is
\[
 (c^a r')^3=c^{1+3a}=1.                                    \tag{E.26i}
\]
There is no solution $a\in\mathbb Z$ or $a\in\mathbb Z_3$.
The other two defining relations hold for all $a,b$, because
$\sigma$ and $r'\sigma$ invert $c$ and have order two.
This proves the integral and relative pro-$3$ assertions.

Here are the exact completion and quotient comparisons. Since
$F$ is normal in $P$ and $\gamma_3F$, $F^3$, and $F^{(j)}$ are
characteristic in $F$, every displayed quotient is a group
extension over the same $S_3$. The relative pro-$p$ class-two
kernel is obtained by completing the normal form (E.26g) to
coordinates in $\mathbb Z_p$. Conjugation (E.26e) extends
continuously: its formulas use addition, multiplication, and
$\binom a2$, the continuous integer-valued polynomial on
$\mathbb Z_p$. The finite cosets of $F$ give the corresponding
continuous extension over $S_3$.
The quotient (E.26c) receives (E.26a) by reducing its three
kernel coordinates modulo three. It therefore also receives
the profinite completion of (E.26a), continuously. A section
of that completion would give a section of (E.26c).

To construct (E.26c) completely, over $\mathbb F_3$ the law
(E.26g) is the Heisenberg group of order $27$ and exponent
three. It is exactly $F/(\gamma_3F F^3)$: collection bounds
its order by $27$, and the matrix model realizes all $27$
normal forms. The quotient has six such cosets, hence order
$162$. With states $x^a y^b c^k\rho^i\sigma^j$, right
multiplication by $\sigma$ toggles $j$. Right multiplication
by $\rho$ increments $i$ when $j=0$, and when $j=1$ it
multiplies the kernel coordinate by $\rho^i x\rho^{-i}$ and
replaces $i$ by $i+2$ modulo three. The three kernel factors are
$x,y,x^{-1}y^{-1}$. These rules construct its regular
permutation representation without a presentation-order guess.
The program \texttt{code/e\_splitting.py} checks closure of the
$162$ normal-form permutations under both generators, all
$4\,251\,528$ associativity triples, and all $27^2=729$
pairs of possible section-generator images. There are no
sections. Its abelian quotient of order $54$ has exactly nine
sections; all $81$ pairs lifting these have their
$\bar\rho$ image of order nine. The chosen complement has
preimage $\langle r',\sigma\rangle=D_{18}$, with
$(r')^3=c$. The code also checks all $216$ two-cocycle
identities and all $243$ normalized one-cochains modulo three:
none has coboundary $f$.

For $j\ge2$, the commutator identities
$[\gamma_aF,\gamma_bF]\subset\gamma_{a+b}F$ give
$F^{(j)}\subset\gamma_{2^j}F\subset\gamma_3F$. Consequently
the identity on generators gives the exact quotient map
\[
 P/F^{(j)}\twoheadrightarrow P/\gamma_3F,
 \quad\ker=\gamma_3F/F^{(j)}.                               \tag{E.26j}
\]
A section of the source extension would compose to a section
of (E.26a). For the profinite version, use the continuous map
to the finite group (E.26c); its kernel contains the closed
$j$th derived subgroup. The same composition gives the
contradiction. This establishes the transfer to the requested
third-derived quotient by an actual surjection.

Finally, the normalized lift calculation (E.26i) works with
central coordinates in $\mathbb Z_p$, $\mathbb Q$, or
$\mathbb Q_\ell$. It has the unique solution $a=-1/3$ exactly
in the coefficient cases specified in the statement. The
parameter $b$ remains arbitrary. Any kernel conjugation
preserving the normalized abelian complement has zero image
in $M$, by $H^0(S_3,M)=0$, and is therefore $c^u$.
It changes $b$ to $b+2u$. Thus the conjugacy classes are
$\mathbb Z_p/2\mathbb Z_p$ for $p\ne3$, and a singleton
over the indicated characteristic-zero fields. Formula
(E.26d) gives their specified representative. In the
cochain description, $-i/3$ cancels (E.26h), so the comparison
to rational coefficients explicitly kills this integral
obstruction.
The rational group here is the class-two relative unipotent
quotient over $S_3$: its six cosets are copies of the
Heisenberg group (E.26g) over the field. To verify the universal
comparison, a homomorphism from $F$ to a class-two unipotent
group, with images $X,Y$, extends uniquely by
$x^a y^b c^k\mapsto X^aY^b[X,Y]^k$. Rational powers are given
by the finite exponential and logarithm polynomials.
The finite-coset relations then give the unique extension
of the original homomorphism from $P$ over $S_3$.
The program also checks the normalized central equations
$1+3a=0$ and the conjugacy relation $b\sim b+2u$ modulo
$p^e$ for $p=2,3,5,7$ and $1\le e\le4$.

We also record precisely how the marked-depth construction
in E.2 compares. There is a surjective Lie homomorphism
\[
 L/D^3L\longrightarrow L/\gamma_3L,
 \qquad \ker=\gamma_3L/D^3L,
\]
because a monomial with at least three occurrences of $y$
has total degree at least three. The specific full
permutation action does not descend to that marked quotient:
the swap $x\leftrightarrow y$ sends
$[y,[y,[y,x]]]\in D^3L$ to
$[x,[x,[x,y]]]$, a nonzero depth-one element modulo $D^3L$.
Thus (E.26j) is an $S_3$-extension comparison for the derived
quotients, whereas the displayed marked-depth map has the
stated Lie source and target. The rational section (E.26d)
and the integral obstruction (E.26i) retain these exact
domains and coefficient rings.
\end{proof}

\subsection{Boundary types and the scope of the covering criterion}

\begin{statement}{E.27}{Proved}{The four boundary types and their actual comparison maps}
Over an algebraically closed field of characteristic zero the geometric
fundamental groups of types $(0,0),(0,1),(0,2),(1,0)$ are, respectively,
\[
 1,\qquad1,\qquad\widehat{\mathbb Z}(1),\qquad
 T E\simeq\widehat{\mathbb Z}^{,2},
\]
where a chosen origin gives the elliptic curve $E$ in the last case.
The last two groups are their own centers and their own abelian and
metabelian quotients. The assertions about nontrivial comparison kernels
in E.7--E.9 have the stated rank-two free-group source. For the discrete
group $F_1=\mathbb Z$ and its subgroup $d\mathbb Z$, the abelian,
metabelian and fixed nilpotent comparisons to their correctly descended
kernels are isomorphisms for every $d$. The rational completion map
is multiplication by $d$ on $\mathbb Q$, hence is also an isomorphism.
For $d=p^a m$ with $(m,p)=1$, the pro-$p$ comparison is the isomorphism
\[
 \mathbb Z_p\longrightarrow p^a\mathbb Z_p,\qquad u\longmapsto du.
\]
The same statements hold for a finite-index lattice in $\mathbb Z^2$,
with its inclusion matrix in place of $d$.

For the corresponding nontrivial finite covers of $B\mathbb Z$ or
$B\mathbb Z^2$, intrinsic-fibration membership for the five localizers
of E.5 still holds only for the $p$-group covers in the pro-$p$ case.
Thus the equivalence between the two criteria in E.10 is proved for its
specified source $F_2$; the displayed lattice comparisons compute what
happens at these boundary sources.
\end{statement}
\begin{proof}
For the projective line a connected degree-$d$ finite etale cover has
$2g-2=-2d$ by Riemann--Hurwitz, forcing $d=1$. For the affine line,
extend a finite etale cover to smooth projective curves. Ramification
can occur only over infinity, and its total contribution is at most
$d-1$. Hence $2g-2\le-d-1$, again forcing $d=1$. For the twice-punctured
line, ramification is allowed at two points. The same formula forces
$g=0$ and total ramification over each of them. Choosing coordinates
at their unique preimages makes the covering $z\mapsto z^d$.
The cofinal inverse system of these maps has deck groups $\mu_d$,
which proves the asserted Tate-module description and its action.

A finite etale cover of a genus-one curve has genus one by the same
formula. Choose a point over the fixed origin. The resulting pointed
map of elliptic curves is an isogeny: its map to the target Jacobian
is the pointed curve map under the degree-zero divisor identification.
Its kernel is finite; its dual isogeny shows that a multiplication
map dominates it. The multiplication coverings are therefore cofinal,
and their inverse deck system is $TE$. In characteristic zero
$E[n]\simeq(\mathbb Z/n)^2$, giving the displayed geometric group.
For the cubic over $\mathbb Q$ the specified origin is a rational
coordinate cusp; E.24 computes its required arithmetic abelian
cohomology. We impose no rational-point conclusion on a genus-one
torsor without such an origin.

For the rank-one discrete group, every verbal commutator subgroup is
zero, so the covering inclusion itself is exactly the displayed
kernel map. Rationalization gives multiplication by $d$.
In $\mathbb Z_p$, multiplication by $m$ is a unit and multiplication
by $p^a$ maps bijectively onto the closed subgroup $p^a\mathbb Z_p$.
For a rank-two lattice, integral Smith normal form gives the same
calculation in two coordinates. These arguments compute both the
image and zero kernel, rather than applying the free-group rank
formula outside its hypotheses. The program \texttt{e\_boundary.py}
checks the rank-one formulas for $d\le24$ and the inverse unit
coordinates modulo $p^k$, for $p=2,3,5,7$ and $k\le4$.

Intrinsic-fibration membership is preserved by pullback, directly
from its universal base-change definition. Pull a nontrivial lattice
cover back by $BF_2\to B\mathbb Z$ or $BF_2\to B\mathbb Z^2$,
using a surjection onto the lattice. Its monodromy is still the
given finite quotient, so E.10 excludes the pulled-back cover in
all the claimed cases. The general $p$-group sufficiency in E.10
proves the other cases. This supplies the comparison between the
two criteria at every boundary source under discussion.

Real descent also retains its arithmetic data. The plus conic
$X_2:x_0^2+x_1^2+x_\infty^2=0$ has no real point, but its
geometric projective fundamental group is trivial, so its arithmetic
fundamental group maps isomorphically to $G_{\mathbb Q}$ and has
its unique section. Deleting the coordinate divisors gives the
nontrivial real obstruction treated by E.25. The exact open-to-proper
quotient kills that obstruction. The elliptic and multiplicative
centers above, the retained central $C_2$ in A.14, and the real
orientation-reversing extension in E.18 are consequently part of
the calculations rather than omitted exceptional cases.
\end{proof}

\subsection{A cubic Kummer obstruction at the prime 109}

\begin{statement}{E.28}{Proved}{A nonempty cubic abelian fiber with no class-two lifts}
In the pro-$3$ version of E.24, consider the toric section class
\[
 (\kappa(6),\kappa(109))
 \in H^1(\mathbb Q,\mathbb Z_3(1)^2).
 \tag{E28.1}
\]
Its fiber of geometrically abelian sections is nonempty and is
a torsor under
$\mu_3^3\times\mathcal R_{\mathcal I}(\mathbb Z_3)\times\mathbb Z_3$.
No point of this fiber lifts to a section of the curve's own
class-two pro-$3$ quotient. No such lift exists for the intrinsic
metabelian quotient, the intrinsic marked-depth quotient of E.25,
or the inherited class-two pro-$3$ quotient either.
\end{statement}
\begin{proof}
Let $F_3$ be the free pro-$3$ group on the base loops $x,y$,
and put $\mathcal H=F_3/\gamma_3F_3$.
In its class-two Lie coordinates over $\mathbb Z_3$ write
$X=\log x$, $Y=\log y$, $C=[X,Y]$.
The Baker--Campbell--Hausdorff formula has only the denominator
two, which is a unit in $\mathbb Z_3$.
The image $H$ of $N_3$ consists exactly of
\[
 \exp(3aX+3bY+cC),\qquad a,b,c\in\mathbb Z_3.
 \tag{E28.2}
\]
The exponent-sum condition proves containment, and the elements
$x^3,y^3,[x,y]\in N_3$ generate the full lattice in (E28.2).
In coordinates $(a,b,c)$ its multiplication is
\[
 (a,b,c)(a',b',c')=
 \left(a+a',b+b',
 c+c'+\frac92(ab'-ba')\right).
 \tag{E28.3}
\]
In particular the commutator is $(0,0,9(ab'-ba'))$.
The coordinates $a,b$ are the two coordinate-unit homology
coordinates of $U_3$; the covering pushforward to base
homology is $(3a,3b)$.

The normalized Galois action sends $x$ to $x^\chi$ and
$y$ to $f_g^{-1}y^\chi f_g$, with $f_g\in F_3'$.
In $\mathcal H$ this conjugator is central, so the action on
(E28.2) is exactly
\[
 (a,b,c)\longmapsto(\chi(g)a,\chi(g)b,\chi(g)^2c).
 \tag{E28.4}
\]
The unit tangent at $t=0$ lifts to the plus cubic by
$w=-x/z$, $t=w^3$ and $y/z=-(1-w^3)^{1/3}$.
Its toric leading coefficients are $(-1,-1)$.
Thus the base rational tangential section lifts to the cubic
and the inherited extension is $H\rtimes G_{\mathbb Q}$.
The pro-$3$ Kummer classes of both these signs, and of the
sign at the origin used in E.24, are zero. Changing between
these two tangential origins therefore fixes the toric
coordinates of every abelian class, so it leaves (E28.1)
and its entire fiber unchanged.

Reduce all three coordinates in (E28.3) modulo $27$.
This gives a Galois-stable finite quotient $H_{27}$ of order
$27^3$, of exponent $27$, whose commutator of lifts of
$(1,0)$ and $(0,1)$ is $(0,0,9)\ne1$.
At the local field $\mathbb Q_{109}$ the action (E28.4) on
$H_{27}$ is trivial: the $3$-adic cyclotomic character is
unramified there, with Frobenius $109\equiv1\pmod {27}$.
The integer $6$ has exact order $108$ in
$\mathbb F_{109}^{\times}$. Consequently its Kummer class
modulo $27$ is unramified with Frobenius value a generator of
$\mu_{27}$, while the Kummer class of $109$ has tame-inertia
value a generator. Choose the common generator of $\mu_{27}$
so that the Frobenius value for $6$ is $1$; a generator of
tame inertia can then be chosen with the value for $109$
also equal to $1$. We may choose Frobenius fixing the selected
roots of the uniformizer. The two projected values are thus
$(1,0)$ and $(0,1)$.

A section would therefore give a homomorphism
$G_{\mathbb Q_{109}}\to H_{27}$ with those two projections.
Wild inertia is pro-$109$ and maps trivially to this finite
$3$-group. In the tame quotient, choose Frobenius $\sigma$
and tame inertia $\tau$ using roots of the uniformizer.
Frobenius acts on their roots of unity by the $109$th power,
so
$\sigma\tau\sigma^{-1}=\tau^{109}$.
Since $H_{27}$ has exponent $27$, this relation forces the
two images to commute. Equation (E28.3) gives their commutator
as $(0,0,9)$, independently of their central coordinates,
a contradiction.
The program \texttt{e\_cubic\_class2.py} checks the order of
$6$ modulo $109$, the exponent and the full multiplication
formula, and all $27^2$ choices of the two central lifts.

Finally the map $N_3\to H_{27}$ kills $\gamma_3N_3$,
$N_3''$ and $\gamma_3 I_3$.
These are the actual factorizations through the quotients in
the statement. The explicit toric projection is $(a,b)$, so
every hypothetical lift of any point of the E.24 fiber would
give the forbidden homomorphism. E.24 proves the asserted
nonemptiness and size of the abelian fiber independently of
this obstruction. There is no reduction map from the rational
group $H_{\mathbb Q_3}$ to $H_{27}$ extending this reduction on
its integral lattice: every element of $H_{\mathbb Q_3}$ is a
$27$th power, whereas every $27$th power in $H_{27}$ is one.
Thus any such homomorphism is trivial.
The finite obstruction is used only for the indicated integral
pro-$3$ quotients. The result does not classify the remaining
cubic class-two sections whose toric pairs pass this local test.
\end{proof}

\IfFileExists{sections/E_rational_pairs.tex}{\subsection{Every rational toric pair in the inherited cubic class-two quotient}

\begin{statement}{E.29}{Proved}{The exact tame-residue criterion and the three classes in each nonempty fibre}
Let $F$ be the free pro-$3$ geometric fundamental group of
$\mathbb P^1-\{0,1,\infty\}$, based at the rational tangent
$t=0$, $dt=1$, and let $N_3=\ker(F\to(\mathbb Z/3)^2)$.
The image of $N_3$ in $F/\overline{\gamma_3F}$ is the pro-$3$
group $H$ with coordinates $(u,v,w)\in\mathbb Z_3^3$, law
\[
 (u,v,w)(u',v',w')
 =\left(u+u',v+v',w+w'+\frac92(uv'-vu')\right),              \tag{E.29a}
\]
and arithmetic action
\[
 g(u,v,w)=(\chi_3(g)u,\chi_3(g)v,\chi_3(g)^2w).
\]
It is the inherited class-two quotient of the plus cubic;
its arithmetic extension, with the lifted tangent, is
$H\rtimes G_{\mathbb Q}$.

For $a,b\in\mathbb Q^*$ let $\mathscr F(a,b)$ be the fibre over
$(\kappa_a,\kappa_b)$ of the continuous section-class map
\[
 H^1_{\mathrm{cts}}(G_{\mathbb Q},H)
 \longrightarrow H^1_{\mathrm{cts}}
       (G_{\mathbb Q},\mathbb Z_3(1)^2).
\]
For a rational prime $q$, put
\[
 \theta_q(a,b)=
 (-1)^{v_q(a)v_q(b)}
 \overline{a^{v_q(b)}/b^{v_q(a)}}\ \in\mathbb F_q^*.
                                                               \tag{E.29b}
\]
The expression under the bar is a $q$-adic unit. Let
$o_3(\theta)$ denote the largest power of three dividing its
multiplicative order. Then the complete answer on rational
toric pairs is
\[
 |\mathscr F(a,b)|=
 \begin{cases}
 3,&o_3(\theta_q(a,b))\le9\text{ for every prime }q,\\
 0,&o_3(\theta_q(a,b))>9\text{ for at least one prime }q.
 \end{cases}                                                   \tag{E.29c}
\]
Only primes dividing a numerator or denominator of $a$ or $b$
need to be tested. For each nonempty fibre, choose one continuous
cocycle $z(g)\in H$. Its three distinct classes are represented by
\[
 z_j(g)=\left(0,0,j\frac{\chi_3(g)^2-1}{3}\right)z(g),
 \qquad j=0,1,2.                                               \tag{E.29d}
\]
The choice of $z$ is the stated ambiguity in this parametrization.
The obstruction itself is exactly
$9\kappa_a\cup\kappa_b\in
H^2_{\mathrm{cts}}(G_{\mathbb Q},\mathbb Z_3(2))$.

For example, $(6,109)$ and $(2,163)$ have empty fibres,
whereas $(6^3,109)$ and $(2^9,163)$ each have three classes.
The source of this classification is the displayed inherited
quotient. The intrinsic rank-ten class-two quotient maps onto
it by
\[
 N_3/\overline{\gamma_3N_3}\twoheadrightarrow H,
 \qquad
 \ker=(N_3\cap\overline{\gamma_3F})/
                     \overline{\gamma_3N_3}.                  \tag{E.29e}
\]
Thus its section classes map to the fibres computed here.

\emph{Manuscript: \S51.16, folios 733v--736v, scans 172--178,
(371)--(379); the Kummer coordinates of the minimum quotient
are retained, and their next inherited integral obstruction is
computed. This statement treats every rational toric pair.}
\end{statement}
\begin{proof}
We first fix the comparison and its factors of three. Put
$X=\log x$, $Y=\log y$, and $C=[X,Y]$ in the integral
class-two pro-$3$ group; the denominator $2$ in its
Baker--Campbell--Hausdorff law is a unit in $\mathbb Z_3$.
The image of $N_3$ has the lattice
\[
 A=3X,\qquad B=3Y,\qquad C=[X,Y],\qquad[A,B]=9C.
\]
Indeed its abelian coordinates are divisible by three,
and $x^3,y^3,[x,y]\in N_3$ generate all these coordinates.
This proves (E.29a) and the exact surjection (E.29e).
For the normalized arithmetic action,
$x\mapsto x^{\chi_3(g)}$ and
$y\mapsto f_g^{-1}y^{\chi_3(g)}f_g$, where $f_g\in F'$.
The element $f_g$ is central after quotienting by
$\overline{\gamma_3F}$, giving the displayed diagonal action.

On the plus cubic put $u=x_0/x_\infty$, $v=x_1/x_\infty$.
The covering is $t=-u^3$, $1-t=-v^3$.
At its rational cusp above $t=0$, the parameter $s=-u$ gives
$t=s^3$ and $v=-(1-s^3)^{1/3}$ with constant coefficient $-1$.
The tangential Galois section at $s=0$, $ds=1$ maps to the
chosen base section: the corresponding embedding of Puiseux
fields sends $t^{1/m}$ to $s^{3/m}$ and commutes with
coefficientwise $G_{\mathbb Q}$-action. This specifies the
comparison at the ramified cusp.
Both leading toric units are $-1$, whose pro-$3$ Kummer
classes are zero, since $-1$ has compatible rational
$3^n$th roots. Thus a toric cocycle $(\kappa_a,\kappa_b)$
has inherited coordinates $\kappa_a A+\kappa_b B$ exactly.
The comparison to base toric coordinates is multiplication
by three, coming from $t=-u^3$, $1-t=-v^3$.

Write $V=\mathbb Z_3(1)^2$, $Z=\mathbb Z_3(2)$, so that
$1\to Z\to H\to V\to1$ is a central extension of pro-$3$
groups with $G_{\mathbb Q}$-action. For a continuous
$V$-cocycle $v_g=(\alpha_g,\beta_g)$, use its continuous
set-theoretic lift $(\alpha_g,\beta_g,0)$. Its obstruction
two-cocycle is
\[
 \omega(g,h)=\frac92\left(
 \alpha_g\chi_3(g)\beta_h-
 \beta_g\chi_3(g)\alpha_h\right).                           \tag{E.29f}
\]
The symmetric sum of the two cup cocycles is the negative
coboundary of $g\mapsto\alpha_g\beta_g$. Hence the class
of (E.29f) is $9[\alpha\cup\beta]$.
Vanishing of this class means precisely that a continuous
central one-cochain cancels (E.29f), which is the cocycle
condition for a lift to $H$.

We now prove the criterion in continuous cohomology itself.
The pairing
\[
 (a,b)\longmapsto\kappa_a\cup\kappa_b
 \quad\text{in }H^2_{\mathrm{cts}}(G_{\mathbb Q},Z)
\]
is bilinear and satisfies the integral Steinberg relation.
For the latter, the rational point $t=a$ of
$\mathbb P^1-\{0,1,\infty\}$, for $a\ne0,1$, supplies
an actual arithmetic section. Its class-two pro-$3$
quotient has toric cocycle $(\kappa_a,\kappa_{1-a})$ and
commutator coefficient one. The calculation (E.29f) with
$9$ replaced by $1$ therefore gives
$\kappa_a\cup\kappa_{1-a}=0$ in continuous integral
cohomology. Bilinearity now defines a homomorphism
\[
 \beta:K_2^M(\mathbb Q)\longrightarrow
 H^2_{\mathrm{cts}}(G_{\mathbb Q},Z),\qquad
 \{a,b\}\longmapsto\kappa_a\cup\kappa_b.                    \tag{E.29g}
\]
This constructs (E.29g) before taking any inverse limit of
degree-two cohomology groups.

We use Matsumoto's presentation and the localization
calculation in the opened source
\cite[III.6.1, III.6.5--III.6.5.1, III.5.2.2]{E_Kbook}.
For $\mathbb Q$ these give the exact sequence
\[
 0\longrightarrow\mathbb Z/2\longrightarrow K_2^M(\mathbb Q)
 \xrightarrow{(\partial_q)_q}
 \bigoplus_q\mathbb F_q^*\longrightarrow0.                  \tag{E.29h}
\]
The kernel is generated by $\{-1,-1\}$. In the convention
of that source, $\partial_q\{a,b\}=\theta_q(a,b)^{-1}$;
coordinatewise inversion is the explicit automorphism of
the right-hand side relating the two conventions.
Every element of the middle group is torsion: its finitely
supported image has finite order, and its remaining kernel
is killed by two. In particular (E.29h) identifies its
$3$-primary subgroup with
$\bigoplus_q(\mathbb F_q^*)\{3\}$.
The target of (E.29g) is a $\mathbb Z_3$-module, so
$\beta$ kills all primary components at primes different
from three.

We give the local detection argument to prove injectivity
on the remaining $3$-primary subgroup. Let $q\ne3$.
For $T_n=\mu_{3^n}^{\otimes2}$, wild inertia has no positive
cohomology with these coefficients, and the $3$-part of
tame inertia is $\mathbb Z_3(1)$. The quotient Galois group
of the finite residue field and this tame $3$-part have
$3$-cohomological dimension one. Their two-term procyclic
resolutions, followed by Hochschild--Serre, give
\[
 H^2(\mathbb Q_q,T_n)
  \simeq H^1(\mathbb F_q,\mu_{3^n})
  \simeq (\mathbb Z/3^n)/(q-1).                              \tag{E.29i}
\]
In particular the degree-one finite-coefficient local
groups are finite, so their first derived inverse limit
vanishes. Passing to the inverse limit of (E.29i) therefore
gives an actual continuous-cohomology isomorphism
\[
 \operatorname{res}^{\mathrm{tame}}_q:
 H^2_{\mathrm{cts}}(\mathbb Q_q,Z)
   \xrightarrow{\sim}
 H^1_{\mathrm{cts}}(\mathbb F_q,\mathbb Z_3(1))
   \simeq\mathbb Z_3/(q-1).                                 \tag{E.29j}
\]

Here is the cup-product normalization of this map.
Choose arithmetic Frobenius $\phi$ and a tame generator
$\tau$ with $\phi\tau\phi^{-1}=\tau^q$; these are the
operators on compatible roots of a uniformizer. The
conjugation formula is also given in the opened source
\cite[(8.2.2)]{E_BrinonConrad}.
For a unit $u$, its Kummer class is unramified. Write
$a=\kappa_u(\phi)$ and choose $\phi$ fixing the chosen
uniformizer roots, so $\kappa_q(\phi)=0$ and
$\kappa_q(\tau)=1$. The extension two-cocycle
$\kappa_u\cup\kappa_q$ has the relation defect $qa$
between $\phi\tau\phi^{-1}$ and $\tau^q$.
Changing the central lift of $\tau$ by $b$ changes that
defect by $q(q-1)b$. Thus (E.29j), which divides the
defect by $q$, sends this cup product to the unit's
residue-field Kummer class. This specifies the map and its
sign, and verifies the edge-map calculation directly.
Two unit Kummer classes have zero cup product because
they factor through the procyclic residue Galois group.
Also $\kappa_q\cup\kappa_q=0$, since graded commutativity
kills it by two and two is a unit in $\mathbb Z_3$.
Writing $a=q^\alpha u$, $b=q^\beta v$ now yields
\[
 \operatorname{res}^{\mathrm{tame}}_q
    (\kappa_a\cup\kappa_b)
    =\kappa_{\overline{u}^{\beta}/\overline{v}^{\alpha}}
    =\kappa_{\theta_q(a,b)}.                                 \tag{E.29k}
\]
The sign in (E.29b) has zero pro-$3$ Kummer class.
Finite-field Kummer theory identifies the image of an
element of $\mathbb F_q^*$ in (E.29j) with its
$3$-primary component, preserving its order.

Every nonzero $3$-primary element of $K_2^M(\mathbb Q)$
has a nonzero tame component at some $q\ne3$, by
(E.29h), since $\mathbb F_3^*$ has order two.
Equation (E.29k) detects that component under $\beta$.
Consequently $\beta$ is injective on this subgroup.
It follows that
\[
 9\kappa_a\cup\kappa_b=0
 \quad\Longleftrightarrow\quad
 9\{a,b\}_{\{3\}}=0
 \quad\Longleftrightarrow\quad
 o_3(\theta_q(a,b))\le9\quad\text{for every }q.               \tag{E.29l}
\]
This proves both necessity and sufficiency in continuous
cohomology, including sufficiency for the existence of a
continuous central cochain. No global inverse-limit
injectivity assertion was used.

It remains to compute each nonempty fibre. We first compute
\[
 H^1_{\mathrm{cts}}(G_{\mathbb Q},\mathbb Z_3(2))
      \simeq\mathbb Z/3.                                     \tag{E.29m}
\]
A continuous cocycle with values in $\mathbb Q_3(2)$
restricts on inertia at $q\ne3$ to a homomorphism.
Wild inertia contributes zero. If $a$ is its value on a
tame $3$-generator, Frobenius equivariance gives
$q^2a=qa$, so $(q-1)a=0$, hence $a=0$.
It therefore vanishes on the closed normal subgroup
generated by these inertia groups and factors through
$G_{\mathbb Q,\{3\}}$.
Theorem 7.2 and Theorem B.1 of the opened
Hain--Matsumoto paper \cite{E_HainMatsumoto} give
\[
 \dim_{\mathbb Q_3}H^1_{\mathrm{cts}}
     (G_{\mathbb Q,\{3\}},\mathbb Q_3(2))=r_2(\mathbb Q)=0.
\]
This proves the vanishing over the full $G_{\mathbb Q}$.
The exact coefficient sequence
$0\to\mathbb Z_3(2)\to\mathbb Q_3(2)
\to\mathbb Q_3/\mathbb Z_3(2)\to0$ thus identifies (E.29m)
with $H^0(G_{\mathbb Q},\mathbb Q_3/\mathbb Z_3(2))$.
The image of $\chi_3$ is $\mathbb Z_3^*$; the ideal generated
by $u^2-1$ for all units $u$ is $3\mathbb Z_3$.
The invariants are consequently
$(\tfrac13\mathbb Z_3)/\mathbb Z_3$, and the boundary of
$j/3$ is exactly the central cocycle in (E.29d).

To see that this calculation gives the whole fibre, choose
one lifted cocycle. The ratio of two lifts with the same
toric cocycle is a central one-cocycle, and every such
cocycle gives another lift. A kernel conjugation which
preserves the exact toric cocycle has projection in
$H^0(G_{\mathbb Q},V)=0$: complex conjugation acts by
$-1$ on the torsion-free group $V$. Such a conjugation
is therefore central and changes the ratio only by a
central coboundary. When the toric cocycles initially
differ by a coboundary, the surjection $H\to V$ supplies
the conjugation which makes them equal. Thus the fibre
is a torsor under the explicitly computed group (E.29m),
with exactly the representatives (E.29d).

The program \texttt{code/e\_rational\_pairs.py} computes
(E.29b) by exact rational arithmetic. Among all $40\,000$
ordered positive integer pairs up to $200$, it finds
$39\,378$ nonempty and $622$ empty fibres. The histogram
of the largest $3$-primary residue order is
\[
\begin{array}{c|rrrrr}
 \text{order}&1&3&9&27&81\\\hline
 \text{number of pairs}&19036&15198&5144&364&258.
\end{array}
\]
It also checks all $101\,124$ ordered pairs of the $318$
nonzero signed reduced rationals of numerator and
denominator height at most $16$, and all $317$ applicable
Steinberg pairs in that set. It checks (E.29d) as a
cocycle modulo $3^n$ for $1\le n\le4$, and the three
twist-two invariants and absence of twist-one invariants
at those finite levels. All output and every obstructed
integer pair are retained in the accompanying JSON file.
\end{proof}
}{}
\IfFileExists{sections/E_proconstant.tex}{\subsection{Arithmetic completion over the proconstant base}

\begin{statement}{E.30}{Proved}{Arithmetic completion over the proconstant base kills the geometric group}
Put $U=\mathbb P^1_{\mathbb Q}-\{0,1,\infty\}$ and
$\mathcal E=\pi_1([U/S_3])$, with the positive rational tangential base point lifted to $U$.
The quotient by its geometric free subgroup $\widehat F_2$ is
$S_3\times G_{\mathbb Q}$.
For $m\geq1$, let $K_m=\mathbb Q(\mu_m)$,
$C_m=\operatorname{Gal}(K_m/\mathbb Q)$ and
$B_m=S_3\times C_m$.
There are two precise versions of the following construction:
abstract relative completion over $k=\mathbb Q$, and continuous relative completion over $k=\mathbb Q_\ell$.
Use the corresponding version throughout each formula, and let
$\mathcal A_{m,k}$ be the relative completion of $G_{\mathbb Q}\to C_m$.
The relative completion of $\mathcal E\to B_m$ is exactly
\[
 \mathcal R_{m,k}=S_3\times\mathcal A_{m,k},\qquad
 1\longrightarrow(G_{K_m})^{\rm un}_k
 \longrightarrow\mathcal A_{m,k}\longrightarrow C_m
 \longrightarrow1.                                      \tag{E.30.1}
\]
For $m\mid n$ the transition maps are induced by the identity on
$G_{\mathbb Q}$; on the displayed kernels they are induced by
$G_{K_n}\hookrightarrow G_{K_m}$.
Consequently the completion over the proconstant group scheme
$B=S_3\times\widehat{\mathbb Z}^{\times}_k$ is the specified inverse limit
\[
 \mathcal R_{B,k}=\varprojlim_{m}\mathcal R_{m,k},\qquad
 1\longrightarrow\varprojlim_m(G_{K_m})^{\rm un}_k
 \longrightarrow\mathcal R_{B,k}\longrightarrow B
 \longrightarrow1.                                      \tag{E.30.2}
\]
The natural homomorphism $\mathcal E\to\mathcal R_{B,k}(k)$ kills
the whole geometric subgroup $\widehat F_2$.
In particular its restriction to every Fermat subgroup is trivial;
it retains none of the geometric homology computed in E.14.

There is a different, explicit completion retaining that geometry.
Let $\mathcal A_\ell$ be the negatively weighted completion of
$G_{\mathbb Q}\xrightarrow{\chi_\ell}\mathbb G_m(\mathbb Q_\ell)$,
with central cocharacter $t\mapsto t^{-2}$.
For the arithmetic group $\pi_1(U)=\widehat F_2\rtimes G_{\mathbb Q}$
with the stated tangential section, its negatively weighted completion is
\[
 (\widehat F_2)^{\rm un}_{\mathbb Q_\ell}\rtimes\mathcal A_\ell.
                                                               \tag{E.30.3}
\]
The geometric factor is the free prounipotent group on two generators.
No algebraic morphism from the proconstant $B$ to $\mathbb G_m$
induces $\chi_\ell$ on its cyclotomic factor.
\end{statement}
\begin{proof}
The morphism $[U/S_3]\to[\operatorname{Spec}\mathbb Q/S_3]$
gives the stated quotient $S_3\times G_{\mathbb Q}$, since the
permutation action is defined over $\mathbb Q$. Its geometric
kernel is $\widehat F_2$. Thus the kernel $N_m$ of
$\mathcal E\to B_m$ has the tangentially split sequence
\[
 1\longrightarrow\widehat F_2\longrightarrow N_m
 \longrightarrow G_{K_m}\longrightarrow1.
\]
On $\widehat F_2^{\rm ab}=\widehat{\mathbb Z}^2$ the arithmetic
action is scalar multiplication by the full cyclotomic character.

First work continuously over $\mathbb Q_\ell$.
Let $r:N_m\to U_0(\mathbb Q_\ell)$ be any continuous homomorphism
to a unipotent algebraic group. Let $A$ be the Zariski closure of
$r(\widehat F_2)$. It is a connected unipotent normal subgroup
of the Zariski closure of $r(N_m)$; unipotent groups in characteristic
zero have no nontrivial finite algebraic quotients.
The map to $A^{\rm ab}$ factors through
$\mathbb Z_\ell^2$ and spans its Lie algebra over $\mathbb Q_\ell$.
Choose $g\in G_{K_m}$ with $\chi_\ell(g)\ne1$, possible because
the cyclotomic image of this finite-index subgroup is infinite.
Conjugation by $r(g)$ on $\operatorname{Lie}(A^{\rm ab})$
is both multiplication by $\chi_\ell(g)$ and a unipotent linear
operator. The latter assertion follows, for example, by putting
$U_0$ in upper triangular unipotent matrices and considering its
conjugation action on its Lie algebra and invariant subquotients.
It follows that $A^{\rm ab}=1$. A nonzero nilpotent Lie algebra
has nonzero abelianization, so $A=1$. Every such $r$ therefore
kills $\widehat F_2$.

For the abstract rational version, the required abelianization
is still exactly $\widehat{\mathbb Z}^2$: the abstract commutator
subgroup of a finitely generated profinite group is closed by
\cite[Theorem1.4]{dNikolovSegal2007}.
Choose $u\in\widehat{\mathbb Z}^{\times}$ with $u\equiv1\pmod m$
and $u-1=dv$, where $d\ne0$ is an integer and
$v\in\widehat{\mathbb Z}^{\times}$, as in D.3.
Explicitly use $u_p=2$ for odd $p\nmid m$,
$u_p=1+p^{v_p(m)}$ for odd $p\mid m$, and
$u_2=1+2^{\max(1,v_2(m))}$.
Cyclotomic surjectivity supplies $g\in G_{K_m}$ with $\chi(g)=u$.
For an abstract homomorphism $r:N_m\to U_0(\mathbb Q)$, there is
an equivariant surjection
\[
 \widehat{\mathbb Z}^2\otimes_{\mathbb Z}\mathbb Q
 \longrightarrow\operatorname{Lie}(A^{\rm ab}).
\]
On the source $u-1=dv$ is invertible. Therefore
$\operatorname{Ad}(r(g))-1$ on the target is surjective. It is
also nilpotent, since $r(g)$ is unipotent. The target is zero,
and the preceding nilpotent Lie algebra argument again gives $A=1$.
This proves the rational assertion without assuming continuity
of a homomorphism to $U_0(\mathbb Q)$.

For a group mapping onto a finite group $C$, its relative
completion is the extension of $C$ by the unipotent completion
of the kernel. One can construct it directly by choosing coset
representatives: multiplication on the finite disjoint union of
copies of the completed kernel is given by the induced conjugation
automorphisms and the original factor-set cocycle. This is the
construction and universal-property proof used in E.14; it applies
to the abstract and continuous categories just considered.
The result above identifies $N_m^{\rm un}$ with
$(G_{K_m})^{\rm un}$. Hence every representation defining
$\mathcal R_{m,k}$ kills $\widehat F_2$ and factors uniquely through
$S_3\times G_{\mathbb Q}$. Completing this direct product gives
$S_3\times\mathcal A_{m,k}$: the finite $S_3$ image commutes with
the arithmetic image, and these two homomorphisms have exactly
the product universal property. This proves (E.30.1) and the
specified transition maps.

An algebraic homomorphism from a proconstant profinite group
scheme to a finite-type affine group scheme factors through a
finite quotient. Indeed its finitely many coordinate generators
belong to one finite stage of the coordinate-ring direct limit;
the multiplication identities also hold at a sufficiently large
stage. Every finite-dimensional algebraic representation of $B$
therefore factors through some $B_m$. Defining relative completion
by representations whose semisimplifications are algebraic
representations of $B$ gives precisely the inverse system
(E.30.2): every such finite-dimensional representation occurs at
one of its stages. This also proves its universal property.
The transition maps on the unipotent kernels are faithfully flat:
a finite-index subgroup has Zariski-dense image in the connected
unipotent closure of the larger group's image. Inverse limits
therefore give the exact proalgebraic sequence displayed in
(E.30.2). The geometric image was already killed at every stage.

For (E.30.3), the genuine action of $G_{\mathbb Q}$ on the
geometric prounipotent group has scalar $\chi_\ell$ on its
abelianization. Its kernel on that abelianization is prounipotent,
with negative weights for $t\mapsto t^{-2}$: the degree-$d$
part of its IA Lie algebra maps a degree-one generator to a
degree-$d$ word and has weight $-2(d-1)<0$ for $d\geq2$.
Thus the universal
property of weighted completion extends this action uniquely to
$\mathcal A_\ell$; the weighted conventions and construction are
those of \cite[\S\S2--4]{E_HainMatsumoto}.
The semidirect product in (E.30.3) receives the actual arithmetic
group. Given any lift of that group into a negatively weighted
extension of $\mathbb G_m$, its restriction to $\widehat F_2$
extends uniquely by continuous unipotent completion, and its
restriction to the tangential $G_{\mathbb Q}$ extends uniquely by
weighted completion. Their equivariance holds on the dense
original images, hence algebraically. Thus these two extensions
define a unique homomorphism from the semidirect product, proving
its weighted universal property. Its geometric kernel is visibly
the free prounipotent group, proving its injection as well.
Finally an algebraic map $B\to\mathbb G_m$ has finite image by
the finite-stage argument, whereas $\chi_\ell(G_{\mathbb Q})$
is infinite. This proves the claimed obstruction to that
specific base-change comparison. We do not identify an abstract
rational completion after scalar extension with a continuous
$\ell$-adic completion; the two constructions above specify
their source categories independently.
\end{proof}
}{}
\IfFileExists{sections/E_intrinsic_obstruction.tex}{\subsection{An intrinsic cubic obstruction inside the zero toric fiber}

\begin{statement}{E.31}{Proved}{Two abelian classes in one toric fiber have different liftability}
Use the pro-$3$ version of E.24, with its unit rational tangent
at $O=[1:-1:0]$. Put $K=\mathbb Q(\zeta)$,
$\lambda=1-\zeta$ and $\pi=3+\zeta$, so
$N_{K/\mathbb Q}(\pi)=7$.
In the frame coordinates of (E24.4) take the rational torus point
\[
 t=(7,7,1,1,\pi)\in T_D(\mathbb Q)\subset G(\mathbb Q).
 \tag{E31.1}
\]
Its Kummer class
$\alpha\in H^1(G_{\mathbb Q},T_3G)$ has zero image in the
toric minimum. Let $\mathcal K$ denote the kernel of that toric
map on section classes. Every class in the nonempty coset
$\alpha+3\mathcal K$ fails to lift to the cubic's own class-two
pro-$3$ arithmetic quotient, its metabelian quotient, and its
marked-depth quotient of E.25. The zero class in the same toric
fiber lifts to a full tangential section.

The obstruction is visible in an explicit quotient of the
intrinsic based Galois action over $K$.
Let $H_3$ be the group of order $27$ with coordinates
$(a,b,c)\in\mathbb F_3^3$ and multiplication
\[
 (a,b,c)(a',b',c')
 =\left(a+a',b+b',
 c+c'+\tfrac12(ab'-ba')\right).
 \tag{E31.2}
\]
There is a surjection from the cubic geometric group to $H_3$
whose Galois action, with the tangent at $O$ as origin, is
\[
 g:(a,b,c)\longmapsto
       (a,b,c+\kappa_3(\zeta)(g)b),
       \qquad g\in G_K .
 \tag{E31.3}
\]
Thus (E31.3) gives the arithmetic cross term explicitly, as
well as a finite obstruction to the indicated section classes.
\end{statement}
\begin{proof}
Equation (E24.4) sends (E31.1) to
$(7/N(\pi),7/N(\pi))=(1,1)$.
Kummer functoriality therefore puts $\alpha$ in $\mathcal K$.
This is a genuine geometrically abelian section class by E.24.
It is nonzero modulo three: the valuation at the prime $(\pi)$
in the third $K^\times$ factor of (E24.7) is zero on its
displayed $\lambda$ relation, and takes the value one on $t$.
It extends to an integral homomorphism on that rational-point
presentation, so $t$ is not three times a rational point of $G'$.

Over $K$ use the actual morphism
\[
 h=\frac{x+y}{x+\zeta y}:
 U_3\longrightarrow\mathbb P^1\setminus\{0,1,\infty\},
 \qquad
 1-h=\frac{(\zeta-1)y}{x+\zeta y}.
 \tag{E31.4}
\]
The equations $h=0$ and $h=\infty$ force $z=0$ on the cubic;
$h=1$ forces $y=0$. All these points have been removed from
$U_3$, so the displayed target and morphism are exact.
Write $C_i=[1:-\zeta^i:0]$ and $B_i=[-\zeta^i:0:1]$.
The numerator and denominator tangent lines give
\[
 \operatorname{div}(h)=3O-3C_2,\qquad
 \operatorname{div}(1-h)=B_0+B_1+B_2-3C_2 .
 \tag{E31.5}
\]
The frame at $C_2$ is $\bar\pi$, not $\pi$.
Formula (E24.3) consequently gives the exact character values
\[
 \chi_h(t)=\bar\pi^{-3},\qquad
 \chi_{1-h}(t)=7\bar\pi^{-3}.
 \tag{E31.6}
\]

The based arithmetic map is obtained by applying (E31.4) to
the tangential fiber functor at $O$ and choosing a geometric
path to the standard unit tangent at zero in the target.
Its affine effect can be computed without a path convention.
At $O$ put $s=z/x$; then
$y/x=-(1+s^3)^{1/3}$, and hence
\[
 h=-\frac{s^3}{3(1-\zeta)}+O(s^6)
   =\zeta\lambda^{-3}s^3+O(s^6),
 \qquad 1-h=1+O(s^3).
 \tag{E31.7}
\]
The identity in the second expression is
$\lambda^2=-3\zeta$.
The map of Puiseux fiber functors under (E31.7) has the
constant Kummer factors $(\zeta\lambda^{-3},1)$.
Thus the image of the tangential section at $O$ has, modulo
three, abelian coordinates
$(\kappa_3(\zeta),0)$ relative to the standard target origin.
Combining (E31.6) and (E31.7), the abelian section with class
$\alpha$ has target coordinates
\[
 \bigl(\kappa_3(\zeta\lambda^{-3}\bar\pi^{-3}),
        \kappa_3(7\bar\pi^{-3})\bigr)
   =(\kappa_3(\zeta),\kappa_3(7)).
 \tag{E31.8}
\]
This retains the nonzero affine origin term explicitly.

Take the quotient of the target free pro-$3$ group in which
the two generators and their central commutator have order
three; it is (E31.2). Its standard based Galois action over
$K$ is trivial. Indeed $\mu_3\subset K$, and the normalized
action of E.20 has cyclotomic exponent one modulo three and
a conjugator with central image.
Let $\phi_h:N_3\to H_3$ be the geometric map induced by
(E31.4). It is surjective. The two geometric Kummer characters
$h,1-h$ are independent modulo cubes: the valuation at a
$B_i$ forces the coefficient of $1-h$ to be zero modulo
three; $h$ is not a cube, since such a cube root would have
divisor $O-C_2$, a degree-one map from an elliptic curve to
$\mathbb P^1$. Their independent abelian images generate
$H_3/[H_3,H_3]$, and their commutator generates its central
subgroup, proving surjectivity.

Let $q_g\in H_3$ be the first component of the image of the
tangential section at $O$ in
$H_3\times G_K$. Equivariance of the arithmetic map gives
$\phi_h(gng^{-1})=q_g\phi_h(n)q_g^{-1}$.
Its abelian coordinate is
$(\kappa_3(\zeta)(g),0)$ by (E31.7); its central coordinate
does not affect conjugation. Conjugation in (E31.2) is therefore
exactly (E31.3).
The character $\kappa_3(\zeta)$ is nonzero: a cube root of
$\zeta$ would be a unit in $\mathbb Z[\zeta]$, whose unit group
$\mu_6$ has only $\pm1$ as cubes.
This proves the asserted nonzero arithmetic cross term.

Now pass to the prime $(\pi)$ over $7$.
Its residue field is $\mathbb F_7$ and
$\zeta\equiv-3\equiv4\pmod\pi$, so
$K_{(\pi)}=\mathbb Q_7$.
The character $\kappa_3(\zeta)$ is unramified and nonzero:
choose $\zeta_9^3=\zeta$ in the unramified cubic extension.
Arithmetic Frobenius sends $\zeta_9$ to $\zeta_9^7$, whose
ratio to $\zeta_9$ is $\zeta^2$.
The Kummer character of the uniformizer $7$ has tame-inertia
value $\zeta$ after a choice of generator, and Frobenius can
be chosen to fix its selected roots.
Consequently (E31.8) sends Frobenius and tame inertia to
$(2,0)$ and $(0,1)$ in $H_3^{\mathrm{ab}}$.
Every pair of lifts has commutator $(0,0,2)\ne1$.

An arithmetic section lifting $\alpha$ would yield a
homomorphism $G_{\mathbb Q_7}\to H_3$, because the standard
target action is trivial. Wild inertia is pro-$7$ and maps
trivially to this finite $3$-group. The tame relation is
$\sigma\tau\sigma^{-1}=\tau^7$.
Every element of $H_3$ has exponent three, so this relation
forces its two images to commute, a contradiction.
The exact field identities, the local leading constants,
all $27$ group elements and all nine choices of the two
central lifts are checked by
\texttt{e\_intrinsic\_obstruction.py}.

The map $\phi_h$ kills $\gamma_3N_3$, $N_3''$ and
$\gamma_3I_3$. Thus it factors through each intrinsic quotient
named in the statement. These factorizations also apply to
the arithmetic maps after restriction to $G_K$, which suffices
to rule out a section over $\mathbb Q$.
There is also a single map over $\mathbb Q$ displaying both
conjugate comparisons: the adjoint of (E31.4) is
\[
 U_3\longrightarrow
 \operatorname{Res}_{K/\mathbb Q}
      (\mathbb P^1\setminus\{0,1,\infty\}).
 \tag{E31.9}
\]
After base change to $\overline{\mathbb Q}$ it is
$(h,\bar h)$, giving the finite geometric image inside
$H_3\times H_3$. The arithmetic map to this image is defined
over $\mathbb Q$ and its geometric kernel contains
$\gamma_3N_3$. Its restriction has the first projection
used above. This supplies an actual global finite quotient
comparison as well.

Adding any element of $3\mathcal K$ to $\alpha$ leaves
(E31.8) unchanged, since the target abelian coefficient group
has exponent three. The same obstruction applies to the
whole coset $\alpha+3\mathcal K$. The zero class is the original
rational tangential section and lifts to the full arithmetic
fundamental group. Both lie over the zero pro-$3$ toric class,
as asserted. In particular this proof distinguishes
liftability inside one specified fiber, using a quotient that
detects an intrinsic direction beyond the inherited toric
comparison of E.28.
\end{proof}
}{}

\IfFileExists{sections/E_completed_pairs.tex}{\subsection{All completed toric pairs in the inherited cubic class-two quotient}

\begin{statement}{E.32}{Proved}{The complete convergence criterion, including the invisible inverse-limit obstruction}
Retain the arithmetic group $H\rtimes G_{\mathbb Q}$ of E.29.
For a set $I$ write
\[
 \mathcal R_I=\{(z_i)\in\textstyle\prod_{i\in I}\mathbb Z_3:
       \text{for every }n,\ z_i\in3^n\mathbb Z_3
       \text{ outside a finite subset of }I\}.
\]
Let $\mathcal P$ be the rational primes and
$\mathcal P^*=\mathcal P-\{3\}$. The Kummer isomorphism is
\[
 H^1_{\rm cts}(G_{\mathbb Q},\mathbb Z_3(1))
       =\mathcal R_{\mathcal P},\qquad
 (a_p)_p\longmapsto\sum_p a_p\kappa_p.                 \tag{E.32a}
\]
Choose a compatible system of roots for each $\kappa_p$;
the sum in (E.32a) converges uniformly as a cocycle.

Here is an explicit criterion for every pair
$a=(a_p),b=(b_p)\in\mathcal R_{\mathcal P}$.
For $q\ne3$ put $e_q=v_3(q-1)$ and choose a primitive
$3^{e_q}$th root $\eta_q\in\mathbb F_q^*$ when $e_q>0$.
For $p\ne q$ let $r_{q,p}\in\{0,\ldots,3^{e_q}-1\}$ be
specified by
\[
 \overline p^{(q-1)/3^{e_q}}=\eta_q^{r_{q,p}}.
\]
When $e_q=0$, put $r_{q,p}=0$. Define convergent sums and a
central correction by
\[
 A_q=\sum_{p\ne q}a_p r_{q,p},\qquad
 B_q=\sum_{p\ne q}b_p r_{q,p},\qquad
 c_q=-\frac{9(A_qb_q-B_qa_q)}{q-1}\in\mathbb Q_3.
                                                               \tag{E.32b}
\]
The fibre of $H^1_{\rm cts}(G_{\mathbb Q},H)$ over the pair
in (E.32a) has exactly three classes precisely when
\[
 c_q\in\mathbb Z_3\quad(q\ne3),\qquad
             (c_q)_{q\ne3}\in\mathcal R_{\mathcal P^*};       \tag{E.32c}
\]
in all other cases that fibre is empty. The criterion is
independent of the roots and generators used above.
For a nonempty fibre its three representatives are given by
(E.29d), starting with any one lift.

The continuous degree-two group, and the morphism used to
compute this obstruction, are
\[
 \begin{aligned}
 H^2_{\rm cts}(G_{\mathbb Q},\mathbb Z_3(2))
   &\xrightarrow[\sim]{\operatorname{loc}_{\rm cts}}
      \mathcal R_{\mathcal P^*}/d\mathcal R_{\mathcal P^*},\\
 9\Big(\sum_pa_p\kappa_p\Big)\cup
       \Big(\sum_pb_p\kappa_p\Big)
   &\longmapsto [\,9(A_qb_q-B_qa_q)\,]_q .
 \end{aligned}                                                 \tag{E.32d}
\]
Here $d(z)_q=(q-1)z_q$.
These are isomorphisms of $\mathbb Z_3$-modules; the denominator
is the actual image of $d$, with no closure added.
The map from this quotient to the product of local groups is
coordinatewise reduction modulo $q-1$. Its exact kernel is
\[
 \frac{\{z\in\mathcal R_{\mathcal P^*}:
                   z_q\in(q-1)\mathbb Z_3\text{ for all }q\}}
      {d\mathcal R_{\mathcal P^*}}.                            \tag{E.32e}
\]

This kernel contains an obstruction of an actual completed
toric pair. Choose recursively increasing primes $q_j$, $j\ge3$,
with
\[
 q_j\equiv1\pmod{3^j},\qquad
             3^{(q_j-1)/3}\not\equiv1\pmod{q_j},
\]
and write $r_j=r_{q_j,3}\in\mathbb Z_3^*$.
Take $a_3=1$, $a_p=0$ for $p\ne3$, and
\[
 b_{q_j}=\frac{q_j-1}{9r_j},\qquad b_p=0
              \text{ at the other primes}.                    \tag{E.32f}
\]
Then $b\in\mathcal R_{\mathcal P}$ and every local obstruction
and every finite-coefficient global obstruction vanishes.
There is nevertheless no continuous section with this completed
toric pair: its corrections are $c_{q_j}=-1$ for all $j$.
Every finite-support partial pair obtained from (E.32f) does
have three continuous section classes.

The classification concerns the inherited quotient $H$.
The map from the cubic's intrinsic quotient is exactly (E.29e);
the additional intrinsic obstruction E.31 remains present.
\emph{Manuscript: \S51.16, folios 733v--736v, scans 172--178,
(371)--(379), with its completed Kummer parameters retained.}
\end{statement}
\begin{proof}
We first calculate the continuous cohomology group and its
localization morphism, retaining the derived inverse limit.
Put
\[
 U=\operatorname{Spec}\mathbb Z[1/3],\qquad
 G_S=G_{\mathbb Q,\{3,\infty\}},\qquad
 T_n=\mathbb Z/3^n(2),\qquad T=\mathbb Z_3(2).
\]
Milne's comparison for locally constant sheaves
\cite[II.2.9]{E_MilneADT} identifies
$R\Gamma(U,T_n)$ with $R\Gamma(G_S,T_n)$.
We need the following actual vanishing:
\[
 H^2(G_S,T_n)=H^3(G_S,T_n)=0,
 \quad H^1(G_S,T_n)\text{ finite},\qquad n\ge1.                \tag{E.32g}
\]
For $n=1$, $T_1=\mathbb F_3$ with trivial action, since
$\chi_3^2\equiv1\pmod3$.
The Kronecker--Weber theorem
\cite[I.4.16]{E_MilneCFT} gives
$H^1(G_S,\mathbb F_3)=\mathbb F_3$.
More explicitly, an abelian extension unramified away from
three is contained in a cyclotomic field, and its quotient
kills the inertia factors at all other primes. Its cubic
quotient therefore factors through
$\operatorname{Gal}(\mathbb Q(\zeta_{3^m})/\mathbb Q)$,
which has a unique quotient of order three. The nontrivial
quotient already occurs in $\mathbb Q(\zeta_9)^+$.
Both $H^0(G_S,\mathbb F_3)$ and $H^1(G_S,\mathbb F_3)$
have order three. The global Euler characteristic
\cite[I.5.1]{E_MilneADT} equals one here: its real factor is
$|H^0(G_{\mathbb R},\mathbb F_3)|/3=1$.
Consequently $H^2(G_S,\mathbb F_3)=0$.
Filter $T_n$ by its powers of three. Each successive quotient
is the same trivial $\mathbb F_3$, so the long exact
cohomology sequence proves $H^2(G_S,T_n)=0$ for all $n$.
Milne I.4.10(c) identifies every degree at least three with
real-place cohomology, which is zero for a module of odd order.
Finiteness is I.4.15. This proves (E.32g).

For a finite set $J\subset\mathcal P^*$, apply the
localization triangle to
$U_J=U-\{q:q\in J\}$, using
\cite[II.0.1(f), II.1.7(b)]{E_MilneADT}.
Here the local purity identification can also be checked
directly: over the strict Henselization at $q$, a locally
constant $T_n$ is constant; the fraction field has tame
$3$-inertia $\mathbb Z_3(1)$, while its wild inertia is
pro-$q$. Its cohomology is $T_n$ in degree zero,
$T_n(-1)$ in degree one, and zero above degree one.
The local support complex is therefore
$T_n(-1)[-2]$, with its residue-field action.
Thus there is the compatible distinguished triangle
\[
 R\Gamma(U,T_n)\longrightarrow R\Gamma(U_J,T_n)
 \longrightarrow\bigoplus_{q\in J}
       R\Gamma(\mathbb F_q,\mathbb Z/3^n(1))[-1]
 \longrightarrow R\Gamma(U,T_n)[1].                         \tag{E.32h}
\]
First take the filtered direct limit over finite $J$.
Filtered direct limits of abelian groups are exact. On the
generic-field term this gives $R\Gamma(G_{\mathbb Q},T_n)$:
every continuous cochain with finite coefficients factors
through a finite quotient of $G_{\mathbb Q}$, whose number
field is ramified at only finitely many primes.
In particular the third term becomes a direct sum at each
fixed coefficient level $n$; no infinite product is substituted.

Next take $R\varprojlim_n$ of these triangles. The comparison
with continuous cochains is Jannsen's theorem
\cite[(2.2), together with (1.6)]{E_JannsenContinuous}.
For the first term (E.32g), the finiteness of the degree-one
groups kills their first derived inverse limit. It follows that
\[
 H^2(R\varprojlim_nR\Gamma(U,T_n))
 =H^3(R\varprojlim_nR\Gamma(U,T_n))=0.                       \tag{E.32i}
\]
For the support term, the procyclic residue Galois group has
the compatible two-term cochain model
\[
 R\Gamma(\mathbb F_q,\mathbb Z/3^n(1))
  \simeq[\,\mathbb Z/3^n\xrightarrow{q-1}\mathbb Z/3^n\,]
                 \quad\text{in degrees }0,1.                \tag{E.32j}
\]
Choose arithmetic Frobenius for this model. After taking the
direct sum over $q\ne3$, the transition in each of its two
terms is the surjective reduction
$\bigoplus_q\mathbb Z/3^{n+1}\to\bigoplus_q\mathbb Z/3^n$.
The inverse systems of terms are therefore acyclic for
$\varprojlim$. Their derived inverse limit is the literal
complex
\[
 [\,\mathcal R_{\mathcal P^*}
         \xrightarrow{d}\mathcal R_{\mathcal P^*}\,].          \tag{E.32k}
\]
Indeed a compatible sequence of finitely supported residues
is exactly a family of $3$-adic integers tending to zero.
The degree-two segment of the resulting triangle, using
(E.32i), proves the first isomorphism in (E.32d).
This calculation takes the inverse limit of complexes. It
does not replace it by the inverse limit of their cohomology.

We next compute the image of the actual obstruction under
this morphism. The Kummer sequence and
$\mathbb Q^*=\{\pm1\}\times\bigoplus_p p^{\mathbb Z}$
give (E.32a); the sign factor has trivial pro-$3$ completion.
Multiplying a fixed cocycle $\kappa_p$ by $a_p$ and summing
is legitimate on continuous cochains: for each $n$, all but
finitely many summands vanish modulo $3^n$, uniformly on
$G_{\mathbb Q}$. Write these sums as $\alpha,\beta$.
Their obstruction is the continuous cocycle
\[
 \omega(g,h)=\frac92\{
    \alpha(g)\chi_3(g)\beta(h)
       -\beta(g)\chi_3(g)\alpha(h)\},                        \tag{E.32l}
\]
representing $9\alpha\cup\beta$ as in E.29.

Fix $q\ne3$ and a decomposition group. Choose tame inertia
$\tau_q$ and arithmetic Frobenius $\phi_q$ such that
\[
 \phi_q\tau_q\phi_q^{-1}=\tau_q^q,
 \qquad \kappa_q(\tau_q)=1,
 \qquad \kappa_q(\phi_q)=0.                                 \tag{E.32m}
\]
These choices use compatible roots of a uniformizer, as in
E.29. Put
$\alpha(\tau_q)=a_q$, $\beta(\tau_q)=b_q$ and
$\alpha(\phi_q)=A'_q$, $\beta(\phi_q)=B'_q$.
On the full inertia group the two cocycles are proportional
to the same tame character; wild inertia maps to zero.
Consequently (E.32l) is identically zero on two inertia
arguments, giving its canonical central-zero inertia splitting.
In the pullback of $H\rtimes G_{\mathbb Q}$, let the central
coordinates of the two lifts be $w$ at $\tau_q$ and $z$ at
$\phi_q$. Direct multiplication in (E.29a) gives the defect
between the two sides of (E.32m) as
\[
 q\{(q-1)w+9(A'_qb_q-B'_qa_q)\}.                            \tag{E.32n}
\]
It is independent of $z$. In particular the central-zero
inertia splitting represents the residue cochain with value
$9(A'_qb_q-B'_qa_q)$ at Frobenius. The division by $q$ in
(E.32n) is the identification
$\operatorname{Hom}(\mathbb Z_3(1),T)=T(-1)$:
the conjugation defect at $\tau_q^q$ is pulled back to its
value at $\tau_q$. Changing the inertia splitting by $w$
adds precisely $(q-1)w$, the differential in (E.32j).
This is the chain-level localization morphism in (E.32h):
it records the failure of an inertia splitting to be invariant
under residue-field Frobenius.

At each finite coefficient level these inertia splittings
agree with an unramified splitting outside finitely many
primes, since $\alpha,\beta$ factor through a finite quotient
modulo $3^n$. They therefore form a compatible representative
in the direct-sum complexes used in (E.32h)--(E.32k).
The resulting family tends to zero, since it contains a
factor $a_q$ or $b_q$ in each summand.
This establishes the localization representative itself,
including its class in the quotient (E.32k).

For a unit $p$ at $q$, finite-field Kummer theory gives
\[
 \kappa_p(\phi_q)\equiv r_{q,p}\pmod{q-1}.
\]
Choose the place over $q$ to send the selected global
$3^{e_q}$th root of unity to $\eta_q$.
The formula follows by applying Frobenius to a
$3^{e_q}$th root of $p$: its quotient is
$\overline p^{(q-1)/3^{e_q}}$.
Write
$\kappa_p(\phi_q)-r_{q,p}=(q-1)t_{q,p}$ with
$t_{q,p}\in\mathbb Z_3$.
The sums of $a_pt_{q,p}$ and $b_pt_{q,p}$ converge and are
uniformly bounded in $\mathbb Z_3$.
Hence replacing $A'_q,B'_q$ by $A_q,B_q$ changes the
representative in (E.32d) by $d$ of a family tending to zero:
the change divided by $q-1$ has a factor $a_q$ or $b_q$.
This proves the second formula in (E.32d).
The same argument proves independence of root choices.
Changing a residue generator scales the corresponding
coordinate by a $3$-adic unit and adds a term of this form,
so it also preserves (E.32c).

For completeness, changing global cocycles by
$\alpha\mapsto\alpha+du$, $\beta\mapsto\beta+dv$, with
$du(g)=(\chi_3(g)-1)u$, changes the residue representative by
\[
 (q-1)\,9(ub_q-va_q).
\]
The family in the second factor lies in
$\mathcal R_{\mathcal P^*}$.
The corresponding central inertia correction changes by
$-9(ub_q-va_q)$, exactly as the change
$z_g\mapsto h^{-1}z_g g(h)$ with $h=(u,v,0)$ gives.
Also replacing $\phi_q$ by
$\tau_q^r\phi_q$ leaves $A'_qb_q-B'_qa_q$ unchanged.
These formulas specify both the coboundary sign and the
actual coefficient comparisons.

The obstruction vanishes precisely when its representative
lies in $d\mathcal R_{\mathcal P^*}$. Since $q-1\ne0$ in
$\mathbb Z_3$, its unique possible central correction is
(E.32b). This proves (E.32c) and the empty-fibre assertion.
The central group has
$H^1_{\rm cts}(G_{\mathbb Q},T)=\mathbb Z/3$ by E.29,
and $H^0(G_{\mathbb Q},\mathbb Z_3(1)^2)=0$.
The exact nonabelian central-extension sequence therefore
makes each nonempty fibre a torsor under this group.
Its three explicit twists are (E.29d).

We finally verify the inverse-limit example. The field
$\mathbb Q(\sqrt[3]{3})$ cannot be contained in a cyclotomic
field: a subfield of an abelian Galois extension is Galois,
whereas this irreducible real cubic does not contain the two
nonreal roots of $X^3-3$.
Thus
$\mathbb Q(\zeta_{3^j},\sqrt[3]{3})/
 \mathbb Q(\zeta_{3^j})$
has degree three. A nonidentity element of this subgroup has
a conjugacy class restricting trivially to the cyclotomic
field. Chebotarev
\cite[V.3.23]{E_MilneCFT} provides infinitely many rational
primes with that class, giving the two requirements on $q_j$.
Choose each $q_j$ to be the least such prime larger than the
previous choice. Noncubicity says $r_j$ is a $3$-adic unit,
so $v_3(b_{q_j})=e_{q_j}-2\ge j-2$.
This proves that (E.32f) is an actual completed Kummer pair.
Its representative in (E.32d) is
\[
 z_{q_j}=q_j-1,\qquad z_q=0\text{ otherwise}.                \tag{E.32o}
\]
It is a restricted family, but its only preimage under
coordinatewise multiplication by $q-1$ is the family with
value one at every $q_j$, which is not restricted.
Thus its class is nonzero in (E.32d), despite lying in the
kernel (E.32e). At every finite level its class is zero in
the cohomology of the direct-sum complex: only finitely
many entries of (E.32o) survive modulo $3^n$, and they
are obtained by applying $d$ to the finitely supported
family with value one at those entries.
The finite-level global obstruction is zero by
(E.32g)--(E.32h). At $q\ne3$ the local obstruction is zero
by coordinatewise divisibility. At three, finite local
duality \cite[I.2.3]{E_MilneADT} gives
$H^2(\mathbb Q_3,T_n)=H^0(\mathbb Q_3,\mathbb Z/3^n(-1))^*=0$:
the local cyclotomic character contains $-1$, whose action
has no fixed vector on this odd-order module.
Finite local degree-one groups then give
$H^2_{\rm cts}(\mathbb Q_3,T)=0$ as well.
Real cohomology with these coefficients is zero.
Every finite-support partial pair in (E.32f) has finitely
many corrections equal to $-1$, so it satisfies (E.32c).
All assertions follow.
\end{proof}
}{}
\subsection{The continuous arithmetic kernel and its trace maps}

\begin{statement}{E.33}{Proved}{The continuous arithmetic kernels are free, with explicit generators as a Galois module}
Fix a prime $\ell$ and put $K_m=\mathbb Q(\mu_m)$,
$d_m=[K_m:\mathbb Q]$ and $r_m=r_2(K_m)+1$.
The continuous unipotent completion occurring in E.30 is
\[
 (G_{K_m})^{\rm un}_{\mathbb Q_\ell}
 \simeq \exp\widehat{\mathbb L}(V_m),\qquad
 V_m=\frac{K_m\otimes_{\mathbb Q}\mathbb Q_\ell}
 {\log_\ell(\mathcal O_{K_m}^{\times}\otimes_{\mathbb Z}\mathbb Q_\ell)}.
                                                        \tag{E.33.1}
\]
Here $\widehat{\mathbb L}$ is completion of the free Lie algebra by
bracket length. The vector space $V_m$ is the canonical abelianization
of this Lie algebra; choosing lifts of a basis gives the displayed
noncanonical free presentation. It has dimension $r_m$.
For $m\geq3$, complex conjugation $c$ gives the explicit
$\operatorname{Gal}(K_m/\mathbb Q)$-equivariant identification
\[
 V_m\xrightarrow{\sim}\mathbb Q_\ell\oplus
        (K_m\otimes\mathbb Q_\ell)^-,\qquad
 [z]\longmapsto
 \left(\frac{\operatorname{Tr}_{K_m/\mathbb Q}(z)}{d_m},
                         \frac{z-cz}{2}\right).           \tag{E.33.2}
\]
For $m=1,2$ it is $V_m=\mathbb Q_\ell$.
For $m\mid n$, the actual homomorphism of completed groups is induced
by $G_{K_n}\hookrightarrow G_{K_m}$, as in E.30, and its map on
abelianizations is
\[
 [z]\longmapsto[\operatorname{Tr}_{K_n/K_m}z].             \tag{E.33.3}
\]
In (E.33.2) this is $(a,w)\mapsto([K_n:K_m]a,
\operatorname{Tr}_{K_n/K_m}w)$. Thus every transition is faithfully
flat. These groups and these specified homomorphisms compute the
continuous arithmetic kernel in (E.30.2); no identification of
the abstract rational completion with this continuous completion is
used.
More explicitly, along the cofinal conductors $m_j=4j!$, free
generators can be chosen compatibly so that each transition keeps
the first $r_{m_j}$ generators and sends every additional generator
to $1$. Thus that kernel is an inverse limit of free prounipotent
groups of the explicit ranks $r_{m_j}$ with these deletion maps.
\end{statement}
\begin{proof}
Write $K=K_m$, and let $S$ consist of the places above $\ell$ and
the infinite places. Every continuous homomorphism
$\rho:G_K\to U(\mathbb Q_\ell)$, for $U$ unipotent algebraic,
is unramified outside $S$. Indeed its compact image is pro-$\ell$:
use a central series of $U$ with additive quotients and the fact
that a compact additive subgroup of a finite-dimensional
$\mathbb Q_\ell$-space is pro-$\ell$. At a finite place $v\nmid\ell$,
wild inertia therefore dies, and tame inertia maps through its
$\mathbb Z_\ell(1)$ quotient. If $N$ is the logarithm of the image
of a tame generator and $F_v$ is arithmetic Frobenius, the exact
relation $F_v\tau F_v^{-1}=\tau^{q_v}$ gives
\[
 \operatorname{Ad}(\rho(F_v))N=q_v N.
\]
The left-hand operator is unipotent and $q_v\ne1$ in
$\mathbb Q_\ell$. Hence $N=0$. The normal closed subgroup generated
by these inertia groups is the kernel of $G_K\to G_{K,S}$,
which proves the claimed factorization, with its universal property.

The class-field-theoretic exact sequence of local principal units,
global units and the maximal abelian pro-$\ell$ extension unramified
outside $S$ gives
\[
 H_1(G_{K,S},\mathbb Q_\ell)
 =\operatorname{coker}\left(
 \mathcal O_K^{\times}\otimes\mathbb Q_\ell
 \xrightarrow{\log_\ell}
 \bigoplus_{v\mid\ell}K_v\right).                       \tag{E.33.4}
\]
The finite ideal class group disappears after tensoring. The exact
sequence, including its proof from reciprocity, is
\cite[Theorem1.5.4 and Theorem1.5.7]{E_SharifiIwasawa}.
The diagonal logarithm is injective here: $K_m/\mathbb Q$ is
abelian, so the proved abelian-field case of Leopoldt applies
\cite[Theorems1.5.17 and1.5.21]{E_SharifiIwasawa}.
To specify the reduction, the unit representation of the maximal
real subfield is its augmentation representation. A rational
normal generator gives conjugate units whose logarithms have just
the norm relation over $\mathbb Q$; Brumer's linear-independence
theorem retains independence over algebraic coefficients.
Diagonalizing the abelian group matrix by characters makes every
nontrivial character eigenvalue nonzero. Thus the logarithm has
the full Dirichlet rank. This is precisely the proof of the cited
Theorem1.5.21, and applies to every $\ell$, including $2$.
Consequently (E.33.4) has dimension
$[K:\mathbb Q]-(r_1(K)+r_2(K)-1)=r_2(K)+1$.

We also need degree two, rather than infer freeness from this
dimension. Apply the finite-coefficient Euler characteristic
\cite[TheoremI.5.1]{E_MilneADT} to $\mathbb Z/\ell^a$ with
trivial action. Its formula is exactly
\[
 \frac{|H^0(G_{K,S},\mathbb Z/\ell^a)|
       |H^2(G_{K,S},\mathbb Z/\ell^a)|}
      {|H^1(G_{K,S},\mathbb Z/\ell^a)|}
       =\ell^{-a r_2(K)}.                               \tag{E.33.5}
\]
All these finite-coefficient groups are finite, by
\cite[CorollaryI.4.15]{E_MilneADT}. The class-field computation says
that $H^1$ has order $\ell^{a(r_2+1)}$ times a bounded factor,
because its abelian pro-$\ell$ source is finitely generated with
finite torsion. Since $|H^0|=\ell^a$, (E.33.5) bounds the orders
of $H^2$ independently of $a$. Inverse limits are exact on these
finite systems; the resulting $H^2(G_{K,S},\mathbb Z_\ell)$
is torsion, and tensoring with $\mathbb Q_\ell$ gives
\[
 H^2_{\rm cts}(G_{K,S},\mathbb Q_\ell)=0.                \tag{E.33.6}
\]
Continuous cochains with $\mathbb Q_\ell$ coefficients are the
union of the cochain complexes with values in
$\ell^{-b}\mathbb Z_\ell$, by compactness of their domain; this
justifies the last tensoring step. The real places at $\ell=2$
cause no exception: (E.33.5) is exactly Milne's degree-zero-through-two
formula, whose footnote explicitly treats this case.

For completeness, the universal unipotent completion has the same
$H^1$ and injects on $H^2$ into continuous group cohomology.
The first assertion follows from homomorphisms to the additive
group. For the second, an algebraic continuous Lie-cohomology
class factors through a finite-dimensional quotient. Represent it
by a central extension by $\mathbb G_a$; exponential coordinates
give a continuous set-theoretic section of the projection.
Its pullback along the dense continuous
group homomorphism is a continuous group extension. A zero pullback
class supplies a continuous lift; the universal property of
unipotent completion extends that lift algebraically. Its composite
with the projection is the identity on the dense original group,
hence the identity, so it splits the original extension. This
proves injectivity, with no comparison theorem left implicit.
Equation (E.33.6) therefore makes the second cohomology of the
completed Lie algebra zero.

Choose lifts of a basis of its abelianization $V_m$. They give a
surjection from $\widehat{\mathbb L}(V_m)$, with a closed kernel
$R$ contained in its commutator ideal. The free-Lie five-term
sequence identifies the continuous dual of
$R/\overline{[\widehat{\mathbb L}(V_m),R]}$ with degree-two cohomology.
A nonzero $R$ has a nonzero term of least bracket degree, which
cannot lie in $\overline{[\widehat{\mathbb L}(V_m),R]}$. Thus $R=0$,
proving (E.33.1).

For the more explicit module in (E.33.2), a unit in a CM field has
unit/conjugate-unit a root of unity. Its logarithm is therefore
fixed by $c$ and has trace zero. The injectivity already proved
and the equal dimensions $d_m/2-1$ show that the logarithmic image
is exactly the trace-zero subspace of
$(K_m\otimes\mathbb Q_\ell)^+$. The kernel and image of the
displayed formula are now as asserted. The scalar $1/2$ is
legitimate also over $\mathbb Q_2$.

Finally reciprocity sends the inclusion $G_{K_n}\to G_{K_m}$
to the local norm maps. Logarithms send norms to traces and send
norms of global units to the logarithms of global units. This
proves (E.33.3), including its compatibility with every tower.
Trace is onto on the displayed quotient: on an embedded element
of $K_m$ it is multiplication by the nonzero degree. A morphism
of prounipotent groups surjective on this finite-dimensional
abelianization is surjective, successively on the lower central
quotients. The transition maps are consequently faithfully flat.
To obtain the final presentation, choose generators at $m_1$.
At the next stage lift their logarithms through the surjection.
Choose a basis of the remaining kernel on abelianizations and
lift it into the Lie kernel: an arbitrary lift can be corrected
by a commutator, because the map on commutator ideals is onto.
Together these elements form a basis modulo commutators, and hence
a free set of topological Lie generators by the freeness already
proved. Iterating gives exactly the claimed deletion maps.

The accompanying exact cyclotomic-field program verifies the
conjugation projections, all divisor trace maps and their tower
compositions through conductor $24$, and prints the corresponding
free-Lie ranks through degree $8$. This checks the explicit maps,
not merely their dimensions.
\end{proof}

\IfFileExists{sections/E_intrinsic_positive.tex}{\subsection{A complete intrinsic family with full geometric pro-\texorpdfstring{$3$}{3} lifts}

\begin{statement}{E.34}{Proved}{Every cubic abelian class unramified outside three lifts fully}
Let $N$ be the free pro-$3$ geometric fundamental group of the
plus Fermat cubic with its nine cusps removed, of rank ten.
Use the rational tangential section above the base tangent
$t=0$, $dt=1$, specified by
$s=-x/z$, $t=s^3$, $y/z=-(1-s^3)^{1/3}$.
Write $V=N^{\mathrm{ab}}=T_3J_{E,D}$ and
$G_S=G_{\mathbb Q,\{3,\infty\}}$.
The genuine action on $N$ factors through $G_S$.
Here a section is called unramified outside three when its
geometric cocycle relative to this tangent factors through $G_S$.

The entire group of such geometrically abelian section classes is
\[
 H^1_{\mathrm{cts}}(G_S,V)
       \simeq \mathbb Z_3^6\oplus(\mathbb Z/3)^3 .
 \tag{E34.1}
\]
Every one of these classes lifts to a section of
$N\rtimes G_{\mathbb Q}$ whose geometric cocycle also factors
through $G_S$. In particular all of them lift to the intrinsic
metabelian, intrinsic class-two and marked-depth quotients.
The actual coordinate-unit map sends (E34.1) onto
\[
 \{(a\kappa_3,b\kappa_3):a,b\in\mathbb Z_3\}
       \simeq\mathbb Z_3^2
 \tag{E34.2}
\]
in the toric minimum. Its kernel is
$\mathbb Z_3^4\oplus(\mathbb Z/3)^3$, and every abelian class
in each of these fibers has a full geometric pro-$3$ lift.
These are assertions about the specified unramified family,
not about all the ramified classes of E.24.

\emph{Manuscript: \S51.16, folios 733v--736v, scans 172--178,
(371)--(379); this is an intrinsic rank-ten geometric pro-$3$
family compared through the actual coordinate-unit map.}
\end{statement}
\begin{proof}
\emph{The actual action and its relative pro-$3$ quotient.}
The projective Fermat cubic over $\mathbb Z[1/3]$ is smooth:
its three partial derivatives cannot vanish simultaneously.
Each of its three coordinate-cusp schemes is finite etale of
degree three over that ring, and they are pairwise disjoint.
The specified cusp and parameter are integral away from three.
For the based comparison, the displayed Puiseux formulas send
$t^{1/m}$ to $s^{3/m}$ and commute with coefficientwise
Galois action. The constant cubic roots are $-1$ and have
compatible rational $3$-power roots. Thus $N$ is the actual
index-nine subgroup of the base free pro-$3$ group $F$ and
its action is the restriction of the standard based action
on $F$, with no additional pro-$3$ Kummer twist.

The based unramifiedness used here is stated explicitly in
Hain--Matsumoto, \S8.5, printed p.~28 of the opened version
\cite{E_HainMatsumoto}: take $X_N=\mathbb P^1-\{0,\mu_N,\infty\}$,
$N=1$ and $\ell=3$ there. Smooth base change gives the
\emph{genuine} action on the fundamental group with tangent
$\overrightarrow{01}$ over $G_S$. Restricting to the above
subgroup proves the asserted unramifiedness for $N$.

Put $K=\mathbb Q(\zeta)$ and let $P$ be the maximal pro-$3$
quotient of $G_{K,\{3\}}$. Its defining kernel is stable under
$G_S$, so the relative quotient
\[
 \Gamma=G_S/\ker(G_{K,\{3\}}\longrightarrow P)
 \quad\text{fits into}\quad
 1\longrightarrow P\longrightarrow\Gamma
       \longrightarrow C_2\longrightarrow1 .
 \tag{E34.3}
\]
Complex conjugation splits this sequence.
Over $K$ the action on $F/\Phi(F)\simeq\mathbb F_3^2$ is
trivial, because $\mu_3\subset K$.
The kernel of
$\operatorname{Aut}(F)\to\operatorname{GL}(F/\Phi(F))$
is pro-$3$. One way to see this is to use the characteristic
lower exponent-$3$ series: successive automorphism kernels
embed in homomorphism groups between its finite
$\mathbb F_3$-graded pieces and are $3$-groups; the finite
Burnside basis theorem starts the induction.
The inverse limit is the asserted pro-$3$ kernel.
The action of $G_{K,\{3\}}$ on $F$, and hence on $N$, therefore
factors through $P$. The full action factors through (E34.3).

The same factorization holds for every $V$-cocycle on $G_S$.
Indeed the map $g\mapsto(\alpha(g),\rho(g))$ takes
$G_{K,\{3\}}$ into the pro-$3$ group
$V\rtimes\rho(G_{K,\{3\}})$. By the defining universal property
of $P$ it kills the kernel in (E34.3).
This proves that both the action and all the cocycles under
consideration factor through the same actual quotient.

\emph{The free Galois kernel and vanishing of obstructions.}
The norm-Euclidean ring $\mathbb Z[\zeta]$ is a PID, its units
are $\mu_6$, and it has the single prime
$(\lambda)$, $\lambda=1-\zeta$, above three.
Kummer theory gives
\[
 H^1(G_{K,\{3\}},\mathbb F_3)
  =\langle\zeta,\lambda\rangle\simeq\mathbb F_3^2 .
 \tag{E34.4}
\]
For completeness, an element defining a cubic extension
unramified away from three has valuation divisible by three
at every other prime. Unique factorization then writes its
class as $\zeta^a\lambda^b$ modulo cubes.
Conversely these classes define extensions unramified away
from three. Their independence follows first from the
$\lambda$ valuation and then from the unit group.

The global Euler characteristic of
Milne, I.5.1 \cite{E_MilneADT}, for this imaginary quadratic
field and the trivial module $\mathbb F_3$ is $1/3$.
The groups in degrees zero and one have orders three and
nine, respectively, by (E34.4). Therefore
\[
 H^2(G_{K,\{3\}},\mathbb F_3)=0 .
 \tag{E34.5}
\]
Inflation embeds $H^2(P,\mathbb F_3)$ into this group.
To verify that point, let $J$ be the kernel defining $P$.
An invariant nonzero character $J\to\mathbb F_3$ would
produce a quotient of $G_{K,\{3\}}$ which is an extension
of $P$ by a central group of order three. That quotient is
pro-$3$, contradicting maximality of $P$.
The five-term sequence therefore proves the claimed injection.

These calculations prove that $P$ is free pro-$3$ of rank two.
Here is a direct explanation of the criterion being used.
The dimension in (E34.4) gives a minimal surjection from a
free pro-$3$ group on two generators to $P$.
The vanishing of $H^2(P,\mathbb F_3)$ gives vanishing in
degree two for every finite $3$-primary $P$-module:
each such module has a filtration with trivial
$\mathbb F_3$ factors, since its action image is a finite
$3$-group. Every finite embedding problem with $3$-group
kernel is then soluble. Indeed use a nonzero characteristic
elementary abelian subgroup of the center of its kernel,
solve the quotient problem by induction, and split the
remaining abelian extension using this $H^2$ vanishing.
For a characteristic open subgroup $U$ of the minimal free
group $F_0$ and its surjection $f:F_0\to P$, the precise
finite-kernel problem is the pullback of
$F_0/U\to P/f(U)$ along $P\to P/f(U)$.
Its solution gives the required homomorphism $P\to F_0/U$.
Apply this to all those finite quotients.
There are finitely many homomorphisms to each finite target,
because $P$ has two generators; compactness gives a compatible
section of the minimal surjection. Its resulting idempotent
endomorphism of the free group is the identity on the
Frattini quotient, hence is an automorphism by the pro-$3$
basis theorem. An idempotent automorphism is the identity,
so the minimal surjection is an isomorphism.

In particular $P$ has $3$-cohomological dimension one.
The Hochschild--Serre sequence for (E34.3), with averaging
over $C_2$ (two is invertible on every $3$-primary module),
gives
\[
 H^2(\Gamma,M)=0
 \quad\text{for every finite $3$-primary $\Gamma$-module }M.
 \tag{E34.6}
\]
The same conclusion holds over $G_S$ for every finite
module whose restriction to $G_{K,\{3\}}$ has $3$-group
action image: filter the restriction by trivial
$\mathbb F_3$ factors, apply (E34.5), and again take
$C_2$ invariants.

Now take any cocycle $\alpha:\Gamma\to V$.
Choose a decreasing characteristic series of $N$ with
finite $3$-group quotients $N_j$ and intersection one.
At each level, pull the surjection
$N_j\rtimes\Gamma\to N_j^{\mathrm{ab}}\rtimes\Gamma$
back along the reduction of $\alpha$.
Its kernel is a finite $3$-group. The same central-subgroup
induction using (E34.6) supplies a section.
There are finitely many choices at every level:
$\Gamma$ has the two generators of $P$ and complex
conjugation, and the geometric target is finite.
The compatible-section sets form a nonempty inverse limit.
Explicitly, every finite collection of levels is supplied
by a section at their largest common refinement; compactness
of the product of the finite sets gives a coherent choice.
The limit is a continuous $N$-cocycle lifting $\alpha$.
Inflating along $G_{\mathbb Q}\to G_S\to\Gamma$ gives
the asserted section of $N\rtimes G_{\mathbb Q}$.
This constructs lifts in the full geometric pro-$3$ group,
not merely compatible names for degree-two obstruction classes.

\emph{Counting the entire unramified abelian family.}
The genuine action on $V/3^nV$ has $3$-group image over
$G_{K,\{3\}}$, by the action factorization already proved.
The last assertion following (E34.6) gives
$H^2(G_S,V/3^nV)=0$ for every $n$.
These finite-coefficient degree-one groups are finite
(Milne I.4.15); hence the inverse-limit sequence also gives
$H^2_{\mathrm{cts}}(G_S,V)=0$.
Consequently
\[
 H^1_{\mathrm{cts}}(G_S,V)/3^n
       \simeq H^1(G_S,V/3^nV).
 \tag{E34.7}
\]

Complex conjugation has four plus eigenvalues and six minus
eigenvalues on $V$. To compute them, it fixes the three
rational cusps and interchanges three pairs among the other
six cusps. Its permutation action on degree-zero cusp
divisors therefore has multiplicities five and three.
The toric Tate twist reverses their signs, giving
multiplicities three and five on the rank-eight toric part.
On $T_3E$, its order is two and its determinant is $-1$ by
the Weil pairing, so its two eigenvalues are $1,-1$.
Adding the exact toric and elliptic pieces gives the claimed
multiplicities. Since two is a unit in $\mathbb Z_3$, the
same ranks hold for the integral eigensublattices.

By E.24 the rational $3$-primary torsion of $G=J_{E,D}$ is
$(\mathbb Z/3)^3$. Thus
$H^0(G_S,V/3^nV)=G(\mathbb Q)[3^n]$ has order $27$
for every $n\ge1$.
Milne's Euler formula I.5.1, now over $\mathbb Q$, has
real factor $3^{4n}/3^{10n}=3^{-6n}$.
Together with the preceding $H^2$ vanishing it yields
\[
 |H^1(G_S,V/3^nV)|=3^{6n+3}.
 \tag{E34.8}
\]
The continuous degree-one group is a finitely generated
$\mathbb Z_3$-module: its cocycles factor through the
three-generated group $\Gamma$ and form a closed
submodule of $V^3$, modulo the closed coboundary image.
Equations (E34.7)--(E34.8) give rank six and torsion of
order $27$ and exponent three. Equivalently that torsion is
$H^0(G_{\mathbb Q},V\otimes\mathbb Q_3/V)
=G(\mathbb Q)[3^\infty]$.
This proves (E34.1).

Finally the algebraic retraction of E.24 gives a split map
of these Galois modules to $\mathbb Z_3(1)^2$.
Kummer theory over $\mathbb Z[1/3]$ gives
$H^1(G_S,\mathbb Z_3(1))=\mathbb Z_3\kappa_3$.
Its rank-two target is torsion-free, and the splitting
therefore proves (E34.2) and the stated rank-four kernel
with its three torsion factors.
The change from the origin at $O$ in E.24 to the present
rational cusp is also unramified outside three. The framed
cusp-difference relation (E24.6) involves only the units
$\lambda$, powers of $3$ and signs, so this framed point
extends over $\mathbb Z[1/3]$ and its Kummer cocycle is
unramified there. Equivalently both integral tangential
fiber functors give the same unramified comparison.
Both origins have zero toric pro-$3$ class.
Thus changing origins preserves the unramified family
and the displayed linear toric map.
The matrix and order calculations are recorded in
\texttt{e\_intrinsic\_positive.py}; the Galois vanishing
and lifting arguments above establish the infinite assertions.
\end{proof}
}{}
\IfFileExists{sections/E_intrinsic_finite.tex}{\subsection{All sections of the intrinsic Heisenberg quotient}

\begin{statement}{E.35}{Proved}{A norm equation and every conjugacy class in the finite intrinsic quotient}
Let $K=\mathbb Q(\zeta)$ and retain the quotient $H_3$ with
the based $G_K$-action (E31.3).
Write $\mathcal A=K^\times/K^{\times3}$ multiplicatively.
For an abelian section pair $(a,b)\in\mathcal A^2$ put
$A=\zeta a$, choose nonzero representatives $A,b\in K$,
and form the etale cubic algebra
$L_A=K[t]/(t^3-A)$.
The fiber of section conjugacy classes of
$H_3\rtimes G_K$ over $(a,b)$ is nonempty exactly when
the explicit equation
\[
 X^3+A Y^3+A^2 Z^3-3AXYZ=b
                 \quad\text{has a solution }X,Y,Z\in K .
 \tag{E35.1}
\]
Every nonempty fiber is a torsor under the explicitly
specified group
\[
 \mathcal A/\langle A,b\rangle
       =K^\times/(K^{\times3}\langle A,b\rangle).
 \tag{E35.2}
\]
The subgroup in the denominator has rank zero, one or two
over $\mathbb F_3$, according to the span of the two
indicated cube classes.
There are concrete coordinates for this group. Choose one
generator $\pi_{\mathfrak p}$ for each nonzero prime ideal of
$\mathbb Z[\zeta]$, and put
$J=\{\zeta\}\amalg\{\mathfrak p\}$.
Unique factorization gives the explicit isomorphism
\[
 \mathcal A\longrightarrow\bigoplus_{j\in J}\mathbb F_3,
 \qquad
 [\zeta^e\prod_{\mathfrak p}\pi_{\mathfrak p}^{e_{\mathfrak p}}]
       \longmapsto(e,(e_{\mathfrak p}))\pmod3 .
 \tag{E35.2a}
\]
The sign in the factorization is a cube and is omitted.
Thus (E35.2) is obtained by Gaussian elimination of the
two finite-support vectors $[A],[b]$: choose zero, one or
two pivot coordinates according to their rank, subtract
the corresponding row-reduced basis vectors, and retain
all remaining coordinates. This is an explicit
parameterization by the direct sum of $\mathbb F_3$ over
the coordinates in $J$ outside those pivots.
After choosing the lift $n_0$ constructed from a norm solution,
all classes are represented by
$n_d(g)=(0,0,\kappa_3(d)(g))n_0(g)$ with
$d\in K^\times/(K^{\times3}\langle A,b\rangle)$.

All dependent cases are included. If $A$ is a cube,
(E35.1) is always soluble and the fiber group is
$\mathcal A/\langle b\rangle$.
If $A$ is not a cube and $b$ is in its span, (E35.1) is
always soluble and the fiber group is
$\mathcal A/\langle A\rangle$.
If both classes are trivial the group is $\mathcal A$.
For two independent classes, a solution of (E35.1)
constructs an actual Heisenberg Galois extension as in the proof.

The quotient is the one obtained from the actual map (E31.4)
and the based arithmetic comparison of E.31.
For example the pair $(a,b)=(1,7)$ has empty fiber,
while $(a,b)=(2\zeta^2,3)$ has nonempty fiber, with
norm solution $1+\sqrt[3]{2}$.
No assertion is made that every pair here is the image of
a class in the full intrinsic rank-ten abelian section group.

\emph{Manuscript: \S51.16, folios 733v--736v, scans 172--178,
(371)--(379); this computes a specified finite intrinsic
quotient reached by the additional morphism (E31.4).}
\end{statement}
\begin{proof}
Choose the path in E.31 within the formal punctured disc
at the target zero. Equation (E31.7) then gives the image
of the source tangential section as
\[
 q_g=(\kappa_3(\zeta)(g),0,0)\in H_3 .
 \tag{E35.3}
\]
To see the zero central coordinate directly, write
$h=\zeta\lambda^{-3}s^3u(s)$ with $u(0)=1$.
Its chosen cube root is
$(\zeta_9/\lambda)s\,u(s)^{1/3}$.
The last factor has coefficients in $K$, so the only
Galois change is the inertia power supplied by
$\zeta_9$. The formal puncture has only that one
geometric generator and gives no commutator term.
Thus $q:G_K\to H_3$ is the homomorphism in (E35.3),
and the intrinsic action is $\operatorname{Ad}(q_g)$.

There is an explicit isomorphism of arithmetic groups
\[
 H_3\rtimes_{\operatorname{Ad}q}G_K
       \longrightarrow H_3\times G_K,\qquad
 (n,g)\longmapsto(nq_g,g).
 \tag{E35.4}
\]
Its inverse sends $(h,g)$ to $(hq_g^{-1},g)$.
It is the identity on the geometric group.
Consequently it identifies section classes with continuous
homomorphisms $G_K\to H_3$ modulo conjugation by $H_3$,
and sends the intrinsic abelian pair $(a,b)$ to the
ordinary Kummer pair $(A,b)=(\zeta a,b)$.

Write these latter Kummer characters as $u,v:G_K\to\mathbb F_3$.
A homomorphism with those abelian coordinates has the form
$(u_g,v_g,w_g)$, where its central coordinate solves
\[
 d w(g,h)=-\tfrac12(u_gv_h-v_gu_h).
 \tag{E35.5}
\]
The two-cocycle on the right, with sign reversed, differs
from $u_gv_h$ by
$d(u_gv_g/2)$; hence its obstruction is exactly
$\kappa_3(A)\cup\kappa_3(b)$.
The Kummer exact sequence and Hilbert 90 embed this group
of classes into $\operatorname{Br}(K)[3]$.

We spell out the resulting norm condition.
Suppose first that $A$ is not a cube, let
$\alpha^3=A$ and let $\sigma(\alpha)=\zeta\alpha$ generate
$\operatorname{Gal}(L_A/K)$.
The cyclic algebra is
\[
 \mathcal C=L_A\oplus L_A U\oplus L_A U^2,\qquad
 U\ell=\sigma(\ell)U,\quad U^3=b .
 \tag{E35.6}
\]
Its factor set is
$b^{\lfloor(i+j)/3\rfloor}$ on $\sigma^i,\sigma^j$.
Choose $\beta^3=b$ and represent the values of $u$ by
$0,1,2$. The multiplicative coboundary of
$g\mapsto\beta^{u_g}$ is
\[
 b^{\lfloor(u_g+u_h)/3\rfloor}\zeta^{v_gu_h}.
 \tag{E35.7}
\]
The identity for $d(u_gv_g)$ in (E35.5) shows that the
class of this factor set is precisely
$\kappa_3(A)\cup\kappa_3(b)$.
The crossed-product comparison with the Brauer group is
Milne, \emph{Class Field Theory}, IV.3.14
\cite{E_MilneCFT}; its construction uses exactly the
multiplication and factor set just displayed.

The algebra (E35.6) is split exactly when $b$ is a norm.
This equivalence also has a direct proof here.
A norm solution $b=N_{L_A/K}(w)$ represents $\mathcal C$
on the three-dimensional $K$-space $L_A$, with $\ell$
acting by multiplication and $U$ acting as $w\sigma$.
The relations hold and the resulting homomorphism is an
isomorphism with $\operatorname{End}_K(L_A)$, by central
simplicity and the equal dimension nine.
Conversely, a splitting realizes $\mathcal C$ on a
three-dimensional $K$-space which, by its embedded
degree-three field $L_A$, is one-dimensional over $L_A$.
The operator $U$ is necessarily $w\sigma$ in a choice
of $L_A$-basis. Its cube is $N(w)$, so $b=N(w)$.
The norm of $X+Y\alpha+Z\alpha^2$, by its multiplication
matrix on this same basis, is
$X^3+A Y^3+A^2Z^3-3AXYZ$.
This proves (E35.1).

For independent cube classes $A,b$ the norm solution
also yields an explicit finite Galois realization.
Write $w=X+Y\alpha+Z\alpha^2$ with $N(w)=b$ and put
\[
 \beta^3=b,\qquad
 v=w^2\sigma(w),\qquad \gamma^3=v.
 \tag{E35.8}
\]
The field $M=L_A(\beta,\gamma)$ has the automorphisms
\[
 \begin{array}{c|ccc}
       &\alpha&\beta&\gamma\\ \hline
 \widetilde\sigma&\zeta\alpha&\beta&\beta\gamma/w\\
 \tau&\alpha&\zeta\beta&\gamma .
 \end{array}
 \tag{E35.9}
\]
Indeed $\sigma(v)/v=b/w^3$, so these formulas preserve
the defining equations, and
$\widetilde\sigma^3(\gamma)
=b\gamma/N(w)=\gamma$.
Both displayed automorphisms have order three.
The root $\gamma$ does not already lie in $L_A(\beta)$.
Otherwise the fact that its cube is in $L_A$ would write
$\gamma=c\beta^j$ for $c\in L_A$.
The equation $\sigma(v)/v=b/w^3$ would then imply that
$b$ is a cube in $L_A$, contradicting the independence
of $A,b$.
Thus $[M:K]=27$, the formulas make $M/K$ Galois, and
their commutator acts on $\gamma$ as multiplication by
$\zeta^{-1}$. Taking that automorphism as the central
generator identifies its Galois group with (E31.2),
with the two stated Kummer characters. Composing the
Galois restriction map with this identification and then
(E35.4)$^{-1}$ gives an actual arithmetic section.

The dependent cases require no degree-27 field.
If $A$ is a cube, replace its representative by one.
Then $L_A\simeq K^3$ and the norm is the product of its
three coordinates, so it is surjective. In (E35.1) the
choice
$X=(b+2)/3$, $Y=Z=(b-1)/3$ gives norm $b$.
The homomorphism $g\mapsto(0,\kappa_3(b)(g),0)$
provides a lift.
If $A$ is nontrivial and $b=A^j c^3$ for $j=0,1,2$,
the element $w=c\alpha^j$ has norm $b$.
The homomorphism
$g\mapsto(\kappa_3(A)(g),j\kappa_3(A)(g),0)$
has the required coordinates; its alternating term is
zero. This includes $b=1$ and all other dependent
cases. The construction is independent of the selected
representatives modulo cubes.

Finally choose one lift $h_g$ of $(u,v)$.
Every other lift is uniquely $z_g h_g$ with
$z\in H^1(G_K,\mathbb F_3)=\mathcal A$.
Conjugating by an element with abelian coordinates
$(r,s)$ adds the central character
$r v-s u$. Central conjugations do nothing.
The exact subgroup of differences generated by all
geometric conjugations is therefore
$\langle\kappa_3(A),\kappa_3(b)\rangle$.
This proves the entire torsor description (E35.2),
including each of the dependent cases just treated.
For completeness, the coordinate computation (E35.2a)
uses the Euclidean norm
$N(a+b\zeta)=a^2-ab+b^2$ on $\mathbb Z[\zeta]$.
Rounding the two rational coefficients of any quotient
to nearest integers leaves errors $r,s$ of absolute
value at most $1/2$, and
$r^2-rs+s^2\leq3/4<1$. The division algorithm proves
that this ring is a principal ideal domain and a unique
factorization domain. Its norm-one units are
$\{\pm1,\pm\zeta,\pm\zeta^2\}$, so their quotient by
cubes is generated freely over $\mathbb F_3$ by $\zeta$.
Factoring the numerator and denominator of an element
of $K^\times$ now proves both the injectivity and
surjectivity of (E35.2a). Row reduction gives an actual
isomorphism from its quotient by $[A],[b]$ to the stated
direct sum: normalize the pivot entries to one and all
other pivot entries to zero, then send a vector to its
unique representative with every pivot entry zero.
This also describes the inverse by inclusion of those
nonpivot coordinates.

At $(a,b)=(1,7)$, equation (E35.1) would assert that
$7$ is a norm from $K(\zeta_9)$, contradicted by the
local calculation at $(3+\zeta)$ in E.31.
For $A=2,b=3$, the element $1+\alpha$ has norm
$1+\alpha^3=3$, proving the positive example.
The program \texttt{e\_intrinsic\_finite.py} checks the
norm polynomial, this solution and the relations in
(E35.8)--(E35.9), and enumerates every conjugation
orbit for all two-generator assignments in $H_3$.
The finite intrinsic quotient map of E.31 remains the
comparison to the full cubic problem; no inverse to
that quotient map is asserted.
\end{proof}
}{}
\subsection{Every fiber of the cuspidally central cubic quotient}

\begin{statement}{E.36}{Proved}{A class-two quotient retaining all ten abelian coordinates, with every fiber evaluated}
Let $E=X_3^+$, let $D$ be its nine coordinate cusps, and put
$O=[1:-1:0]$. Use the tangential base point at $O$ with
parameter $s=z/x$ and tangent one, both on $U_3=E\setminus D$
and on $E\setminus\{O\}$.
Write
\[
 N=\pi_1(U_{3,\overline{\mathbb Q}})^{(3)},\qquad
 V=N^{\rm ab}=T_3J_{E,D},\qquad T=T_3E,
 \qquad I=\ker(N\longrightarrow T).
\]
Closures are understood in commutator subgroups. Set
\[
 N^{\rm cc}=N/[N,I],\qquad
 C=\pi_1((E\setminus\{O\})_{\overline{\mathbb Q}})^{(3)}
                   /\gamma_3 .
\]
The maps induced by the open immersion
$U_3\hookrightarrow E\setminus\{O\}$ and by abelianization
give an actual, based, $G_{\mathbb Q}$-equivariant isomorphism
\[
 N^{\rm cc}\ \xrightarrow{\ \sim\ }\ C\times_T V,
       \qquad n\longmapsto(q_O(n),n^{\rm ab}).
 \tag{E36.1}
\]
In particular the quotient retains all ten coordinates of
$V$ and its nonzero elliptic commutator. The exact quotient
from the full intrinsic class-two group is
\[
 1\longrightarrow
 \ker\!\left(\bigwedge\nolimits^2 V
           \longrightarrow\bigwedge\nolimits^2 T\right)
 \longrightarrow N/\gamma_3N
 \longrightarrow N^{\rm cc}\longrightarrow1 .
 \tag{E36.2}
\]
The displayed kernel is a free $\mathbb Z_3$-module of rank
$44$, with the Galois action induced on the exterior square;
the remaining central commutator module is
$\bigwedge^2T\simeq\mathbb Z_3(1)$.

Every geometric-abelian section class
$\alpha\in H^1(\mathbb Q,V)$ lifts to a section of
$N^{\rm cc}\rtimes G_{\mathbb Q}$.
For every such $\alpha$, the complete fiber of geometric
conjugacy classes is a torsor under
\[
 \mathcal R_{\mathbb P}(\mathbb Z_3)
 =\left\{(a_p)_{p\ {
m rational\ prime}}\in
                \prod_p\mathbb Z_3:
       \substack{\text{for every }m,\ a_p\in3^m\mathbb Z_3\\
                         \text{ for all but finitely many }p}\right\}.
 \tag{E36.3}
\]
The parameter $a=(a_p)$ acts by the central Kummer cocycle
of the compatible classes $\prod_p p^{a_p}$ modulo
$\mathbb Q^{\times3^m}$. These finite products modulo each
$3^m$ define the indicated element of the completion for
multiplication by powers of three.
Thus this is a computation of every fiber over the entire
rank-ten abelian section group evaluated in E.24, including
its ramified classes. The classes obstructed in E.31 also
lift to (E36.1), while their further lifting through (E36.2)
is obstructed by the explicit quotient of E.31.

\emph{Manuscript: \S51.16, folios 733v--736v, scans 172--178,
(371)--(379). The quotient maps (E36.1)--(E36.2) specify
exactly the intrinsic cubic quotient evaluated here.}
\end{statement}
\begin{proof}
First compute the geometric maps integrally. An oriented
presentation of the punctured complex elliptic curve is
\[
 N=\langle a,b,c_O,c_1,\ldots,c_8\mid
              [a,b]c_Oc_1\cdots c_8=1\rangle_{\rm pro\text{-}3}.
 \tag{E36.4}
\]
Here $a,b$ map to a basis of $T$; the $c_d$ are the
peripheral loops. Eliminating $c_O$ gives the free pro-three
group of rank ten. The kernel $I$ is normally generated by
these peripheral loops. Passing to $N/[N,I]$ makes all
$c_d$ central and leaves only the indicated relation.
Consequently the geometric cuspidal kernel is
\[
 I/[N,I]\simeq\mathbb Z_3(1)[D],
 \qquad [a,b]=-\sum_{d\in D}c_d,
 \tag{E36.5}
\]
with its actual Galois action: the cusps are permuted and
each inertia generator has the cyclotomic twist. This action
can also be obtained from the torus
$\operatorname{Res}_{D/\mathbb Q}\mathbb G_m$.
The central quotient is the one defined in
Borne--Emsalem--Stix, \S3.2; their Proposition 3.5 identifies
it with the fundamental group of their torus torsor by
the canonical map from the open curve
\cite{E_BorneEmsalemStix}.

For an entirely explicit check of (E36.5), take the
two-step Lie module with basis
$a,b,c_1,\ldots,c_8,k$ and only bracket $[a,b]=k$.
The law $u*v=u+v+[u,v]/2$ is integral over $\mathbb Z_3$.
Sending $c_O$ to $-k-\sum_{i=1}^8c_i$ satisfies (E36.4).
Conversely its generators and brackets give precisely the
relations of $N/[N,I]$. The two maps are inverse, so there
are no further relations or torsion in the nine-dimensional
central lattice of (E36.5).

The map $q_O$ fills the other eight punctures. On (E36.5)
it is projection to the $O$-coordinate. Abelianization
quotients (E36.5) by the diagonal element
$\delta=\sum_{d\in D}c_d$. The maps on the two geometric
cuspidal kernels are therefore
\[
 \mathbb Z_3(1)[D]\longrightarrow
 \mathbb Z_3(1)\oplus
       \bigl(\mathbb Z_3(1)[D]/\mathbb Z_3(1)\delta\bigr),
 \quad t\longmapsto(t_O,\overline t).
 \tag{E36.6}
\]
This is an isomorphism: in coordinates it sends
$(t_O,t_1,\ldots,t_8)$ to
$(t_O,t_1-t_O,\ldots,t_8-t_O)$, with inverse
$(u,v_1,\ldots,v_8)\mapsto(u,v_1+u,\ldots,v_8+u)$.
It is $G_{\mathbb Q}$-equivariant because $O$ is rational
and the diagonal element is invariant apart from its
cyclotomic twist. Both sides of (E36.1) map onto $T$,
and (E36.6) identifies their kernels. This proves (E36.1).
The isomorphism is induced by the stated morphisms of
curves and abelianization, with the same tangent throughout.
It therefore identifies the actual arithmetic extensions,
not just their abstract geometric groups. There is no
change of section origin in (E36.1).

The commutator module of the free class-two pro-three
group on $V$ is $\bigwedge^2V$. The bracket surviving in
(E36.5) is the exterior-square map induced by $V\to T$.
In the basis from (E36.4) it sends $a\wedge b$ to one
generator and sends the other $44$ basic wedges to zero.
Its kernel is therefore exactly $[N,I]/\gamma_3N$,
free of rank $\binom{10}{2}-1=44$; its image has index one.
The orientation identifies the image with
$\bigwedge^2T=\mathbb Z_3(1)$, compatibly with the Weil
pairing. Functoriality proves the asserted arithmetic
action and (E36.2).

We next establish the existence of all the sections in
the statement. Use the rational point $O$ as origin on
the proper elliptic curve. Its full geometric fundamental
group is
$T_\bullet E=\prod_\ell T_\ell E$, and its arithmetic
fundamental group is the corresponding semidirect product
with $G_{\mathbb Q}$. For a given $\alpha$, let
$\beta\in H^1(\mathbb Q,T_3E)$ be its image, and choose
compatible cocycle representatives. There is a specified
full proper-curve section
\[
 g\longmapsto\bigl((\beta(g),0,0,\ldots),g\bigr),
 \tag{E36.7}
\]
where every component at a prime other than three is zero.
The inclusion of the three-primary factor is continuous
and $G_{\mathbb Q}$-equivariant, so (E36.7) is an actual
section of the full arithmetic fundamental group of $E$.

Borne--Emsalem--Stix, Theorem 5.10, states that for a
smooth projective positive-genus curve over $\mathbb Q$
and a boundary which is a union of torsion sub-packets,
every section lifts to its cuspidally central quotient
\cite{E_BorneEmsalemStix}. Apply the theorem to $E$ and
the singleton boundary $\{O\}$. This is a torsion
sub-packet by their Definition 2.3: the only difference
of two points in its support is zero in the Picard group.
All hypotheses are now satisfied. The cuspidal kernel
for a once-punctured elliptic curve is the geometric
commutator subgroup; making it geometrically central
is precisely quotienting by $\gamma_3$. Projecting the
resulting full lift of (E36.7) to its geometric pro-three
quotient gives a section $c_\beta$ of $C\rtimes G_{\mathbb Q}$
above $\beta$.

If this section projects to a representative conjugate
to the chosen $\beta$, conjugate it by a geometric lift
of the relevant element of $T$. Its projection then
equals the chosen representative. Pairing it with the
given abelian section $\alpha$ gives a section of the
fiber product (E36.1). This proves existence for every
$\alpha$, without imposing any ramification condition.

Finally the kernel of $N^{\rm cc}\to V$ is the central
copy of $\mathbb Z_3(1)$ generated by the elliptic
commutator. After choosing one lift of a fixed
representative of $\alpha$, the other lifts differ by
continuous one-cocycles with values in this module.
Central geometric conjugation changes the cocycle by
an ordinary coboundary. There is no additional
identification from geometric elements with nonzero
image in $V$: such an element could fix the prescribed
abelian cocycle only if its image lay in
$H^0(\mathbb Q,V)$, and this group is zero.
Here is a direct verification. In the exact sequence
$0\to T_D\to J_{E,D}\to E\to0$, the rational torsion
of $T_D\simeq\mathbb G_m^2\times
(\operatorname{Res}_{K/\mathbb Q}\mathbb G_m)^3$
is finite, being a product of groups of roots of unity.
The group $E(\mathbb Q)$ has order three by E.24.
Thus $J_{E,D}(\mathbb Q)[3^\infty]$ is finite.
An invariant in $T_3J_{E,D}$ is a compatible sequence
in this finite three-primary group under multiplication
by three. Every such sequence is zero. Hence
$H^0(\mathbb Q,V)=0$ as asserted.

The exact fiber is consequently a torsor under
$H^1(\mathbb Q,\mathbb Z_3(1))$.
Kummer theory identifies this group with
$\varprojlim_m\mathbb Q^\times/\mathbb Q^{\times3^m}$.
Unique factorization in $\mathbb Q^\times$ gives the
isomorphism with (E36.3): the sign is a cube, each
prime valuation is reduced modulo $3^m$, and every
such residue vector has finite support. The inverse
constructs the compatible Kummer classes from precisely
these finite-support residue vectors. This proves the
complete coordinate description of all fibers.

The matrix in (E36.6) has determinant one, and the
exterior-square map in (E36.2) has Smith diagonal one
and kernel rank $44$. The program
\texttt{e\_cuspidally\_central.py} verifies these integral
lattice computations. The existence theorem and the
arithmetic fiber calculation above establish the
infinite section assertions.
\end{proof}


\clearpage
\bibliographystyle{alpha}
\bibliography{LM7_FERMAT_TOWER_references}
\end{document}
